\documentclass{amsart}
\usepackage[margin=1.5in]{geometry} 
\usepackage{amsmath}
\usepackage{tcolorbox}
\usepackage{amssymb}
\usepackage{amsthm}
\usepackage{lastpage}
\usepackage{fancyhdr}
\usepackage{accents}
\usepackage{xcolor}
\usepackage{color}
\usepackage{euscript}
\usepackage[T1]{fontenc}
\usepackage[utf8]{inputenc}
\usepackage{microtype}
\emergencystretch=1em
\usepackage[bbgreekl]{mathbbol}
\DeclareSymbolFontAlphabet{\mathbb}{AMSb} % to ensure \mathbb does not change
\DeclareMathAlphabet{\mathdutchcal}{U}{dutchcal}{m}{n} 
\DeclareSymbolFontAlphabet{\mathbbl}{bbold}
% Fields
\newcommand{\A}{\mathbb{A}}
\newcommand{\CC}{\mathbb{C}}
\newcommand{\EE}{\mathbb{E}}
\newcommand{\FF}{\mathbb{F}}
\newcommand{\RR}{\mathbb{R}}
\newcommand{\QQ}{\mathbb{Q}}
\newcommand{\ZZ}{\mathbb{Z}}
\newcommand{\NN}{\mathbb{N}}
\newcommand{\LL}{\mathbb{L}}

% mathcal letters
\newcommand{\Acal}{\mathcal{A}}
\newcommand{\Bcal}{\mathcal{B}}
\newcommand{\Ccal}{\mathcal{C}}
\newcommand{\Dcal}{\mathcal{D}}
\newcommand{\Ecal}{\mathcal{E}}
\newcommand{\Fcal}{\mathcal{F}}
\newcommand{\Gcal}{\mathcal{G}}
\newcommand{\Hcal}{\mathcal{H}}
\newcommand{\Ical}{\mathcal{I}}
\newcommand{\Jcal}{\mathcal{J}}
\newcommand{\Kcal}{\mathcal{K}}
\newcommand{\Lcal}{\mathcal{L}}
\newcommand{\Mcal}{\mathcal{M}}
\newcommand{\Ncal}{\mathcal{N}}
\newcommand{\Ocal}{\mathcal{O}}
\newcommand{\Pcal}{\mathcal{P}}
\newcommand{\Qcal}{\mathcal{Q}}
\newcommand{\Rcal}{\mathcal{R}}
\newcommand{\Scal}{\mathcal{S}}
\newcommand{\Tcal}{\mathcal{T}}
\newcommand{\Ucal}{\mathcal{U}}
\newcommand{\Vcal}{\mathcal{V}}
\newcommand{\Wcal}{\mathcal{W}}
\newcommand{\Xcal}{\mathcal{X}}
\newcommand{\Ycal}{\mathcal{Y}}
\newcommand{\Zcal}{\mathcal{Z}}

% abstract categories
\newcommand{\Asf}{\mathsf{A}}
\newcommand{\Bsf}{\mathsf{B}}
\newcommand{\Csf}{\mathsf{C}}
\newcommand{\Dsf}{\mathsf{D}}
\newcommand{\Fsf}{\mathsf{F}}
\newcommand{\hsf}{\mathsf{h}}
\newcommand{\Rsf}{\mathsf{R}}
\newcommand{\Ssf}{\mathsf{S}}
\newcommand{\Tsf}{\mathsf{T}}
\newcommand{\Lsf}{\mathsf{L}}
\newcommand{\ksf}{\mathsf{k}}
\newcommand{\ysf}{\mathsf{y}}

% infinity categories
\newcommand{\Ascr}{\EuScript{A}}
\newcommand{\Bscr}{\EuScript{B}}
\newcommand{\Cscr}{\EuScript{C}}
\newcommand{\Dscr}{\EuScript{D}}
\newcommand{\Escr}{\EuScript{E}}
\newcommand{\Iscr}{\EuScript{I}}
\newcommand{\Kscr}{\EuScript{K}}
\newcommand{\Lscr}{\EuScript{L}}
\newcommand{\Mscr}{\EuScript{M}}
\newcommand{\Nscr}{\EuScript{N}}
\newcommand{\Vscr}{\EuScript{V}}
\newcommand{\Xscr}{\EuScript{X}}

% algebraic geometry
\newcommand{\spec}{\operatorname{Spec}}
\newcommand{\proj}{\operatorname{Proj}}

% categories 
\newcommand{\id}{\mathrm{id}}
\newcommand{\Obj}{\mathrm{Obj}}
\newcommand{\Mor}{\mathrm{Mor}}
\newcommand{\Hom}{\mathrm{Hom}}
\newcommand{\Aut}{\mathrm{Aut}}
\newcommand{\Sets}{\mathsf{Sets}}
\newcommand{\SSets}{\mathsf{SSets}}
\newcommand{\kVec}{\mathsf{Vec}_{k}}
\newcommand{\Alg}{\mathsf{Alg}}
\newcommand{\Ring}{\mathsf{Ring}}
\newcommand{\Mod}{\mathsf{Mod}}
\newcommand{\Grp}{\mathsf{Grp}}
\newcommand{\AbGrp}{\mathsf{AbGrp}}
\newcommand{\PSh}{\mathsf{PShv}}
\newcommand{\Sh}{\mathsf{Shv}}
\newcommand{\PSch}{\mathsf{PSch}}
\newcommand{\Sch}{\mathsf{Sch}}
\newcommand{\Top}{\mathsf{Top}}
\newcommand{\Com}{\mathsf{Com}}
\newcommand{\Coh}{\mathsf{Coh}}
\newcommand{\QCoh}{\mathsf{QCoh}}
\newcommand{\Opens}{\mathsf{Opens}}
\newcommand{\Opp}{\mathsf{Opp}}
\newcommand{\Cat}{\mathsf{Cat}}
\DeclareMathOperator*{\colim}{colim}
\newcommand{\Grpd}{\mathsf{Grpd}}
\newcommand{\Fun}{\mathrm{Fun}}
\newcommand{\cofib}{\mathrm{cofib}}
\newcommand{\fib}{\mathrm{fib}}

% simplicial sets
\newcommand{\DDelta}{\Updelta}
\newcommand{\Sing}{\operatorname{Sing}}

% condensed math
\newcommand{\LCA}{\mathsf{LCA}}
\newcommand{\Cond}{\mathsf{Cond}}
\newcommand{\dom}{\operatorname{dom}}
\newcommand{\Cov}{\operatorname{Cov}}
\newcommand{\proet}{\mathsf{pro}\mathsf{\acute{e}}\mathsf{t}}
\newcommand{\et}{\mathsf{\acute{e}t}}
\newcommand{\Ab}{\mathsf{Ab}}
\newcommand{\ProFin}{\mathsf{ProFin}}
\newcommand{\Prof}{\mathsf{Prof}}
\newcommand{\PFin}{\mathsf{PFin}}
\newcommand{\PAni}{\mathsf{PAni}}
\newcommand{\Pro}{\mathsf{Pro}}
\newcommand{\Fin}{\mathsf{Fin}}
\newcommand{\ExtDiscHaus}{\mathsf{ExtDiscHaus}}
\newcommand{\CHaus}{\mathsf{CHaus}}

% K-theory
\newcommand{\St}{\mathsf{St}}
\newcommand{\Ind}{\mathsf{Ind}}
\newcommand{\Sp}{\mathsf{Sp}}
\newcommand{\Loc}{\mathsf{Loc}}
\newcommand{\ev}{\mathrm{ev}}
\newcommand{\coev}{\mathrm{coev}}
\newcommand{\dual}{\mathsf{dual}}


\newcommand{\Aff}{\mathsf{Aff}}
\newcommand{\Ani}{\mathsf{Ani}}
\newcommand{\rex}{\mathsf{rex}}
\newcommand{\lex}{\mathsf{lex}}
\newcommand{\st}{\mathsf{st}}
\newcommand{\PSp}{\mathsf{PSp}}
\newcommand{\const}{\mathrm{const}}
\newcommand{\Fil}{\mathsf{Fil}}
\newcommand{\gr}{\mathrm{gr}}
\newcommand{\Ch}{\mathrm{Ch}}
\newcommand{\Dec}{\mathrm{D\acute{e}c}}
\newcommand{\img}{\mathrm{im}}
\newcommand{\intHom}{\underline{\Hom}}
\newcommand{\PrL}{\mathsf{Pr}^{\mathsf{L}}}
\newcommand{\CAlg}{\mathsf{CAlg}}
\newcommand{\triv}{\mathsf{triv}}
\newcommand{\Sph}{\mathbb{S}}
\newcommand{\Cont}{\mathrm{Cont}}
\newcommand{\RMod}{\mathsf{RMod}}
\newcommand{\LMod}{\mathsf{LMod}}
\newcommand{\An}{\mathsf{An}}
\newcommand{\static}{\mathsf{static}}
\newcommand{\PriL}{\mathsf{Pr}^{\mathsf{iL}}}
\newcommand{\iL}{\mathsf{iL}}
\newcommand{\cn}{\mathsf{cn}}
\newcommand{\tF}{\mathrm{F}}
\newcommand{\tH}{\mathrm{H}}
\newcommand{\infPi}{\mathbbl{\Pi}}
\newcommand{\Pairs}{\mathsf{Pairs}}
\newcommand{\End}{\mathrm{End}}
\newcommand{\Std}{\mathrm{Std}}
\newcommand{\ex}{\mathsf{ex}}
\newcommand{\Huber}{\mathsf{Huber}}
\newcommand{\DMod}{\mathsf{DMod}}
\newcommand{\Tor}{\mathrm{Tor}}
\newcommand{\AnSpec}{\mathrm{AnSpec}}
\newcommand{\AnStk}{\mathsf{AnStk}}
\newcommand{\Stk}{\mathsf{Stk}}
\newcommand{\PCoh}{\mathsf{PCoh}}
\newcommand{\Perf}{\mathsf{Perf}}
\newcommand{\Dbl}{\mathsf{Dbl}}
\newcommand{\BNuc}{\mathsf{BNucl}}
\newcommand{\Nuc}{\mathsf{Nucl}}
\newcommand{\WH}{\text{WH}}
\newcommand{\Gr}{\mathsf{Gr}}
\newcommand{\coker}{\mathrm{coker}}
\newcommand{\Idem}{\mathsf{Idem}}
\newcommand{\Tow}{\mathrm{Tow}}
\newcommand{\Corr}{\mathrm{Corr}}
\newcommand{\trcl}{\mathsf{trcl}}
\newcommand{\tsld}{\scalebox{0.4}{$\square$}}
\newcommand{\sld}{\scalebox{0.5}{$\square$}}
\newcommand{\msld}{\scalebox{0.7}{$\square$}}
\newcommand{\mperp}{\scalebox{0.7}{$\perp$}}
\newcommand{\tri}{\scalebox{0.5}{$\rhd$}}
\newcommand{\pr}{\mathrm{pr}}
\newcommand{\quot}{/ \! /}
\newcommand{\ord}{\mathrm{ord}}
\newcommand{\Rig}{\mathrm{Rig}}
\newcommand{\Ret}{\mathsf{Ret}}
\newcommand{\PD}{\mathrm{PD}}
\setlength{\headheight}{40pt}
\setcounter{tocdepth}{2}
 
\definecolor{hot}{RGB}{65,105,225}
\usepackage{amsmath, amssymb, tikz, amsthm, csquotes, multicol, footnote, tablefootnote, biblatex, wrapfig, float, quiver, mathrsfs, enumitem, upgreek, todonotes}
\usepackage{thmtools}
\usepackage[colorlinks, allcolors=hot, hypertexnames=false]{hyperref}
\usepackage{cleveref}
\usepackage{epigraph}
\addbibresource{refs.bib}
 
 
\theoremstyle{definition}
\declaretheorem[numberwithin=section,name=Theorem]{theorem}
\declaretheorem[sibling=theorem,name=Lemma]{lemma}
\declaretheorem[sibling=theorem,name=Corollary]{corollary}
\declaretheorem[sibling=theorem,name=Example]{example}
\declaretheorem[sibling=theorem,name=Proposition]{proposition}
\declaretheorem[sibling=theorem,name=Remark]{remark}
\declaretheorem[sibling=theorem,name=Definition]{definition}
\declaretheorem[numbered=no,style=definition,name=Theorem]{theorem*}
\declaretheorem[numbered=no,style=definition,name=Lemma]{lemma*}
\declaretheorem[numbered=no,style=definition,name=Proposition]{proposition*}
\declaretheorem[numbered=no,style=definition,name=Definition]{definition*}
\numberwithin{equation}{section}

\setuptodonotes{color=blue!20, size=tiny}
 
 
\begin{document}
\large
\title[A six-functor formalism for solid sheaves in spectral algebraic geometry]{A six-functor formalism for solid sheaves in spectral algebraic geometry}
\author{Wern Juin Gabriel Ong}
\address{Universit\"{a}t Bonn, Bonn, D-53113}
\email{wgabrielong@uni-bonn.de}
\urladdr{https://wgabrielong.github.io/}
\begin{abstract}
  We develop a theory of analytic rings relative to connective $\EE_{\infty}$-algebras in condensed spectra and apply it to the construction of a six-functor formalism for solid quasicoherent sheaves on spectral algebraic stacks. This extends work of Clausen--Scholze in the case of static rings, and of Hansen--Mann, L. Mann, and Rodr\'{i}guez Camargo in the animated setting.
\end{abstract}
\maketitle
\tableofcontents
\section{Introduction}\label{sec: introduction}
\subsection{Results}\label{subsec: intro results}
This work concerns the construction of a six-functor formalism for solid quasicoherent sheaves in spectral algebraic geometry. The geometric objects of this setting, the spectral algebraic stacks, carry a natural theory of quasicoherent sheaves and their cohomology. Our main result, stated here with some imprecision, is that once quasicoherent sheaves are replaced by their solid analogues from condensed mathematics, this cohomology theory upgrades to a six-functor formalism.
\begin{theorem*}[\ref{thm: main six-functor formalism}]
    Let $\kappa$ be an uncountable regular cardinal. There is a six-functor formalism $X\mapsto\QCoh_{\sld}(X)$ on the $(\infty,1)$-category $\Stk_{\kappa}$ of $\kappa$-small spectral algebraic stacks for the faithfully flat and finitely presented topology, assigning to each $X$ its closed symmetric monoidal $(\infty,1)$-category $\QCoh_{\sld}(X)$ of solid quasicoherent sheaves, with the following properties:
    \begin{enumerate}
        \item For every morphism $f: X\to Y$ there is an adjoint pair
        $$f^{*}:\QCoh_{\sld}(Y)\rightleftarrows\QCoh_{\sld}(X): f_{*}$$
        with $f^{*}$ symmetric monoidal.
        \item For every morphism $f: X\to Y$ that locally in the faithfully flat and finitely presented topology factors as an integral morphism followed by a finitely presented morphism, there is a further adjoint pair
        $$f_{!}:\QCoh_{\sld}(X)\rightleftarrows\QCoh_{\sld}(Y): f^{!}$$
        satisfying proper base change and the projection formula.
        \item For every morphism $f: X\to Y$ that locally in the faithfully flat and finitely presented topology factors as $\spec(C)\to\spec(B)\to\spec(A)$ with $A\to B$ flat and almost of finite presentation, and $C$ perfect over $B$, there is a natural equivalence
        $$f^{!}(-)\simeq\omega_{f}\otimes_{X}^{\sld}f^{*}(-):\QCoh_{\sld}(Y)\to\QCoh_{\sld}(X),$$
        where $\omega_{f}:=f^{!}1_{Y}\in\QCoh_{\sld}(X)$ is the dualizing sheaf of $f$.
        \begin{itemize}
            \item[(3.1)] If $f$ is a fiber-smooth map of $\kappa$-small affine spectral schemes, then $\omega_{f}$ is $\otimes_{X}^{\sld}$-invertible.
            \item[(3.2)] If $f$ is \'{e}tale-locally an \'{e}tale morphism, then $\omega_{f}\simeq1_{X}$.
        \end{itemize}
        \item If $f$ is a proper morphism of $\kappa$-small affine spectral schemes, then there is a natural equivalence
        $$f_{!}(-)\simeq f_{*}(-):\QCoh_{\sld}(X)\to\QCoh_{\sld}(Y).$$
        \item If $f: X\to Y$ is as in (3), then the following variant of Grothendieck--Verdier duality holds.
        \begin{itemize}
            \item[(5.1)] The natural morphism $\intHom_{Y}(f_{!}M,N)\to f_{*}\intHom_{X}(M,f^{!}N)$ is an equivalence in $\QCoh_{\sld}(Y)$ for all $M\in\QCoh_{\sld}(X),N\in\QCoh_{\sld}(Y)$.
            \item[(5.2)] The functor given on objects by $M\mapsto\intHom_{X}(M,\omega_{f})$ is an anti-autoequivalence of the full subcategory $\Coh_{\sld}(X,f)$ of $f$-coherent sheaves.
        \end{itemize}
    \end{enumerate}
\end{theorem*}
The theorem sits at the confluence of several areas of mathematics: stable homotopy theory and spectral algebraic geometry, the abstract theory of six-functor formalisms, and condensed mathematics.

In stable homotopy theory the Abelian groups of ordinary algebra give way to the $(\infty,1)$-category $\Sp$ of spectra, the modules over the sphere spectrum $\Sph$, with connective $\EE_{\infty}$-algebras in $\Sp$ taking the place of commutative rings. Following Lurie \cite{SAG}, one builds algebraic geometry over such rings, obtaining spectral analogues of the objects of algebraic geometry -- among them the spectral algebraic stacks and quasicoherent sheaves above.

A sheaf theory is often organized by a six-functor formalism, which records the pushforward and pullback functors, their exceptional counterparts, and the coherences among them. An abstract $(\infty,1)$-categorical account of such formalisms, independent of any one geometric context -- \'{e}tale sheaves, $\Dcal$-modules, and the like -- has emerged in recent work of Gaitsgory--Rozenblyum \cite{GR-DAG}, Heyer--L. Mann \cite{HeyerMann6FF}, Liu--Zheng \cite{LiuZheng6FF}, L. Mann \cite{MannRigid6FF}, and Scholze \cite{Scholze6FF}. One would like the quasicoherent sheaves of spectral algebraic geometry to fit into such a formalism. As in the classical setting, they do not: the exceptional pushforward $f_{!}$ is not defined beyond the case of proper morphisms.

The remedy is due to the condensed mathematics of Clausen and Scholze. They observe that the obstruction can be overcome by enlarging the category of sheaves, replacing quasicoherent sheaves by their topologically complete, or solid, analogues. Within this enlargement, the exceptional pushforward exists for morphisms of finite type \cite[Thm. 8.13]{Condensed} and the six operations assemble into a full six-functor formalism.

The present work carries this out in the setting of spectral algebraic geometry, generalizing the solid six-functor formalism of Clausen and Scholze in the static case \cite{Condensed,Complex} and of Hansen--L. Mann \cite[\S 2.7]{HansenMannCLL}, L. Mann \cite[\S 2.9]{MannRigid6FF}, and Rodr\'{i}guez Camargo \cite{CamargoSolid} in the animated setting. The passage to connective $\EE_{\infty}$ rings is a generalization of the static and animated frameworks, with the forgetful functor from animated rings to connective $\EE_{\infty}$-algebras over $\ZZ$ exhibiting the animated theory as a special case. It is also where the difficulty lies, for the structural conveniences of animated algebra are unavailable over the sphere: the free $\EE_{\infty}$-algebra over the sphere $\Sph\{T\}\simeq\bigoplus_{n\geq0}\Sigma^{\infty}_{+}B\Sigma_{n}$ is of infinite $\Tor$-amplitude, so key factorizations underlying the animated arguments do not transfer.
\subsection{Leitfaden}\label{subsec: intro leitfaden}
We now describe the contents in more detail. The foundation is a theory of analytic rings relative to connective $\EE_{\infty}$-algebras in condensed spectra, generalizing the analytic rings of Clausen--Scholze and their animated refinements, and its construction occupies much of the work. With this theory in place we specialize to the solid setting, building the six operations first on affine spectral schemes and then on stacks, where a form of Grothendieck--Verdier duality is recovered.
\subsubsection{Condensed spectra and solid spectra}\label{subsubsec: intro condensed}
\Cref{sec: condensed spectra and solid spectra} recalls the light condensed mathematics adopted throughout. Unlike in the animated setting, where the theory can be built by applying universal constructions such as animation to the static case, we are forced to work with $(\infty,1)$-categorical constructions from the start. The $(\infty,1)$-categories of condensed animae $\Cond\Ani$ and condensed spectra $\Cond\Sp$ are constructed in \Cref{subsec: condensed animae,subsec: condensed spectra} as hypercomplete sheaves on the site of light profinite sets, and their basic properties are established. Solid spectra are defined in \Cref{subsec: solid spectra}, where the definition of Wagner \cite[\S B.6]{WagnerHHab} through summability of nullsequences is shown to agree with that of Clausen \cite[Def. 13.2]{ClausenLie} through homotopy groups (\Cref{thm: properties of solid spectra}).
\subsubsection{Analytic \texorpdfstring{$\EE_{\infty}$}{E-infinity}-rings}\label{subsubsec: intro analytic}
\Cref{sec: analytic rings} develops the theory of analytic $\EE_{\infty}$-rings. \Cref{subsec: basic definitions} introduces the central object, a condensed $\EE_{\infty}$-ring together with a distinguished subcategory of its modules.
\begin{definition*}[\ref{def: analytic E-infinity ring}]
    Let $A^{\tri}$ be a connective $\EE_{\infty}$-algebra object in $\Cond\Sp$. An \emph{analytic $\EE_{\infty}$-ring} on $A^{\tri}$-modules is a pair $A:=(A^{\tri},\Mod_{A})$ where $\Mod_{A}\subseteq\Mod_{A^{\tri}}$ is a full subcategory of $A^{\tri}$-modules in condensed spectra such that:
    \begin{enumerate}[label=(\roman*)]
        \item $\Mod_{A}$ is closed under small limits and colimits in $\Mod_{A^{\tri}}$.
        \item $\Mod_{A}$ is $\Cond\Sp$-linear: $\intHom_{A^{\tri}}(M,N)\in\Mod_{A}$ for $M\in\Mod_{A^{\tri}},N\in\Mod_{A}$.
        \item The left adjoint $L^{A}:\Mod_{A^{\tri}}\to\Mod_{A}$ to the inclusion is right $t$-exact.
    \end{enumerate}
\end{definition*}
Properties of $A$-analytic modules are studied in \Cref{subsec: analytic modules}, and the functoriality of the $(\infty,1)$-category $\An\CAlg$ of analytic $\EE_{\infty}$-rings, together with its colimits, in \Cref{subsec: functoriality and colimits}. \Cref{subsec: completion} constructs completions (\Cref{thm: completions exist}), the passage to the full subcategory $\An\CAlg^{\wedge}$ of complete analytic $\EE_{\infty}$-rings where $A^{\tri}\in\Mod_{A}\subseteq\Mod_{A^{\tri}}$. This is simpler here than in the animated setting: the completion of an analytic $\EE_{\infty}$-ring is again an analytic $\EE_{\infty}$-ring, whereas endowing the completion of an animated analytic ring with an animated structure is delicate. \Cref{subsec: topological invariance} establishes a topological invariance phenomenon, showing that an analytic $\EE_{\infty}$-ring structure is detected on the heart and so recorded by a purely static datum over $\pi_{0}(A^{\tri})$ (\Cref{thm: heart recognition for analytic rings}). \Cref{subsec: mereology of analytic modules} then turns to the internal finiteness structure of $\Mod_{A}$, relating nuclearity and dualizability (\Cref{prop: dualizable iff nuclear and compact}) as a special case of the discussion in \Cref{app: nuclearity}.
\subsubsection{Morphisms of analytic \texorpdfstring{$\EE_{\infty}$}{E-infinity}-rings}\label{subsubsec: intro morphisms}
\Cref{sec: morphisms of analytic rings} takes up morphisms in earnest. \Cref{subsec: constructing morphisms} records the $\EE_{\infty}$-analogue of a construction principle of Scholze, promoting a morphism of the underlying rings to one of analytic $\EE_{\infty}$-rings under a compatibility on $\pi_{0}$ (\Cref{prop: constructing morphisms of analytic rings}). \Cref{subsec: steady morphisms} treats the steady morphisms.
\begin{definition*}[\ref{def: steady morphism}]
    A morphism $(A^{\tri},\Mod_{A})\to(B^{\tri},\Mod_{B})$ in $\An\CAlg^{\wedge}$ is \emph{steady} if for every morphism $(A^{\tri},\Mod_{A})\to(C^{\tri},\Mod_{C})$ and coCartesian square with pushout $(D^{\tri},\Mod_{D})$ the natural morphism
    $$B\otimes_{A}(M|_{A^{\tri}})\longrightarrow(D\otimes_{C}M)|_{B^{\tri}}$$
    is an equivalence in $\Mod_{B}$ for all $M\in\Mod_{C}$.
\end{definition*}
For a morphism of analytic $\EE_{\infty}$-rings, the scalar extension $\Mod_{A}\to\Mod_{B}$ need not be tensoring with $B^{\tri}$ regarded as an algebra in $\Mod_{A}$, so base change is not automatic. 

\Cref{subsec: open immersions} treats open immersions (\Cref{def: open immersion}) and, in particular, a restricted form of an expectation of Peter Scholze, relayed to us by L. Mann, that open immersions are detected on the heart. The recognition reduces the question to the underlying static analytic rings once the analytic $\Tor$-amplitude is at most one.
\begin{proposition*}[\ref{prop: open immersion detected on static rings}]
    Let $(A^{\tri},\Mod_{A})\to(B^{\tri},\Mod_{B})$ be a morphism of complete analytic $\EE_{\infty}$-rings with $B\otimes_{A}\pi_{0}(A^{\tri})\simeq\pi_{0}(B^{\tri})$ and such that:
    \begin{enumerate}[label=(\roman*)]
        \item The right adjoint $(-)|_{A^{\tri}}:\Mod_{B}\to\Mod_{A}$ to scalar extension is fully faithful.
        \item The functor $\pi_{0}(B)\otimes_{\pi_{0}(A)}-$ is of analytic $\Tor$-amplitude $\leq 1$.
        \item The $t$-structure on $\Mod_{A}$ is compatible with small products, and the induced morphism $(\pi_{0}(A^{\tri}),\Mod_{\pi_{0}(A)})\to(\pi_{0}(B^{\tri}),\Mod_{\pi_{0}(B)})$ is an open immersion.
    \end{enumerate}
    Then $(A^{\tri},\Mod_{A})\to(B^{\tri},\Mod_{B})$ is an open immersion.
\end{proposition*}
\Cref{subsec: analytically proper morphisms} treats the complementary analytically proper morphisms and their base change and projection formulae. Together the open immersions and analytically proper morphisms form a suitable decomposition of a geometric setup on $\An\Aff:=(\An\CAlg^{\wedge})^{\Opp}$ (\Cref{lem: geometric setup on analytic E-infinity rings}), from which \Cref{subsec: six operations on analytic rings} constructs a six-functor formalism on analytic $\EE_{\infty}$-rings (\Cref{prop: six-functor formalism on analytic E-infinity rings}), following the framework of \cite[\S 2-3]{HeyerMann6FF}. \Cref{subsec: descent} closes with a discussion of descent for analytic modules (\Cref{thm: descent for descendable morphisms}), building on work of Mikami \cite[\S 2]{MikamiFFDescent} and L. Mann's formalism of steady endofunctors \cite[\S 2.5-2.6]{MannRigid6FF}.
\subsubsection{Discrete solid \texorpdfstring{$\EE_{\infty}$}{E-infinity}-rings}\label{subsubsec: intro solid}
\Cref{sec: solid commutative rings} begins the study of discrete solid $\EE_{\infty}$-rings. The solid spectra of \Cref{subsec: solid spectra} assemble into an analytic $\EE_{\infty}$-ring over the sphere, the solid sphere $\Sph_{\sld}$ of \Cref{subsec: solid sphere}, whose generators are internally compact projective rather than merely $\omega_{1}$-compact (\Cref{prop: generation of solid spectra}, cf. \Cref{prop: A-analytic tensor product}). The same construction, carried out on compact projective algebras and extended by sifted colimits, produces the discrete solid $\EE_{\infty}$-ring $A_{\sld}$ of an arbitrary connective $\EE_{\infty}$-algebra $A\in\CAlg^{\cn}$ in \Cref{subsec: solid E-infinity rings} (\Cref{prop: solid rings general}). \Cref{subsec: spectral Huber pairs} refines these to discrete spectral Huber pairs.
\begin{definition*}[\ref{def: discrete spectral Huber pair}]
    Let $A\in\CAlg^{\cn}$. A \emph{discrete spectral Huber pair} $(A,A^{+})$ is a morphism $A^{+}\to A$ in $\CAlg^{\cn}$ such that $\pi_{0}(A^{+})\subseteq\pi_{0}(A)$ is an integrally closed subring of $\pi_{0}(A)$. A \emph{morphism of discrete spectral Huber pairs} $f:(A,A^{+})\to(B,B^{+})$ is a morphism $f: A\to B$ in $\CAlg^{\cn}$ with $\pi_{0}(f)(\pi_{0}(A^{+}))\subseteq\pi_{0}(B^{+})$.
\end{definition*}
In contrast to the extant literature on the animated theory, the plus-part $A^{+}$ is permitted to be non-static, the subring condition being imposed only on $\pi_{0}$. Such pairs give rise to complete analytic $\EE_{\infty}$-rings by a fully faithful functor (\Cref{lem: Huber pairs to solid rings,lem: functoriality of Huber rings}), and afford both finer control of the solid structure and a convenient factorization of morphisms through an analytically proper one (\Cref{prop: factorization of Huber ring maps}). \Cref{subsec: mereology of solid modules} shows that over a discrete solid Huber ring the nuclear objects are exactly the discrete ones, so every dualizable object is discrete and $A_{\sld}$ is Fredholm (\Cref{prop: nuclear objects in solid Huber rings,prop: solid rings are Fredholm}).

The exceptional pushforwards rest on the factorization discrete spectral Huber pairs supply. The morphisms carrying them are those of finite expansion (\Cref{def: finite expansion Einfty}), the morphisms $A\to B$ in $\CAlg^{\cn}$ admitting a factorization $A\to A\{T_{1},\dots,T_{n}\}\to B$ through a free $\EE_{\infty}$-algebra with $\pi_{0}(A\{T_{1},\dots,T_{n}\})\to\pi_{0}(B)$ integral. The construction of $f_{!}\dashv f^{!}$ for such a morphism reduces to the free-algebra map $\Sph_{\sld}\to\Sph\{T\}_{\sld}$, treated in \Cref{subsec: solid affine line,subsec: exceptional functors finite expansion}.
\begin{proposition*}[\ref{prop: exceptional functors for the affine line}]
    Let $f:\Sph_{\sld}\to\Sph\{T\}_{\sld}$ be the morphism of discrete solid $\EE_{\infty}$-rings to the free $\EE_{\infty}$-algebra.
    \begin{enumerate}
        \item $f$ admits a factorization $\Sph_{\sld}\xrightarrow{p}(\Sph\{T\},\Sph)_{\sld}\xrightarrow{j}\Sph\{T\}_{\sld}$ in which $p$ is analytically proper and $j$ an open immersion.
        \item The functor $f_{!}:=p_{*}j_{!}:\Mod_{\Sph\{T\}_{\sld}}\to\Mod_{\Sph_{\sld}}$ preserves small colimits and satisfies the projection formula.
    \end{enumerate}
\end{proposition*}
That $j$ is an open immersion is shown in two ways, either through the idempotent algebra of \Cref{subsec: solid affine line} or by the recognition criterion of \Cref{prop: open immersion detected on static rings} described above. Base change gives the free algebra in finitely many variables over an arbitrary ring. Noting that integral morphisms are analytically proper, the two factors yield the exceptional functors of any finite-expansion morphism (\Cref{thm: exceptional functor for finite expansion}). \Cref{subsec: flat morphisms of solid rings} develops the structure theory of the flat, almost of finite presentation morphisms after Mikami, on which \Cref{subsec: descent for solid rings} shows discrete solid $\EE_{\infty}$-rings satisfy descent for the faithfully flat and finitely presented topology.
\begin{theorem*}[\ref{thm: solid fppf descent}]
    Let $A_{\sld}\to B_{\sld}$ be a finitely presented and faithfully flat morphism of discrete solid $\EE_{\infty}$-rings. Then the functor
    $$\Mod_{A_{\sld}}\longrightarrow\lim_{[n]\in\Delta}\Mod_{B_{\tsld}^{\otimes_{A_{\tsld}}(n+1)}}$$
    is an equivalence.
\end{theorem*}
\subsubsection{Six operations for solid quasicoherent sheaves}\label{subsubsec: intro six operations}
\Cref{sec: six functors for solid sheaves} constructs the six-functor formalism for solid quasicoherent sheaves. On $\kappa$-small affine spectral schemes it is obtained in \Cref{subsec: six functors on affines} by restricting the formalism on complete analytic $\EE_{\infty}$-rings along the embedding of discrete solid $\EE_{\infty}$-rings (\Cref{prop: six-functor formalism affine spectral schemes}). 

\Cref{subsec: affine prim and suave} studies the prim, suave, and smooth morphisms of Heyer--L. Mann in this setting, where the divergence from the animated theory is sharpest. In the animated setting, a morphism of solid affine schemes is suave if and only if it is almost of finite presentation and of finite $\Tor$-amplitude. The proof there rests on a structure theorem factoring such a morphism as $A\to A[T_{1},\dots,T_{n}]\to B$ over the polynomial algebra, and no such theorem is available over the sphere. We isolate instead the pseudo-lisse morphisms -- those factoring as $A\to B\to C$ with $A\to B$ flat and almost of finite presentation and $C$ perfect over $B$ -- and prove they are $\msld$-suave (\Cref{prop: affine suave morphisms}), though it is not clear that all $\msld$-suave morphisms are of this form (\Cref{rmk: comparison to lurie jiang}). 

\Cref{subsec: affine perfect and coherent} introduces the perfect and $f$-coherent sheaves and their duality. \Cref{subsec: construction on stacks} globalizes the affine formalism to spectral algebraic stacks by descent for the faithfully flat and finitely presented topology (\Cref{thm: six-functor formalism on stacks}), using that such morphisms are $\msld$-suave (\Cref{lem: ffp is universal cover}) to reduce many questions to the affine case. \Cref{subsec: duality theory for spectral algebraic stacks} then develops Grothendieck--Verdier duality (\Cref{thm: Grothendieck Verdier duality stacks}) and restates the main theorem.
\subsubsection{Appendix A: \texorpdfstring{$t$}{t}-bistable categories and their localizations}\label{subsubsec: intro appendix A}
\Cref{app: categories} collects results in $(\infty,1)$-category theory that we could not find in the literature. Several statements usually proved for analytic rings already hold for stable presentable categories carrying compatible $t$- and symmetric monoidal structures -- the $t$-bistable categories. We work at this level of generality, both to clarify the hypotheses on which the statements rest, and because we expect them to be of independent interest. The localizations of such categories, and the recognition of symmetric monoidal localizations with smashing kernel, imply many of the topological invariance results of \Cref{subsec: topological invariance} and underlie the previously described open-immersion criterion of \Cref{subsec: open immersions}.
\subsubsection{Appendix B: Nuclearity in the \texorpdfstring{$\kappa$}{kappa}-presentable setting}\label{subsubsec: intro appendix B}
\Cref{app: nuclearity} treats trace class morphisms and nuclear objects in stable $\kappa$-presentably symmetric monoidal $(\infty,1)$-categories with compact unit, showing that the compact-generation hypotheses of \cite[Lectures VIII \& IX]{Complex} are unnecessary for the categories of interest. The nuclear objects (\Cref{def: basic nuclear in category}) form a full subcategory closed under small colimits, and an object is dualizable exactly when it is compact and nuclear (\Cref{prop: dualizable iff nuclear and compact category}).
\subsubsection{Appendix C: Commutative algebra of \texorpdfstring{$\EE_{\infty}$}{E-infinity}-rings}\label{subsubsec: intro appendix C}
\Cref{app: commutative algebra} records auxiliary results in the commutative algebra of connective $\EE_{\infty}$-algebras in $\Sp$. Finiteness properties of modules and morphisms are recalled in \Cref{subsec: finiteness properties}. We then introduce in \Cref{subsec: morphisms of finite expansion} the morphisms of finite expansion, due to Hamacher \cite{FiniteExpansion} in the static case. We adapt this to the $\EE_{\infty}$ setting and verify its permanence properties (\Cref{prop: finite expansion permanence}). \Cref{subsec: pseudo-lisse morphisms} concludes with the pseudo-lisse morphisms $A\to C$ of connective $\EE_{\infty}$-algebras in $\Sp$, those admitting a factorization $A\to B\to C$ with $B$ flat and almost of finite presentation over $A$ and $C$ perfect as a $B$-module, alongside the elementary pseudo-lisse morphisms, where $B$ is required to be the polynomial algebra over $A$ in finitely many variables (\Cref{def: pseudo-lisse appendix}), and their basic properties (\Cref{prop: properties of pseudo-lisse morphisms}).
\subsection{Notation and conventions}\label{subsec: notation and conventions}
In the rest of this text, we will work with $\infty$-categories and $\infty$-categorical algebra throughout following \cite{HTT,HA,kerodon}, referring to $(\infty,1)$-categories as categories and distinguishing ordinary categories as needed. We will adopt the following: 
\begin{itemize}
  \item $\Ani$ and $\Sp$ will refer to the categories of animae (also known as spaces) and spectra, respectively. 
  \item $\PrL$ will refer to Lurie's category of presentable categories and left adjoint functors. For a regular cardinal $\kappa$, $\PrL_{\kappa}$ will refer to Lurie's category of $\kappa$-presentable categories -- the non-full subcategory of $\PrL$ spanned by the $\kappa$-presentable categories and functors preserving $\kappa$-compact objects. $\PrL$ admits a symmetric monoidal structure via the Lurie tensor product which restricts to a symmetric monoidal structure on $\PrL_{\kappa}$.
  \item We say a category $\Cscr$ is presentably symmetric monoidal if $\Cscr$ is symmetric monoidal and the tensor product preserves small colimits separately in each argument. This implies that $\Cscr$ is closed symmetric monoidal. For a regular cardinal $\kappa$, we say $\Cscr$ is $\kappa$-presentably symmetric monoidal if $\Cscr\in\CAlg(\PrL_{\kappa})$: that is, the unit $1_{\Cscr}$ is $\kappa$-compact and the tensor product preserves $\kappa$-compact objects. 
\end{itemize}
Additionally, we will make use of the theory of (light) condensed mathematics as developed by Clausen and Scholze, the language of enriched categories (ie. enriched $\infty$-categories), as well as the axiomatic perspective on six-functor formalisms. We refer the reader to \cite{Condensed,Analytic,Complex,AnalyticStacks,CamargoSolid} as references for (light) condensed mathematics, \cite{MannRigid6FF,HeyerMann6FF,RamziLocRigid,RamziDualizable,ReutterZettoEnriched} for enriched categories, and \cite{MannRigid6FF,HeyerMann6FF,ConstructionCLL} for abstract six-functor formalisms. We will adopt the following: 
\begin{itemize}
  \item Condensed objects are assumed to be light, that is, hypersheaves on the category $\PFin_{\NN}$ of profinite sets arising as sequential limits of finite sets and continuous maps between them.
  \item An object is said to be static if it is 0-truncated as an object (eg. if it lies in the heart of a $t$-structured stable category), reserving discreteness for condensed objects satisfying a certain property. 
  \item For a category $\Cscr$ enriched in a category $\Vscr$, we will denote by $\Mor_{\Cscr}(-,-)$ and $\Hom_{\Cscr}^{\Vscr}(-,-)$ the $\Ani$-valued and $\Vscr$-valued morphisms, respectively. When $\Cscr$ is stable and $\Vscr=\Sp$, we write $\Hom_{\Cscr}^{\Sp}(-,-)$ as $\Hom_{\Cscr}(-,-)$. If $\Cscr$ is closed symmetric monoidal (eg. presentably symmetric monoidal) we will use $\intHom_{\Cscr}(-,-)$ to refer to the mapping object internal to $\Cscr$.
\end{itemize}
In this text, we have at times elected to deviate our nomenclature and notation from the extant literature. We have done so only where, in our judgment, the alternative conventions render the underlying ideas more transparent -- first to ourselves, and, we hope, to the reader. In each such instance, the justification for our choice is given in the remark immediately succeeding the definition in which the deviation is introduced. We apologize in advance for any unnecessary confusion these departures may cause the reader already accustomed to the established terminology.
\subsection{Acknowledgements}\label{subsec: acknowledgements}
I would like to thank my advisors: Dr. Lucas Mann, for his most felicitous choice of topic for my master's thesis, for his generosity in discussion and critique, and for his unwavering support throughout the year; and Prof. Dr. Peter Scholze, for the inspiring courses he taught as well as for helpful conversations and correspondence. I have been fortunate to discuss the content of this thesis, in part or in whole, with various mathematicians. I thank Benjamin Antieau, Daniel Arone, Sam Brinkerhoff, Dustin Clausen, Thorger Gei{\ss}, Niklas Kipp, Jordan Levin, Jiacheng Liang, Niko Naumann, Marius Nielsen, Lucas Piessevaux, Phil P\"{u}tzst\"{u}ck, Maxime Ramzi, Gabriel Ribeiro, Florian Riedel, Maria Stroe, Ferdinand Wagner, and Yuchen Wu for conversations that have shaped this work, even if not always in its content. I am particularly grateful to Maria Stroe, who read a draft with great care.
\section{Condensed spectra and solid spectra}\label{sec: condensed spectra and solid spectra}
Lurie's theory of presentable categories produces, from a small site $\Cscr$ and a presentable category $\Vscr$, the hypercomplete $\Vscr$-valued sheaves on $\Cscr$ as the Lurie tensor product of the hypercomplete topos of $\Cscr$ with $\Vscr$. The category of condensed animae is such a hypercomplete topos, and furnishes the basic building block of homotopy-coherent condensed mathematics, and of condensed spectra in particular.

We recall its construction after Clausen \cite[\S 4]{ClausenLie}, building on earlier work of Barwick--Glasman--Haine \cite[\S 13]{Exodromy} and Wolf \cite{WolfPyknotic}, before forming the stable category of condensed spectra by tensoring with $\Sp$. Solid spectra are then isolated within condensed spectra as the analogue of the derived category of solid condensed Abelian groups, and the descriptions of Clausen \cite[Def. 13.2]{ClausenLie} and Wagner \cite[\S B.6]{WagnerHHab} are shown to agree.
\subsection{Condensed animae}\label{subsec: condensed animae}
It was observed by Johnstone that the category $\CHaus$ of compact Hausdorff spaces is monadic over $\Sets$ through the ultrafilter monad \cite[\S III.2.4]{JohnstoneStone}, indicating that a better-behaved topological algebra should follow from the algebraic character of such spaces. This expectation is realized in part by the condensed mathematics of Clausen--Scholze \cite{Condensed}, built from set-valued sheaves on profinite sets for the Grothendieck pretopology of finite jointly surjective covers. 

For most purposes it suffices to restrict to those profinite sets arising as sequential limits of finite discrete spaces, which form an essentially small category and yield the light condensed mathematics used throughout this work. These profinite sets organize into a Grothendieck site whose hypercomplete sheaves are the (light) condensed animae. 
\begin{definition}\label{def: light profinite sets}
    The category of \emph{light profinite sets} $\PFin_{\NN}$ is the sequential $\Pro$-category of finite sets. 
\end{definition}
\begin{lemma}\label{lem: Grothendieck topology on light profinite sets}
    Let $S\in\PFin_{\NN}$ and say a family $\{S_{i}\to S\}_{1\leq i\leq n}$ in $(\PFin_{\NN})_{/S}$ is a covering family if $\coprod_{1\leq i\leq n}S_{i}\to S$ is a surjection. The collection of covering families defines a Grothendieck pretopology on $\PFin_{\NN}$. We refer to the corresponding Grothendieck topology as the finitary Grothendieck topology. 
\end{lemma}
\begin{proof}
    This follows from the discussion in \cite[\S 2.2]{CamargoSolid}. More explicitly, the covering sieves for the finitary Grothendieck topology on $\PFin_{\NN}$ are those sieves that contain a sieve generated by a covering family. 
\end{proof}
\begin{definition}\label{def: condensed animae}
    The category of \emph{condensed animae} $\Cond\Ani$ is the category of hypercomplete $\Ani$-valued sheaves on $\PFin_{\NN}$ with respect to the finitary Grothendieck topology of \Cref{lem: Grothendieck topology on light profinite sets}. 
\end{definition}
Two properties of condensed animae recur throughout. 
\begin{proposition}\label{prop: properties of condensed animae}
    Let $\Cond\Ani$ be the category of condensed animae. 
    \begin{enumerate}
        \item $\Cond\Ani$ is an $\omega_{1}$-presentable topos. 
        \item $\Cond\Ani$ is a replete topos: if 
        $$X:=\lim_{n}\left(\dots\to X_{2}\to X_{1}\to X_{0}\right)$$
        is a sequential diagram in $\Cond\Ani$ with transition maps that are effective epimorphisms then each $X\to X_{n}$ is an effective epimorphism. 
    \end{enumerate}
\end{proposition}
\begin{proof}[Proof of (1)]
    That $\Cond\Ani$ is a topos is immediate from construction. For $\omega_{1}$-presentability, note that the presheaf category $\PSh(\PFin_{\NN};\Ani)$ is $\omega$-presentable so by \cite[Lem. 5.5.1.4]{HTT} it suffices to show the inclusion $\Cond\Ani\hookrightarrow\PSh(\PFin_{\NN};\Ani)$, right adjoint to the formation of hypersheaves, preserves $\omega_{1}$-filtered colimits. Noting that $\omega_{1}$-filtered colimits commute with countable limits in $\Ani$ \cite[Prop. 5.3.3.3]{HTT}, the inclusion preserves them, and the claim follows. 
\end{proof}
\begin{proof}[Proof of (2)]
    This is \cite[Lem. 4.9]{ClausenLie}. 
\end{proof}
\subsection{Condensed spectra}\label{subsec: condensed spectra}
Condensed animae form a presentable topos, and it is by tensoring with spectra that the stable homotopy theory underlying the rest of this work is obtained. The construction rests on the following general identification. 
\begin{proposition}\label{prop: hypersheaves as tensor}
    Let $\Cscr$ be a site with associated hypercomplete topos $\Xscr$ and $\Vscr$ a presentable category. Then precomposition with the functor $\Cscr\to\Xscr$ induces an equivalence of categories $\Xscr\otimes\Vscr\xrightarrow{\sim}\widehat{\Sh}(\Cscr;\Vscr)$. 
\end{proposition}
\begin{proof}
    This follows by combining \cite[Rmk. 1.3.1.6]{SAG}, which identifies $\Xscr\otimes\Vscr$ with $\Vscr$-valued sheaves on $\Xscr$, with \cite[Lem. A.4.22]{HeyerMann6FF}. 
\end{proof}
The derived category of condensed Abelian groups occupies the central position in animated analytic geometry, and condensed spectra take over that role here as its spectral, homotopy-coherent counterpart. 
\begin{definition}\label{def: condensed spectra}
    The category of \emph{condensed spectra} $\Cond\Sp$ is the category of $\Sp$-valued hypersheaves on $\PFin_{\NN}$ with respect to the finitary Grothendieck topology of \Cref{lem: Grothendieck topology on light profinite sets}. 
\end{definition}
\begin{remark}\label{rmk: characterizations of CondSp}
    By \Cref{prop: hypersheaves as tensor} and \cite[Rmk. 1.3.2.2]{SAG}, the category of condensed spectra is defined equivalently as $\Cond\Ani\otimes\Sp$ or as spectrum objects in condensed animae. 
\end{remark}
For the following, recall the notions of an $\omega_{1}$-presentably symmetric monoidal category from \Cref{subsec: notation and conventions} and of a compact projective object from \Cref{def: compact projective objects} (1). With this terminology, the elementary properties of condensed spectra are as follows. 
\begin{theorem}\label{thm: properties of condensed spectra}
    Let $\Cond\Sp$ be the category of condensed spectra. 
    \begin{enumerate}
        \item $\Cond\Sp$ is a stable $\omega_{1}$-presentably symmetric monoidal category with compact projective unit $\Sph$. 
        \item $\Cond\Sp$ admits an accessible $t$-structure $(\Cond\Sp_{\geq0},\Cond\Sp_{\leq0})$ with the following properties: 
        \begin{enumerate}[label=(\roman*)]
            \item $\Cond\Sp_{\geq0},\Cond\Sp_{\leq0}$ are spanned by those condensed spectra $M$ such that the homotopy sheaves $\pi_{k}(M)$ vanish for $k<0$ and $k>0$, respectively. The heart of this $t$-structure is equivalent to the category of (static) condensed Abelian groups $\Cond\Ab$. 
            \item $\Cond\Sp_{\geq0}$ contains the unit in $\Cond\Sp$ and is closed under tensor products in $\Cond\Sp$. 
            \item $\Cond\Sp_{\leq0}$ is closed under filtered colimits in $\Cond\Sp$. 
            \item $\Cond\Sp_{\geq0}$ is closed under countable products in $\Cond\Sp$. 
            \item The $t$-structure on $\Cond\Sp$ is both left and right complete. 
        \end{enumerate}
        \item There is a $t$-exact symmetric monoidal equivalence of categories 
        $$\Mod_{\pi_{0}(\Sph)}(\Cond\Sp)\simeq\Dscr(\Cond\Ab).$$ 
    \end{enumerate}
\end{theorem}
\begin{proof}[Proof of (1)]
    Defining $\Cond\Sp$ as the Lurie tensor product $\Cond\Ani\otimes\Sp$, the statement on $\omega_{1}$-presentability of the underlying category follows readily from the fact that the Lurie tensor product on presentable categories restricts to one on $\PrL_{\kappa}$ for all regular cardinals $\kappa$ \cite[Lem. 5.3.2.11]{HA}. Stability follows from the identification $\Cond\Sp\simeq\Sp(\Cond\Ani)$ in \Cref{rmk: characterizations of CondSp}. The category admits a symmetric monoidal structure by \cite[\S 1.3.4]{SAG}.
    
    For compactness of the unit, we note that every hypercover of the singleton splits, so $\Hom_{\Cond\Sp}(\Sph,-):\Cond\Sp\to\Sp$ preserves filtered colimits which, in conjunction with the finite colimit preservation furnished by exactness, shows the functor preserves small colimits and thus gives compactness of $\Sph$. For projectivity, we show that $\Hom_{\Cond\Sp}(\Sph,-):\Cond\Sp\to\Sp$ given on objects by $M\mapsto M(*)$ is $t$-exact. For $k\in\ZZ$, $\pi_{k}(M(*))\simeq\pi_{k}(M)(*)$ as sheafification of the pointwise homotopy presheaf does not change values at the singleton. Since the $t$-structure on $\Cond\Sp$ is characterized by homotopy sheaves, the functor is $t$-exact. 

    The unit is compact, in particular $\omega_{1}$-compact, so it remains to show that the tensor product preserves $\omega_{1}$-compact objects. The tensor product on $\Cond\Sp$ is exact and preserves small colimits separately in each argument, so using symmetric monoidality of $\Sigma^{\infty}_{+}:\Cond\Ani\to\Cond\Sp$, we have
    $$\Sph[S_{1}]\otimes_{\Cond\Sp}\Sph[S_{2}]\simeq\Sigma^{\infty}_{+}(S_{1})\otimes_{\Cond\Sp}\Sigma^{\infty}_{+}(S_{2})\simeq\Sigma^{\infty}_{+}(S_{1}\times S_{2})\simeq\Sph[S_{1}\times S_{2}]$$
    for $S_{1},S_{2}\in\PFin_{\NN}$. $\PFin_{\NN}$ admits countable limits \cite[Prop. 2.1.4]{CamargoSolid}, so $S_{1}\times S_{2}$ is light profinite and $\Sph[S_{1}\times S_{2}]$ is $\omega_{1}$-compact. The free condensed spectra $\Sph[S]$ for $S\in\PFin_{\NN}$ are $\omega_{1}$-compact and generate $\Cond\Sp$ under small colimits and shifts. Hence an object of $\Cond\Sp$ is $\omega_{1}$-compact if and only if it is contained in the minimal full subcategory of $\Cond\Sp$ closed under countable colimits and shifts and containing the free condensed spectra on light profinite sets -- such a category is already idempotent complete by \cite[Cor. 4.4.5.15]{HTT}. Since the tensor product is exact and preserves colimits separately in each variable, the computation on generators above shows that binary tensors of $\omega_{1}$-compact objects are $\omega_{1}$-compact. Thus $\Cond\Sp$ is $\omega_{1}$-presentably symmetric monoidal. 
\end{proof}
\begin{proof}[Proof of (2)]
    \cite[Prop. 1.3.2.7]{SAG} gives all properties but the condition on countable products and left completeness. 

    \cite[Lem. 4.10 (4)]{ClausenLie} implies that countable products in $\Cond\Sp$ are $t$-exact. In particular, they preserve connective objects, so $\Cond\Sp_{\geq0}$ is closed under countable products. 

    For left completeness, note that the category of condensed animae is a Postnikov complete topos \cite[Thm. A]{PostnikovReplete} as its underlying 1-topos of condensed sets is replete \cite[Prop. 2.1.4 \& Rmk. 2.2.2]{CamargoSolid}. The $t$-structure on $\Sp$-valued sheaves is then left complete by \cite[Cor. 3.26]{PostnikovReplete}. 
\end{proof}
\begin{proof}[Proof of (3)]
    This is \cite[Cor. 2.6.7 \& 2.6.9]{BrinkCondensed}. 
\end{proof}
\begin{remark}\label{rmk: CondSp is pres sym mon t plus bistable}
    The conditions of \Cref{thm: properties of condensed spectra} (1) and (2) show that $\Cond\Sp$ is a presentably symmetric monoidal $t^{+}$-bistable category in the sense of \Cref{def: symmetric monoidal t-bistable category} (2), which situates condensed spectra within the framework of the $t^{+}$-bistable categories discussed in \Cref{app: categories}. 
\end{remark}
\subsection{Solid spectra}\label{subsec: solid spectra}
A solid condensed Abelian group is one in which nullsequences are summable, a condition encoded by the shift operator on the free condensed Abelian group on the nullsequence. The same construction applies to condensed spectra: a condensed spectrum is solid when this shift condition holds for $\Sph[\NN_{\infty}]$, following Wagner \cite[\S B.6]{WagnerHHab}. This coincides with the formulation of Clausen \cite[Def. 13.2]{ClausenLie}, under which solidity is tested on homotopy groups.
\begin{definition}\label{def: solid spectra}
    Let $\NN\cup\{\infty\}$ be the one-point compactification of the natural numbers as a light profinite set and $\Sph[\NN_{\infty}]:=\cofib(\Sph[\{\infty\}]\to\Sph[\NN\cup\{\infty\}])$ be the cofiber of the morphism of condensed spectra induced by the inclusion $\{\infty\}\hookrightarrow\NN\cup\{\infty\}$ in $\PFin_{\NN}\subseteq\Cond\Ani$. 
    \begin{enumerate}
        \item A condensed spectrum $M$ is \emph{solid} if the natural morphism 
        $$\id-\sigma^{*}:\intHom_{\Cond\Sp}(\Sph[\NN_{\infty}],M)\to\intHom_{\Cond\Sp}(\Sph[\NN_{\infty}],M)$$
        induced by the shift operator $\sigma$ given by $n\mapsto n+1$ on $\NN\cup\{\infty\}$ is an equivalence. 
        \item The category of \emph{solid spectra} $\Cond\Sp^{\sld}$ is the full subcategory of $\Cond\Sp$ spanned by the solid spectra. 
    \end{enumerate}
\end{definition}
The following records the colimit behaviour of the internal hom out of the nullsequence quotient, reducing it to Brink's computation for internal cohomology of free condensed spectra on light profinite sets \cite[Lem. 2.7.1]{BrinkCondensed}.
\begin{lemma}\label{lem: Brink internal cohomology}
    Let $\Sph[\NN_{\infty}]:=\cofib(\Sph[\{\infty\}]\to\Sph[\NN\cup\{\infty\}])$ be the object of \Cref{def: solid spectra}. Then the functor $\intHom_{\Cond\Sp}(\Sph[\NN_{\infty}],-):\Cond\Sp\to\Cond\Sp$ commutes with direct sums of uniformly bounded above objects: 
    $$\intHom_{\Cond\Sp}\left(\Sph[\NN_{\infty}],\bigoplus_{i\in I}M_{i}\right)\simeq\bigoplus_{i\in I}\intHom_{\Cond\Sp}\left(\Sph[\NN_{\infty}],M_{i}\right)$$
    for every $k\geq0$ and every family $\{M_{i}\}_{i\in I}\subseteq\Cond\Sp_{\leq k}$. 
\end{lemma}
\begin{proof}
    Applying $\intHom_{\Cond\Sp}(-,-)$ to the (co)fiber sequence $\Sph[\{\infty\}]\to\Sph[\NN\cup\{\infty\}]\to\Sph[\NN_{\infty}]$ in the first argument gives a (co)fiber sequence 
    $$\intHom_{\Cond\Sp}(\Sph[\NN_{\infty}],-)\to\intHom_{\Cond\Sp}(\Sph[\NN\cup\{\infty\}],-)\to\intHom_{\Cond\Sp}(\Sph[\{\infty\}],-)$$
    of endofunctors of $\Cond\Sp$. By \cite[Lem. 2.7.1]{BrinkCondensed}, the latter two functors commute with filtered colimits of uniformly bounded above objects. Since limits and colimits in functor categories are computed pointwise, and finite limits commute with filtered colimits in $\Cond\Sp$, $\intHom_{\Cond\Sp}(\Sph[\NN_{\infty}],-)$ preserves filtered colimits of uniformly bounded above objects. 

    Now in the case of direct sums, we may write $\bigoplus_{i\in I}M_{i}$ as a filtered colimit of its finite subsums. Since $\intHom_{\Cond\Sp}(\Sph[\NN_{\infty}],-)$ is exact it preserves finite direct sums, and additionally preserves filtered colimits, giving the claim. 
\end{proof}
\begin{proposition}\label{prop: nullsequence quotient is internally compact projective}
    Let $\Sph[\NN_{\infty}]:=\cofib(\Sph[\{\infty\}]\to\Sph[\NN\cup\{\infty\}])$ be the object of \Cref{def: solid spectra}. The object $\Sph[\NN_{\infty}]$ is an internally compact projective object of $\Cond\Sp$ in the sense of \Cref{def: internally compact projective}: $\intHom_{\Cond\Sp}(\Sph[\NN_{\infty}],-):\Cond\Sp\to\Cond\Sp$ is $t$-exact and preserves small colimits. 
\end{proposition}
\begin{proof}
    We first treat $t$-exactness. Note $\Sph[\NN_{\infty}]$ is connective, being the cofiber of connective objects, so $\intHom_{\Cond\Sp}(\Sph[\NN_{\infty}],-)$ is exact, limit-preserving, and left $t$-exact as an endofunctor of $\Cond\Sp$ as it is right adjoint to the right $t$-exact left adjoint $\Sph[\NN_{\infty}]\otimes_{\Cond\Sp}-$ \cite[Prop. 1.3.17 (ii)]{Perverse}. The equivalence of \Cref{thm: properties of condensed spectra} (3), being $t$-exact, identifies the respective hearts. Under this identification, \Cref{lem: modules inherit bistability} (2) and \Cref{prop: relationship with heart} (1) identify the restriction of $\intHom_{\Cond\Sp}(\Sph[\NN_{\infty}],-)$ to the heart with $\intHom_{\Dscr(\Cond\Ab)}(\ZZ[\NN_{\infty}],-)|_{\Sph}$. By internal projectivity of $\ZZ[\NN_{\infty}]$ in $\Cond\Ab$ \cite[Thm. 2.3.3]{CamargoSolid}, the latter agrees with $\intHom_{\Cond\Ab}(\ZZ[\NN_{\infty}],-)|_{\Sph}$ \cite[Rmk. 3.2.2]{CamargoSolid} and hence preserves the heart. Furthermore, $\Cond\Sp$ is $\lim^{1}$-epi-acyclic by \Cref{rmk: 1-localic is lim1 epi acyclic}, its heart being Abelian sheaves on the replete 1-topos of condensed sets \cite[Prop. 2.1.4 \& Rmk. 2.2.2]{CamargoSolid}, so the hypotheses of \Cref{prop: lim1 acyclic t-exact} are satisfied and $\intHom_{\Cond\Sp}(\Sph[\NN_{\infty}],-)$ is $t$-exact.

    By exactness of $\intHom_{\Cond\Sp}(\Sph[\NN_{\infty}],-)$, preservation of small colimits can be reduced to the preservation of direct sums. Note that $\Cond\Sp_{\leq k}$ is closed under arbitrary direct sums in $\Cond\Sp$. Indeed, every direct sum is the filtered colimit of its finite subsums, the latter belong to $\Cond\Sp_{\leq k}$, and $\Cond\Sp_{\leq k}$ is closed under filtered colimits by compatibility with filtered colimits \Cref{thm: properties of condensed spectra} (2) (iii). Hence direct sums of objects of $\Cond\Sp_{\leq k}$ agree with those in $\Cond\Sp$. For a direct sum $\bigoplus_{i\in I}M_{i}$, we have $\tau_{\leq k}\bigoplus_{i\in I}M_{i}\simeq\bigoplus_{i\in I}\tau_{\leq k}M_{i}$ by commuting Postnikov truncation, a left adjoint, with the colimit so $\bigoplus_{i\in I}M_{i}\simeq\lim_{k}\tau_{\leq k}\bigoplus_{i\in I}M_{i}\simeq\lim_{k}\bigoplus_{i\in I}\tau_{\leq k}M_{i}$ by left completeness. We then compute 
    \begin{align*}
        \intHom_{\Cond\Sp}\left(\Sph[\NN_{\infty}],\bigoplus_{i\in I}M_{i}\right) &\simeq \lim_{k}\intHom_{\Cond\Sp}\left(\Sph[\NN_{\infty}],\bigoplus_{i\in I}\tau_{\leq k}M_{i}\right) \\
        &\simeq\lim_{k}\bigoplus_{i\in I}\intHom_{\Cond\Sp}\left(\Sph[\NN_{\infty}],\tau_{\leq k}M_{i}\right) && \text{\Cref{lem: Brink internal cohomology}} \\
        &\simeq \lim_{k}\bigoplus_{i\in I}\tau_{\leq k}\intHom_{\Cond\Sp}\left(\Sph[\NN_{\infty}],M_{i}\right) && t\text{-exactness} \\
        &\simeq \lim_{k}\tau_{\leq k}\bigoplus_{i\in I}\intHom_{\Cond\Sp}\left(\Sph[\NN_{\infty}],M_{i}\right) && \tau_{\leq k}\text{ left adj.} \\
        &\simeq \bigoplus_{i\in I}\intHom_{\Cond\Sp}(\Sph[\NN_{\infty}],M_{i}) && \text{left compl.}
    \end{align*}
    yielding the claim. 
\end{proof}
With this key computation in hand, the closure properties and the homotopy group criterion follow directly.
\begin{theorem}\label{thm: properties of solid spectra}
    Let $\Cond\Sp^{\sld}$ be the category of solid spectra. 
    \begin{enumerate}
        \item $\Cond\Sp^{\sld}$ is closed under small limits, small colimits, and extensions in $\Cond\Sp$. 
        \item If $M\in\Cond\Sp$ then $M\in\Cond\Sp^{\sld}$ if and only if $\pi_{k}(M)\in\Cond\Ab^{\sld}$ for all $k\in\ZZ$. 
    \end{enumerate}
\end{theorem}
\begin{proof}[Proof of (1)]
    The functor $\intHom_{\Cond\Sp}(\Sph[\NN_{\infty}],-)$ is a right adjoint, hence preserves small limits. Closure of $\Cond\Sp^{\sld}$ under small colimits follows from internal compact projectivity shown in \Cref{prop: nullsequence quotient is internally compact projective}. For closure under extensions, we let $M_{1}\to M_{2}\to M_{3}$ be a (co)fiber sequence in $\Cond\Sp$ with $M_{1},M_{3}\in\Cond\Sp^{\sld}$. Contemplate the diagram 
    $$% https://q.uiver.app/#q=WzAsNixbMCwwLCJNX3sxfSJdLFswLDEsIjAiXSxbMiwwLCJNX3syfSJdLFsyLDEsIk1fezN9Il0sWzQsMSwiTV97MX1bMV0iXSxbNCwwLCIwIl0sWzAsMl0sWzIsNV0sWzUsNF0sWzIsM10sWzMsNF0sWzAsMV0sWzEsM11d
    \begin{tikzcd}
        {M_{1}} && {M_{2}} && 0 \\
        0 && {M_{3}} && {M_{1}[1]}
        \arrow[from=1-1, to=1-3]
        \arrow[from=1-1, to=2-1]
        \arrow[from=1-3, to=1-5]
        \arrow[from=1-3, to=2-3]
        \arrow[from=1-5, to=2-5]
        \arrow[from=2-1, to=2-3]
        \arrow[from=2-3, to=2-5]
    \end{tikzcd}$$
    where both squares are biCartesian. Since $M_{1}[1],M_{3}\in\Cond\Sp^{\sld}$ so is $M_{2}$, being the fiber of $M_{3}\to M_{1}[1]$. 
\end{proof}
\begin{proof}[Proof of (2)]
    Suppose $M\in\Cond\Sp^{\sld}$. By $t$-exactness of $\intHom_{\Cond\Sp}(\Sph[\NN_{\infty}],-)$ we have 
    $$\pi_{k}\left(\intHom_{\Cond\Sp}(\Sph[\NN_{\infty}],M)\right)\simeq\intHom_{\Cond\Ab}(\ZZ[\NN_{\infty}],\pi_{k}(M))$$
    where the equivalence follows from \Cref{lem: modules inherit bistability} (2) and internal projectivity of $\ZZ[\NN_{\infty}]$ in $\Cond\Ab$, which identifies the derived internal hom with the underived internal hom \cite[Thm. 2.3.3 \& Rmk. 3.2.2]{CamargoSolid}. This equivalence is natural in $\id-\sigma^{*}$, so solidity of $M$ implies that $\pi_{k}(M)\in\Cond\Ab^{\sld}$ for every $k$. The same computation applied to an object $A\in\Cond\Ab\simeq\Cond\Sp^{\heartsuit}$ shows that $A\in\Cond\Ab^{\sld}$ if and only if $A$, regarded as a condensed spectrum, belongs to $\Cond\Sp^{\sld}$.

    Conversely suppose $\pi_{k}(M)\in\Cond\Ab^{\sld}$ for all $k\in \ZZ$. By right completeness and closure of $\Cond\Sp^{\sld}$ under small colimits from (1), it suffices to treat $M$ bounded below, and after shifting we may assume that $M$ is connective. By left completeness, we write $M\simeq\lim_{n}\tau_{\leq n}M$. Since $\Cond\Sp^{\sld}$ is closed under limits, it suffices to show that each term of the Postnikov tower lies in $\Cond\Sp^{\sld}$. For $n\geq 1$ the terms of the Postnikov tower form (co)fiber sequences 
    $$\pi_{n}(M)[n]\to\tau_{\leq n}M\to\tau_{\leq n-1}M$$
    from which the desideratum is obtained by induction: for $n=1$ the (co)fiber sequence is $\pi_{1}(M)[1]\to\tau_{\leq 1}M\to\pi_{0}(M)$ so $\tau_{\leq 1}M\in\Cond\Sp^{\sld}$ by closure under extensions. For larger $n$ we argue similarly, for the leftmost term is in $\Cond\Sp^{\sld}$ by assumption and the rightmost by the induction hypothesis. 
\end{proof}
\begin{remark}\label{rmk: definitions of solid spectra agree}
    \Cref{thm: properties of solid spectra} (2) shows that the definitions of solid spectra given by Wagner in \cite[\S B.6]{WagnerHHab} and Clausen in \cite[Def. 13.2]{ClausenLie} coincide. More generally, this is a manifestation of the fact that sufficiently nice localizations of stable presentable categories are determined by their heart. See \Cref{subsec: localizations of stable categories} for variations on this theme. 
\end{remark}
\section{Analytic \texorpdfstring{$\EE_{\infty}$}{E-infinity}-rings}\label{sec: analytic rings}
The category of all modules over a topological ring retains little of its topology. One often works instead inside a subcategory of complete modules -- those on which the topology is appropriately reflected -- and replaces the algebraic tensor product by a completed variant. 

In condensed mathematics, this is axiomatized by Clausen and Scholze \cite[Lecture VII]{Condensed} through the notion of an analytic ring: a condensed ring together with a full subcategory of its modules -- the analytic modules -- satisfying good closure properties in the ambient category and admitting an analytic refinement of the tensor product. 

Here we treat the fully commutative, ie. $\EE_{\infty}$, case, taking as base a connective $\EE_{\infty}$-algebra in condensed spectra.
\subsection{Basic definitions}\label{subsec: basic definitions}
The $\EE_{\infty}$-algebras in condensed spectra form a category $\Cond\CAlg:=\CAlg(\Cond\Sp)$ with full subcategory $\Cond\CAlg^{\cn}$ of connective objects. Since the forgetful functor $\CAlg\to\Sp$ preserves and reflects limits, we have 
$$\Cond\CAlg\simeq\widehat{\Sh}(\PFin_{\NN};\CAlg)\text{ and }\Cond\CAlg^{\cn}\simeq\widehat{\Sh}(\PFin_{\NN};\CAlg^{\cn}).$$

An analytic $\EE_{\infty}$-ring equips an object of $\Cond\CAlg^{\cn}$ with a distinguished subcategory of its modules. 
\begin{definition}\label{def: analytic E-infinity ring}
    Let $A^{\tri}\in\Cond\CAlg^{\cn}$. 
    \begin{enumerate}
        \item An \emph{analytic $\EE_{\infty}$-ring} on $A^{\tri}$-modules is the data of a pair $A:=(A^{\tri},\Mod_{A})$ where $\Mod_{A}\subseteq\Mod_{A^{\tri}}$ is a full subcategory of $A^{\tri}$-modules in condensed spectra satisfying the following additional conditions: 
        \begin{enumerate}[label=(\roman*)]
            \item $\Mod_{A}$ is closed under small limits and colimits in $\Mod_{A^{\tri}}$. 
            \item $\Mod_{A}$ is a $\Cond\Sp$-linear category: for $M\in\Mod_{A^{\tri}},N\in\Mod_{A}$ we have $\intHom_{A^{\tri}}(M,N)\in\Mod_{A}$. 
            \item The left adjoint $L^{A}:\Mod_{A^{\tri}}\to\Mod_{A}$ to the inclusion $\Mod_{A}\hookrightarrow\Mod_{A^{\tri}}$ is right $t$-exact (as an endofunctor of $\Mod_{A^{\tri}}$). 
        \end{enumerate} 
        We refer to the subcategory $\Mod_{A}$ of $\Mod_{A^{\tri}}$ as the subcategory of \emph{$A$-analytic modules} and the functor $L^{A}:\Mod_{A^{\tri}}\to\Mod_{A}$ as \emph{$A$-analytification}. 
        \item An analytic $\EE_{\infty}$-ring is \emph{complete} if $A^{\tri}\in\Mod_{A}$. 
        \item A \emph{morphism of analytic $\EE_{\infty}$-rings} is a morphism $A^{\tri}\to B^{\tri}$ in $\Cond\CAlg^{\cn}$ such that restriction of scalars along this morphism takes $B$-analytic modules to $A$-analytic modules. 
        \item Denote by $\An\CAlg$ the category with objects analytic $\EE_{\infty}$-rings as in (1) and morphisms those of analytic $\EE_{\infty}$-rings in the sense of (3). Denote by $\An\CAlg^{\wedge}$ the full subcategory of $\An\CAlg$ spanned by the complete analytic $\EE_{\infty}$-rings. 
    \end{enumerate}
\end{definition}
\begin{remark}\label{rmk: analytic v complete modules}
    Our nomenclature here is slightly nonstandard. In the existing literature, one typically refers to an analytic $\EE_{\infty}$-ring as a pre-analytic or uncompleted analytic ring, to a complete analytic $\EE_{\infty}$-ring as an analytic ring, and to $\Mod_{A}\subseteq\Mod_{A^{\tri}}$ as the $A$-complete modules. We have adopted our conventions to avoid confusion with the notion of completeness of analytic $\EE_{\infty}$-rings. 
\end{remark}
\begin{remark}\label{rmk: conditions on analytic modules}
    The condition that $\Mod_{A}$ is a $\Cond\Sp$-linear category is equivalent to the tensor product $-\otimes_{A^{\tri}}-$ on $\Mod_{A^{\tri}}$ preserving $L^{A}$-equivalences \cite[Lem. 7.5]{RagimovReflection}. This implies that $\Mod_{A}$ admits the structure of a symmetric monoidal category with respect to which $A$-analytification is a symmetric monoidal functor (cf. \Cref{prop: sym mon localizations}). 
\end{remark}
\begin{remark}\label{rmk: Ek and nonconnective}
    One might instead take as base an $\EE_{k}$-algebra in condensed spectra for $1\leq k<\infty$ without connectivity assumptions. We expect that an analytic theory over connective $\EE_{k}$-algebras can be developed along the same lines, with the module categories carrying $\EE_{k-1}$-monoidal structures rather than symmetric monoidal ones by Dunn additivity \cite[Cor. 5.1.2.6]{HA}, in analogy to Francis' less commutative algebraic geometry \cite{FrancisThesis}. The fully commutative case is adopted only because it suffices for our intended applications and affords a symmetric monoidal theory. Connectivity, on the other hand, is used essentially, as several results in the subsequent text rely crucially on it, and working in the nonconnective setting would require substantial modifications. 
\end{remark}
\subsection{\texorpdfstring{$A$}{A}-analytic modules}\label{subsec: analytic modules}
This subsection records the basic structure of the module category of a fixed analytic $\EE_{\infty}$-ring $(A^{\tri},\Mod_{A})$: stability, generation, the analytic tensor product, and the $t$-structure. For a complete analytic $\EE_{\infty}$-ring the module category is closely tied to its discrete modules, namely the ordinary $A^{\tri}(*)$-modules in spectra.
\begin{proposition}\label{prop: stability of A-analytic modules}\label{prop: A-analytic tensor product}
    Let $(A^{\tri},\Mod_{A})$ be an analytic $\EE_{\infty}$-ring. 
    \begin{enumerate}
        \item $\Mod_{A^{\tri}}$ is a stable $\omega_{1}$-presentably symmetric monoidal category, generated under small colimits and shifts by the $\omega_{1}$-compact objects $A^{\tri}[S]:=A^{\tri}\otimes_{\Sph}\Sph[S]$ for $S\in\PFin_{\NN}$. 
        \item $\Mod_{A}$ is a stable $\omega_{1}$-presentably symmetric monoidal category for the tensor product $M\otimes_{A}N:=L^{A}(M\otimes_{A^{\tri}}N)$, generated under small colimits and shifts by the $\omega_{1}$-compact objects $A[S]:=L^{A}(A^{\tri}[S])$, with respect to which $L^{A}$ is symmetric monoidal and the inclusion $\Mod_{A}\hookrightarrow\Mod_{A^{\tri}}$ lax symmetric monoidal. The unit is compact projective in the sense of \Cref{def: compact projective objects} (1), and internal homs are computed by $\intHom_{A}(M,N)\simeq\intHom_{A^{\tri}}(M,N)$. 
    \end{enumerate}
\end{proposition}
\begin{proof}[Proof of (1)]
    The category $\Cond\Sp$ is stable and $\omega_{1}$-presentably symmetric monoidal with compact projective unit and $\omega_{1}$-compact generators the $\Sph[S]$ by \Cref{thm: properties of condensed spectra} (1), so $\Mod_{A^{\tri}}$ is stable \cite[Prop. 7.1.1.4]{HA} and $\omega_{1}$-presentably symmetric monoidal \cite[Lem. 2.2]{AokiHigherLimits}. The forgetful functor $\Mod_{A^{\tri}}\to\Cond\Sp$ is conservative and preserves small limits and colimits \cite[Cor. 4.2.3.5]{HA}, so its left adjoint $A^{\tri}\otimes_{\Sph}-$ preserves $\omega_{1}$-compact objects and the $A^{\tri}[S]$ are $\omega_{1}$-compact. 
    
    By adjunction and conservativity of scalar restriction, joint conservativity of $\Sph[S]$ in $\Cond\Sp$ imply the joint conservativity of $A^{\tri}[S]$ in $\Mod_{A^{\tri}}$. 
\end{proof}
\begin{proof}[Proof of (2)]
    By $\Cond\Sp$-linearity of \Cref{def: analytic E-infinity ring} (1) (ii), the localization $L^{A}$ is compatible with the symmetric monoidal structure of $\Mod_{A^{\tri}}$ \cite[Lem. 7.4 \& 7.5]{RagimovReflection}, yielding the stated symmetric monoidal structure with $L^{A}$ symmetric monoidal and the inclusion lax symmetric monoidal. The formula for internal homs being \cite[Cor. 7.6]{RagimovReflection}. 
    
    Closure of $\Mod_{A}$ under small limits, colimits, and shifts in $\Mod_{A^{\tri}}$ gives stability, and since the inclusion preserves all small colimits, $L^{A}$ preserves $\omega_{1}$-compact objects, so the $A[S]$ are $\omega_{1}$-compact and generate $\Mod_{A}$ under small colimits and shifts by (1). Since $L^{A}$ is symmetric monoidal and preserves $\omega_{1}$-compact objects and generators, $\Mod_{A}$ is $\omega_{1}$-presentably symmetric monoidal as well. 
    
    Finally, the unit $L^{A}(A^{\tri})$ is connective by right $t$-exactness of $L^{A}$, and by adjunction $\Hom_{A}(L^{A}(A^{\tri}),-)\simeq\Hom_{\Cond\Sp}(\Sph,(-)|_{\Sph})$ where $(-)|_{\Sph}:\Mod_{A}\to\Cond\Sp$ is $t$-exact and preserves small colimits, so compact projectivity of the unit reduces to that of $\Sph$ established in \Cref{thm: properties of condensed spectra} (1). 
\end{proof}
Membership in $\Mod_{A}$ is detected on the generators of \Cref{prop: stability of A-analytic modules}.
\begin{corollary}\label{cor: testing by free on profinites}
    Let $(A^{\tri},\Mod_{A})$ be an analytic $\EE_{\infty}$-ring and $M\in\Mod_{A^{\tri}}$. Then $M\in\Mod_{A}$ if and only if the natural morphism 
    \begin{equation}\label{eqn: analytic by free objects}
        \intHom_{A^{\tri}}(A[S],M)\longrightarrow\intHom_{A^{\tri}}(A^{\tri}[S],M)
    \end{equation}
    is an equivalence for all $S\in\PFin_{\NN}$. 
\end{corollary}
\begin{proof}
    $(\Rightarrow)$ This is \cite[Cor. 7.6]{RagimovReflection}. 

    $(\Leftarrow)$ Suppose the natural morphism (\ref{eqn: analytic by free objects}) is an equivalence for all $S\in\PFin_{\NN}$. By \Cref{prop: stability of A-analytic modules} (2) the $A^{\tri}[S]$ generate $\Mod_{A^{\tri}}$ under colimits and shifts, so $M$ is a colimit of shifts of the $A^{\tri}[S]$ and $L^{A}(M)$ the colimit of the corresponding shifts of the $A[S]$, with $\eta_{M}:M\to L^{A}(M)$ their colimit. Since $\intHom_{A^{\tri}}(-,M)$ carries colimits to limits and shifts to shifts up to sign, applying it to $\eta_{M}$ returns a limit of shifts of (\ref{eqn: analytic by free objects}), so $\intHom_{A^{\tri}}(\eta_{M},M):\intHom_{A^{\tri}}(L^{A}(M),M)\to\intHom_{A^{\tri}}(M,M)$ is an equivalence. The preimage $r:L^{A}(M)\to M$ of $\id_{M}$ then satisfies $r\circ\eta_{M}\simeq\id_{M}$, exhibiting $M$ as a retract of $L^{A}(M)$ in $\Mod_{A^{\tri}}$. Now $\Mod_{A}$ is closed under retracts in $\Mod_{A^{\tri}}$ by \Cref{prop: properties of pi stably reflective localizations} (1), so $M\in\Mod_{A}$. 
\end{proof}
As for algebras in spectra, both module categories carry a $t$-structure for which the forgetful functors to $\Cond\Sp$ are $t$-exact.
\begin{proposition}\label{prop: t-structure on A-analytic modules}
    Let $(A^{\tri},\Mod_{A})$ be an analytic $\EE_{\infty}$-ring. The category $\Mod_{A}$ admits an accessible $t$-structure $(\Mod_{A,\geq0},\Mod_{A,\leq0})$ with the following properties:
    \begin{enumerate}[label=(\roman*)]
        \item $\Mod_{A,\geq0},\Mod_{A,\leq0}$ are spanned by those $A$-analytic modules $M$ such that the homotopy sheaves $\pi_{k}(M)$ vanish for $k<0$ and $k>0$, respectively.
        \item $\Mod_{A,\geq0}$ contains the unit in $\Mod_{A}$ and is closed under the $A$-analytic tensor product in $\Mod_{A}$. 
        \item $\Mod_{A,\leq0}$ is closed under filtered colimits in $\Mod_{A}$. 
        \item $\Mod_{A,\geq0}$ is closed under countable products in $\Mod_{A}$.
        \item The $t$-structure on $\Mod_{A}$ is both left and right complete. 
        \item An $A^{\tri}$-module $M$ in condensed spectra is $A$-analytic if and only if $\pi_{k}(M)\in\Mod_{A}^{\heartsuit}$ for all $k\in\ZZ$. 
    \end{enumerate}
\end{proposition} 
\begin{proof}
    All but (iv) and (vi) hold for $\Mod_{A}=\Mod_{A^{\tri}}$ by \cite[Prop. 2.1.1.1, War. 2.1.1.2, \& Rmk. 2.1.2.1]{SAG}. For (iv), the forgetful functor $\Mod_{A^{\tri}}\to\Cond\Sp$ preserves countable products and reflects connectivity, so the claim reduces to closure of $\Cond\Sp_{\geq0}$ under countable products, which holds by \Cref{thm: properties of condensed spectra} (2) (iv). 

    When $\Mod_{A}\subsetneq\Mod_{A^{\tri}}$, conditions (i) and (iii) of \Cref{def: analytic E-infinity ring} exhibit $\Mod_{A}$ as a $\pi$-stably reflective subcategory (\Cref{def: pi stably reflective}) of the $t$-bistable $\Mod_{A^{\tri}}$, so \Cref{prop: properties of pi stably reflective localizations} and \Cref{prop: compatible localization implies pres sym mon t-bistable} apply to give the conclusion for (vi). 
\end{proof}
When an analytic $\EE_{\infty}$-ring is complete, the underlying ring $A^{\tri}$ is itself the unit of $\Mod_{A}$, and its global sections generate a fully faithful copy of $\Mod_{A^{\tri}(*)}(\Sp)$ inside $\Mod_{A}$.
\begin{proposition}\label{prop: discrete modules}
    Let $(A^{\tri},\Mod_{A})$ be a complete analytic $\EE_{\infty}$-ring. 
    \begin{enumerate}
        \item There is a fully faithful functor $\Cond_{A^{\tri}(*)}:\Mod_{A^{\tri}(*)}(\Sp)\hookrightarrow\Mod_{A}$. 
        \item An object of $\Mod_{A}$ lies in the essential image of the embedding in (1) if and only if it lies in the subcategory of $\Mod_{A}$ generated by $A^{\tri}$ under colimits and shifts. 
    \end{enumerate} 
\end{proposition}
\begin{proof}[Proof of (1)]
    By completeness $A^{\tri}$ is the unit of $\Mod_{A}$ which is compact by \Cref{prop: A-analytic tensor product}. Let $\Escr\subseteq\Mod_{A}$ be the smallest full subcategory containing $A^{\tri}$ and closed under colimits and shifts. Then $\Escr$ is stable and presentable and $A^{\tri}$ is a compact generator of it, so the Schwede--Shipley theorem \cite[Prop. 7.1.2.7]{HA} yields an equivalence $\Escr\simeq\Mod_{A^{\tri}(*)}(\Sp)$ where $\Hom_{A}(A^{\tri},A^{\tri})\simeq A^{\tri}(*)$ as $\EE_{\infty}$-rings, identifying the endomorphisms of the unit with its global sections. The desired embedding is the composite $\Mod_{A^{\tri}(*)}(\Sp)\simeq\Escr\hookrightarrow\Mod_{A}$, fully faithful as $\Escr$ is a full subcategory of $\Mod_{A}$. 
\end{proof}
\begin{proof}[Proof of (2)]
    The essential image of the embedding of (1) is $\Escr$, which by construction is the subcategory of $\Mod_{A}$ generated by $A^{\tri}$ under colimits and shifts. 
\end{proof}
The objects so embedded are the discrete ones.
\begin{definition}\label{def: discrete modules}
    Let $(A^{\tri},\Mod_{A})$ be a complete analytic $\EE_{\infty}$-ring. An object $M\in\Mod_{A}$ is \emph{discrete} if it lies in the essential image of the fully faithful embedding $\Mod_{A^{\tri}(*)}(\Sp)\hookrightarrow\Mod_{A}$. 
\end{definition}
The relationship between discreteness and finiteness of analytic modules is taken up in \Cref{subsec: mereology of analytic modules}.
\subsection{Functoriality and colimits in \texorpdfstring{$\An\CAlg$}{AnCAlg}}\label{subsec: functoriality and colimits}
The previous subsection fixed an analytic $\EE_{\infty}$-ring. We now let the ring vary and record the basic functoriality of $\An\CAlg$, before establishing that $\An\CAlg$ admits small colimits.

The category of $A^{\tri}$-modules in $\Cond\Sp$ always defines an analytic structure and is complete since $A^{\tri}\in\Mod_{A^{\tri}}$.
\begin{definition}\label{def: trivial analytic ring structure}
    Let $(A^{\tri},\Mod_{A})$ be an analytic $\EE_{\infty}$-ring. $(A^{\tri},\Mod_{A})$ is \emph{trivial} if $\Mod_{A}=\Mod_{A^{\tri}}$. 
\end{definition}
The trivial structure is functorial in the underlying ring and left adjoint to the forgetful functor.
\begin{proposition}\label{prop: trivial is left adjoint}
    The functor $\Cond\CAlg^{\cn}\to\An\CAlg$ given on objects by $A^{\tri}\mapsto(A^{\tri},\Mod_{A^{\tri}})$ is fully faithful with essential image in $\An\CAlg^{\wedge}$ and left adjoint to the forgetful functor given on objects by $(A^{\tri},\Mod_{A})\mapsto A^{\tri}$. 
\end{proposition}
\begin{proof}
    For any analytic $\EE_{\infty}$-ring $(B^{\tri},\Mod_{B})$ we have an equivalence of animae 
    $$\Mor_{\An\CAlg}\left((A^{\tri},\Mod_{A^{\tri}}),(B^{\tri},\Mod_{B})\right)\simeq\Mor_{\Cond\CAlg^{\cn}}\left(A^{\tri},B^{\tri}\right)$$
    as the condition that scalar restriction takes $B$-analytic modules to $A^{\tri}$-modules is tautologically satisfied. This proves the desired adjunction and that it is fully faithful. The description of the essential image is clear as $A^{\tri}\in\Mod_{A^{\tri}}$. 
\end{proof}
A morphism of analytic $\EE_{\infty}$-rings induces a symmetric monoidal scalar extension functor. 
\begin{proposition}\label{prop: scalar extension functor is symmetric monoidal}
    Let $(A^{\tri},\Mod_{A})\to(B^{\tri},\Mod_{B})$ be a morphism of analytic $\EE_{\infty}$-rings. Scalar restriction $\Mod_{B}\to\Mod_{A}$ admits a symmetric monoidal left adjoint $B\otimes_{A}-:\Mod_{A}\to\Mod_{B}$ for which the square
    \begin{equation}\label{eqn: scalar extension is sym mon}
        % https://q.uiver.app/#q=WzAsNCxbMCwwLCJcXE1vZF97QV57XFx0cml9fSJdLFsyLDAsIlxcTW9kX3tCXntcXHRyaX19Il0sWzAsMSwiXFxNb2Rfe0F9Il0sWzIsMSwiXFxNb2Rfe0J9Il0sWzEsMywiTF57Qn0iXSxbMCwxLCJCXntcXHRyaX1cXG90aW1lc197QV57XFx0cml9fS0iXSxbMCwyLCJMXntBfSIsMl0sWzIsMywiQlxcb3RpbWVzX3tBfS0iLDJdXQ==
        \begin{tikzcd}
            {\Mod_{A^{\tri}}} && {\Mod_{B^{\tri}}} \\
            {\Mod_{A}} && {\Mod_{B}}
            \arrow["{B^{\tri}\otimes_{A^{\tri}}-}", from=1-1, to=1-3]
            \arrow["{L^{A}}"', from=1-1, to=2-1]
            \arrow["{L^{B}}", from=1-3, to=2-3]
            \arrow["{B\otimes_{A}-}"', from=2-1, to=2-3]
        \end{tikzcd}
    \end{equation}
    is commutative. 
\end{proposition}
\begin{proof}
    The forgetful functor $\Mod_{B}\to\Mod_{A}$ is a functor between presentable categories preserving small limits and colimits, so the adjoint functor theorem \cite[Cor. 5.5.2.9]{HTT} furnishes a left adjoint $B\otimes_{A}-$. The composite $L^{B}\circ(B^{\tri}\otimes_{A^{\tri}}-):\Mod_{A^{\tri}}\to\Mod_{B}$ is symmetric monoidal, being a composite of such, and we claim it annihilates $\ker(L^{A})$. Let $K\in\ker(L^{A})\subseteq\Mod_{A^{\tri}}$ and $M\in\Mod_{B}$. Since the map is a morphism of analytic rings, $M|_{A^{\tri}}\in\Mod_{A}$, and we compute
    \begin{align*}
        \Hom_{B^{\tri}}(L^{B}(B^{\tri}\otimes_{A^{\tri}}K),M) &\simeq \Hom_{B^{\tri}}(B^{\tri}\otimes_{A^{\tri}}K,M) && L^{B}\dashv\iota_{B}\\
        &\simeq \Hom_{A^{\tri}}(K,M|_{A^{\tri}}) && B^{\tri}\otimes_{A^{\tri}}(-)\dashv(-)|_{A^{\tri}}\\
        &\simeq \Hom_{A^{\tri}}(L^{A}(K),M|_{A^{\tri}}) && L^{A}\dashv\iota_{A},\ M|_{A^{\tri}}\in\Mod_{A}\\
        &\simeq 0 && K\in\ker(L^{A})
    \end{align*}
    where $\iota_{A},\iota_{B}$ denote the inclusions of the analytic modules. So $L^{B}(B^{\tri}\otimes_{A^{\tri}}K)\simeq0$. By \cite[Thm. I.3.6 \& App. S]{TC}, the symmetric monoidal functor $L^{B}\circ B^{\tri}\otimes_{A^{\tri}}-:\Mod_{A^{\tri}}\to\Mod_{B}$ descends to a symmetric monoidal functor $B\otimes_{A}-:\Mod_{A}\to\Mod_{B}$ left adjoint to $(-)|_{A^{\tri}}:\Mod_{B}\to\Mod_{A}$.
\end{proof}
The small colimits of $\An\CAlg$ are assembled from pushouts along trivial structures, for which we first isolate the analytic structure a ring map induces on its target.
\begin{proposition}\label{prop: induced analytic ring structure}
    Let $(A^{\tri},\Mod_{A})$ be an analytic $\EE_{\infty}$-ring and $A^{\tri}\to B^{\tri}$ a morphism in $\Cond\CAlg^{\cn}$. 
    \begin{enumerate}
        \item The full subcategory of $\Mod_{B^{\tri}}$ spanned by those objects $M$ such that $M|_{A^{\tri}}\in\Mod_{A}$ defines an analytic $\EE_{\infty}$-ring $(B^{\tri},\Mod_{B^{\tri}/A})$. 
        \item The analytic $\EE_{\infty}$-ring $(B^{\tri},\Mod_{B^{\tri}/A})$ of (1) is the pushout of 
        $$(A^{\tri},\Mod_{A})\leftarrow(A^{\tri},\Mod_{A^{\tri}})\to(B^{\tri},\Mod_{B^{\tri}})$$
        in $\An\CAlg$. 
    \end{enumerate}
\end{proposition}
\begin{proof}[Proof of (1)]
    The forgetful functor $(-)|_{A^{\tri}}:\Mod_{B^{\tri}}\to\Mod_{A^{\tri}}$ is conservative, $t$-exact, and preserves small limits and colimits so $\Mod_{B^{\tri}/A}$ admits the induced $t$-structure from $\Mod_{B^{\tri}}$, $\Mod_{B^{\tri}/A}$ is closed under small limits and colimits in $\Mod_{B^{\tri}}$ and there is an adjunction $L^{B^{\tri}/A}:\Mod_{B^{\tri}}\rightleftarrows\Mod_{B^{\tri}/A}:\iota_{B^{\tri}/A}$ by the adjoint functor theorem \cite[Cor. 5.5.2.9]{HTT}. By adjunction between $B^{\tri}\otimes_{A^{\tri}}-:\Mod_{A^{\tri}}\rightleftarrows\Mod_{B^{\tri}}:(-)|_{A^{\tri}}$ we have 
    $$\intHom_{B^{\tri}}(B^{\tri}[S],N)|_{A^{\tri}}\simeq\intHom_{A^{\tri}}(A^{\tri}[S],N|_{A^{\tri}})$$
    for all $N\in\Mod_{B^{\tri}/A}$, where the right-hand side lies in $\Mod_{A}$ by \Cref{def: analytic E-infinity ring} (1) (ii). A general $M\in\Mod_{B^{\tri}}$ is a colimit of shifts of the $B^{\tri}[S]$ by \Cref{prop: stability of A-analytic modules} (2), and $\intHom_{B^{\tri}}(-,N)|_{A^{\tri}}$ carries colimits to limits and shifts to shifts up to sign, so $\intHom_{B^{\tri}}(M,N)|_{A^{\tri}}\in\Mod_{A}$ by closure of $\Mod_{A}$ under small limits and shifts. Thus $\Mod_{B^{\tri}/A}$ is $\Cond\Sp$-linear. It remains to show $L^{B^{\tri}/A}$ is right $t$-exact. For $N\in\Mod_{B^{\tri}/A,\leq -1},M\in\Mod_{B^{\tri},\geq0}$ the mapping spectra $\Hom_{B^{\tri}}(M, N)\simeq\Hom_{B^{\tri}}(L^{B^{\tri}/A}(M),N)$ have trivial 0th homotopy groups so $L^{B^{\tri}/A}$ is right $t$-exact. This verifies the conditions of \Cref{def: analytic E-infinity ring} (1) showing $(B^{\tri},\Mod_{B^{\tri}/A})$ is an analytic $\EE_{\infty}$-ring. 
\end{proof}
\begin{proof}[Proof of (2)]
    Let $(C^{\tri},\Mod_{C})\in\An\CAlg$. By the Yoneda lemma it suffices to show that the square of animae
    $$
    \begin{tikzcd}
        {\Mor_{\An\CAlg}\left((B^{\tri},\Mod_{B^{\tri}/A}),(C^{\tri},\Mod_{C})\right)} && {\Mor_{\Cond\CAlg^{\cn}}\left(B^{\tri},C^{\tri}\right)} \\
        {\Mor_{\An\CAlg}\left((A^{\tri},\Mod_{A}),(C^{\tri},\Mod_{C})\right)} && {\Mor_{\Cond\CAlg^{\cn}}\left(A^{\tri},C^{\tri}\right)}
        \arrow[from=1-1, to=1-3]
        \arrow[from=1-1, to=2-1]
        \arrow[from=1-3, to=2-3]
        \arrow[from=2-1, to=2-3]
    \end{tikzcd}
    $$
    is Cartesian, where the horizontal maps send an analytic morphism to its underlying ring map and are, by \Cref{def: analytic E-infinity ring} (3), the inclusions of the subanimae cut out by the scalar-restriction condition. The fibered product of the bottom and right morphisms is the anima of ring maps $B^{\tri}\to C^{\tri}$ whose restriction along $A^{\tri}\to B^{\tri}\to C^{\tri}$ sends $\Mod_{C}$ into $\Mod_{A}$. Since $\Mod_{B^{\tri}/A}$ consists precisely of those $B^{\tri}$-modules whose underlying $A^{\tri}$-module is $A$-analytic, this is exactly the anima of ring maps whose restriction along $B^{\tri}\to C^{\tri}$ sends $\Mod_{C}$ into $\Mod_{B^{\tri}/A}$, namely the top-left corner. So the square is Cartesian. 
\end{proof}
\begin{definition}\label{def: induced analytic ring}
    Let $(A^{\tri},\Mod_{A})\to(B^{\tri},\Mod_{B})$ be a morphism of analytic $\EE_{\infty}$-rings. $(B^{\tri},\Mod_{B})$ has the \emph{induced analytic $\EE_{\infty}$-ring structure} from $(A^{\tri},\Mod_{A})$ if $\Mod_{B}=\Mod_{B^{\tri}/A}$. 
\end{definition}
Colimits along a fixed underlying ring are computed as intersections of module categories.
\begin{lemma}\label{lem: ringwise colimits}
    Let $\{(A^{\tri},\Mod_{A_{i}})\}_{i\in I}$ be a diagram of analytic $\EE_{\infty}$-ring structures on $A^{\tri}$-modules for a fixed $A^{\tri}\in\Cond\CAlg^{\cn}$. 
    \begin{enumerate}
        \item The colimit $\colim_{i\in I}(A^{\tri},\Mod_{A_{i}})$ in analytic $\EE_{\infty}$-rings under $(A^{\tri},\Mod_{A^{\tri}})$ exists and is given by 
        $$\colim_{i\in I}(A^{\tri},\Mod_{A_{i}})\simeq\left(A^{\tri},\bigcap_{i\in I}\Mod_{A_{i}}\right).$$
        \item If all analytic ring structures in (1) are assumed to be complete, then the colimit is complete as well. 
    \end{enumerate}
\end{lemma}
\begin{proof}[Proof of (1)]
    Note that $(A^{\tri},\Mod_{A_{i}})\to(A^{\tri},\Mod_{A_{j}})$ is a morphism of analytic $\EE_{\infty}$-rings if and only if $\Mod_{A_{j}}$ is a subcategory of $\Mod_{A_{i}}$. Observe \Cref{def: analytic E-infinity ring} (1) (i) and (ii) are preserved under intersections in $\Mod_{A^{\tri}}$. By \cite[Lem. 5.5.4.18]{HTT}, the intersection $\bigcap_{i\in I}\Mod_{A_{i}}$ is the objects local for the union of small collections generating the $L^{A_{i}}$-equivalences. Since each $\Mod_{A_{i}}$ is closed under Postnikov and Whitehead truncation in $\Mod_{A^{\tri}}$ by \Cref{prop: properties of pi stably reflective localizations} (2), so is the intersection. Let $L^{A}$ be the left adjoint to the inclusion $\Mod_{A}:=\bigcap_{i\in I}\Mod_{A_{i}}\hookrightarrow\Mod_{A^{\tri}}$. Assume $M\in\Mod_{A^{\tri},\geq0}$ and $N\in\Mod_{A,\leq-1}$. By adjunction, $\Hom_{A^{\tri}}(L^{A}(M),N)\simeq\Hom_{A^{\tri}}(M,N)$ has trivial $\pi_{0}$. So the truncation map $L^{A}(M)\to\tau_{\leq-1}L^{A}(M)$ is zero and therefore $\tau_{\leq-1}L^{A}(M)\simeq0$. Thus $L^{A}(M)$ is connective and $L^{A}$ is right $t$-exact, verifying \Cref{def: analytic E-infinity ring} (1) (iii). This shows that $(A^{\tri},\bigcap_{i\in I}\Mod_{A_{i}})$ is an analytic $\EE_{\infty}$-ring. 
    
    For the universal property, fix an analytic $\EE_{\infty}$-ring $(B^{\tri},\Mod_{B})$ under $(A^{\tri},\Mod_{A^{\tri}})$. For any $i$, the map $(A^{\tri},\Mod_{A^{\tri}})\to(B^{\tri},\Mod_{B})$ descends to one $(A^{\tri},\Mod_{A_{i}})\to(B^{\tri},\Mod_{B})$ if and only if scalar restriction takes $B$-analytic modules to $A_{i}$-analytic ones so the map descends to one $\left(A^{\tri},\bigcap_{i\in I}\Mod_{A_{i}}\right)\to(B^{\tri},\Mod_{B})$ if and only if scalar restriction takes $B$-analytic modules to $A_{i}$-analytic ones for all $i\in I$, which is precisely the universal property of the colimit (cf. \cite[Lem. 4.1.10]{CamargoSolid}). 
\end{proof}
\begin{proof}[Proof of (2)]
    Membership of $A^{\tri}$ in the subcategory $\Mod_{A_{i}}$ is preserved under intersection as well. 
\end{proof}
The general case follows from a result on colimits in coCartesian fibrations over a cocomplete base.
\begin{proposition}\label{prop: colimits in AnCAlg}
   Let $\{(A_{i}^{\tri},\Mod_{A_{i}})\}_{i\in I}$ be a small diagram in $\An\CAlg$. Then the colimit $\colim_{i\in I}(A_{i}^{\tri},\Mod_{A_{i}})$ exists and is given by 
   $$\colim_{i\in I}(A_{i}^{\tri},\Mod_{A_{i}})\simeq(B^{\tri},\Mod_{B}):=\left(B^{\tri},\bigcap_{i\in I}\Mod_{B^{\tri}/A_{i}}\right)$$
   where $B^{\tri}$ is the colimit $\colim_{i\in I}A^{\tri}_{i}$ in $\Cond\CAlg^{\cn}$. 
\end{proposition}
\begin{proof}
    The forgetful functor $\An\CAlg\to\Cond\CAlg^{\cn}$, $(A^{\tri},\Mod_{A})\mapsto A^{\tri}$, of \Cref{prop: trivial is left adjoint} is a coCartesian fibration: for a morphism $A^{\tri}\to B^{\tri}$ and an analytic structure $\Mod_{A}$ over $A^{\tri}$, the induced structure of \Cref{prop: induced analytic ring structure} supplies a coCartesian lift, its universal property being the pushout of \Cref{prop: induced analytic ring structure} (2). Fibre transport is thus $\Mod_{A}\mapsto\Mod_{B^{\tri}/A}$, which preserves colimits as $\Mod_{B^{\tri}/\bigcap_{i}A_{i}}=\bigcap_{i}\Mod_{B^{\tri}/A_{i}}$. 
    
    Now $\Cond\CAlg^{\cn}$ is presentable, hence cocomplete, and each fibre is cocomplete by \Cref{lem: ringwise colimits}, so \cite[Lem. 10.7]{LaxFree} gives that $\An\CAlg$ is cocomplete, the colimit being computed by transporting each $\Mod_{A_{i}}$ to the fibre over $B^{\tri}=\colim_{i\in I}A_{i}^{\tri}$ and intersecting, namely $(B^{\tri},\bigcap_{i\in I}\Mod_{B^{\tri}/A_{i}})$. 
\end{proof}
On the free objects a sifted colimit is computed termwise, which we record for later use.
\begin{lemma}\label{lem: sifted colimits in AnCAlg}
    Let $\{(A_{i}^{\tri},\Mod_{A_{i}})\}_{i\in I}$ be a sifted diagram in $\An\CAlg$ with colimit $(B^{\tri},\Mod_{B})$, and let $S\in\PFin_{\NN}$. The colimit $\colim_{i\in I}A_{i}[S]$ of underlying condensed spectra carries a natural $B^{\tri}$-module structure with respect to which it is $B$-analytic and equivalent to $B[S]$.
\end{lemma}
\begin{proof}
    Since $I$ is sifted, $\colim_{i\in I}A_{i}^{\tri}\simeq B^{\tri}$ acts on $M:=\colim_{i\in I}A_{i}[S]$, the colimit of underlying condensed spectra, making it a $B^{\tri}$-module: the diagonal $I\to I\times I$ is cofinal, so the action maps assemble into $B^{\tri}\otimes_{\Sph}M\simeq\colim_{i\in I}(A_{i}^{\tri}\otimes_{\Sph}A_{i}[S])\to\colim_{i\in I}A_{i}[S]=M$, and moreover $M\simeq\colim_{i\in I}(B^{\tri}\otimes_{A_{i}^{\tri}}A_{i}[S])$ in $\Mod_{B^{\tri}}$.

    Since $L^{B}$ preserves colimits, $L^{B}(M)\simeq\colim_{i\in I}L^{B}(B^{\tri}\otimes_{A_{i}^{\tri}}A_{i}[S])$. The factorization $L^{B}\circ(B^{\tri}\otimes_{A_{i}^{\tri}}-)\simeq(B\otimes_{A_{i}}-)\circ L^{A_{i}}$ of \Cref{prop: scalar extension functor is symmetric monoidal}, applied to the analytification unit $\eta_{i}: A_{i}^{\tri}[S]\to A_{i}[S]$ inverted by $L^{A_{i}}$, shows that $L^{B}$ inverts $B^{\tri}\otimes_{A_{i}^{\tri}}\eta_{i}: B^{\tri}[S]\to B^{\tri}\otimes_{A_{i}^{\tri}}A_{i}[S]$, where $B^{\tri}\otimes_{A_{i}^{\tri}}A_{i}^{\tri}[S]\simeq B^{\tri}[S]$. The diagram $i\mapsto B^{\tri}\otimes_{A_{i}^{\tri}}A_{i}^{\tri}[S]$ is thus constant at $B^{\tri}[S]$, and the $\eta_{i}$ are natural in $i$, so $L^{B}(M)\simeq\colim_{i\in I}B[S]\simeq B[S]$. It therefore suffices to show that $M$ is $B$-analytic, in which case the unit $M\to L^{B}(M)$ is an equivalence.

    Fix $i\in I$. We claim the forgetful map $I_{i/}\to I$ is cofinal. Siftedness gives the weak contractibility of $I\times_{I\times I}(I\times I)_{(i,j)/}$ for every $(i,j)\in I\times I$, and $(I\times I)_{(i,j)/}\simeq I_{i/}\times I_{j/}$ identifies with $I_{i/}\times_{I}I_{j/}$ via the diagonal. Quillen's Theorem A \cite[Thm. 4.1.3.1]{HTT} applied to $I_{i/}\to I$ proves the claim. Restriction of scalars makes $(\alpha: i\to j)\mapsto A_{j}[S]|_{A_{i}^{\tri}}$ a diagram $I_{i/}\to\Mod_{A_{i}^{\tri}}$ with values in $\Mod_{A_{i}}$, as each morphism is one of analytic $\EE_{\infty}$-rings by \Cref{def: analytic E-infinity ring} (3). Denote its colimit by $N$, which lies in $\Mod_{A_{i}}$ by closure under colimits. The natural maps $A_{j}[S]\to M|_{A_{j}^{\tri}}$ restrict to $A_{i}^{\tri}$-linear maps inducing equivalences $N\to M|_{A_{i}^{\tri}}$: the forgetful functor $\Mod_{A_{i}^{\tri}}\to\Cond\Sp$ is conservative and preserves small colimits \cite[Cor. 4.2.3.5]{HA}, and the underlying map is the comparison $\colim_{(i\to j)\in I_{i/}}A_{j}[S]\to M$, invertible by cofinality \cite[Prop. 4.1.1.8]{HTT}. Hence $M|_{A_{i}^{\tri}}\in\Mod_{A_{i}}$. Since $i\in I$ was arbitrary, $M\in\bigcap_{i\in I}\Mod_{B^{\tri}/A_{i}}=\Mod_{B}$ by \Cref{prop: colimits in AnCAlg}, whence $M\simeq L^{B}(M)\simeq B[S]$. 
\end{proof}
\subsection{Completion}\label{subsec: completion}
We produce a left adjoint to the inclusion $\An\CAlg^{\wedge}\hookrightarrow\An\CAlg$, so that every analytic $\EE_{\infty}$-ring admits a completion, obtained by applying the analytification functor to its underlying ring.
\begin{theorem}\label{thm: completions exist}
    The full inclusion $\An\CAlg^{\wedge}\hookrightarrow\An\CAlg$ admits a left adjoint $(-)^{\wedge}:\An\CAlg\to\An\CAlg^{\wedge}$ given on objects by $(A^{\tri},\Mod_{A})\mapsto (A^{\tri,=},\Mod_{A})$ where $A^{\tri,=}:=L^{A}(A^{\tri})$ and $L^{A}$ the $A$-analytification functor. 
\end{theorem}
\begin{proof}
    The following proof is due to Longke Tang and Yixiao Li provided to us by Peter Scholze in private correspondence \cite{ScholzePC}.

    By the functor $L^{A}$ exhibits $\Mod_{A}$ as a symmetric monoidal localization of $\Mod_{A^{\tri}}$, with symmetric monoidal unit $A^{\tri,=}:=L^{A}(A^{\tri})$. The unit of the localization $\eta:A^{\tri}\to A^{\tri,=}$ is a morphism in $\Cond\CAlg^{\cn}$. By \cite[Const. 5.23]{NilpotenceDescent}, the forgetful functor $R:\Mod_{A}\hookrightarrow\Mod_{A^{\tri}}$ factors as 
    $$\Mod_{A}\xrightarrow{\overline{R}}\Mod_{A^{\tri,=}}\xrightarrow{(-)|_{A^{\tri}}}\Mod_{A^{\tri}}$$ 
    and there is an induced left adjoint $\overline{L}:\Mod_{A^{\tri,=}}\to\Mod_{A}$ which identifies with $L^{A}$. The counit of the adjunction $\overline{L}\dashv\overline{R}$ is the canonical equivalence $L^{A}(N)\simeq N$ for $N\in\Mod_{A}$ so $\overline{R}$ is fully faithful, and the unit is an equivalence precisely on those objects of $\Mod_{A^{\tri,=}}$ whose underlying $A^{\tri}$-modules are $A$-analytic. So $\overline{L}$ is the localization for the induced structure $\Mod_{A^{\tri,=}/A}$ and conservativity of the forgetful functor $(-)|_{A^{\tri}}$ defines an equivalence of categories $\Mod_{A^{\tri,=}/A}\simeq\Mod_{A}$ identifying the induced structure with $(A^{\tri,=},\Mod_{A})$ which is complete as $A^{\tri,=}|_{A^{\tri}}\in\Mod_{A}$ by construction. 

    Finally, let $(B^{\tri},\Mod_{B})\in\An\CAlg^{\wedge}$ and let $f:(A^{\tri},\Mod_{A})\to(B^{\tri},\Mod_{B})$ be a morphism of analytic $\EE_{\infty}$-rings. Completeness gives $B^{\tri}\in\Mod_{B}$ and analyticity of $f$ then gives $B^{\tri}|_{A^{\tri}}\in\Mod_{A}$, so $B^{\tri}$, regarded as a commutative algebra object of $\Mod_{A^{\tri}}$, lies in the essential image of the fully faithful functor $\CAlg(\Mod_{A})\to\CAlg(\Mod_{A^{\tri}})$, under which $A^{\tri,=}$ corresponds to the unit of $\Mod_{A}$. Since the unit is initial in $\CAlg(\Mod_{A})$ \cite[Cor. 3.2.1.9]{HA}, the anima of morphisms of condensed $\EE_{\infty}$-rings $A^{\tri,=}\to B^{\tri}$ under $A^{\tri}$ (ie. factorizations of $f$ over $\eta$) is contractible. The unique factorization $g:A^{\tri,=}\to B^{\tri}$ is a morphism of analytic rings $(A^{\tri,=},\Mod_{A})\to(B^{\tri},\Mod_{B})$, since $N|_{A^{\tri}}\in\Mod_{A}$ for $N\in\Mod_{B}$. Conversely, composition with the canonical morphism $(A^{\tri},\Mod_{A})\to(A^{\tri,=},\Mod_{A})$ preserves analyticity. This gives an equivalence 
    $$\Mor_{\An\CAlg^{\wedge}}\left((A^{\tri,=},\Mod_{A}),(B^{\tri},\Mod_{B})\right)\simeq\Mor_{\An\CAlg}\left((A^{\tri},\Mod_{A}),(B^{\tri},\Mod_{B})\right),$$
    naturally in $(B^{\tri},\Mod_{B})\in\An\CAlg^{\wedge}$ giving the claim. 
\end{proof}
Reflectivity transports the colimits of \Cref{prop: colimits in AnCAlg} to $\An\CAlg^{\wedge}$.
\begin{corollary}\label{cor: complete AnCAlg has colimits}
    The category $\An\CAlg^{\wedge}$ of complete analytic $\EE_{\infty}$-rings admits small colimits, given as completions of colimits in $\An\CAlg$.  
\end{corollary}
\begin{proof}
    The adjunction exhibits $\An\CAlg^{\wedge}$ as a reflective subcategory in $\An\CAlg$ so the existence of small colimits in the latter -- as shown in \Cref{prop: colimits in AnCAlg} -- implies their existence in the former, as well as the formula for their computation \cite[\href{https://kerodon.net/tag/04JW}{Tag 04JW}]{kerodon}. 
\end{proof}
In the case of sifted colimits, no completion of the result is needed. 
\begin{proposition}\label{prop: sifted colimits in complete}
    The category $\An\CAlg^{\wedge}$ of complete analytic $\EE_{\infty}$-rings is closed under sifted colimits computed in $\An\CAlg$. 
\end{proposition}
\begin{proof}
    Let $\{(A_{i}^{\tri},\Mod_{A_{i}})\}_{i\in I}$ be a sifted diagram in $\An\CAlg^{\wedge}$ with colimit $(B^{\tri},\Mod_{B})$ in $\An\CAlg$. By completeness $A_{i}^{\tri}\simeq A_{i}[*]$, so $B^{\tri}\simeq\colim_{i\in I}A_{i}^{\tri}\simeq\colim_{i\in I}A_{i}[*]$, which is $B$-analytic by \Cref{lem: sifted colimits in AnCAlg} at $S=*$. Hence $B^{\tri}\in\Mod_{B}$, so $(B^{\tri},\Mod_{B})$ is complete. 
\end{proof}
\subsection{Topological invariance}\label{subsec: topological invariance}
The completion of \Cref{subsec: completion} replaced the underlying ring by its analytification while leaving the analytic modules fixed. We now show that the analytic structure is insensitive to the higher homotopy of the underlying condensed $\EE_{\infty}$-ring altogether: an analytic $\EE_{\infty}$-ring structure on $A^{\tri}$-modules is detected on the heart.
\begin{theorem}\label{thm: heart recognition for analytic rings}
    Let $A^{\tri}\in\Cond\CAlg^{\cn}$. The analytic $\EE_{\infty}$-ring structures on $A^{\tri}$-modules are in bijection with the full subcategories $\Asf$ of the Grothendieck Abelian category $\Mod_{\pi_{0}(A^{\tri})}^{\heartsuit}$ of $\pi_{0}(A^{\tri})$-modules in static condensed Abelian groups satisfying the following properties:
     \begin{enumerate}[label=(\roman*)]
        \item $\Asf$ is closed under small limits, colimits, and extensions in $\Mod_{\pi_{0}(A^{\tri})}^{\heartsuit}$. 
        \item $\Asf$ is closed under arbitrary higher derived products $\prod^{(n)}_{I}$. 
        \item For all $S\in\PFin_{\NN}$ and $N\in\Asf$, the homotopy groups 
        $$\pi_{-k}(\intHom_{A^{\tri}}(A^{\tri}[S],N))\simeq\pi_{-k}(\intHom_{\Cond\Sp}(\Sph[S],N))$$ 
        lie in $\Asf$ for all $k\geq0$. 
    \end{enumerate}
    The bijection sends $\Mod_{A}$ to $\Mod_{A}^{\heartsuit}$ and $\Asf$ to the full subcategory of $\Mod_{A^{\tri}}$ spanned by those objects $M$ such that $\pi_{k}(M)\in\Asf$ for all $k\in\ZZ$. 
\end{theorem}
\begin{proof}
    We follow \cite[Prop. 4.4.2 \& Thm. 4.4.1]{CamargoSolid}, noting that the identification in condition (iii) is \Cref{lem: modules inherit bistability} (2) using $t$-exactness of the forgetful functor to $\Cond\Sp$.  

    Given an analytic $\EE_{\infty}$-ring structure $\Mod_{A}$, set $\Asf:=\Mod_{A}^{\heartsuit}$. By \Cref{prop: properties of pi stably reflective localizations} (2) and (3), $\Mod_{A}$ carries the induced $t$-structure and membership in $\Mod_{A}$ is detected on homotopy groups. Regarding short exact sequences in $\Asf$ as (co)fiber sequences in $\Mod_{A}$, $\Asf$ is closed under extensions, kernels, and cokernels. Colimits in $\Asf$ are computed as $\pi_{0}$ of colimits computed in $\Mod_{A}$ hence in $\Asf$ by closure under homotopy. In the case of limits, we note that $\Asf$ is closed under higher derived products, which are the homotopy groups $\pi_{-k}(\prod_{i\in I}M_{i})$ of products formed in $\Mod_{A}$, and $\Asf$ is closed under small limits by writing these as kernels of derived products. This shows (i) and (ii). For (iii), the internal hom $\intHom_{A^{\tri}}(A^{\tri}[S],N)\in\Mod_{A}$ for $N\in\Asf\subseteq\Mod_{A}$ by \Cref{def: analytic E-infinity ring} (1) (ii) so its homotopy groups lie in $\Asf$ by homotopy detection.  

    Conversely, given $\Asf$ satisfying (i)-(iii), let $\Mod_{A}$ be spanned by those $M$ with all $\pi_{k}(M)\in\Asf$ for all $k\in\ZZ$. The long exact sequence of homotopy groups and exactness of direct sums give closure under fibers, cofibers, and small colimits. By the same argument using the biCartesian square in the proof of \Cref{thm: properties of solid spectra} (1), $\Mod_{A}$ is closed under extensions as well. To show closure of $\Mod_{A}$ under small limits, it suffices to show closure under small products. For a family $\{M_{i}\}_{i\in I}$ in $\Mod_{A}$ and $k\in\ZZ$, the third term of the fiber sequence $\prod_{i\in I}\tau_{\geq r}M_{i}\to\prod_{i\in I}M_{i}\to\prod_{i\in I}\tau_{\leq r-1}M_{i}$ is in $\Mod_{A^{\tri},\leq r-1}$ as coconnective objects are closed under products, so $\pi_{k}(\prod_{i\in I}M_{i})\simeq\pi_{k}(\prod_{i\in I}\tau_{\geq r}M_{i})$ for $r<k$. The truncations are uniformly bounded below, so left completeness and commutation of products with limits give $\prod_{i\in I}\tau_{\geq r}M_{i}\simeq\lim_{m}\prod_{i\in I}\tau_{[r,m]}M_{i}$, and induction on $m$ exhibits each $\prod_{i\in I}\tau_{[r,m]}M_{i}$ as an iterated extension of the $\prod_{i\in I}\pi_{n}(M_{i})[n]$, whose homotopy groups are higher derived products, so lie in $\Asf$. The Milnor sequence of \Cref{prop: Milnor sequence in t-plus bistable} then exhibits $\pi_{k}(\prod_{i\in I}\tau_{\geq r}M_{i})$ as an extension of $\lim_{m}\pi_{k}(\prod_{i\in I}\tau_{[r,m]}M_{i})$ by $\lim^{1}_{m}\pi_{k+1}(\prod_{i\in I}\tau_{[r,m]}M_{i})$. Both terms are kernels and cokernels of maps between countable products of objects of $\Asf$ by the construction of the Milnor sequence, so $\pi_{k}(\prod_{i\in I}M_{i})\in\Asf$ by (i). This shows \Cref{def: analytic E-infinity ring} (1) (i). 
    
    The inclusion admits a left adjoint $L^{A}$ by \cite[Thm. 1.1]{RagimovReflection} which is right $t$-exact by the same proof as \Cref{lem: ringwise colimits} (1) using adjunction. This establishes \Cref{def: analytic E-infinity ring} (1) (iii). For the final condition of \Cref{def: analytic E-infinity ring}, note that a general $M\in\Mod_{A^{\tri}}$ is a colimit of shifts of the $A^{\tri}[S]$ by \Cref{prop: stability of A-analytic modules} (2) which $\intHom_{A^{\tri}}(-,N)$ carries to a limit of shifts in $\Mod_{A}$ so it suffices to show $\intHom_{A^{\tri}}(A^{\tri}[S],N)\in\Mod_{A}$ for all $N\in\Mod_{A}$. Using the Postnikov tower and left completeness, we can further reduce to $N$ coconnective. In this case $\intHom_{A^{\tri}}(A^{\tri}[S],-)$ commutes with Whitehead towers of bounded above objects by \Cref{lem: internal hom and towers} (1) so using the (co)fiber sequences $\tau_{\geq k}N\to\tau_{\geq k-1}N\to\pi_{k-1}(N)[k-1]$ we use induction and closure of $\Mod_{A}$ under shifts and extensions to reduce to the case of $N\in\Asf$ which was assumed in (iii) above. This shows that $(A^{\tri},\Mod_{A})$ is an analytic $\EE_{\infty}$-ring structure. 

    The assignments are mutually inverse: the heart of the subcategory associated to $\Asf$ is $\Asf$, and an analytic $\EE_{\infty}$-ring structure is recovered from its heart by \Cref{prop: properties of pi stably reflective localizations} (3). 
\end{proof}
\begin{remark}\label{rmk: higher derived products}
    Closure under higher derived products is absent in \cite[Prop. 12.19]{Analytic} due to exactness of products in the non-light setting. In the light condensed setting, only countable products of connective objects remain connective by \Cref{thm: properties of condensed spectra} (2) (iv), and closure under higher derived products is necessary, cf. \cite[Rmk. 4.4.3]{CamargoSolid}. 
\end{remark}
In particular the data classifying analytic $\EE_{\infty}$-ring structures on $A^{\tri}$-modules and $\pi_{0}(A^{\tri})$-modules agree. 
\begin{proposition}\label{prop: functorial heart invariance}
    Let $A^{\tri}\in\Cond\CAlg^{\cn}$. The natural morphism $A^{\tri}\to\pi_{0}(A^{\tri})$ induces a functorial bijection of analytic $\EE_{\infty}$-ring structures on the categories of $A^{\tri}$- and $\pi_{0}(A^{\tri})$-modules in condensed spectra taking an analytic $\EE_{\infty}$-ring structure on $A^{\tri}$-modules to the induced analytic $\EE_{\infty}$-ring structure on $\pi_{0}(A^{\tri})$-modules. 
\end{proposition}
\begin{proof}
    The hearts of $\Mod_{A^{\tri}}$ and $\Mod_{\pi_{0}(A^{\tri})}$ coincide by \Cref{prop: relationship with heart} (1), so the sets of subcategories of the heart satisfying conditions (i)-(iii) of \Cref{thm: heart recognition for analytic rings} are identical. The bijection is realized by the induced analytic $\EE_{\infty}$-ring structure (\Cref{def: induced analytic ring}) as its analytic modules $\Mod_{\pi_{0}(A^{\tri})/A}$ -- those $N$ with $N|_{A^{\tri}}\in\Mod_{A}$ -- having heart $\Mod_{A}^{\heartsuit}$ by \Cref{prop: properties of pi stably reflective localizations} (3). Functoriality is provided by that of the induced structure. 
\end{proof}
Consequently the analytic structure depends only on the static truncation. 
\begin{corollary}\label{cor: induced for equivalence on pi0}
    Let $A^{\tri}\to B^{\tri}$ be a morphism in $\Cond\CAlg^{\cn}$ which induces an equivalence $\pi_{0}(A^{\tri})\to\pi_{0}(B^{\tri})$. Then there is a functorial bijection of analytic $\EE_{\infty}$-ring structures on $A^{\tri}$-modules and analytic $\EE_{\infty}$-ring structures on $B^{\tri}$-modules via the induced analytic $\EE_{\infty}$-ring structure. 
\end{corollary}
\begin{proof}
    Under the identification $\pi_{0}(A^{\tri})\simeq\pi_{0}(B^{\tri})$, any analytic $\EE_{\infty}$-ring structure $(A^{\tri},\Mod_{A})$ on $A^{\tri}$-modules gives a unique analytic $\EE_{\infty}$-ring structure on $\pi_{0}(B^{\tri})$-modules which lifts to a unique analytic $\EE_{\infty}$-ring structure on $B^{\tri}$-modules by \Cref{prop: functorial heart invariance}. 
\end{proof}
Analytic $\EE_{\infty}$-ring structures are moreover invariant under nilpotent thickenings. 
\begin{corollary}\label{cor: analytic structures invariant under nilpotent thickenings}
    Let $A^{\tri}\to B^{\tri}$ be a morphism in $\Cond\CAlg^{\cn}$ such that $\pi_{0}(A^{\tri})\to\pi_{0}(B^{\tri})$ is surjective with nilpotent kernel. Then there is a functorial bijection of analytic $\EE_{\infty}$-ring structures on $A^{\tri}$-modules and analytic $\EE_{\infty}$-ring structures on $B^{\tri}$-modules via the induced analytic $\EE_{\infty}$-ring structure. 
\end{corollary}
\begin{proof}
    By \Cref{cor: induced for equivalence on pi0} applied to the truncations, we may assume $A^{\tri}$ and $B^{\tri}$ static, where the claim is \cite[Prop. 4.4.5]{CamargoSolid}. 
\end{proof}
\subsection{Mereology of \texorpdfstring{$A$}{A}-analytic modules}\label{subsec: mereology of analytic modules}
As in static algebra and algebraic geometry, it is often fruitful to single out subcategories of modules that are finitary in a suitable sense. We consider several notions of finiteness for the modules over a complete analytic $\EE_{\infty}$-ring $(A^{\tri},\Mod_{A})$ and relate them to finiteness in the category $\Mod_{A^{\tri}(*)}(\Sp)$ of $A^{\tri}(*)$-modules in spectra. 

Recall that $\Mod_{A}$ is $\omega_{1}$-presentably symmetric monoidal with compact unit by \Cref{prop: A-analytic tensor product}, so the theory of \Cref{app: nuclearity} applies for $\kappa=\omega_{1}$.

We first specialize dualizability to this setting.
\begin{definition}\label{def: dualizability}
    Let $(A^{\tri},\Mod_{A})$ be a complete analytic $\EE_{\infty}$-ring.
    \begin{enumerate}
        \item Let $M\in\Mod_{A}$. The \emph{pre-dual} of $M$ is the object $M^{\vee}:=\intHom_{A}(M, A^{\tri})$. Let $\ev_{M}:M\otimes_{A}M^{\vee}\to A^{\tri}$ be the counit morphism of $M\otimes_{A}(-)\dashv\intHom_{A}(M, -)$ at $A^{\tri}$.  
        \item An object $M\in\Mod_{A}$ is \emph{dualizable} if there exists a morphism $\coev_{M}:A^{\tri}\to M^{\vee}\otimes_{A}M$ such that the composites 
        $$M\xrightarrow{\id_{M}\otimes_{A}\coev_{M}}M\otimes_{A}M^{\vee}\otimes_{A}M\xrightarrow{\ev_{M}\otimes_{A}\id_{M}}M$$
        $$M^{\vee}\xrightarrow{\coev_{M}\otimes_{A}\id_{M^{\vee}}}M^{\vee}\otimes_{A}M\otimes_{A}M^{\vee}\xrightarrow{\id_{M^{\vee}}\otimes_{A}\ev_{M}}M^{\vee}$$
        are identities. In this case, we say the pre-dual $M^{\vee}$ is the \emph{dual} of $M$. 
    \end{enumerate}
\end{definition}
The nuclear objects are built from the compact ones along trace class maps.
\begin{definition}\label{def: nuclearity and trace class morphisms}
    Let $(A^{\tri},\Mod_{A})$ be a complete analytic $\EE_{\infty}$-ring. 
    \begin{enumerate}
        \item A morphism $f:M\to N$ in $\Mod_{A}$ is \emph{trace class} if the morphism $A^{\tri}\to\intHom_{A}(M, N)$ classifying $f$ admits a factorization as $A^{\tri}\to M^{\vee}\otimes_{A}N\to\intHom_{A}(M,N)$. 
        \item An object of $\Mod_{A}$ is \emph{basic nuclear} if it is the colimit of a $\QQ_{\geq0}$-indexed diagram of $\omega_{1}$-compact objects along trace class morphisms. 
        \item An object of $\Mod_{A}$ is \emph{nuclear} if it lies in the subcategory of $\Mod_{A}$ generated under small colimits by the basic nuclear objects. 
    \end{enumerate}
    Denote the full subcategory of nuclear objects by $\Nuc_{A}\subseteq\Mod_{A}$. 
\end{definition}
\begin{remark}\label{rmk: nuclear objects}
    The generators $A[S]$ of $\Mod_{A}$ are only $\omega_{1}$-compact by \Cref{prop: stability of A-analytic modules} (3), so $\Mod_{A}$ need not be $\omega$-presentable and the results of \cite[Lecture VIII]{Complex} do not apply. However, $\Mod_{A}$ is $\omega_{1}$-presentably symmetric monoidal with compact unit by \Cref{prop: A-analytic tensor product} (2) bringing $\Mod_{A}$ within the scope of \Cref{app: nuclearity}, and allowing us to produce a satisfactory theory of nuclear objects.
\end{remark}
The collection of nuclear modules satisfies the following permanence properties. 
\begin{proposition}[{\cite[Prop. 2.3.22]{MannRigid6FF}}]\label{prop: permanence properties of nuclear objects}
    Let $(A^{\tri},\Mod_{A})\to(B^{\tri},\Mod_{B})$ be a morphism in $\An\CAlg^{\wedge}$.
    \begin{enumerate}
        \item $\Nuc_{A}$ is closed under small colimits computed in $\Mod_{A}$. 
        \item The scalar extension functor $B\otimes_{A}-:\Mod_{A}\to\Mod_{B}$ preserves nuclear objects: if $M\in\Nuc_{A}$ then $B\otimes_{A}M\in\Nuc_{B}\subseteq\Mod_{B}$. 
    \end{enumerate}
\end{proposition}
\begin{proof}
    (1) is \Cref{lem: comparison functor from rigidification} (3) and \Cref{prop: nuclearity and rigidity}, while (2) is \Cref{prop: preservation of nuclearity} using that scalar extension is a symmetric monoidal functor which preserves $\omega_{1}$-compact generators.  
\end{proof}
Discrete modules provide a convenient source of nuclear modules.
\begin{corollary}\label{cor: discrete implies nuclear}
    Let $(A^{\tri},\Mod_{A})$ be a complete analytic $\EE_{\infty}$-ring. Any discrete object of $\Mod_{A}$ is nuclear.
\end{corollary}
\begin{proof}
    By \Cref{prop: discrete modules}, the discrete objects of $\Mod_{A}$ are generated under colimits and shifts by $A^{\tri}$. But $A^{\tri}$ is dualizable, hence nuclear, and $\Nuc_{A}$ is closed under colimits and shifts in $\Mod_{A}$, thus in particular contains the discrete objects.  
\end{proof}
In a closed symmetric monoidal category with compact unit, any dualizable object is necessarily compact. However, compactness is often not sufficient. The precise criterion is nuclearity together with compactness.
\begin{proposition}[{\cite[Prop. 9.3]{Complex}}]\label{prop: dualizable iff nuclear and compact}
    Let $(A^{\tri},\Mod_{A})$ be a complete analytic $\EE_{\infty}$-ring. An object $M\in\Mod_{A}$ is dualizable if and only if it is nuclear and compact. 
\end{proposition}
\begin{proof}
    This is \Cref{prop: dualizable iff nuclear and compact category}. 
\end{proof}
\begin{remark}\label{rmk: beyond the omega-presentable case}
    \cite[Prop. 9.3]{Complex} imposes the condition that the ambient category is $\omega$-presentable with compact unit. The statement in \Cref{app: nuclearity} shows this still holds when the assumption of $\omega$-presentability on the category is removed but the unit is still compact. 
\end{remark}
As the preceding discussion suggests, the discrete compact objects of $\Mod_{A}$ are dualizable, and we introduce the following name for when the embedding of \Cref{prop: discrete modules} restricts to an equivalence on dualizable objects. 
\begin{definition}[{\cite[Def. 3.3.1]{dRFF}}]\label{def: Fredholm rings}
    A complete analytic $\EE_{\infty}$-ring $(A^{\tri},\Mod_{A})$ is \emph{Fredholm} if all dualizable objects of $\Mod_{A}$ are discrete. 
\end{definition}
\begin{remark}\label{rmk: Fredholm rings}
    The notion of a Fredholm compactly generated closed symmetric monoidal category with compact unit first appeared as \cite[Def. 9.7]{Complex}, of which \Cref{def: Fredholm rings} can be seen as special case under additional assumptions on the category of analytic modules. 
\end{remark}
\section{Morphisms of analytic \texorpdfstring{$\EE_{\infty}$}{E-infinity}-rings}\label{sec: morphisms of analytic rings}
Having built the category $\An\CAlg^{\wedge}$ of complete analytic $\EE_{\infty}$-rings in \Cref{sec: analytic rings}, we isolate the morphisms suited to the six operations. The governing requirement is compatibility with base change: scalar extension need not be tensoring with an algebra object, so restriction of scalars commutes with base change only for the steady morphisms. 

Open immersions and analytically proper morphisms then furnish a suitable decomposition of a geometric setup on the opposite category of complete analytic $\EE_{\infty}$-rings $\An\Aff$, yielding the six-functor formalism. We conclude with descent of analytic modules, following L. Mann's theory of steady endofunctors \cite[\S 2.5-2.6]{MannRigid6FF}.
\subsection{Constructing morphisms of analytic \texorpdfstring{$\EE_{\infty}$}{E-infinity}-rings}\label{subsec: constructing morphisms}
We begin with the $\EE_{\infty}$-analogue of a construction principle of Scholze, by which a morphism of the underlying connective condensed $\EE_{\infty}$-rings promotes to a morphism of analytic $\EE_{\infty}$-rings once a compatibility on $\pi_{0}$ is imposed. The reduction to the static case is immediate.
\begin{proposition}[{\cite[Prop. 4.5.1]{CamargoSolid}}]\label{prop: constructing morphisms of analytic rings}
    Let $(A^{\tri},\Mod_{A}),(B^{\tri},\Mod_{B})\in\An\CAlg^{\wedge}$ and let $f: A^{\tri}\to B^{\tri}$ be a morphism in $\Cond\CAlg^{\cn}$. Suppose for all $S\in\PFin_{\NN}$ there are:
    \begin{enumerate}[label=(\roman*)]
        \item Morphisms $\gamma_{S}:A[S]\to B[S]|_{A^{\tri}}$ natural in $S$ such that $S\to \Omega^{\infty}B[S]$ factors as $S\to\Omega^{\infty}A[S]\to\Omega^{\infty}B[S]$ on underlying condensed animae. 
        \item For every morphism of condensed animae $g:S\to \Omega^{\infty}A^{\tri}$ inducing morphisms $\alpha:A[S]\to A^{\tri},\beta:B[S]\to B^{\tri}$ adjoint to $g,f\circ g$ the diagram 
        $$% https://q.uiver.app/#q=WzAsNCxbMCwwLCJBW1NdIl0sWzAsMSwiQltTXXxfe0Fee1xcdHJpfX0iXSxbMiwwLCJBXntcXHRyaX0iXSxbMiwxLCJCXntcXHRyaX18X3tBXntcXHRyaX19Il0sWzAsMiwiXFxhbHBoYSJdLFsyLDMsImYiXSxbMCwxLCJcXGdhbW1hX3tTfSIsMl0sWzEsMywiXFxiZXRhIiwyXV0=
        \begin{tikzcd}
            {A[S]} && {A^{\tri}} \\
            {B[S]|_{A^{\tri}}} && {B^{\tri}|_{A^{\tri}}}
            \arrow["\alpha", from=1-1, to=1-3]
            \arrow["{\gamma_{S}}"', from=1-1, to=2-1]
            \arrow["f", from=1-3, to=2-3]
            \arrow["\beta"', from=2-1, to=2-3]
        \end{tikzcd}$$
        in $\Mod_{A^{\tri}}$ is commutative on $\pi_{0}$.
    \end{enumerate}
    Then $f$ promotes to a morphism $(A^{\tri},\Mod_{A})\to(B^{\tri},\Mod_{B})$ in $\An\CAlg^{\wedge}$.
\end{proposition}
\begin{proof}
    The morphism $f$ promotes to a morphism of analytic $\EE_{\infty}$-rings precisely when restriction of scalars $(-)|_{A^{\tri}}:\Mod_{B}\to\Mod_{A^{\tri}}$ lands in $\Mod_{A}$. The objects $B[S]$ for $S\in\PFin_{\NN}$ generate $\Mod_{B}$ under colimits and shifts by \Cref{prop: stability of A-analytic modules} (3), restriction to $\Mod_{A^{\tri}}$ preserves colimits and shifts, and $\Mod_{A}$ is closed under them, so it suffices to show $B[S]|_{A^{\tri}}\in\Mod_{A}$ for all $S\in\PFin_{\NN}$.

    By \Cref{prop: t-structure on A-analytic modules} (vi) this is checked on homotopy groups: $B[S]|_{A^{\tri}}\in\Mod_{A}$ if and only if $\pi_{k}(B[S]|_{A^{\tri}})\in\Mod_{A}^{\heartsuit}$ for all $k\in\ZZ$. \Cref{prop: functorial heart invariance} identifies $\Mod_{A}^{\heartsuit}$ with the heart of the induced static analytic ring on $\pi_{0}(A^{\tri})$-modules, reducing the claim to the underlying morphism $\pi_{0}(A^{\tri})\to\pi_{0}(B^{\tri})$ of static rings, where it is the criterion of \cite[Prop. 7.12]{Condensed}, the hypotheses of which are precisely the conditions of this proposition.
\end{proof}
\subsection{Steady morphisms}\label{subsec: steady morphisms}
Steadiness concerns the interaction of restriction of scalars with base change, so we first record the explicit description of pushouts in $\An\CAlg^{\wedge}$, using the colimit construction of \Cref{subsec: functoriality and colimits,subsec: completion}. Throughout, $B\otimes_{A}-:\Mod_{A}\to\Mod_{B}$ denotes the scalar extension along a morphism $(A^{\tri},\Mod_{A})\to(B^{\tri},\Mod_{B})$.
\begin{corollary}\label{cor: pushout description}
    Let $(B^{\tri},\Mod_{B})\leftarrow(A^{\tri},\Mod_{A})\to(C^{\tri},\Mod_{C})$ be a diagram in $\An\CAlg^{\wedge}$ with pushout $(D^{\tri},\Mod_{D})$ in $\An\CAlg^{\wedge}$ as in 
    \begin{equation}\label{eqn: pushout recognition}
        % https://q.uiver.app/#q=WzAsNCxbMCwwLCIoQV57XFxyaGR9LFxcTW9kX3tBfSkiXSxbMCwxLCIoQ157XFxyaGR9LFxcTW9kX3tDfSkiXSxbMiwwLCIoQl57XFxyaGR9LFxcTW9kX3tCfSkiXSxbMiwxLCIoRF57XFxyaGR9LFxcTW9kX3tEfSkiXSxbMCwyXSxbMiwzXSxbMCwxXSxbMSwzXV0=
        \begin{tikzcd}
            {(A^{\tri},\Mod_{A})} && {(B^{\tri},\Mod_{B})} \\
            {(C^{\tri},\Mod_{C})} && {(D^{\tri},\Mod_{D}).}
            \arrow[from=1-1, to=1-3]
            \arrow[from=1-1, to=2-1]
            \arrow[from=1-3, to=2-3]
            \arrow[from=2-1, to=2-3]
        \end{tikzcd}
    \end{equation}
    Write $\widetilde{D}^{\tri}:=B^{\tri}\otimes_{A^{\tri}}C^{\tri}$ for the pushout of the underlying rings in $\Cond\CAlg^{\cn}$. Then 
    $$\Mod_{D}=\Mod_{\widetilde{D}^{\tri}/B}\cap\Mod_{\widetilde{D}^{\tri}/C}$$
    as subcategories of $\Mod_{\widetilde{D}^{\tri}}$, and $D^{\tri}$ is the analytification of $\widetilde{D}^{\tri}$ with respect to this analytic structure.
\end{corollary}
\begin{proof}
    This follows from \Cref{prop: colimits in AnCAlg} and \Cref{cor: complete AnCAlg has colimits}. Explicitly, the pushout in $\An\CAlg$ is $(\widetilde{D}^{\tri},\Mod_{D})$ which has underlying connective condensed $\EE_{\infty}$-algebra $\widetilde{D}^{\tri}\simeq B^{\tri}\otimes_{A^{\tri}}C^{\tri}$ and module category $\Mod_{D}=\Mod_{\widetilde{D}^{\tri}/B}\cap\Mod_{\widetilde{D}^{\tri}/C}$. The colimit in $\An\CAlg^{\wedge}$ is its completion, which is $(D^{\tri},\Mod_{D})$. 
\end{proof}
We now introduce steady morphisms.
\begin{definition}\label{def: steady morphism}
    A morphism $(A^{\tri},\Mod_{A})\to(B^{\tri},\Mod_{B})$ in $\An\CAlg^{\wedge}$ is \emph{steady} if for all morphisms $(A^{\tri},\Mod_{A})\to(C^{\tri},\Mod_{C})$ with pushout as in (\ref{eqn: pushout recognition}) the morphism 
    $$B\otimes_{A}(M|_{A^{\tri}})\longrightarrow(D\otimes_{C}M)|_{B^{\tri}}$$
    is an equivalence in $\Mod_{B}$ for all $M\in\Mod_{C}$. 
\end{definition}
\begin{remark}\label{rmk: definition of steady maps}
    The notion of a steady morphism was introduced in \cite[Def. 12.13]{Analytic} and the conventions we have adopted in \Cref{def: steady morphism} are in line with those in the literature \cite[Def. 2.3.16 \& Rmk. 2.3.17]{MannRigid6FF} (cf. \cite[Def. 1.17 \& Rmk. 1.18]{MikamiFFDescent}). We refer the reader to \cite[p. 102]{Analytic} for an example of a non-steady morphism. 
\end{remark}
The collection of steady morphisms enjoys the expected permanence properties.
\begin{proposition}[{\cite[Prop. 2.3.19]{MannRigid6FF}}]\label{prop: permanence properties of steady morphisms}
    The collection of steady morphisms satisfies the following permanence properties: 
    \begin{enumerate}
        \item (Composition) If $(A^{\tri},\Mod_{A})\to (B^{\tri},\Mod_{B}), (B^{\tri},\Mod_{B})\to (C^{\tri},\Mod_{C})$ are steady morphisms in $\An\CAlg^{\wedge}$ then so is the composite $(A^{\tri},\Mod_{A})\to (B^{\tri},\Mod_{B})\to (C^{\tri},\Mod_{C})$. 
        \item (Base Change) If $(A^{\tri},\Mod_{A})\to (B^{\tri},\Mod_{B})$ is a steady morphism and $(A^{\tri},\Mod_{A})\to (C^{\tri},\Mod_{C})$ another morphism in $\An\CAlg^{\wedge}$ the induced morphism from $(C^{\tri},\Mod_{C})$ to the pushout is steady. 
        \item (Left Cancellation) If a composite $(A^{\tri},\Mod_{A})\to (B^{\tri},\Mod_{B})\to (C^{\tri},\Mod_{C})$ is steady and $(A^{\tri},\Mod_{A})\to (B^{\tri},\Mod_{B})$ is steady then so is $(B^{\tri},\Mod_{B})\to (C^{\tri},\Mod_{C})$. 
        \item (Colimits) If $\{(A_{i}^{\tri},\Mod_{A_{i}})\to (B_{i}^{\tri},\Mod_{B_{i}})\}_{i\in I}$ is a small diagram of steady morphisms in $\Fun([1],\An\CAlg^{\wedge})$ then so is the colimit. 
    \end{enumerate}
\end{proposition}
\begin{proof}
    The permanence properties are \cite[Prop. 2.3.19]{MannRigid6FF} via the characterization of \cite[Prop. 2.3.18]{MannRigid6FF}. Neither of the proofs uses the animated structure, so they carry over to the light condensed $\EE_{\infty}$-setting.
\end{proof}
\subsection{Open immersions}\label{subsec: open immersions}
Open immersions are the localizations of categories of analytic modules with smashing kernel -- equivalently, those whose scalar extension admits a fully faithful left adjoint satisfying the projection formula \cite[Prop. 6.5]{Complex}. We record their permanence and discuss a criterion detecting an open immersion on the induced morphism of static analytic rings under additional hypotheses.
\begin{definition}[{\cite[Lecture VI]{Complex}}]\label{def: open immersion}
    A morphism $j:(A^{\tri},\Mod_{A})\to(B^{\tri},\Mod_{B})$ in $\An\CAlg^{\wedge}$ whose scalar extension $j^{*}:\Mod_{A}\to\Mod_{B}$ has fully faithful right adjoint $j_{*}$ is an \emph{open immersion} if it satisfies either of the following equivalent conditions:
    \begin{enumerate}[label=(\alph*)]
        \item $j^{*}$ admits a fully faithful left adjoint $j_{\sharp}$ that satisfies the projection formula: for all $M\in\Mod_{A},N\in\Mod_{B}$ the natural morphism $j_{\sharp}(j^{*}M\otimes_{B}N)\to M\otimes_{A}j_{\sharp}N$ is an equivalence in $\Mod_{A}$. 
        \item The kernel $\ker(j^{*}):=\{K\in\Mod_{A}:j^{*}K\simeq0\}\subseteq\Mod_{A}$ is equivalent to $\Mod_{Q}(\Mod_{A})$ for some idempotent algebra $Q\in\Idem_{A}$.
    \end{enumerate}
\end{definition}
\begin{remark}\label{rmk: open immersion as categorical open immersion}
    \Cref{def: open immersion} is the special case of \Cref{def: categorical open immersion} for the symmetric monoidal functor $B\otimes_{A}-$ in $\CAlg(\PrL_{\St})$.
\end{remark}
Using that the kernel of an open immersion is equivalent to the category of modules over an idempotent algebra, we show that open immersions are preserved by base change in $\An\CAlg^{\wedge}$. 
\begin{proposition}\label{prop: pushouts along open immersions}
    Let $j:(A^{\tri},\Mod_{A})\to(B^{\tri},\Mod_{B})$ be an open immersion in $\An\CAlg^{\wedge}$ with $\ker(j^{*})=\Mod_{Q}(\Mod_{A})$, let $g:(A^{\tri},\Mod_{A})\to(C^{\tri},\Mod_{C})$ be a morphism in $\An\CAlg^{\wedge}$ with pushout $(D^{\tri},\Mod_{D})$ as in 
    $$% https://q.uiver.app/#q=WzAsNCxbMCwwLCIoQV57XFx0cml9LFxcTW9kX3tBfSkiXSxbMiwwLCIoQl57XFx0cml9LFxcTW9kX3tCfSkiXSxbMCwxLCIoQ157XFx0cml9LFxcTW9kX3tDfSkiXSxbMiwxLCIoRF57XFx0cml9LFxcTW9kX3tEfSkiXSxbMCwxLCJqIl0sWzEsMywiXFx3aWRldGlsZGV7Z30iXSxbMCwyLCJnIiwyXSxbMiwzLCJcXHdpZGV0aWxkZXtqfSIsMl1d
    \begin{tikzcd}
        {(A^{\tri},\Mod_{A})} && {(B^{\tri},\Mod_{B})} \\
        {(C^{\tri},\Mod_{C})} && {(D^{\tri},\Mod_{D})}
        \arrow["j", from=1-1, to=1-3]
        \arrow["g"', from=1-1, to=2-1]
        \arrow["{\widetilde{g}}", from=1-3, to=2-3]
        \arrow["{\widetilde{j}}"', from=2-1, to=2-3]
    \end{tikzcd}$$
    Then $\widetilde{j}$ is an open immersion with $\ker(\widetilde{j}^{*})\simeq\Mod_{C\otimes_{A}Q}(\Mod_{C})$. 
\end{proposition}
\begin{proof}
    Since $C\otimes_{A}-$ is symmetric monoidal, $\widetilde{Q}:=g^{*}Q=C\otimes_{A}Q$ is an idempotent algebra in $\Mod_{C}$ and commutativity of the pushout diagram and stability imply that $\widetilde{g}^{*}j^{*}Q\simeq\widetilde{j}^{*}g^{*}Q\simeq0$ in $\Mod_{D}$, where $\widetilde{g}:(B^{\tri},\Mod_{B})\to(D^{\tri},\Mod_{D})$ is the induced morphism. Write $I:=\fib_{A}(A^{\tri}\to Q)$ and $\widetilde{I}:=\fib_{C}(C^{\tri}\to\widetilde{Q})\simeq g^{*}I$.

    The right adjoint $j_{*}$ of $j^{*}$ is scalar restriction, so it preserves small limits and colimits. We claim the essential image of $j_{*}$ is $\Escr:=\{N\in\Mod_{A}:\intHom_{A}(Q,N)\simeq0\}\subseteq\Mod_{A}$. The natural map $N\to\intHom_{A}(I,N)\simeq j_{*}j^{*}N$ of \cite[Rmk. 5.1.5 (2)]{CamargoSolid} has fiber $\intHom_{A}(Q,N)$, so for $N\in\Escr$ it is an equivalence and $N$ lies in the essential image. Conversely if $N$ is in the essential image, the unit map $N\to j_{*}j^{*}N\simeq\intHom_{A}(I,N)$ is an equivalence, so the fiber $\intHom_{A}(Q,N)$ vanishes and $N\in\Escr$. The computation $\intHom_{A^{\tri}}(Q,\intHom_{A^{\tri}}(M,N))\simeq\intHom_{A^{\tri}}(M,\intHom_{A}(Q,N))\simeq0$ for $N\in\Escr$ and $M\in\Mod_{A^{\tri}}$ shows $\Escr$ is moreover closed under internal homs from arbitrary $A^{\tri}$-modules, and the localization $L^{\Escr}$ is right $t$-exact, being the composite $j_{*}j^{*}\circ L^{A}$ of right $t$-exact functors. This defines an analytic $\EE_{\infty}$-ring structure $(A^{\tri},\Escr)$ on $A^{\tri}$ with completion $(B^{\tri},\Mod_{B})$. \Cref{cor: complete AnCAlg has colimits} shows $\An\CAlg^{\wedge}$ is a reflective subcategory of $\An\CAlg$ so the pushout $(D^{\tri},\Mod_{D})$ above is also the completion of the pushout of $(C^{\tri},\Mod_{C})\leftarrow(A^{\tri},\Mod_{A})\to(A^{\tri},\Escr)$ taken in $\An\CAlg$. \Cref{cor: pushout description} then identifies $\Mod_{D}$ with the full subcategory of $M\in\Mod_{C}$ with $g_{*}M\in \Escr$. For every $M\in\Mod_{C}$ the adjunction gives $g_{*}\intHom_{C}(\widetilde{Q},M)\simeq\intHom_{A}(Q,g_{*}M)$, so by the description of the essential image and conservativity of scalar restriction \cite[Cor. 4.2.3.2]{HA},
    $$\Mod_{D}=\{M\in\Mod_{C}:\intHom_{C}(\widetilde{Q},M)\simeq0\}\subseteq\Mod_{C},$$
    showing in particular that the right adjoint $\widetilde{j}_{*}$ is fully faithful.


    By idempotence of $\widetilde{Q}$, an object $K\in\Mod_{C}$ is a $\widetilde{Q}$-module if and only if the natural map $K\to\widetilde{Q}\otimes_{C}K$ is an equivalence. By symmetric monoidality of $\widetilde{j}^{*}$, we have 
    $$\widetilde{j}^{*}(K)\simeq\widetilde{j}^{*}(\widetilde{Q})\otimes_{D}\widetilde{j}^{*}(K)\simeq 0\otimes_{D}\widetilde{j}^{*}(K)\simeq0$$
    showing $\Mod_{\widetilde{Q}}(\Mod_{C})\subseteq\ker(\widetilde{j}^{*})$. Conversely, suppose $\widetilde{j}^{*}(K)\simeq0$ so $\Hom_{C}(K,N)\simeq0$ for all $N\in\Mod_{D}\subseteq\Mod_{C}$. For $P\in\Mod_{C}$, $Q\otimes_{A}I\simeq0$ implies $\widetilde{Q}\otimes_{C}\widetilde{I}\simeq0$ and we have equivalences
    $$\intHom_{C}\left(\widetilde{Q},\intHom_{C}\left(\widetilde{I},P\right)\right)\simeq\intHom_{C}\left(\widetilde{Q}\otimes_{C}\widetilde{I},P\right)\simeq0$$
    and $\intHom_{C}(\widetilde{I},P)\in\Mod_{D}\subseteq\Mod_{C}$ so adjunction gives 
    $$\Hom_{C}\left(K\otimes_{C}\widetilde{I},P\right)\simeq\Hom_{C}\left(K,\intHom_{C}\left(\widetilde{I},P\right)\right)\simeq0.$$
    Taking $P=K\otimes_{C}\widetilde{I}$ we get that the identity $\id_{K\otimes_{C}\widetilde{I}}$ is zero and $K\otimes_{C}\widetilde{I}\simeq0$. Thus tensoring $K$ with the (co)fiber sequence $\widetilde{I}\to C^{\tri}\to\widetilde{Q}$ we conclude that $K\to\widetilde{Q}\otimes_{C}K$ is an equivalence, and $\ker(\widetilde{j}^{*})\subseteq\Mod_{\widetilde{Q}}(\Mod_{C})$. 
\end{proof}
The permanence of open immersions under base change follows from the result above, while permanence under composition and left cancellation are formal from the categorical properties. 
\begin{proposition}\label{prop: permanence of open immersions}
    The collection of open immersions satisfies the following permanence properties: 
    \begin{enumerate}
        \item (Composition) If $(A^{\tri},\Mod_{A})\to (B^{\tri},\Mod_{B}), (B^{\tri},\Mod_{B})\to (C^{\tri},\Mod_{C})$ are open immersions in $\An\CAlg^{\wedge}$ then so is the composite $(A^{\tri},\Mod_{A})\to (B^{\tri},\Mod_{B})\to (C^{\tri},\Mod_{C})$. 
        \item (Base Change) If $(A^{\tri},\Mod_{A})\to (B^{\tri},\Mod_{B})$ is an open immersion and $(A^{\tri},\Mod_{A})\to (C^{\tri},\Mod_{C})$ another morphism in $\An\CAlg^{\wedge}$ the induced morphism from $(C^{\tri},\Mod_{C})$ to the pushout is an open immersion. 
        \item (Left Cancellation) If a composite $(A^{\tri},\Mod_{A})\to (B^{\tri},\Mod_{B})\to (C^{\tri},\Mod_{C})$ is an open immersion and $(A^{\tri},\Mod_{A})\to (B^{\tri},\Mod_{B})$ is an open immersion then so is $(B^{\tri},\Mod_{B})\to (C^{\tri},\Mod_{C})$. 
    \end{enumerate}
\end{proposition}
\begin{proof}
    Permanence under base change is established in \Cref{prop: pushouts along open immersions}. The forgetful functor from $\Cond\Sp$-linear categories preserves composition, so permanence under composition and left cancellation can be checked on underlying stable presentably symmetric monoidal categories, and thus follows from \cite[Prop. 6.5 \& Cor. 6.6]{Complex}. 
\end{proof}
Beyond permanence, the open immersions satisfy the base change that the six-functor formalism requires: the exceptional left adjoint $j_{\sharp}$ commutes with scalar extension along an arbitrary morphism.
\begin{proposition}\label{prop: base change for open immersions}
    Let $j:(A^{\tri},\Mod_{A})\to(B^{\tri},\Mod_{B})$ be an open immersion in $\An\CAlg^{\wedge}$. For every $g:(A^{\tri},\Mod_{A})\to(C^{\tri},\Mod_{C})$ with pushout $(D^{\tri},\Mod_{D})$ in $\An\CAlg^{\wedge}$, the base change morphism $\widetilde{j}_{\sharp}\,\widetilde{g}^{*}\to g^{*}j_{\sharp}$ of functors $\Mod_{B}\to\Mod_{C}$ associated to
    \begin{equation}\label{eqn: open immersion base change}
        % https://q.uiver.app/#q=WzAsNCxbMCwwLCJcXE1vZF97QX0iXSxbMiwwLCJcXE1vZF97Qn0iXSxbMCwxLCJcXE1vZF97Q30iXSxbMiwxLCJcXE1vZF97RH0iXSxbMSwwLCJqX3tcXHNoYXJwfVxcZGFzaHYgal57Kn0iLDIseyJvZmZzZXQiOjF9XSxbMCwxLCIiLDAseyJvZmZzZXQiOjF9XSxbMywyLCIiLDIseyJvZmZzZXQiOjF9XSxbMiwzLCJcXHdpZGV0aWxkZXtqfV97XFxzaGFycH1cXGRhc2h2XFx3aWRldGlsZGV7an1eeyp9IiwyLHsib2Zmc2V0IjoxfV0sWzAsMiwiZ157Kn0iLDJdLFsxLDMsIlxcd2lkZXRpbGRle2d9XnsqfSJdXQ==
        \begin{tikzcd}
            {\Mod_{A}} && {\Mod_{B}} \\
            {\Mod_{C}} && {\Mod_{D}}
            \arrow[shift right, from=1-1, to=1-3]
            \arrow["{g^{*}}"', from=1-1, to=2-1]
            \arrow["{j_{\sharp}\dashv j^{*}}"', shift right, from=1-3, to=1-1]
            \arrow["{\widetilde{g}^{*}}", from=1-3, to=2-3]
            \arrow["{\widetilde{j}_{\sharp}\dashv\widetilde{j}^{*}}"', shift right, from=2-1, to=2-3]
            \arrow[shift right, from=2-3, to=2-1]
        \end{tikzcd}
    \end{equation}
    is an equivalence.
\end{proposition}
\begin{proof}
    From \Cref{prop: pushouts along open immersions}, $\widetilde{j}$ is an open immersion and there exists an idempotent algebra $Q\in\Idem_{A}$ such that $\ker(j^{*})\simeq\Mod_{Q}(\Mod_{A})$ and $\ker(\widetilde{j}^{*})\simeq\Mod_{\widetilde{Q}}(\Mod_{C})$, where $\widetilde{Q}:=g^{*}Q$. Write $I:=\fib_{A}(A^{\tri}\to Q)$ and $\widetilde{I}:=\fib_{C}(C^{\tri}\to\widetilde{Q})\simeq g^{*}I$.

    Since $j^{*}:\Mod_{A}\to\Mod_{B}$ is essentially surjective, it suffices to show the morphism $\widetilde{j}_{\sharp}\widetilde{g}^{*}j^{*}\to g^{*}j_{\sharp}j^{*}$ of functors $\Mod_{A}\to\Mod_{C}$ is an equivalence. Using $\widetilde{g}^{*}j^{*}\simeq\widetilde{j}^{*}g^{*}$, this is the morphism $\widetilde{j}_{\sharp}\widetilde{j}^{*}g^{*}\to g^{*}j_{\sharp}j^{*}$, which the natural equivalences $j_{\sharp}j^{*}\simeq I\otimes_{A}-$ and $\widetilde{j}_{\sharp}\widetilde{j}^{*}\simeq\widetilde{I}\otimes_{C}-$ of \cite[Rmk. 5.1.5 (2)]{CamargoSolid} carry to the canonical equivalence $\widetilde{I}\otimes_{C}g^{*}(-)\simeq g^{*}(I\otimes_{A}-)$ supplied by symmetric monoidality of $g^{*}$ and $\widetilde{I}\simeq g^{*}I$, giving the claim.
\end{proof}
It follows that open immersions are steady morphisms. 
\begin{corollary}\label{cor: open immersions are steady}
    If $j:(A^{\tri},\Mod_{A})\to(B^{\tri},\Mod_{B})$ is an open immersion, then $j$ is steady. 
\end{corollary}
\begin{proof}
    Form the pushout as in (\ref{eqn: open immersion base change}). The morphism $j$ is steady if and only if $j^{*}g_{*}\to\widetilde{g}_{*}\widetilde{j}^{*}$ is an equivalence, which follows by passing to the adjoint morphism of the equivalence $\widetilde{j}_{\sharp}\widetilde{g}^{*}\xrightarrow{\sim}g^{*}j_{\sharp}$. 
\end{proof}
It remains to recognize open immersions in practice. We prove a restricted form of an expectation of Peter Scholze, relayed to us by Lucas Mann, that whether a morphism is an open immersion should be detected on the heart. The recognition criterion of \Cref{subsec: detecting smashing kernel} reduces the question to the underlying static analytic rings once the analytic $\Tor$-amplitude is at most one.
\begin{definition}\label{def: finite analytic tor-amplitude}
    A morphism $(A^{\tri},\Mod_{A})\to (B^{\tri},\Mod_{B})$ in $\An\CAlg$ is said to be \emph{of analytic $\Tor$-amplitude $\leq n$} if the functor $B\otimes_{A}-:\Mod_{A}\to\Mod_{B}$ takes $\Mod_{A}^{\heartsuit}$ to $\Mod_{B,\leq n}$. 
\end{definition}
\begin{remark}\label{rmk: analytic tor-amplitude}
    If $A^{\tri},B^{\tri}$ are discrete, the notion of analytic $\Tor$-amplitude differs from the $\Tor$-amplitude of the morphism of discrete rings $A^{\tri}\to B^{\tri}$ in $\Sp$ in the sense of \Cref{def: finiteness conditions of morphisms} (5). In particular, if $A^{\tri}\to B^{\tri}$ is a flat morphism (ie. of $\Tor$-amplitude $\leq0$) of animated rings, it need only be of analytic $\Tor$-amplitude $\leq 1$ \cite[Lem. 4.10]{MikamiFPPF}. 
\end{remark}
\begin{lemma}[{\cite[Lem. 4.8]{MikamiFPPF}}]\label{lem: tor amplitude and limits}
    Let $(A^{\tri},\Mod_{A})\to (B^{\tri},\Mod_{B})$ be a morphism in $\An\CAlg^{\wedge}$ such that $B\otimes_{A}\pi_{0}(A^{\tri})\simeq\pi_{0}(B^{\tri})$. Denote by $(\pi_{0}(A^{\tri}),\Mod_{\pi_{0}(A)})\to(\pi_{0}(B^{\tri}),\Mod_{\pi_{0}(B)})$ the induced morphism of static analytic rings and suppose the functor $\pi_{0}(B)\otimes_{\pi_{0}(A)}-$ is of analytic $\Tor$-amplitude $\leq n$.
    \begin{enumerate}
        \item Suppose $\pi_{0}(B)\otimes_{\pi_{0}(A)}-:\Mod_{\pi_{0}(A)}\to\Mod_{\pi_{0}(B)}$ preserves countable limits. Then $B\otimes_{A}-:\Mod_{A}\to\Mod_{B}$ is of analytic $\Tor$-amplitude $\leq n$ and preserves countable limits. 
        \item Assume further the $t$-structure on $\Mod_{A}$ is compatible with small products and $\pi_{0}(B)\otimes_{\pi_{0}(A)}-$ preserves small limits. Then $B\otimes_{A}-$ is of analytic $\Tor$-amplitude $\leq n$ and preserves small limits. 
    \end{enumerate}
\end{lemma}
\begin{proof}[Proof of (1)]
    We first show the claim concerning analytic $\Tor$-amplitude: namely, we seek to prove that $B\otimes_{A}M\in\Mod_{B,\leq n}$ for $M\in\Mod_{A}^{\heartsuit}$. Indeed such $M$ admits the structure of a $\pi_{0}(A^{\tri})$-module and we have an equivalence $\pi_{0}(B)\otimes_{\pi_{0}(A)}M\simeq B\otimes_{A}M$ and the former lies in $\Mod_{\pi_{0}(B),\leq n}$ by assumption. Now by $t$-exactness of the forgetful functor $\Mod_{\pi_{0}(B)}\to\Mod_{B}$, we have $B\otimes_{A}M\in\Mod_{B,\leq n}$ as well. By induction on the (co)fiber sequences $\pi_{k}(M)[k]\to\tau_{\leq k}M\to\tau_{\leq k-1}M$ we obtain $B\otimes_{A}-$ sends $\Mod_{A,[0,r]}$ to $\Mod_{B,[0,n+r]}$. 

    We now show preservation of countable limits. The functor $B\otimes_{A}-$ is an exact and right $t$-exact left adjoint. Using exactness, we can reduce to the case of $B\otimes_{A}-$ preserving countable products, and by compatibility of the $t$-structures on $\Mod_{A},\Mod_{B}$ with countable products and the Whitehead tower in $\Mod_{A}$, it suffices to show that 
    \begin{equation}\label{eqn: countable products}
        B\otimes_{A}\left(\prod_{i\in I}M_{i}\right)\longrightarrow\prod_{i\in I}\left(B\otimes_{A}M_{i}\right)
    \end{equation} 
    is an equivalence in $\Mod_{B,\geq0}$ for countably indexed collections $\{M_{i}\}_{i\in I}\subseteq\Mod_{A,\geq0}$. Passing to Postnikov towers in $\Mod_{B}$ and applying Whitehead's theorem there, it suffices to verify 
    $$\pi_{k}\left(\tau_{\leq m}\left(B\otimes_{A}\left(\prod_{i\in I}M_{i}\right)\right)\right)\to\pi_{k}\left(\tau_{\leq m}\left(\prod_{i\in I}\left(B\otimes_{A}M_{i}\right)\right)\right)$$
    is an equivalence for all $m\geq0$ and $0\leq k\leq m$. Since $B\otimes_{A}-$ is right $t$-exact and the $t$-structures on $\Mod_{A},\Mod_{B}$ are compatible with countable products, the truncations $\tau_{\leq m}\left(B\otimes_{A}\left(\prod_{i\in I}M_{i}\right)\right),\tau_{\leq m}\left(\prod_{i\in I}\left(B\otimes_{A}M_{i}\right)\right)$ depend only on the truncations $\tau_{\leq m}M_{i}$: the fibers of the comparisons with the truncated system are $B\otimes_{A}\left(\prod_{i\in I}\tau_{\geq m+1}M_{i}\right)$ and $\prod_{i\in I}B\otimes_{A}\left(\tau_{\geq m+1}M_{i}\right)$ respectively, both lying in $\Mod_{B,\geq m+1}$. We may therefore assume $\{M_{i}\}_{i\in I}\subseteq\Mod_{A,[0,m]}$. Now writing $\prod_{i\in I}M_{i}$ as its finite Postnikov tower in $\Mod_{A}$, which is preserved by $B\otimes_{A}-$, we can reduce to verifying the claim for $\{M_{i}\}_{i\in I}\subseteq\Mod_{A}^{\heartsuit}$ which was assumed to hold.
\end{proof}
\begin{proof}[Proof of (2)]
    The arguments given in (1) hold, replacing countable products by small ones. More precisely, the claim concerning analytic $\Tor$-amplitude follows from an identical argument to that given in (1). For the second claim, we note that it suffices to show preservation of small products. By the same reductions as in (1) for small products in place of countable ones, we can reduce to the case of $\pi_{0}(B)\otimes_{\pi_{0}(A)}-$ preserving small limits, which was assumed. 
\end{proof}
The lemma in hand, we obtain the recognition result. 
\begin{proposition}\label{prop: open immersion detected on static rings}
    Let $(A^{\tri},\Mod_{A})\to(B^{\tri},\Mod_{B})$ be a morphism in $\An\CAlg^{\wedge}$ with $B\otimes_{A}\pi_{0}(A^{\tri})\simeq\pi_{0}(B^{\tri})$ and such that:
    \begin{enumerate}[label=(\roman*)]
        \item The right adjoint $(-)|_{A^{\tri}}:\Mod_{B}\to\Mod_{A}$ to scalar extension is fully faithful.
        \item The functor $\pi_{0}(B)\otimes_{\pi_{0}(A)}-$ is of analytic $\Tor$-amplitude $\leq 1$.
        \item The $t$-structure on $\Mod_{A}$ is compatible with small products, and the induced morphism $(\pi_{0}(A^{\tri}),\Mod_{\pi_{0}(A)})\to(\pi_{0}(B^{\tri}),\Mod_{\pi_{0}(B)})$ is an open immersion.
    \end{enumerate}
    Then $(A^{\tri},\Mod_{A})\to(B^{\tri},\Mod_{B})$ is an open immersion.
\end{proposition}
\begin{proof}
    We need to show $\ker(B\otimes_{A}-)\subseteq\Mod_{A}$ is equivalent to the category of modules over some idempotent algebra. 

    By (i) the restriction $(-)|_{A^{\tri}}$ is a fully faithful right adjoint to the symmetric monoidal right $t$-exact functor $B\otimes_{A}-$, so $L:=(B\otimes_{A}-)|_{A^{\tri}}$ is an $\otimes_{A}$-compatible $\pi$-stably reflective localization of the presentably symmetric monoidal $t^{+}$-bistable category $\Mod_{A}$ (cf. \Cref{def: compatible localization}). The induced morphism on static rings in (iii) is an open immersion making $\pi_{0}(B)\otimes_{\pi_{0}(A)}-$ preserve small limits, so \Cref{lem: tor amplitude and limits} (2) applies with $n=1$ by conditions (ii) and (iii) above. It follows that $B\otimes_{A}-$ preserves small limits and is of analytic $\Tor$-amplitude $\leq 1$ (cf. \Cref{def: t-amplitude}). \Cref{cor: recognition at t-amplitude 1} then reduces the claim to showing $\intHom_{A}(M, K)\in\ker(L)$ for all $M\in\Mod_{A}^{\heartsuit}$ and $K\in\ker(L)\cap\Mod_{A}^{\heartsuit}$. 

    Write $L^{0}:=(\pi_{0}(B)\otimes_{\pi_{0}(A)}-)|_{\pi_{0}(A^{\tri})}$ for the induced localization. By \Cref{prop: relationship with heart} (1) and \Cref{prop: properties of pi stably reflective localizations} (3), the conservative $t$-exact scalar restriction functor induces an equivalence on the hearts $\Mod_{\pi_{0}(A)}^{\heartsuit}\to\Mod_{A}^{\heartsuit}$. A static object $N\in\Mod_{A}^{\heartsuit}$ admits the structure of a $\pi_{0}(A^{\tri})$-module for which $B\otimes_{A}N\simeq\pi_{0}(B)\otimes_{\pi_{0}(A)}N$ as in the proof of \Cref{lem: tor amplitude and limits} (1). This implies that there is an equivalence $\ker(L)\cap\Mod_{A}^{\heartsuit}\simeq\ker(L^{0})\cap\Mod_{\pi_{0}(A)}^{\heartsuit}$. Both $L, L^{0}$ are exact, right $t$-exact, preserve small colimits, and are of $t$-amplitude $\leq 1$ -- the former by the argument above and the latter by assumption. So $\ker(L)$ and $\ker(L^{0})$ admit the induced $t$-structure by \Cref{lem: homotopy detected kernels of t-amplitude less than 1 functors}. Both localizations preserve small limits, so \Cref{lem: kernel structure L limit preserving} (2) applies to each: an object of $\Mod_{A}$ lies in $\ker(L)$ if and only if each of its homotopy groups lies in $\ker(L)\cap\Mod_{A}^{\heartsuit}$, and an object of $\Mod_{\pi_{0}(A)}$ lies in $\ker(L^{0})$ if and only if each of its homotopy groups lies in $\ker(L^{0})\cap\Mod_{\pi_{0}(A)}^{\heartsuit}$.

    Fix $M, K$ as above. \Cref{lem: modules inherit bistability} (2) gives an equivalence 
    $$\intHom_{A}(M, K)\simeq\intHom_{\pi_{0}(A)}\left(\pi_{0}(A^{\tri})\otimes_{A}M, K\right)|_{A^{\tri}}.$$
    Since $\ker(L^{0})\simeq\Mod_{Q_{0}}(\Mod_{\pi_{0}(A)})$ for some idempotent algebra $Q_{0}\in\Idem_{\pi_{0}(A)}$ by (iii), and $K$ lies in $\ker(L^{0})\cap\Mod_{\pi_{0}(A)}^{\heartsuit}$ under the agreement of kernels on the heart, \Cref{prop: idempotent recognition} (c) gives $\intHom_{\pi_{0}(A)}(\pi_{0}(A^{\tri})\otimes_{A}M, K)\in\ker(L^{0})$. Homotopy detection of $\ker(L^{0})$ places each homotopy group of $\intHom_{\pi_{0}(A)}(\pi_{0}(A^{\tri})\otimes_{A}M, K)$ in $\ker(L^{0})\cap\Mod_{\pi_{0}(A)}^{\heartsuit}$, and $t$-exactness of $(-)|_{A^{\tri}}$ then gives $\pi_{k}(\intHom_{A}(M, K))\in\ker(L)\cap\Mod_{A}^{\heartsuit}$ for all $k\in\ZZ$, so $\intHom_{A}(M, K)\in\ker(L)$, as desired. 
\end{proof}
\subsection{Analytically proper morphisms}\label{subsec: analytically proper morphisms}
Complementary to the open immersions are the analytically proper morphisms, along which the target carries the analytic structure induced from the source. We establish their permanence and show that the right adjoint to scalar extension satisfies the projection formula and base change.
\begin{definition}\label{def: analytically proper}
    A morphism $f:(A^{\tri},\Mod_{A})\to(B^{\tri},\Mod_{B})$ in $\An\CAlg^{\wedge}$ is \emph{analytically proper} if $(B^{\tri},\Mod_{B})$ carries the induced analytic $\EE_{\infty}$-ring structure from $(A^{\tri},\Mod_{A})$ along $f$ in the sense of \Cref{def: induced analytic ring}, ie. $\Mod_{B}=\Mod_{B^{\tri}/A}$.
\end{definition}
\begin{remark}\label{rmk: analytically proper naming}
    Analytically proper morphisms are elsewhere called proper. We retain the qualifier to distinguish them from the proper morphisms of spectral algebraic geometry \cite[Def. 5.1.2.1]{SAG} and the proper morphisms of a six-functor formalism \cite[Def. 4.6.1]{HeyerMann6FF}.
\end{remark}
\begin{proposition}\label{prop: permanence of analytically proper}
    The collection of analytically proper morphisms satisfies the following permanence properties: 
    \begin{enumerate}
        \item (Composition) If $(A^{\tri},\Mod_{A})\to (B^{\tri},\Mod_{B}), (B^{\tri},\Mod_{B})\to (C^{\tri},\Mod_{C})$ are analytically proper morphisms in $\An\CAlg^{\wedge}$ then so is the composite $(A^{\tri},\Mod_{A})\to (B^{\tri},\Mod_{B})\to (C^{\tri},\Mod_{C})$. 
        \item (Base Change) If $(A^{\tri},\Mod_{A})\to (B^{\tri},\Mod_{B})$ is an analytically proper morphism and $(A^{\tri},\Mod_{A})\to (C^{\tri},\Mod_{C})$ another morphism in $\An\CAlg^{\wedge}$ the induced morphism from $(C^{\tri},\Mod_{C})$ to the pushout is analytically proper.
        \item (Left Cancellation) If a composite $(A^{\tri},\Mod_{A})\to (B^{\tri},\Mod_{B})\to (C^{\tri},\Mod_{C})$ is analytically proper and $(A^{\tri},\Mod_{A})\to (B^{\tri},\Mod_{B})$ is analytically proper then so is $(B^{\tri},\Mod_{B})\to (C^{\tri},\Mod_{C})$. 
    \end{enumerate}
\end{proposition}
\begin{proof}[Proof of (1)]
    By induced structures for the constituent morphisms of the composite, a $C^{\tri}$-module is $C$-analytic if and only if its underlying $B^{\tri}$-module is $B$-analytic if and only if its underlying $A^{\tri}$-module is $A$-analytic, so $(C^{\tri},\Mod_{C})$ has the induced structure from $(A^{\tri},\Mod_{A})$ along this composite as well. 
\end{proof}
\begin{proof}[Proof of (2)]
    Since $p:(A^{\tri},\Mod_{A})\to (B^{\tri},\Mod_{B})$ is analytically proper, it is the pushout of $(A^{\tri},\Mod_{A^{\tri}})\to(B^{\tri},\Mod_{B^{\tri}})$ along $(A^{\tri},\Mod_{A^{\tri}})\to(A^{\tri},\Mod_{A})$ as in \Cref{prop: induced analytic ring structure} (2). Now let $(D^{\tri},\Mod_{D})$ be the pushout of $(C^{\tri},\Mod_{C})\leftarrow(A^{\tri},\Mod_{A})\xrightarrow{p}(B^{\tri},\Mod_{B})$. By the pasting lemma, this assembles into a diagram 
    $$% https://q.uiver.app/#q=WzAsNixbMCwwLCIoQV57XFx0cml9LFxcTW9kX3tBXntcXHRyaX19KSJdLFsyLDAsIihCXntcXHRyaX0sXFxNb2Rfe0Jee1xcdHJpfX0pIl0sWzAsMSwiKEFee1xcdHJpfSxcXE1vZF97QX0pIl0sWzIsMSwiKEJee1xcdHJpfSxcXE1vZF97Qn0pIl0sWzAsMiwiKENee1xcdHJpfSxcXE1vZF97Q30pIl0sWzIsMiwiKERee1xcdHJpfSxcXE1vZF97RH0pIl0sWzAsMV0sWzEsM10sWzAsMl0sWzIsM10sWzIsNF0sWzQsNV0sWzMsNV1d
    \begin{tikzcd}
        {(A^{\tri},\Mod_{A^{\tri}})} && {(B^{\tri},\Mod_{B^{\tri}})} \\
        {(A^{\tri},\Mod_{A})} && {(B^{\tri},\Mod_{B})} \\
        {(C^{\tri},\Mod_{C})} && {(D^{\tri},\Mod_{D})}
        \arrow[from=1-1, to=1-3]
        \arrow[from=1-1, to=2-1]
        \arrow[from=1-3, to=2-3]
        \arrow[from=2-1, to=2-3]
        \arrow[from=2-1, to=3-1]
        \arrow[from=2-3, to=3-3]
        \arrow[from=3-1, to=3-3]
    \end{tikzcd}$$
    where the outer rectangle is coCartesian. By \Cref{cor: pushout description} applied to the outer rectangle, a $D^{\tri}$-module is $D$-analytic if and only if its underlying $C^{\tri}$-module is $C$-analytic and its underlying $B^{\tri}$-module is $B^{\tri}$-analytic. The condition that its underlying $B^{\tri}$-module is $B^{\tri}$-analytic is vacuously satisfied since $(B^{\tri},\Mod_{B^{\tri}})$ has the trivial analytic $\EE_{\infty}$-ring structure. So $(C^{\tri},\Mod_{C})\to(D^{\tri},\Mod_{D})$ is analytically proper. 
\end{proof}
\begin{proof}[Proof of (3)]
    Since $B\to C$ is a morphism of analytic $\EE_{\infty}$-rings, $\Mod_{C}\subseteq\Mod_{C^{\tri}/B}$. Conversely, let $M\in\Mod_{C^{\tri}}$ with $M|_{B^{\tri}}\in\Mod_{B}$. As $A\to B$ is a morphism, $M|_{A^{\tri}}\simeq(M|_{B^{\tri}})|_{A^{\tri}}\in\Mod_{A}$, and since $A\to C$ is analytically proper $\Mod_{C}=\Mod_{C^{\tri}/A}$, so $M\in\Mod_{C}$. Hence $\Mod_{C}=\Mod_{C^{\tri}/B}$ and $(B^{\tri},\Mod_{B})\to(C^{\tri},\Mod_{C})$ is analytically proper.
\end{proof}
The base change results for analytically proper morphisms rest on an identification of the unit of the pushout.
\begin{lemma}\label{lem: unit of pushout along proper}
    Let $p:(A^{\tri},\Mod_{A})\to(B^{\tri},\Mod_{B})$ be an analytically proper morphism and $g:(A^{\tri},\Mod_{A})\to(C^{\tri},\Mod_{C})$ a morphism in $\An\CAlg^{\wedge}$ with pushout $(D^{\tri},\Mod_{D})$. The natural morphism $g^{*}(B^{\tri}|_{A^{\tri}})\to D^{\tri}|_{C^{\tri}}$ is an equivalence in $\Mod_{C}$.
\end{lemma}
\begin{proof}
    Use the notation of \Cref{cor: pushout description}. Since $p$ is analytically proper, a $\widetilde{D}^{\tri}$-module has $B$-analytic underlying $B^{\tri}$-module if and only if its underlying $A^{\tri}$-module is $A$-analytic, so $\Mod_{\widetilde{D}^{\tri}/C}\subseteq\Mod_{\widetilde{D}^{\tri}/B}=\Mod_{\widetilde{D}^{\tri}/A}$, and \Cref{cor: pushout description} gives $\Mod_{D}=\Mod_{\widetilde{D}^{\tri}/B}\cap\Mod_{\widetilde{D}^{\tri}/C}=\Mod_{\widetilde{D}^{\tri}/C}$. 
    
    This identifies $\Mod_{D}$ with the category of modules in $\Mod_{C}$ over the analytified algebra $L^{C}(\widetilde{D}^{\tri}|_{C^{\tri}})\in\CAlg(\Mod_{C})$, under which the composite $\widetilde{p}_{*}\widetilde{p}^{*}$ is tensoring with $D^{\tri}|_{C^{\tri}}$. The unit of $\Mod_{D}$ has underlying $C^{\tri}$-module $D^{\tri}|_{C^{\tri}}\simeq L^{C}(\widetilde{D}^{\tri}|_{C^{\tri}})$. The factorization (\ref{eqn: scalar extension is sym mon}) for $g$ applied to $B^{\tri}|_{A^{\tri}}$ gives
    $$g^{*}(B^{\tri}|_{A^{\tri}})\simeq L^{C}(C^{\tri}\otimes_{A^{\tri}}(B^{\tri}|_{A^{\tri}}))\simeq L^{C}(\widetilde{D}^{\tri}|_{C^{\tri}}),$$
    whence the claim.
\end{proof}
This identification underlies the proof of the projection formula and base change below, and enters again in the steadiness criterion of \Cref{prop: nuclear implies steady}.
\begin{proposition}\label{prop: analytically proper base change properties}
    Let $p:(A^{\tri},\Mod_{A})\to(B^{\tri},\Mod_{B})$ be an analytically proper morphism in $\An\CAlg^{\wedge}$.
    \begin{enumerate}
        \item The natural morphism $M\otimes_{A}p_{*}N\to p_{*}(p^{*}M\otimes_{B}N)$ is an equivalence for all $M\in\Mod_{A},N\in\Mod_{B}$. 
        \item For every $g:(A^{\tri},\Mod_{A})\to(C^{\tri},\Mod_{C})$ with pushout $(D^{\tri},\Mod_{D})$ in $\An\CAlg^{\wedge}$ the base change morphism $g^{*}p_{*}\to\widetilde{p}_{*}\widetilde{g}^{*}$ of functors $\Mod_{B}\to\Mod_{C}$ associated to
        $$% https://q.uiver.app/#q=WzAsNCxbMCwwLCJcXE1vZF97QX0iXSxbMiwwLCJcXE1vZF97Qn0iXSxbMCwxLCJcXE1vZF97Q30iXSxbMiwxLCJcXE1vZF97RH0iXSxbMCwyLCJnXnsqfSIsMl0sWzEsMywiXFx3aWRldGlsZGV7Z31eeyp9Il0sWzAsMSwicF57Kn1cXGRhc2h2IHBfeyp9IiwwLHsib2Zmc2V0IjotMX1dLFsxLDAsIiIsMSx7Im9mZnNldCI6LTF9XSxbMiwzLCIiLDIseyJvZmZzZXQiOi0xfV0sWzMsMiwiXFx3aWRldGlsZGV7cH1eeyp9XFxkYXNodlxcd2lkZXRpbGRle3B9X3sqfSIsMCx7Im9mZnNldCI6LTF9XV0=
        \begin{tikzcd}
            {\Mod_{A}} && {\Mod_{B}} \\
            {\Mod_{C}} && {\Mod_{D}}
            \arrow["{p^{*}\dashv p_{*}}", shift left, from=1-1, to=1-3]
            \arrow["{g^{*}}"', from=1-1, to=2-1]
            \arrow[shift left, from=1-3, to=1-1]
            \arrow["{\widetilde{g}^{*}}", from=1-3, to=2-3]
            \arrow[shift left, from=2-1, to=2-3]
            \arrow["{\widetilde{p}^{*}\dashv\widetilde{p}_{*}}", shift left, from=2-3, to=2-1]
        \end{tikzcd}$$
        is an equivalence. 
    \end{enumerate}
\end{proposition}
\begin{proof}[Proof of (1)]
    Since $p:(A^{\tri},\Mod_{A})\to(B^{\tri},\Mod_{B})$ is analytically proper, we have $\Mod_{B^{\tri}}(\Mod_{A})\simeq\Mod_{B}$ hence the projection formula is satisfied by \cite[Cor. 4.8.5.21]{HA}. 
\end{proof}
\begin{proof}[Proof of (2)]
    Note that all functors $g^{*},p_{*},\widetilde{p}_{*},\widetilde{g}^{*}$ commute with colimits and $\Mod_{B}$ generated under small colimits by objects of the form $p^{*}M$ for $M\in\Mod_{A}$. Thus to show $g^{*}p_{*}\to\widetilde{p}_{*}\widetilde{g}^{*}$ is an equivalence it suffices to show the morphism $g^{*}p_{*}p^{*}\to\widetilde{p}_{*}\widetilde{g}^{*}p^{*}\simeq\widetilde{p}_{*}\widetilde{p}^{*}g^{*}$ is one. Under the equivalences $p_{*}p^{*}\simeq B^{\tri}|_{A^{\tri}}\otimes_{A}-$ and $\widetilde{p}_{*}\widetilde{p}^{*}\simeq D^{\tri}|_{C^{\tri}}\otimes_{C}-$ obtained from (1), it is carried to the canonical equivalence $g^{*}(B^{\tri}|_{A^{\tri}}\otimes_{A}-)\simeq D^{\tri}|_{C^{\tri}}\otimes_{C}g^{*}(-)$ supplied by symmetric monoidality of $g^{*}$ and \Cref{lem: unit of pushout along proper}, giving the claim.
\end{proof}
An analytically proper morphism need not be steady: scalar extension is tensoring with the algebra $B^{\tri}$, which need not commute with base change. When $B$ is nuclear as an $A$-module, the projection formula of \Cref{prop: preservation of nuclearity} (2), applied to each base change $g$ with $B^{\tri}|_{A^{\tri}}$ nuclear, supplies steadiness.
\begin{proposition}\label{prop: nuclear implies steady}
    Let $p:(A^{\tri},\Mod_{A})\to(B^{\tri},\Mod_{B})$ be an analytically proper morphism such that $B^{\tri}|_{A^{\tri}}$ is nuclear in $\Mod_{A}$. Then $p$ is a steady morphism. 
\end{proposition}
\begin{proof}
    Let $g:(A^{\tri},\Mod_{A})\to(C^{\tri},\Mod_{C})$ be a general morphism in $\An\CAlg^{\wedge}$ and consider the pushout as in the statement of \Cref{prop: analytically proper base change properties} (2). Since $p_{*}$ is conservative, it suffices to show the equivalence $p^{*}g_{*}\to\widetilde{g}_{*}\widetilde{p}^{*}$ after applying $p_{*}$. By \Cref{lem: unit of pushout along proper}, we have $g^{*}(B^{\tri}|_{A^{\tri}})\simeq D^{\tri}|_{C^{\tri}}\in\Mod_{C}$ and $g_{*}$ preserves small colimits so \Cref{prop: preservation of nuclearity} (2) applied to the adjoint pair $g^{*}\dashv g_{*}$ implies that the natural map
    $$B^{\tri}|_{A^{\tri}}\otimes_{A}g_{*}M\longrightarrow g_{*}(D^{\tri}|_{C^{\tri}}\otimes_{C}M)\simeq g_{*}(g^{*}(B^{\tri}|_{A^{\tri}})\otimes_{C}M)$$
    is an equivalence. Under the equivalence $p_{*}p^{*}\simeq(B^{\tri}|_{A^{\tri}})\otimes_{A}-$ of \Cref{prop: analytically proper base change properties} (1), the base change morphism $p^{*}g_{*}\to\widetilde{g}_{*}\widetilde{p}^{*}$ is taken to the projection formula above under $p_{*}$, both being the adjoint morphism of $\widetilde{g}^{*}p^{*}\simeq \widetilde{p}^{*}g^{*}$ under $g^{*}\dashv g_{*}$. In the case of the target, we compute 
    \begin{align*}
        g_{*}(D^{\tri}|_{C^{\tri}}\otimes_{C}M) &\simeq g_{*}\widetilde{p}_{*}\widetilde{p}^{*}M && \widetilde{p}\text{ analytically proper; \Cref{prop: permanence of analytically proper} (2)} \\
        &\simeq p_{*}\widetilde{g}_{*}\widetilde{p}^{*}M && g_{*}\widetilde{p}_{*}\simeq p_{*}\widetilde{g}_{*} \text{ by commutativity}
    \end{align*}
    so $p_{*}p^{*}g_{*}\xrightarrow{\sim}p_{*}\widetilde{g}_{*}\widetilde{p}^{*}$ is an equivalence, showing $p^{*}g_{*}\xrightarrow{\sim}\widetilde{g}_{*}\widetilde{p}^{*}$ is an equivalence by conservativity, so $p$ is steady. 
\end{proof}
Together with \Cref{prop: permanence properties of steady morphisms}, this furnishes a supply of steady morphisms among the analytically proper ones, to which the descent theory for analytic $\EE_{\infty}$-rings will apply. 
\subsection{Six operations on analytic \texorpdfstring{$\EE_{\infty}$}{E-infinity}-rings}\label{subsec: six operations on analytic rings}
The open immersions and analytically proper morphisms carry the exceptional functors of a six-functor formalism on analytic $\EE_{\infty}$-rings. We recall the framework of geometric setups \cite[\S 2]{HeyerMann6FF} and exhibit the two classes as a suitable decomposition, from which the formalism follows.
\begin{definition}\label{def: geometric setup and suitable decomposition}
    Let $\Cscr$ be a category. 
    \begin{enumerate}
        \item A \emph{geometric setup} $(\Cscr,E)$ is a pair consisting of the category $\Cscr$ and a homotopy class $E$ of \emph{exceptional} morphisms satisfying the following properties:
        \begin{enumerate}[label=(\roman*)]
            \item $E$ contains all equivalences. 
            \item $E$ is closed under composition and pullbacks along morphisms in $\Cscr$. 
            \item $E$ is right cancellative in $\Cscr$: if $f:X\to Y, g:Y\to Z$ are morphisms in $\Cscr$ such that $g\circ f, g\in E$ then $f\in E$. 
        \end{enumerate} 
        \item A \emph{morphism of geometric setups} $(\Cscr_{1},E_{1})\to(\Cscr_{2},E_{2})$ is a functor $F:\Cscr_{1}\to\Cscr_{2}$ such that $F(E_{1})\subseteq E_{2}$ and for all diagrams $X\to Z\leftarrow Y$ in $\Cscr_{1}$ with $X\to Z$ in $E_{1}$, there is an equivalence $F(X\times_{Z}Y)\simeq F(X)\times_{F(Z)}F(Y)$ in $\Cscr_{2}$. 
        \item A \emph{suitable decomposition} of a geometric setup $(\Cscr,E)$ is a pair of classes $I,P\subseteq E$ of morphisms such that:
        \begin{enumerate}[label=(\roman*)]
            \item The morphisms in $I$ and $P$ contain all equivalences in $\Cscr$ and are closed under composition and pullbacks along morphisms in $\Cscr$. 
            \item Each morphism $f\in E$ admits a factorization as $f\simeq p\circ j$ for $p\in P, j\in I$. 
            \item $I$ and $P$ are right cancellative in $\Cscr$. 
            \item Each morphism in $I\cap P$ is $n$-truncated for some $n\geq-2$ (possibly depending on the morphism): some finite iterate of the diagonal morphism is an equivalence. 
        \end{enumerate}
    \end{enumerate}
\end{definition}
\begin{remark}\label{rmk: shriek functors and correspondences}
    The notion of a geometric setup $(\Cscr,E)$ is meant to capture the idea of a category of geometric objects $\Cscr$ and a distinguished collection of morphisms $E$ for which exceptional (or $!$-able) functors exist and satisfy additional properties such as projection and base change formulae.
\end{remark}
We construct a suitable decomposition where we take $I$ to be the open immersions and $P$ to be the analytically proper morphisms, allowing us to define exceptional functors for those morphisms in $\An\CAlg^{\wedge}$ that factor as an analytically proper morhism followed by an open immersion. 
\begin{definition}\label{def: exceptional morphisms of analytic rings}
    A morphism $(A^{\tri},\Mod_{A})\to (B^{\tri},\Mod_{B})$ in $\An\CAlg^{\wedge}$ is \emph{exceptional} if the natural map from the induced analytic $\EE_{\infty}$-ring $(B^{\tri},\Mod_{B^{\tri}/A})\to (B^{\tri},\Mod_{B})$ is an open immersion.
\end{definition} 
\begin{lemma}\label{lem: geometric setup on analytic E-infinity rings}
    Let $\An\Aff:=(\An\CAlg^{\wedge})^{\Opp}$ be the opposite category of complete analytic $\EE_{\infty}$-rings. The collection of exceptional morphisms in the sense of \Cref{def: exceptional morphisms of analytic rings} defines a geometric setup $(\An\Aff,E_{!})$ on $\An\Aff$ and the collections of open immersions $I$ and analytically proper morphisms $P$ define a suitable decomposition of the geometric setup. 
\end{lemma}
\begin{proof}
    The collection of exceptional morphisms contains equivalences (as these are both open immersions and analytically proper) and is preserved under base change since the open immersions and analytically proper morphisms are by \Cref{prop: permanence of open immersions,prop: permanence of analytically proper}. To show that $(\An\Aff,E_{!})$ is a geometric setup, it remains to show closure under composition and right cancellation. For $(A^{\tri},\Mod_{A})\to(B^{\tri},\Mod_{B})$ exceptional and $(B^{\tri},\Mod_{B})\to(C^{\tri},\Mod_{C})$ arbitrary, we have a coCartesian square in $\An\CAlg^{\wedge}$
    $$% https://q.uiver.app/#q=WzAsNCxbMSwwLCIoQl57XFx0cml9LFxcTW9kX3tCfSkiXSxbMCwwLCIoQl57XFx0cml9LFxcTW9kX3tCXntcXHRyaX0vQX0pIl0sWzAsMSwiKENee1xcdHJpfSxcXE1vZF97Q157XFx0cml9L0F9KSJdLFsxLDEsIihDXntcXHRyaX0sXFxNb2Rfe0Nee1xcdHJpfS9CfSkiXSxbMSwyLCJQIiwyXSxbMSwwLCJJIl0sWzAsMywiUCJdLFsyLDMsIkkiLDJdXQ==
    \begin{tikzcd}
        {(B^{\tri},\Mod_{B^{\tri}/A})} & {(B^{\tri},\Mod_{B})} \\
        {(C^{\tri},\Mod_{C^{\tri}/A})} & {(C^{\tri},\Mod_{C^{\tri}/B})}
        \arrow["I", from=1-1, to=1-2]
        \arrow["P"', from=1-1, to=2-1]
        \arrow["P", from=1-2, to=2-2]
        \arrow["I"', from=2-1, to=2-2]
    \end{tikzcd}$$
    where the vertical morphisms are analytically proper and horizontal morphisms are open immersions by the respective permanence properties of these morphisms. Now suppose $(B^{\tri},\Mod_{B})\to(C^{\tri},\Mod_{C})$ is also exceptional. We conclude $(C^{\tri},\Mod_{C^{\tri}/A})\to(C^{\tri},\Mod_{C})$ is an open immersion as it factors as the composite of open immersions $(C^{\tri},\Mod_{C^{\tri}/A})\to(C^{\tri},\Mod_{C^{\tri}/B})\to(C^{\tri},\Mod_{C})$ showing exceptional morphisms are preserved by composition. For the left cancellative property in $\An\CAlg^{\wedge}$, let $(A^{\tri},\Mod_{A})\to (B^{\tri},\Mod_{B}),(A^{\tri},\Mod_{A})\to(B^{\tri},\Mod_{B})\to(C^{\tri},\Mod_{C})$ be exceptional. $(B^{\tri},\Mod_{B})\to(C^{\tri},\Mod_{C^{\tri}/B})$ is analytically proper being the base change of an analytically proper morphism and the factorization of the open immersion $(C^{\tri},\Mod_{C^{\tri}/A})\to(C^{\tri},\Mod_{C})$ above implies $(C^{\tri},\Mod_{C^{\tri}/B})\to(C^{\tri},\Mod_{C})$ is an open immersion by cancellation. This shows that the exceptional morphisms are left cancellative in $\An\CAlg^{\wedge}$ (thus right cancellative in $\An\Aff$). 

    For the latter claim, the permanence properties \Cref{prop: permanence of open immersions,prop: permanence of analytically proper} and the definition of exceptional morphisms \Cref{def: exceptional morphisms of analytic rings} readily imply \Cref{def: geometric setup and suitable decomposition} (3) (i)-(iii), and (iv) holds as open immersions are monomorphisms of analytic $\EE_{\infty}$-rings. 
\end{proof}
A geometric setup determines an operad of correspondences on which a six-functor formalism is a lax symmetric monoidal functor. We recall the two constructions in the generality we need, referring to \cite[\S 2]{HeyerMann6FF} for the correspondence operad and to \cite{ConstructionCLL} for its universal property.
\begin{lemma}[{\cite[Prop. 2.3.3]{HeyerMann6FF}} \& {\cite[\S 4 \& 5]{ConstructionCLL}}]\label{lem: existence of correspondence operad}
    Let $(\Cscr,E)$ be a geometric setup. 
    \begin{enumerate}
        \item There exists an operad $\Corr(\Cscr,E)$ with objects those of $\Cscr$. 
        \item Let $(I,P)$ be a suitable decomposition of $(\Cscr,E)$. There exists an 2-operad $\Corr_{I,P}(\Cscr,E)$ with underlying operad $\Corr(\Cscr,E)$. 
    \end{enumerate}
\end{lemma}
A six-functor formalism assigns to each object a symmetric monoidal category and to each correspondence a pair of adjoint functors, encoded as a lax symmetric monoidal functor out of the correspondence operad. We give the definition in both the general and the presentable forms, the latter being the one realized on analytic $\EE_{\infty}$-rings.
\begin{definition}\label{def: six-functor formalism}
    Let $(\Cscr,E)$ be a geometric setup. 
    \begin{enumerate}
        \item A \emph{six-functor formalism} is a lax symmetric monoidal functor $\Dscr:\Corr(\Cscr,E)\to\widehat{\Cat}$, where we regard the very large category of large categories $\widehat{\Cat}$ as equipped with the Cartesian symmetric monoidal structure, satisfying the following additional conditions: 
        \begin{enumerate}[label=(\roman*)]
            \item For each $X\in \Cscr$, $\Dscr(X)$ is a closed symmetric monoidal category. 
            \item For each $f:X\to Y$ in $\Cscr$, there is an adjoint pair $f^{*}:\Dscr(Y)\rightleftarrows\Dscr(X):f_{*}$. 
            \item For each $f:X\to Y$ in $E$, there is an adjoint pair $f_{!}:\Dscr(X)\rightleftarrows\Dscr(Y):f^{!}$. 
        \end{enumerate}
        \item A \emph{presentable six-functor formalism} is a lax symmetric monoidal functor $\Dscr:\Corr(\Cscr,E)\to\PrL$ where we regard $\PrL$ as equipped with the Lurie tensor product of presentable categories.
    \end{enumerate} 
\end{definition}
\begin{remark}\label{rmk: presentable is automatically 6ff}
    Any lax symmetric monoidal functor $\Corr(\Cscr,E)\to\PrL$ pointwise admits right adjoints and automatically satisfies (i)-(iii) of \Cref{def: six-functor formalism} (1) \cite[Lem. 3.2.5]{HeyerMann6FF}. 
\end{remark}
The two base-change results of the preceding subsections allow us to construct a six-functor formalism on complete analytic $\EE_{\infty}$-rings, unique once refined over the decorated correspondence $2$-category of \Cref{lem: existence of correspondence operad} (2). 
\begin{proposition}\label{prop: six-functor formalism on analytic E-infinity rings}
    Let $(\An\Aff,E_{!}, I, P)$ be the geometric setup and suitable decomposition of \Cref{lem: geometric setup on analytic E-infinity rings} and $\mathdutchcal{Pr}^{\Lsf}$ the category $\PrL$ regarded as a 2-category with non-invertible natural transformations. There is a unique lax symmetric monoidal $2$-functor
    $$\Corr_{I,P}(\An\Aff,E_{!})\to\mathdutchcal{Pr}^{\Lsf}$$
    given on objects by $(A^{\tri},\Mod_{A})\mapsto\Mod_{A}$. Its underlying lax symmetric monoidal functor $\Corr(\An\Aff,E_{!})\to\PrL$ is a presentable six-functor formalism.
\end{proposition}
\begin{proof}
    By \Cref{prop: base change for open immersions,prop: analytically proper base change properties}, conditions (a) and (b) of \cite[Prop. 3.3.3]{HeyerMann6FF} are satisfied, and (c) holds by \cite[Cor. 3.3.5]{HeyerMann6FF}, so the hypotheses of \cite[Thm. 5.17]{ConstructionCLL} are met. This gives the unique lax symmetric monoidal 2-functor $\Corr_{I,P}(\An\Aff,E_{!})\to\mathdutchcal{Pr}^{\Lsf}$. Its underlying functor is $\Corr(\An\Aff,E_{!})\to\PrL$, given on objects by $(A^{\tri},\Mod_{A})\mapsto\Mod_{A}$.
\end{proof}
This is the six-functor formalism on complete analytic $\EE_{\infty}$-rings. 
\begin{definition}\label{def: six-functor formalism on analytic E-infinity rings}
    The \emph{six-functor formalism on complete analytic $\EE_{\infty}$-rings} is the six-functor formalism of \Cref{prop: six-functor formalism on analytic E-infinity rings}. 
\end{definition}
We conclude by showing a property of this six-functor formalism. 
\begin{proposition}\label{prop: categorical kunneth for analytic rings}
    The functor $\An\CAlg^{\wedge}\to\CAlg(\PrL_{\Cond\Sp})$ given on objects by $(A^{\tri},\Mod_{A})\mapsto\Mod_{A}$ commutes with small colimits. 
\end{proposition}
\begin{proof}
    This follows from the proof of \cite[Prop. 4.1.14]{CamargoSolid}, which does not make use of the animated structure. 
\end{proof}
This in particular implies that the six-functor formalism on complete analytic $\EE_{\infty}$-rings satisfies the categorical K\"{u}nneth formula. We recall the definition for presentable six-functor formalisms, then deduce the result for complete analytic $\EE_{\infty}$-rings.
\begin{definition}\label{def: categorical Kunneth formula}
    Let $(\Cscr,E)$ be a geometric setup and $\Dscr:\Corr(\Cscr,E)\to\PrL$ a six-functor formalism. The six-functor formalism \emph{satisfies the categorical K\"{u}nneth formula} if the underlying functor $\Dscr:\Cscr^{\Opp}\to\CAlg(\PrL)$ commutes with pushouts. 
\end{definition}
\begin{corollary}\label{cor: six-functor formalism on analytic E-infinity rings is kunneth}
    The six-functor formalism on complete analytic $\EE_{\infty}$-rings satisfies the categorical K\"{u}nneth formula. 
\end{corollary}
\begin{proof}
    Under the equivalences $\PrL_{\Cond\Sp}\simeq\Mod_{\Cond\Sp}(\PrL)$ and $\CAlg(\PrL_{\Cond\Sp})\simeq\CAlg(\PrL)_{\Cond\Sp/}$, the forgetful functor $\CAlg(\PrL)_{\Cond\Sp/}\to\CAlg(\PrL)$ is a left fibration \cite[\href{https://kerodon.net/tag/018F}{Tag 018F}]{kerodon}, and hence preserves and reflects connected colimits \cite[\href{https://kerodon.net/tag/02KT}{Tag 02KT}]{kerodon}. 
\end{proof}
\subsection{Descent}\label{subsec: descent}
We develop a framework for descent of analytic $\EE_{\infty}$-rings, following L. Mann \cite[\S 2.5-2.8]{MannRigid6FF} and Mikami \cite[\S 2]{MikamiFFDescent}, beginning with a mild refinement of the latter's descent criterion.
\begin{proposition}[{\cite[Lem. 2.17]{MikamiFFDescent}}]\label{prop: Mikami descent}
    Let $(A^{\tri},\Mod_{A})\to (B^{\tri},\Mod_{B})$ be a steady morphism in $\An\CAlg^{\wedge}$ such that: 
    \begin{enumerate}[label=(\roman*)]
        \item The scalar extension functor $F:\Mod_{A}\to\Mod_{B}$ is conservative. 
        \item $F$ preserves totalizations of $F$-split cosimplicial objects in $\Mod_{A}$. 
    \end{enumerate}
    Then the functor 
    \begin{equation}\label{eqn: Cech nerve}
        \Mod_{A}\longrightarrow\lim_{[n]\in\Delta}\Mod_{B^{\otimes_{A}(n+1)}}
    \end{equation}
    is an equivalence, where $\Mod_{B^{\otimes_{A}(\bullet+1)}}$ is the module category of the \v{C}ech nerve of $(A^{\tri},\Mod_{A})\to (B^{\tri},\Mod_{B})$ in $\An\CAlg^{\wedge}$. 
\end{proposition}
\begin{proof}
    See \cite[Rmk. 2.18]{MikamiFFDescent}. We show the hypotheses of the dual of \cite[Cor. 4.7.5.3]{HA} are satisfied: $F$ is conservative and $F$ preserves totalizations of $F$-split cosimplicial objects, and for each $\alpha:[m]\to [n]$ in $\Delta$ the diagram 
    $$% https://q.uiver.app/#q=WzAsNCxbMCwwLCJcXE1vZF97Ql57XFxvdGltZXNfe0F9bX19Il0sWzIsMCwiXFxNb2Rfe0Jee1xcb3RpbWVzX3tBfW0rMX19Il0sWzAsMSwiXFxNb2Rfe0Jee1xcb3RpbWVzX3tBfW59fSJdLFsyLDEsIlxcTW9kX3tCXntcXG90aW1lc197QX1uKzF9fSJdLFswLDJdLFswLDEsImReezB9Il0sWzEsM10sWzIsMywiZF57MH0iLDJdXQ==
    \begin{tikzcd}
        {\Mod_{B^{\otimes_{A}(m+1)}}} && {\Mod_{B^{\otimes_{A}(m+2)}}} \\
        {\Mod_{B^{\otimes_{A}(n+1)}}} && {\Mod_{B^{\otimes_{A}(n+2)}}}
        \arrow["{d^{0}}", from=1-1, to=1-3]
        \arrow[from=1-1, to=2-1]
        \arrow[from=1-3, to=2-3]
        \arrow["{d^{0}}"', from=2-1, to=2-3]
    \end{tikzcd}$$
    is right adjointable. Indeed, any such diagram is coCartesian, and all morphisms are steady. The claim follows as right adjointability of such squares is precisely the definition of steadiness. 
\end{proof}
The criterion of \Cref{prop: Mikami descent} is sufficient but not manifestly stable under composition and base change, since the scalar extension $B\otimes_{A}-$ is tensoring with the algebra $B^{\tri}$ only when the morphism is analytically proper. 

L. Mann's remedy transfers the question to the descendability of steady endofunctors of $\Mod_{A}$. We give a treatment adequate for our purposes, simplified by the light condensed setting. As the arguments of \cite[\S 2.5]{MannRigid6FF} make no use of the animated structure, they apply verbatim here.

For a morphism $f:(A^{\tri},\Mod_{A})\to(B^{\tri},\Mod_{B})$ in $\An\CAlg^{\wedge}$, the scalar extension $f^{*}$ exhibits $\Mod_{B}$ as $\Mod_{A}$-enriched. As $\Mod_{B}$ is closed symmetric monoidal, hence enriched over itself, this $\Mod_{A}$-enrichment is $f_{*}\intHom_{B}(-,-)$.
\begin{definition}\label{def: enriched endofunctors}
    Let $f:(A^{\tri},\Mod_{A})\to(B^{\tri},\Mod_{B})$ be a morphism in $\An\CAlg^{\wedge}$. 
    \begin{enumerate}
        \item The category $\Fun_{f}(B,A)$ of \emph{$\Mod_{A}$-enriched exact functors $\Mod_{B}\to\Mod_{A}$} is the full subcategory of $\Mod_{A}$-enriched functors $\Mod_{B}\to\Mod_{A}$ whose underlying functor is an exact functor of stable categories. 
        \item Let $\End_{A}:=\Fun_{\id_{\Mod_{A}}}(A,A)$ be the category of \emph{$\Mod_{A}$-enriched exact endofunctors}. 
    \end{enumerate}
\end{definition}
\begin{remark}[{\cite[Rmk. 2.5.3, Ex. 2.5.4, \& Lem. A.4.5]{MannRigid6FF}}]\label{rmk: enriched endofunctors}
    The category $\Fun_{f}(B, A)$ of $\Mod_{A}$-enriched exact functors $\Mod_{B}\to\Mod_{A}$ is a stable category admitting small limits and small colimits and contains the objects 
    $$f_{*}\intHom_{B}(M, -)\otimes_{A}N$$ 
    for $M\in\Mod_{B},N\in\Mod_{A}$. 
\end{remark}
The formation of these enriched functor categories is functorial in the morphism $f$. 
\begin{proposition}[{\cite[Lem. 2.5.6]{MannRigid6FF}}]\label{prop: endofunctor functoriality}
    Let 
    $$% https://q.uiver.app/#q=WzAsNCxbMCwwLCIoQV57XFxyaGR9LFxcTW9kX3tBfSkiXSxbMiwwLCIoQl57XFxyaGR9LFxcTW9kX3tCfSkiXSxbMCwxLCIoQ157XFxyaGR9LFxcTW9kX3tDfSkiXSxbMiwxLCIoRF57XFxyaGR9LFxcTW9kX3tEfSkiXSxbMCwxLCJmIl0sWzEsMywiXFx3aWRldGlsZGV7Z30iXSxbMCwyLCJnIiwyXSxbMiwzLCJcXHdpZGV0aWxkZXtmfSIsMl1d
    \begin{tikzcd}
        {(A^{\tri},\Mod_{A})} && {(B^{\tri},\Mod_{B})} \\
        {(C^{\tri},\Mod_{C})} && {(D^{\tri},\Mod_{D})}
        \arrow["f", from=1-1, to=1-3]
        \arrow["g"', from=1-1, to=2-1]
        \arrow["{\widetilde{g}}", from=1-3, to=2-3]
        \arrow["{\widetilde{f}}"', from=2-1, to=2-3]
    \end{tikzcd}$$
    be a commutative diagram in $\An\CAlg^{\wedge}$. There is a natural adjunction 
    $$% https://q.uiver.app/#q=WzAsMyxbMCwwLCIoXFx3aWRldGlsZGV7Zn0sZilee1xcbmF0dXJhbH06XFxGdW5fe2d9KEMsIEEpIl0sWzIsMCwiXFxGdW5fe1xcd2lkZXRpbGRle2d9fShELCBCKTooXFx3aWRldGlsZGV7Zn0sZilfe1xcbmF0dXJhbH0iXSxbMSwwLCJcXG1wZXJwIl0sWzAsMSwiIiwyLHsib2Zmc2V0IjotMn1dLFsxLDAsIiIsMix7Im9mZnNldCI6LTJ9XV0=
    \begin{tikzcd}
        {(\widetilde{f},f)^{\natural}:\Fun_{g}(C, A)} & \mperp & {\Fun_{\widetilde{g}}(D, B):(\widetilde{f},f)_{\natural}.}
        \arrow[shift left=2, from=1-1, to=1-3]
        \arrow[shift left=2, from=1-3, to=1-1]
    \end{tikzcd}$$
\end{proposition}
\begin{remark}\label{rmk: single endofunctors}
    For a morphism $f:(A^{\tri},\Mod_{A})\to (B^{\tri},\Mod_{B})$ in $\An\CAlg^{\wedge}$, denote by $f^{\natural}:=(f,f)^{\natural}$ and $f_{\natural}:=(f,f)_{\natural}$ the adjoint pair between $\End_{A}$ and $\End_{B}$. 
\end{remark}
Steadiness isolates the endofunctors compatible with restriction of scalars: those $F$ for which $F\circ f_{*}$ factors again through $f_{*}$ for every $f:(A^{\tri},\Mod_{A})\to(B^{\tri},\Mod_{B})$.
\begin{definition}[{\cite[Def. 2.5.8]{MannRigid6FF}}]\label{def: steady endofunctors}
    Let $(A^{\tri},\Mod_{A})$ be a complete analytic $\EE_{\infty}$-ring. An enriched endofunctor $F\in\End_{A}$ is a \emph{steady endofunctor} if for all $f:(A^{\tri},\Mod_{A})\to(B^{\tri},\Mod_{B})$, the object $(f,\id_{\Mod_{A}})^{\natural}F\in\Fun_{f}(B, A)$ lies in the essential image of $(\id_{\Mod_{B}},f)_{\natural}:\End_{B}\to\Fun_{f}(B, A)$. 
\end{definition}
\begin{proposition}[{\cite[Prop. 2.5.9 \& Lem. 2.5.10]{MannRigid6FF}}]\label{prop: properties of steady endofunctors}
    Let $f:(A^{\tri},\Mod_{A})\to (B^{\tri},\Mod_{B})$ be a steady morphism in $\An\CAlg^{\wedge}$. 
    \begin{enumerate}
        \item $\Std_{A}$ is a subcategory of $\End_{A}$ closed under small colimits and composition, and contains functors of the form $\intHom_{A}(M,-)$ for $M\in\Mod_{A}$. 
        \item There is an adjunction 
        $$% https://q.uiver.app/#q=WzAsMyxbMCwwLCJmXntcXG5hdHVyYWx9OlxcU3RkX3tBfSJdLFsyLDAsIlxcU3RkX3tCfTpmX3tcXG5hdHVyYWx9Il0sWzEsMCwiXFxtcGVycCJdLFswLDEsIiIsMix7Im9mZnNldCI6LTJ9XSxbMSwwLCIiLDIseyJvZmZzZXQiOi0yfV1d
        \begin{tikzcd}
            {f^{\natural}:\Std_{A}} & \mperp & {\Std_{B}:f_{\natural}}
            \arrow[shift left=2, from=1-1, to=1-3]
            \arrow[shift left=2, from=1-3, to=1-1]
        \end{tikzcd}$$
        where $f_{\natural}$ is conservative and preserves small colimits restricting the one of \Cref{prop: endofunctor functoriality} and \Cref{rmk: single endofunctors}. 
        \item For a coCartesian diagram in $\An\CAlg^{\wedge}$
        $$% https://q.uiver.app/#q=WzAsNCxbMCwwLCIoQV57XFxyaGR9LFxcTW9kX3tBfSkiXSxbMiwwLCIoQl57XFxyaGR9LFxcTW9kX3tCfSkiXSxbMCwxLCIoQ157XFxyaGR9LFxcTW9kX3tDfSkiXSxbMiwxLCIoRF57XFxyaGR9LFxcTW9kX3tEfSkiXSxbMCwxLCJmIl0sWzEsMywiXFx3aWRldGlsZGV7Z30iXSxbMCwyLCJnIiwyXSxbMiwzLCJcXHdpZGV0aWxkZXtmfSIsMl1d
        \begin{tikzcd}
            {(A^{\tri},\Mod_{A})} && {(B^{\tri},\Mod_{B})} \\
            {(C^{\tri},\Mod_{C})} && {(D^{\tri},\Mod_{D})}
            \arrow["f", from=1-1, to=1-3]
            \arrow["g"', from=1-1, to=2-1]
            \arrow["{\widetilde{g}}", from=1-3, to=2-3]
            \arrow["{\widetilde{f}}"', from=2-1, to=2-3]
        \end{tikzcd}$$
        the natural morphism $g^{\natural}f_{\natural}\to\widetilde{f}_{\natural}\widetilde{g}^{\natural}$ is an equivalence of functors from $\Std_{B}$ to $\Std_{C}$. 
    \end{enumerate}
\end{proposition}
\begin{remark}\label{rmk: equivalently in enriched endofunctors}
    By \Cref{prop: properties of steady endofunctors} (1), the inclusion $\Std_{A}\hookrightarrow\End_{A}$ is exact and closed under composition and retracts. Thus the subcategory $\langle(B\otimes_{A}-)|_{A^{\tri}}\rangle$ containing $(B\otimes_{A}-)|_{A^{\tri}}$ closed under finite limits, finite colimits, retracts, and composition is formed identically in $\Std_{A}$ and in $\End_{A}$.
\end{remark}
We can now define descendability.
\begin{definition}\label{def: descendable morphisms}
    Let $f:(A^{\tri},\Mod_{A})\to (B^{\tri},\Mod_{B})$ be a steady morphism in $\An\CAlg^{\wedge}$. Let $K_{f}:=\fib(\id_{\Mod_{A}}\to (B\otimes_{A}-)|_{A^{\tri}})$ in $\End_{A}$ and denote by $K_{f}^{n}$ the $n$-fold composition of $K_{f}$ with itself in $\End_{A}$. 
    \begin{enumerate}
        \item $f$ is \emph{descendable} if there exists some $n\geq0$ such that the natural map $K_{f}^{n}\to \id_{\Mod_{A}}$ is zero in $\End_{A}$. 
        \item $f$ is \emph{descendable of index $n$} if $n$ is the smallest natural number such that the map $K_{f}^{n}\to \id_{\Mod_{A}}$ in (1) is zero in $\End_{A}$.
    \end{enumerate}
\end{definition}
\begin{remark}\label{rmk: formation of K in Std and End}
    By \Cref{rmk: equivalently in enriched endofunctors}, the fiber $K_{f}$, its composites $K_{f}^{n}$, and membership in $\langle(B\otimes_{A}-)|_{A^{\tri}}\rangle$ can be tested equivalently in $\Std_{A}$ and $\End_{A}$. We work with $\End_{A}$ to adhere to the conventions of Mikami's work \cite[\S 4]{MikamiFPPF}. 
\end{remark}
This connects to Mann's notion of descendability \cite[Def. 2.6.1]{MannRigid6FF}, itself closer to Mathew's original definition \cite[\S 3]{MatthewGalois}.
\begin{lemma}[{\cite[Prop. 2.6.5]{MannRigid6FF}}]\label{lem: descendable equivalent}
    Let $f:(A^{\tri},\Mod_{A})\to (B^{\tri},\Mod_{B})$ be a steady morphism in $\An\CAlg^{\wedge}$. The following are equivalent:
    \begin{enumerate}[label=(\alph*)]
        \item $f$ is descendable in the sense of \Cref{def: descendable morphisms} (1). 
        \item $\id_{\Mod_{A}}$ is contained in the smallest full subcategory $\langle (B\otimes_{A}-)|_{A^{\tri}}\rangle$ of $\End_{A}$ containing $(B\otimes_{A}-)|_{A^{\tri}}$ closed under finite limits, finite colimits, retracts, and composition. 
    \end{enumerate}
\end{lemma}
Descendability is likewise stable under composition, base change, and cancellation.
\begin{proposition}[{\cite[Lem. 2.6.9 \& 2.6.10]{MannRigid6FF}}]\label{prop: permanence of descendable morphisms}
    The collection of descendable morphisms satisfies the following permanence properties: 
    \begin{enumerate}
        \item (Composition) If $(A^{\tri},\Mod_{A})\to (B^{\tri},\Mod_{B})$ is descendable of index $n$ and $(B^{\tri},\Mod_{B})\to (C^{\tri},\Mod_{C})$ is descendable of index $m$ then the composite $(A^{\tri},\Mod_{A})\to (B^{\tri},\Mod_{B})\to (C^{\tri},\Mod_{C})$ is descendable of index at most $mn$. 
        \item (Base Change) If $(A^{\tri},\Mod_{A})\to (B^{\tri},\Mod_{B})$ is descendable of index $n$ and $(A^{\tri},\Mod_{A})\to (C^{\tri},\Mod_{C})$ another morphism in $\An\CAlg^{\wedge}$ the induced morphism from $(C^{\tri},\Mod_{C})$ to the pushout is descendable of index at most $n$. 
        \item (Left Cancellation) If a composite $(A^{\tri},\Mod_{A})\to (B^{\tri},\Mod_{B})\to (C^{\tri},\Mod_{C})$ is descendable of index $n$ and $(A^{\tri},\Mod_{A})\to (B^{\tri},\Mod_{B})$ is steady then $(A^{\tri},\Mod_{A})\to (B^{\tri},\Mod_{B})$ is descendable of index at most $n$. 
    \end{enumerate}
\end{proposition}
We conclude the general discussion by showing that analytic module categories descend along descendable morphisms.
\begin{theorem}\label{thm: descent for descendable morphisms}
    If $(A^{\tri},\Mod_{A})\to(B^{\tri},\Mod_{B})$ is a descendable morphism in $\An\CAlg^{\wedge}$ then analytic modules admit descent: the functor (\ref{eqn: Cech nerve}) is an equivalence. 
\end{theorem}
\begin{proof}
    We verify the conditions of \Cref{prop: Mikami descent}. Let $F:\Mod_{A}\to\Mod_{B}$ be the scalar extension functor and $G:\Mod_{B}\to\Mod_{A}$ the scalar restriction functor, which is conservative, so $F$ is conservative if and only if $G\circ F$ is. Indeed, by stability, it suffices to show that the endofunctor $G\circ F$ preserves nonzero objects. Let $M\in\Mod_{A}$ be such that $(G\circ F)(M)\simeq0$. Then $H(M)\simeq0$ for all $H\in\langle G\circ F\rangle$. Picking $H$ to be $\id_{\Mod_{A}}$ gives $M\simeq0$. 

    An $F$-split cosimplicial object $M^{\bullet}$ is also $G\circ F$-split, and $(G\circ F)(M^{\bullet})$ is pro-constant, thus $H(M^{\bullet})$ is pro-constant for any $H\in\langle G\circ F\rangle$. Taking $H$ to be $\id_{\Mod_{A}}$ we get that $F$ preserves the totalization of $M^{\bullet}$, as desired (cf. \cite[Ex. 3.11]{MatthewGalois}).
\end{proof}
Finally, we relate the scalar extension $B\otimes_{A}-$ to its induced morphism $\pi_{0}(B)\otimes_{\pi_{0}(A)}-$ on static analytic rings, reducing conservativity to the latter.
\begin{lemma}\label{lem: conservativity of bounded below}
    Let $(A^{\tri},\Mod_{A})\in\An\CAlg^{\wedge}$. If $f:M\to N$ is a morphism of bounded below objects in $\Mod_{A}$ and $f\otimes_{A}\id_{\pi_{0}(A^{\tri})}:M\otimes_{A}\pi_{0}(A^{\tri})\to N\otimes_{A}\pi_{0}(A^{\tri})$ is an equivalence, then so is $f$. 
\end{lemma}
\begin{proof}
    This is a special case of \Cref{prop: relationship with heart} (2). 
\end{proof}
\section{Discrete solid \texorpdfstring{$\EE_{\infty}$}{E-infinity}-rings}\label{sec: solid commutative rings}
The affine building blocks of spectral algebraic geometry are the connective $\EE_{\infty}$-rings of $\Sp$. Their condensed counterparts are the discrete solid $\EE_{\infty}$-rings, analytic $\EE_{\infty}$-rings with discrete underlying rings, whose modules carry the solid structure, one in which nullsequences are summable as for the solid spectra of \Cref{subsec: solid spectra}. This generalizes the theory of Clausen--Scholze for static rings and of Rodr\'{i}guez Camargo and Hansen--Mann in the animated setting. 

We construct the discrete solid $\EE_{\infty}$-ring of a connective $\EE_{\infty}$-algebra from the free case by sifted colimits, and introduce discrete spectral Huber pairs which afford a more refined structure theory of morphisms. This theory yields the exceptional functors $f_{!}\dashv f^{!}$ for morphisms of finite expansion, which, together with descent for the faithfully flat and finitely presented topology after Mikami, provide the foundation on which the six-functor formalism for spectral algebraic stacks is built.
\subsection{The solid sphere}\label{subsec: solid sphere}
The solid spectra of \Cref{subsec: solid spectra} assemble into an analytic $\EE_{\infty}$-ring over the sphere, the base case for the solid structures of this section.
\begin{proposition}\label{prop: solid spectra are complete analytic}
    The pair $(\Sph,\Cond\Sp^{\sld})$ determines a complete analytic $\EE_{\infty}$-ring structure on $\Cond\Sp$. 
\end{proposition}
\begin{proof}
    By the results on solid spectra proven earlier in \Cref{thm: properties of solid spectra}, $\Cond\Sp^{\sld}$ is closed under small limits, small colimits, and extensions in $\Cond\Sp$ and is $\Cond\Sp$-linear by \cite[Thm. 13.3]{ClausenLie}. Writing $L^{\sld}:\Cond\Sp\to\Cond\Sp^{\sld}$ for the left adjoint to the inclusion, we have for $M\in\Cond\Sp_{\geq0},N\in\Cond\Sp^{\sld}_{\leq-1}$ an equivalence 
    $$\Hom_{\Cond\Sp^{\sld}}(L^{\sld}(M),N)\simeq\Hom_{\Cond\Sp}(M,N)\simeq0$$
    showing $L^{\sld}(M)\in\Cond\Sp^{\sld}_{\geq0}$, $L^{\sld}$ is right $t$-exact, and that the conditions of \Cref{def: analytic E-infinity ring} (1) are satisfied. 

    To show that this analytic $\EE_{\infty}$-ring structure is complete, we need to show $\Sph\in\Cond\Sp^{\sld}$. By \Cref{thm: properties of solid spectra} (2) it suffices to show each homotopy group is solid. But by Serre's theorem, the homotopy groups of $\Sph$ are either $\ZZ$ or a finite Abelian group, which are solid as Abelian groups, so $\Sph\in\Cond\Sp^{\sld}$ as desired. 
\end{proof}
This gives the following definition. 
\begin{definition}\label{def: solid spectra ring}
    The \emph{solid sphere} is the analytic $\EE_{\infty}$-ring $\Sph_{\sld}:=(\Sph,\Mod_{\Sph_{\sld}})=(\Sph,\Cond\Sp^{\sld})$ of \Cref{prop: solid spectra are complete analytic}. 
\end{definition}
The generators of $\Mod_{\Sph_{\sld}}$ are internally compact projective, unlike the merely $\omega_{1}$-compact generators of a general analytic $\EE_{\infty}$-ring (cf. \Cref{prop: A-analytic tensor product}). 
\begin{proposition}\label{prop: generation of solid spectra}
    Let $\Sph_{\sld}$ be the solid sphere. 
    \begin{enumerate}
        \item The object $\prod_{\NN}\Sph$ is internally compact projective in $\Mod_{\Sph_{\sld}}$: $$\intHom_{\Sph_{\sld}}\left(\prod_{\NN}\Sph,-\right):\Mod_{\Sph_{\sld}}\to\Mod_{\Sph_{\sld}}$$
        is $t$-exact and preserves small colimits. 
        \item The objects $\Sph_{\sld}[S]$ for varying $S\in\PFin_{\NN}$ generate $\Mod_{\Sph_{\sld}}$ under small colimits and shifts. In particular, $\Mod_{\Sph_{\sld}}$ is $\omega$-presentable. 
        \item Any compact object of $\Mod_{\Sph_{\sld}}$ is also internally compact and the compact objects are preserved by the tensor product. 
        \item The $t$-structure on $\Mod_{\Sph_{\sld}}$ is compatible with small products. 
    \end{enumerate}
\end{proposition}
\begin{proof}[Proof of (1)]
    This is shown in \cite[Thm. 13.3]{ClausenLie}. 
\end{proof}
\begin{proof}[Proof of (2)]
    The objects $\Sph_{\sld}[S]$ generate $\Mod_{\Sph_{\sld}}$ under small colimits and shifts by \Cref{prop: stability of A-analytic modules} (2). Since they are internally compact projective, and the units of $\Cond\Sp,\Mod_{\Sph_{\sld}}$ are compact projective by \Cref{prop: A-analytic tensor product}, the objects $\Sph_{\sld}[S]$ are compact projective by \Cref{lem: permanence of internal compact projectivity} (2) showing that $\Mod_{\Sph_{\sld}}$ is $\omega$-presentable. 
\end{proof}
\begin{proof}[Proof of (3)]
    From \cite[Lem. 13.7]{ClausenLie}, the internally compact projective generators of $\Mod_{\Sph_{\sld}}$ are preserved by the tensor product. Since $\Mod_{\Sph_{\sld}}$ is presentably symmetric monoidal, we can write each compact object as a retract of finite colimits of the compact generators, and these operations are preserved by the tensor product, showing that the compact objects are preserved by the tensor product. These objects are furthermore internally compact by \Cref{lem: compact implies internally compact}. 
\end{proof}
\begin{proof}[Proof of (4)]
    This is a special case of \Cref{prop: compact projectively generated implies product compatibility}. 
\end{proof}
\begin{remark}\label{rmk: generation of solid spectra}
    In fact, $\prod_{\NN}\Sph$ is a compact projective generator of $\Mod_{\Sph_{\sld}}$. The free solid spectra on an infinite light profinite set is equivalent to $\prod_{\NN}\Sph$ and the free solid spectrum on a finite set is a finite direct sum, hence a retract of $\prod_{\NN}\Sph$. 
\end{remark}
The product compatibility of (4) holds more generally for any complete analytic $\EE_{\infty}$-ring admitting a morphism from $\Sph_{\sld}$.
\begin{proposition}\label{prop: product compatibility rings under solid sphere}
    Let $(A^{\tri},\Mod_{A})$ be a complete analytic $\EE_{\infty}$-ring. If there exists a morphism $(\Sph,\Mod_{\Sph_{\sld}})\to (A^{\tri},\Mod_{A})$ in $\An\CAlg^{\wedge}$ then the $t$-structure on $\Mod_{A}$ is compatible with small products. 
\end{proposition}
\begin{proof}
    The functor $(-)|_{\Sph}:\Mod_{A}\to\Mod_{\Sph_{\sld}}$ is $t$-exact, conservative, and preserves small limits. For a collection $\{M_{i}\}_{i\in I}$ of connective objects in $\Mod_{A}$ with product $M$, we have $M|_{\Sph}\simeq\prod_{i\in I}(M_{i}|_{\Sph})$. The latter is connective by compatibility of the $t$-structure on $\Mod_{\Sph_{\sld}}$ with small products by \Cref{prop: generation of solid spectra} (4) and thus so is the former. By the properties of the forgetful functor, it follows that $M$ is connective in $\Mod_{A}$. 
\end{proof}
\subsection{Solid structures on discrete \texorpdfstring{$\EE_{\infty}$}{E-infinity}-rings}\label{subsec: solid E-infinity rings}
We construct a discrete solid $\EE_{\infty}$-ring from any connective $\EE_{\infty}$-algebra in $\Sp$, regarded as a discrete object of $\Cond\CAlg^{\cn}$. As complete analytic $\EE_{\infty}$-rings are closed under sifted colimits by \Cref{prop: sifted colimits in complete} and $\CAlg^{\cn}$ is generated under sifted colimits by retracts of the free algebras by \Cref{prop: approximation by finitely presented algebras} (1), it suffices to carry out the construction on the compact projective algebras and then pass to sifted colimits.
\begin{lemma}\label{lem: solid rings compact projective}
    Let $A\in\CAlg^{\cn}$ be a compact projective object and regard it as a discrete object of $\Cond\CAlg^{\cn}$. Denote its module category by $\Mod_{A}$. 
    \begin{enumerate}
        \item Let $S\in\PFin_{\NN}$ and $S=\lim_{n}S_{n}$ a presentation of $S$ as a sequential limit of finite sets. Define $A_{\sld}[S]:=\lim_{n}A[S_{n}]$. Then $A_{\sld}[S]\simeq\prod_{I}A$ for $|I|\leq|\NN|$. 
        \item The full subcategory $\Mod_{A_{\sld}}\subseteq\Mod_{A}$ spanned by those modules all of whose homotopy groups are $\pi_{0}(A)_{\sld}$-analytic defines a complete analytic $\EE_{\infty}$-ring structure $A_{\sld}$ which is the unique analytic $\EE_{\infty}$-ring structure on $A$ whose base change to $\pi_{0}(A)$ is $\pi_{0}(A)_{\sld}$. Moreover, there is an equivalence $\pi_{0}(A_{\sld}[S])\simeq\pi_{0}(A)_{\sld}[S]$ of static condensed $\pi_{0}(A)$-modules for all $S\in\PFin_{\NN}$. 
    \end{enumerate}
\end{lemma}
\begin{proof}[Proof of (1)]
    Note that 
    \begin{equation}\label{eqn: free solid modules}
        A_{\sld}[S]:=\lim_{n}A[S_{n}]\simeq\lim_{n}\prod_{S_{n}}A\simeq\lim_{n}\bigoplus_{S_{n}}A\simeq\intHom_{A}\left(\colim_{n}\prod_{S_{n}}A,A\right)
    \end{equation} 
    and 
    \begin{equation}\label{eqn: continuous functions}
        A\otimes_{\Sph}\left(\colim_{n}\prod_{S_{n}}\Sph\right)\simeq\colim_{n}\prod_{S_{n}}A\simeq\colim_{n}\bigoplus_{S_{n}}A
    \end{equation} 
    by $\Cond\Sp$ being presentably symmetric monoidal and exactness of the tensor product. (\ref{eqn: continuous functions}) naturally identifies with $\Cont(S,\Sph)\otimes_{\Sph}A\simeq\Cont(S, A)$ (cf. \cite[Def. 3.5.10 ff.]{HeyerMann6FF}). Thus by base-changing the equivalence $\bigoplus_{I}\Sph\xrightarrow{\sim}\Cont(S,\Sph)$ for $|I|\leq|\NN|$ in \cite[Pf. Thm. 13.3]{ClausenLie}, we get that $\bigoplus_{I}A\xrightarrow{\sim}\Cont(S, A)$ is an equivalence. Applying $\intHom_{A}(-,A)$, we obtain the desired equivalence from (\ref{eqn: free solid modules}). 
\end{proof}
\begin{proof}[Proof of (2)]
    Since $A$ is compact projective in $\CAlg^{\cn}$, $\pi_{0}(A)$ is a finitely presented $\ZZ$-algebra and carries the static solid analytic $\EE_{\infty}$-ring structure $\pi_{0}(A)_{\sld}$ \cite[Thm. 5.3.2]{CamargoSolid}. By \Cref{prop: functorial heart invariance} it lifts to a unique analytic $\EE_{\infty}$-ring structure $(A,\Mod_{A_{\sld}})$ on $\Mod_{A}$ which induces $\pi_{0}(A)_{\sld}$ on base change to $\pi_{0}(A)$, spanned by those modules whose homotopy groups are $\pi_{0}(A)_{\sld}$-analytic by \Cref{prop: t-structure on A-analytic modules} (vi). It is complete as each $\pi_{k}(A)$ is a discrete, hence $\pi_{0}(A)$-solid, $\pi_{0}(A)$-module, so the unit $A$ lies in $\Mod_{A_{\sld}}$. 

    For the second claim, (1) identifies $A_{\sld}[S]$ with $\prod_{I}A$ for $|I|\leq|\NN|$. Using exactness of countable products in $\Mod_{A}$, $\pi_{0}(A_{\sld}[S])\simeq\pi_{0}(\prod_{I}A)\simeq\prod_{I}\pi_{0}(A)\simeq\pi_{0}(A)_{\sld}[S]$ by \cite[Thm. 5.3.2]{CamargoSolid} once more.  
\end{proof}
An arbitrary connective $\EE_{\infty}$-algebra is obtained by sifted colimits.
\begin{proposition}\label{prop: solid rings general}
    Let $B\in\CAlg^{\cn}$ and regard it as a discrete object of $\Cond\CAlg^{\cn}$. Denote its module category by $\Mod_{B}$. Consider the diagram 
    \begin{equation}\label{eqn: sifted colimit of cp subalgebras}
        I:= \left\{A_{i}\to B: A_{i}\text{ compact projective in }\CAlg^{\cn}\right\}.
    \end{equation}
    \begin{enumerate}
        \item The diagram $I$ of (\ref{eqn: sifted colimit of cp subalgebras}) is sifted and the natural map 
        $$\colim_{i\in I}A_{i}\to B$$
        is an equivalence in $\CAlg^{\cn}$. 
        \item The colimit in $\An\CAlg^{\wedge}$ 
        $$B_{\sld}:=\colim_{i\in I}A_{i,\sld}$$
        defines a complete analytic $\EE_{\infty}$-ring structure on $\Mod_{B}$. 
        \item The objects 
        $$B_{\sld}[S]\simeq\colim_{i\in I}A_{i,\sld}[S]$$
        for $S\in\PFin_{\NN}$ are internally compact projective in $\Mod_{B_{\sld}}$ and generate $\Mod_{B_{\sld}}$ under small colimits and shifts in $\Mod_{B}$. In particular, $\Mod_{B_{\sld}}$ is $\omega$-presentable. 
    \end{enumerate}
\end{proposition}
\begin{proof}[Proof of (1)]
    Siftedness of the diagram and the equivalence follow from \Cref{prop: approximation by finitely presented algebras} (1).
\end{proof}
\begin{proof}[Proof of (2)]
    Recall from \Cref{prop: sifted colimits in complete} that complete analytic $\EE_{\infty}$-rings are closed under sifted colimits computed in analytic $\EE_{\infty}$-rings. Since each $A_{i,\sld}$ is a complete analytic $\EE_{\infty}$-ring by \Cref{lem: solid rings compact projective}, it follows that $B_{\sld}$ is a complete analytic $\EE_{\infty}$-ring as well. 
\end{proof}
\begin{proof}[Proof of (3)]
    The claimed equivalence is provided by \Cref{lem: sifted colimits in AnCAlg} and $B_{\sld}$ is a complete analytic $\EE_{\infty}$-ring by \Cref{prop: sifted colimits in complete}. By \Cref{prop: stability of A-analytic modules} (3), the objects $B_{\sld}[S]$ generate $\Mod_{B_{\sld}}$ under small colimits and shifts. 

    We have a map of discrete solid $\EE_{\infty}$-rings $\Sph_{\sld}\to B_{\sld}$ by \Cref{prop: constructing morphisms of analytic rings}. Adjunction then furnishes $\intHom_{B_{\sld}}(B_{\sld}[S],-)|_{\Sph}\simeq\intHom_{\Sph_{\sld}}(\Sph_{\sld}[S],(-)|_{\Sph})$. Since $(-)|_{\Sph}:\Mod_{B_{\sld}}\to\Mod_{\Sph_{\sld}}$ is conservative, $t$-exact, and preserves small colimits, internal compact projectivity of $\Sph_{\sld}[S]$ in $\Mod_{\Sph_{\sld}}$ implies the same for $B_{\sld}[S]$ in $\Mod_{B_{\sld}}$. The unit of $\Mod_{B_{\sld}}$ is compact projective by \Cref{prop: A-analytic tensor product}, so \Cref{lem: permanence of internal compact projectivity} (2) exhibits $B_{\sld}[S]$ as compact projective generators of $\Mod_{B_{\sld}}$. 
\end{proof}
This allows us to define discrete solid $\EE_{\infty}$-rings. 
\begin{definition}\label{def: solid E-infinity ring}
    \begin{enumerate}
        \item Let $A\in\CAlg^{\cn}$. The \emph{discrete solid $\EE_{\infty}$-ring} of $A$ is the analytic $\EE_{\infty}$-ring $A_{\sld}:=(A,\Mod_{A_{\sld}})$ of \Cref{prop: solid rings general}.
        \item The \emph{category of discrete solid $\EE_{\infty}$-rings} $\An\CAlg^{\sld}$ is the full subcategory of $\An\CAlg^{\wedge}$ spanned by discrete solid $\EE_{\infty}$-rings of connective $\EE_{\infty}$-algebras in $\Sp$. 
    \end{enumerate}
\end{definition}
The formation of discrete solid $\EE_{\infty}$-rings enjoys the following functoriality. 
\begin{lemma}\label{lem: functoriality of solid rings}
    There is a functor $\CAlg^{\cn}\to\An\CAlg^{\wedge}$ given on objects by $A\mapsto A_{\sld}$ which is fully faithful and preserves nonempty small colimits. 
\end{lemma}
\begin{proof}
    By full faithfulness of $\CAlg^{\cn}\hookrightarrow\Cond\CAlg^{\cn}$ and the adjunction of \Cref{prop: trivial is left adjoint}, \Cref{def: analytic E-infinity ring} (3) exhibits $\Mor_{\An\CAlg^{\wedge}}(A_{\sld},B_{\sld})$ as the subanima of $\Mor_{\CAlg^{\cn}}(A,B)$ spanned by those morphisms along which restriction of scalars carries $\Mod_{B_{\sld}}$ into $\Mod_{A_{\sld}}$. This condition holds for every morphism: $B_{\sld}$ is the colimit of the $C_{\sld}$ over the diagram (\ref{eqn: sifted colimit of cp subalgebras}) of all morphisms $C\to B$ from compact projective objects $C$, computed in $\An\CAlg$ by \Cref{prop: sifted colimits in complete}, so \Cref{prop: colimits in AnCAlg} shows that a $B$-module is $B_{\sld}$-analytic if and only if its restriction along every morphism from a compact projective object is analytic, and the same holds for $A$. For $M\in\Mod_{B_{\sld}}$ and any morphism $C\to A$ from a compact projective $C$, the composite $C\to A\to B$ then gives $(M|_{A})|_{C}\simeq M|_{C}\in\Mod_{C_{\sld}}$, so $M|_{A}\in\Mod_{A_{\sld}}$. The functor is therefore fully faithful. 

    For colimit preservation, note first that the functor preserves sifted colimits: $\CAlg^{\cn}$ is generated under sifted colimits by its compact projective objects by \Cref{prop: approximation by finitely presented algebras} (1), and \Cref{prop: solid rings general} exhibits $A\mapsto A_{\sld}$ as the extension of its values on compact projective objects by sifted colimits, so it preserves sifted colimits \cite[Prop. 5.5.8.15]{HTT}. It remains to show preservations of binary products. Since the coproduct $\otimes_{\Sph}$ in $\CAlg^{\cn}$ preserves sifted colimits separately in each argument, it suffices to show that the functor preserves binary coproducts of compact projective objects. 

    Let $B,C$ be compact projective with coproduct $D:=B\otimes_{\Sph}C$ in $\CAlg^{\cn}$. $D$ is compact projective as $\Mor_{\CAlg^{\cn}}(D,-)\simeq\Mor_{\CAlg^{\cn}}(B,-)\times\Mor_{\CAlg^{\cn}}(C,-)$ preserves sifted colimits. The embedding $\Sp\hookrightarrow\Cond\Sp$ is a symmetric monoidal left adjoint, so $D$ is also the coproduct in $\Cond\CAlg^{\cn}$ and \Cref{prop: colimits in AnCAlg} computes the coproduct of $B_{\sld}$ and $C_{\sld}$ in $\An\CAlg$ as $(D,\Mod_{D/B_{\sld}}\cap\Mod_{D/C_{\sld}})$: those $D$-modules whose underlying $B$- and $C$-modules are $B_{\sld}$- and $C_{\sld}$-analytic. This structure is complete, the homotopy groups of $D$ being discrete and hence solid as $\pi_{0}(B)$- and $\pi_{0}(C)$-modules, so it is also the coproduct in $\An\CAlg^{\wedge}$ by \Cref{cor: complete AnCAlg has colimits}. 

    It remains to identify $(D,\Mod_{D/B_{\sld}}\cap\Mod_{D/C_{\sld}})$ with $D_{\sld}$. Restriction of scalars is $t$-exact and membership in $\Mod_{B_{\sld}}$, $\Mod_{C_{\sld}}$, and $\Mod_{D_{\sld}}$ is detected on homotopy groups by \Cref{lem: solid rings compact projective} (2), so both analytic structures on $\Mod_{D}$ are detected on homotopy groups and \Cref{prop: functorial heart invariance} reduces the identification to the static solid structures on $\pi_{0}(D)\cong\pi_{0}(B)\otimes_{\ZZ}\pi_{0}(C)$. A static condensed $\pi_{0}(D)$-module is $\pi_{0}(D)_{\sld}$-analytic if and only if its restrictions to $\pi_{0}(B)$ and $\pi_{0}(C)$ are $\pi_{0}(B)_{\sld}$- and $\pi_{0}(C)_{\sld}$-analytic, which is the preservation of binary coproducts by the animated solidification functor of \cite[Prop. 5.4.11]{CamargoSolid}. 
\end{proof}
\begin{remark}\label{rmk: solid E-infinity rings is equivalent to connective rings}
    The functor $\CAlg^{\cn}\to\An\CAlg^{\sld}$ given on objects by $A\mapsto A_{\sld}$ is fully faithful and essentially surjective and thus an equivalence of categories. 
\end{remark}
\begin{remark}\label{rmk: morphisms of solid rings}
    Since $\CAlg^{\cn}\to\An\CAlg^{\sld}$ is an equivalence of categories, any property $\Pcal$ of morphisms in $\CAlg^{\cn}$ applies unambiguously to morphisms in $\An\CAlg^{\sld}$: a morphism $A_{\sld}\to B_{\sld}$ is said to be $\Pcal$ if and only if it is the image of a morphism $A\to B$ in $\CAlg^{\cn}$ with property $\Pcal$ under the functor of \Cref{lem: functoriality of solid rings}. We recall some standard notions of morphisms of (connective) $\EE_{\infty}$-algebras in $\Sp$ in \Cref{def: finiteness conditions of morphisms}. 
\end{remark}
\subsection{Discrete spectral Huber pairs}\label{subsec: spectral Huber pairs}
The discrete solid $\EE_{\infty}$-rings of \Cref{subsec: solid E-infinity rings} are refined to discrete solid Huber rings, whose extra plus-part affords finer control of the solid structure on the modules over a fixed $\EE_{\infty}$-ring and, most crucially, a convenient factorization of morphisms.
\begin{definition}\label{def: discrete spectral Huber pair}
    Let $A\in\CAlg^{\cn}$ be a connective $\EE_{\infty}$-algebra in $\Sp$. 
    \begin{enumerate}
        \item A \emph{discrete spectral Huber pair} $(A,A^{+})$ is a morphism $A^{+}\to A$ in $\CAlg^{\cn}$ such that $\pi_{0}(A^{+})\subseteq\pi_{0}(A)$ is an integrally closed subring of the static ring $\pi_{0}(A)$. 
        \item A \emph{morphism of discrete spectral Huber pairs} $f:(A,A^{+})\to (B,B^{+})$ is a morphism $f:A\to B$ in $\CAlg^{\cn}$ such that the induced morphism on static rings $\pi_{0}(f):\pi_{0}(A)\to\pi_{0}(B)$ satisfies $\pi_{0}(f)(\pi_{0}(A^{+}))\subseteq\pi_{0}(B^{+})$. 
        \item Denote by $\Huber$ the \emph{category of discrete spectral Huber pairs} with objects discrete spectral Huber pairs and morphisms those of discrete spectral Huber pairs in the sense of (2). 
    \end{enumerate}
\end{definition}
\begin{remark}\label{rmk: adjectives discrete and spectral}
    In what follows, we will often elide the adjective ``spectral'' when discussing the constructions of the preceding definition. 
\end{remark}
\begin{remark}\label{rmk: comparison with M and RC}
    Compare the definitions of Rodr\'{i}guez Camargo \cite[\S 5.5]{CamargoSolid} and of L. Mann \cite[\S 2.9]{MannRigid6FF} where discrete animated Huber pairs are defined to be pairs consisting of an animated ring together with a distinguished subanima of elements and a distinguished static subring, respectively. 
\end{remark}
Discrete spectral Huber pairs give rise to complete analytic $\EE_{\infty}$-rings.
\begin{lemma}\label{lem: Huber pairs to solid rings}
    Let $(A,A^{+})\in\Huber$ and $(A,\Mod_{(A,A^{+})_{\sld}}):=(A,\Mod_{A/A^{+}_{\sld}})$ be the analytic $\EE_{\infty}$-ring induced from $A^{+}_{\sld}$ along $A^{+}\to A$. 
    \begin{enumerate}
        \item $(A,\Mod_{(A,A^{+})_{\sld}})$ is a complete analytic $\EE_{\infty}$-ring structure on $\Mod_{A}$. 
        \item $\Mod_{(A,A^{+})_{\sld}}$ is generated under colimits and shifts by internally compact projective objects $A\otimes_{A^{+}}A^{+}_{\sld}[S]$ for $S\in\PFin_{\NN}$. In particular, $\Mod_{(A,A^{+})_{\sld}}$ is $\omega$-presentable with compact unit. 
    \end{enumerate}
\end{lemma}
\begin{proof}[Proof of (1)]
    By \Cref{prop: induced analytic ring structure}, $(A,\Mod_{(A,A^{+})_{\sld}})$ is an analytic $\EE_{\infty}$-ring structure. Completeness follows from $A$ being discrete in $\Mod_{A^{+}_{\sld}}$. 
\end{proof}
\begin{proof}[Proof of (2)]
    The objects $A\otimes_{A^{+}}A^{+}_{\sld}[S]$ are internally compact projective and generate $\Mod_{(A,A^{+})_{\sld}}$ under small colimits and shifts by \Cref{prop: scalar extension preserves internal compact projectives}. The unit of $\Mod_{(A,A^{+})_{\sld}}$ is compact projective by \Cref{prop: A-analytic tensor product}, so any internally compact projective $M$ is itself compact projective, its mapping spectrum $\Hom_{(A,A^{+})_{\sld}}(M,-)\simeq\Hom_{(A,A^{+})_{\sld}}(A,\intHom_{(A,A^{+})_{\sld}}(M,-))$ being a composite of $t$-exact functors preserving small colimits. In particular the generators are compact projective, whence $\Mod_{(A,A^{+})_{\sld}}$ is $\omega$-presentable. 
\end{proof}
We refer to these as discrete solid Huber rings. 
\begin{definition}\label{def: solid rings of Huber pairs}
    Let $(A,A^{+})\in\Huber$. The \emph{discrete solid Huber ring} $(A,A^{+})_{\sld}$ of the Huber pair $(A,A^{+})$ is the complete analytic $\EE_{\infty}$-ring of \Cref{lem: Huber pairs to solid rings} (1). 
\end{definition}
\begin{remark}\label{rmk: Huber pairs for arbitrary plus parts}
    For an arbitrary morphism $A^{+}\to A$ in $\CAlg^{\cn}$, we continue to write $(A,A^{+})_{\sld}$ for the analytic $\EE_{\infty}$-ring induced from $A^{+}_{\sld}$ along $A^{+}\to A$. The statement and proof of \Cref{lem: Huber pairs to solid rings} apply verbatim. This convention parallels \cite[Def. 2.9.3]{MannRigid6FF} (c) as the induced structure is invariant under passage to the integral closure of the image. 
\end{remark}
As with discrete solid $\EE_{\infty}$-rings in \Cref{lem: functoriality of solid rings}, the formation of discrete solid Huber rings is likewise functorial. 
\begin{lemma}\label{lem: functoriality of Huber rings}
    There is a functor $\Huber\to\An\CAlg^{\wedge}$ given on objects by $(A,A^{+})\mapsto(A,A^{+})_{\sld}$ which is fully faithful. 
\end{lemma}
\begin{proof}
    Let $f:(A,A^{+})\to(B,B^{+})$ be a morphism in $\Huber$ and note that a $B$-module is $(B,B^{+})_{\sld}$-analytic if and only if all its homotopy groups are $\pi_{0}(B^{+})_{\sld}$-analytic. But this implies they are in particular $\pi_{0}(A^{+})_{\sld}$-analytic by the containment $\pi_{0}(f)(\pi_{0}(A^{+}))\subseteq\pi_{0}(B^{+})$. Conversely, for a map of discrete solid Huber rings $(A,A^{+})_{\sld}\to(B,B^{+})_{\sld}$, we can verify the condition that this comes from a map of discrete spectral Huber pairs on the underlying static rings: \cite[Prop. 13.16]{Analytic} shows that the image of $\pi_{0}(A^{+})$ in $\pi_{0}(B)$ is integral over $\pi_{0}(B^{+})$, and integral closedness of $\pi_{0}(B^{+})$ in $\pi_{0}(B)$ gives the containment $\pi_{0}(f)(\pi_{0}(A^{+}))\subseteq\pi_{0}(B^{+})$. 
\end{proof}
\begin{remark}\label{rmk: Huber to analytic ring preserves colimits}
    In \cite[Prop. 2.9.6]{MannRigid6FF}, it is also shown that the formation of discrete solid animated Huber rings from discrete animated Huber pairs preserves nonempty small colimits. Since we do not need this result, we do not reprove its spectral analogue here.
\end{remark}
When $A^{+}\simeq A$, the discrete solid Huber ring $(A,A^{+})_{\sld}$ recovers the discrete solid $\EE_{\infty}$-ring $A_{\sld}$, and morphisms then factor through an analytically proper one.
\begin{proposition}\label{prop: factorization of Huber ring maps}
    Let $(A,A^{+})_{\sld}\to (B,B^{+})_{\sld}$ be a morphism of discrete solid Huber rings. 
    \begin{enumerate}
        \item The morphism admits a factorization as 
        $$(A,A^{+})_{\sld}\to (B,A^{+})_{\sld}\to (B,B^{+})_{\sld}$$
        where $(A,A^{+})_{\sld}\to (B,A^{+})_{\sld}$ is an analytically proper morphism. 
        \item In particular, any morphism $A_{\sld}\to B_{\sld}$ of discrete solid $\EE_{\infty}$-rings admits a factorization 
        $$A_{\sld}\to (B,A)_{\sld}\to B_{\sld}$$
        where $A_{\sld}\to (B,A)_{\sld}$ is an analytically proper morphism. 
    \end{enumerate}
\end{proposition}
\begin{proof}[Proof of (1)]
    Since $(A,A^{+})_{\sld}\to(B,B^{+})_{\sld}$ is a morphism of discrete solid Huber rings, \Cref{lem: functoriality of Huber rings} supplies the containment $\pi_{0}(f)(\pi_{0}(A^{+}))\subseteq\pi_{0}(B^{+})$. The factorization consists of morphisms of analytic $\EE_{\infty}$-rings with $(B,A^{+})_{\sld}$ following \Cref{rmk: Huber pairs for arbitrary plus parts}. Restriction along $A\to B$ carries $(B,A^{+})_{\sld}$-analytic modules to $(A,A^{+})_{\sld}$-analytic ones, both structures being induced from $A^{+}_{\sld}$ showing the first morphism is analytically proper. 

    To identify the composite, note that we have $L^{(B,B^{+})_{\sld}}\circ L^{(B,A^{+})_{\sld}}\simeq L^{(B,B^{+})_{\sld}}$ by containment $\Mod_{(B,B^{+})_{\sld}}\subseteq\Mod_{(B,A^{+})_{\sld}}$ so the scalar extension functor of the factorization is $L^{(B,B^{+})_{\sld}}\circ L^{(B,A^{+})_{\sld}}\circ (B\otimes_{A}-)\simeq L^{(B,B^{+})_{\sld}}\circ(B\otimes_{A}-)$ which is precisely the scalar extension functor of the original morphism by \Cref{prop: scalar extension functor is symmetric monoidal}.  
\end{proof}
\begin{proof}[Proof of (2)]
    This is immediate from (1). 
\end{proof}
\subsection{Mereology of solid modules}\label{subsec: mereology of solid modules}
In analytic modules over a discrete solid Huber ring, nuclear objects coincide with the discrete ones. 
\begin{proposition}[Cf. {\cite[Prop. 2.9.7]{MannRigid6FF}}]\label{prop: nuclear objects in solid Huber rings}
    Let $\Acal:=(A,A^{+})_{\sld}$ be a discrete solid Huber ring and denote the generators $\Acal[S]:=(A,A^{+})_{\sld}[S]=A\otimes_{A^{+}}A^{+}_{\sld}[S]$ for varying light profinite sets $S\in\PFin_{\NN}$. The following conditions are equivalent for an object $M\in\Mod_{(A,A^{+})_{\sld}}=\Mod_{\Acal}$: 
    \begin{enumerate}[label=(\alph*)]
        \item $M$ is discrete. 
        \item $M$ is nuclear in the sense of \Cref{def: nuclearity and trace class morphisms} (3). 
        \item For all $N\in\Mod_{\Acal}$ and $S\in\PFin_{\NN}$, the natural morphism
        $$\Hom_{\Acal}\left(A, \intHom_{\Acal}\left(\Acal[S],N\right)\otimes_{\Acal}M\right)\to\Hom_{\Acal}\left(\Acal[S], N\otimes_{\Acal}M\right)$$
        is an equivalence in $\Sp$. 
    \end{enumerate}
\end{proposition}
\begin{proof}
    We show (a)$\Rightarrow$(b)$\Rightarrow$(c)$\Rightarrow$(a). 

    (a)$\Rightarrow$(b) is \Cref{cor: discrete implies nuclear}. 

    (b)$\Rightarrow$(c) Since $A$ is compact by \Cref{prop: A-analytic tensor product} and $P:=\Acal[S]$ (internally) compact projective by \Cref{lem: Huber pairs to solid rings} (2), the natural morphism in (c) commutes with small colimits in $M$. We may thus take $M$ to be a $\QQ_{\geq0}$-indexed colimit of compact objects along trace class morphisms $M:=\colim_{q\in\QQ_{\geq0}}M_{q}$. Choosing a factorization
    $$A\xrightarrow{\eta_{q,q'}}M_{q}^{\vee}\otimes_{\Acal}M_{q'}\to \intHom_{\Acal}(M_{q},M_{q'})$$
    of $f_{q,q'}:M_{q}\to M_{q'}$, we get a morphism 
    \begin{align*}
        \Hom_{\Acal}\left(P,N\otimes_{\Acal}M_{q}\right) &\simeq \Hom_{\Acal}\left(A,\intHom_{\Acal}\left(P,N\otimes_{\Acal}M_{q}\right)\right) \\
        &\to \Hom_{\Acal}\left(A,\intHom_{\Acal}\left(P,N\otimes_{\Acal}M_{q}\right)\otimes_{\Acal}M_{q}^{\vee}\otimes_{\Acal}M_{q'}\right) \\
        &\to \Hom_{\Acal}\left(A,\intHom_{\Acal}\left(P,N\otimes_{\Acal}M_{q}\otimes_{\Acal}M_{q}^{\vee}\right)\otimes_{\Acal}M_{q'}\right) \\
        &\to \Hom_{\Acal}\left(A,\intHom_{\Acal}\left(P,N\right)\otimes_{\Acal}M_{q'}\right)
    \end{align*}
    where the morphisms from the second line onwards are, respectively, induced by $\eta_{q,q'}$, the canonical map $\intHom_{\Acal}(P,X)\otimes_{\Acal}Y\to\intHom_{\Acal}(P,X\otimes_{\Acal}Y)$, and the evaluation map $\ev_{M_{q}}:M_{q}\otimes_{\Acal}M_{q}^{\vee}\to A$ of $M_{q}$. This provides a diagonal filler to the diagram 
    $$% https://q.uiver.app/#q=WzAsNCxbMCwxLCJcXEhvbV97XFxBY2FsfVxcbGVmdChQLE5cXG90aW1lc197XFxBY2FsfU1fe3F9XFxyaWdodCkiXSxbMCwwLCJcXEhvbV97XFxBY2FsfVxcbGVmdChBLFxcaW50SG9tX3tcXEFjYWx9XFxsZWZ0KFAsTlxccmlnaHQpXFxvdGltZXNfe1xcQWNhbH1NX3txfVxccmlnaHQpIl0sWzIsMCwiXFxIb21fe1xcQWNhbH1cXGxlZnQoQSxcXGludEhvbV97XFxBY2FsfVxcbGVmdChQLE5cXHJpZ2h0KVxcb3RpbWVzX3tcXEFjYWx9TV97cSd9XFxyaWdodCkiXSxbMiwxLCJcXEhvbV97XFxBY2FsfVxcbGVmdChQLE5cXG90aW1lc197XFxBY2FsfU1fe3EnfVxccmlnaHQpIl0sWzEsMF0sWzAsM10sWzEsMl0sWzIsM10sWzAsMiwiXFxleGlzdHMiLDEseyJzdHlsZSI6eyJib2R5Ijp7Im5hbWUiOiJkYXNoZWQifX19XV0=
    \begin{tikzcd}
        {\Hom_{\Acal}\left(A,\intHom_{\Acal}\left(P,N\right)\otimes_{\Acal}M_{q}\right)} && {\Hom_{\Acal}\left(A,\intHom_{\Acal}\left(P,N\right)\otimes_{\Acal}M_{q'}\right)} \\
        {\Hom_{\Acal}\left(P,N\otimes_{\Acal}M_{q}\right)} && {\Hom_{\Acal}\left(P,N\otimes_{\Acal}M_{q'}\right)}
        \arrow[from=1-1, to=1-3]
        \arrow[from=1-1, to=2-1]
        \arrow[from=1-3, to=2-3]
        \arrow["\exists"{description}, dashed, from=2-1, to=1-3]
        \arrow[from=2-1, to=2-3]
    \end{tikzcd}$$ 
    by its construction from $\eta_{q,q'}$ which induces an equivalence on filtered colimits by \cite[Lem. 2.16]{AokiSchwartz}.

    (c)$\Rightarrow$(a) Taking $N=A$, we have 
    $$\Hom_{\Acal}(A,P^{\vee}\otimes_{\Acal}M)\simeq\Hom_{\Acal}(P,M).$$
    Let $p:A^{+}_{\sld}\to\Acal\simeq(A,A^{+})_{\sld}$ be the natural map of analytic $\EE_{\infty}$-rings, which is analytically proper. For $P^{\vee}:=\intHom_{\Acal}(p^{*}A^{+}_{\sld}[S],A)$ we have $p_{*}P^{\vee}\simeq \intHom_{A^{+}_{\sld}}(A^{+}_{\sld}[S],A|_{A^{+}})$ by \Cref{lem: modules inherit bistability} (2) which by $A^{+}_{\sld}$-analyticity of $A|_{A^{+}}$ and adjunction is equivalent to $\intHom_{A^{+}}(A^{+}[S],A|_{A^{+}})$. Since $A$ is discrete, $p_{*}P^{\vee}\simeq\colim_{i\in I}(A|_{A^{+}})^{\oplus |S_{i}|}$ for a presentation $S\simeq\lim_{i\in I}S_{i}$ as a sequential limit of finite sets. Since $p_{*}$ is conservative and preserves small colimits, it follows $P^{\vee}\simeq\colim_{i\in I}A^{\oplus|S_{i}|}$ and we have 
    \begin{align*}
        \Hom_{\Acal}\left(A,\colim_{i\in I}A^{\oplus|S_{i}|}\otimes_{\Acal}M\right) &\simeq \Hom_{\Acal}\left(A,\colim_{i\in I}M^{\oplus |S_{i}|}\right) \\
        &\simeq \colim_{i\in I}\Hom_{\Acal}\left(A,M^{\oplus |S_{i}|}\right) \\
        &\simeq \colim_{i\in I}\Hom_{\Acal}\left(A^{\oplus |S_{i}|},M\right)
    \end{align*}
    where we have used, respectively, that the tensor product preserves small colimits, that $A$ is compact in $\Mod_{\Acal}=\Mod_{(A,A^{+})_{\sld}}$, and that finite products and coproducts coincide by exactness. By the assumed equivalence and that $\Hom_{\Acal}(A^{\oplus |S_{i}|},M)\simeq M(S_{i})$, we get that $M(S)\simeq\colim_{i\in I}M(S_{i})$. This shows that the value of $M$ on a light profinite set is left Kan-extended from its values on finite sets, hence discrete (see \cite[Prop. 2.1.8 (ii)]{MannRigid6FF} and the surrounding discussion). 
\end{proof}
Since nuclearity is necessary for dualizability, the dualizable objects are discrete as well.
\begin{proposition}\label{prop: solid rings are Fredholm}
    If $(A,A^{+})_{\sld}$ is a discrete solid Huber ring, then $(A,\Mod_{(A,A^{+})_{\sld}})$ is a Fredholm analytic $\EE_{\infty}$-ring in the sense of \Cref{def: Fredholm rings}: any dualizable object is discrete. In particular, there is an equivalence of categories $\Dbl(\Mod_{(A,A^{+})_{\sld}})\simeq\Dbl(\Mod_{A}(\Sp))$. 
\end{proposition}
\begin{proof}
    From \Cref{prop: nuclear objects in solid Huber rings}, the nuclear and discrete objects coincide so using \Cref{prop: dualizable iff nuclear and compact}, an object of $\Mod_{(A,A^{+})_{\sld}}$ is dualizable if and only if it is discrete and compact. By \Cref{lem: comparison functor from rigidification} (3), the embedding of discrete objects identifies with the embedding of nuclear objects, hence induces an equivalence on dualizable objects. 
\end{proof}
\subsection{The solid affine line}\label{subsec: solid affine line}
We construct the solid affine line $\Sph\{T\}_{\sld}$ and factor the morphism $\Sph_{\sld}\to\Sph\{T\}_{\sld}$ into an analytically proper morphism followed by an open immersion of complete analytic $\EE_{\infty}$-rings. The open immersion is recognized through an idempotent algebra, which we construct first.
\begin{lemma}\label{lem: solid nullsequence}
    Let $\Sph[\NN_{\infty}]:=\cofib(\Sph[\{\infty\}]\to\Sph[\NN\cup\{\infty\}])$. 
    \begin{enumerate}
        \item $\Sph[\NN_{\infty}]$ admits the structure of an $\EE_{\infty}$-algebra in $\Cond\Sp$. Denote its solidification (as an $\EE_{\infty}$-algebra) by $\Sph[[q]]$. 
        \item Let $\Sph[[q]]\{T\}:=\Sph[[q]]\otimes_{\Sph_{\sld}}\Sph\{T\}\in\Mod_{(\Sph\{T\},\Sph)_{\sld}}$ be the base change of $\Sph[[q]]$ along the morphism of discrete solid Huber rings $\Sph_{\sld}=:(\Sph,\Sph)_{\sld}\to(\Sph\{T\},\Sph)_{\sld}$. Let $\Qcal:=\cofib_{(\Sph\{T\},\Sph)_{\sld}}(\Sph[[q]]\{T\}\xrightarrow{1-qT}\Sph[[q]]\{T\})$. Then there exists a morphism $e:\Sph\{T\}\to \Qcal$ which exhibits $\Qcal$ as an idempotent $\EE_{\infty}$-algebra of $\Mod_{(\Sph\{T\},\Sph)_{\sld}}$. 
    \end{enumerate} 
\end{lemma}
\begin{proof}[Proof of (1)]
    We use the formalism of Smith ideals (see \cite[\S 3.1]{MaoCrystalline}). The category $\Fun((\Delta^{1})^{\Opp},\Cond\Sp)$ carries the Day convolution symmetric monoidal structure induced by the maximum monoidal structure on $(\Delta^{1})^{\Opp}$, with unit $[0\to\Sph]$, and there is a symmetric monoidal equivalence with $\Fun(\Delta^{1},\Cond\Sp)$ under the pointwise structure carrying $[I\to R]$ to $[R\to\cofib(I\to R)]$ \cite[Prop. 3.1]{MaoCrystalline}. An $\EE_{\infty}$-algebra for the Day convolution structure thus determines a canonical $\EE_{\infty}$-algebra structure on the cofiber together with an $\EE_{\infty}$-algebra morphism $R\to\cofib(I\to R)$ \cite[Cor. 3.3]{MaoCrystalline}. 

    The addition on $\NN$ extends to the compactification by $n+\infty=\infty=\infty+n$, and this operation is continuous, showing the subset $\{\infty\}$ is a monoid ideal making $\{\infty\}\to\NN\cup\{\infty\}$ a commutative algebra object for the Day convolution structure on $\Cond\Ani$. 

    The functor $\Sph[-]:\Cond\Ani\to\Cond\Sp$ is symmetric monoidal and preserves small colimits, so it induces a symmetric monoidal functor on Day convolution categories \cite[\S 3]{StableOperads}, carrying the commutative algebra above to a Day convolution algebra $[\Sph[\{\infty\}]\to\Sph[\NN\cup\{\infty\}]]$ in $\Cond\Sp$. The preceding argument equips $\Sph[\NN_{\infty}]=\cofib(\Sph[\{\infty\}]\to\Sph[\NN\cup\{\infty\}])$ with the structure of an $\EE_{\infty}$-algebra in $\Cond\Sp$. Symmetric monoidality of solidification makes $\Sph[[q]]$ an $\EE_{\infty}$-algebra in $\Cond\Sp^{\sld}$. 
\end{proof}
\begin{proof}[Proof of (2)]
    The self-map of $\Sph[[q]]\{T\}$ given by multiplication by the element $1-qT$ is a map of $\Sph\{T\}$-modules, so we can consider the cofiber $\Qcal$ as an $\Sph\{T\}$-module. Let $e:\Sph\{T\}\to \Qcal$ be the composite of the structure map $\Sph\{T\}\to\Sph[[q]]\{T\}$ of $\Sph[[q]]\{T\}$ as an $\EE_{\infty}$-algebra in $\Mod_{(\Sph\{T\},\Sph)_{\sld}}$ and the canonical map to the cofiber. It suffices by \cite[Rmk. 4.8.2.6 \& 4.8.2.9]{HA} to show that 
    $$\Qcal\xrightarrow{\id_{\Qcal}\otimes_{(\Sph\{T\},\Sph)_{\sld}}e}\Qcal\otimes_{(\Sph\{T\},\Sph)_{\sld}}\Qcal$$ 
    is an equivalence, where all objects are connective by compatibility of the $t$-structure with the symmetric monoidal structure. Note that the functor $\ZZ[T]\otimes_{(\Sph\{T\},\Sph)_{\sld}}-:\Mod_{(\Sph\{T\},\Sph)_{\sld}}\to\Mod_{(\ZZ[T],\ZZ)_{\sld}}$ is conservative on bounded below objects by \Cref{lem: conservativity of bounded below} so the above morphism is an equivalence if and only if it is so after base change to $(\ZZ[T],\ZZ)_{\sld}$-modules. In this case, exactness of the functor identifies $\ZZ[T]\otimes_{(\Sph\{T\},\Sph)_{\sld}}\Qcal$ with $\ZZ((T^{-1})):=\cofib_{(\ZZ[T],\ZZ)_{\sld}}(\ZZ[[q]][T]\xrightarrow{1-qT}\ZZ[[q]][T])$ which is an idempotent algebra along the base change of $e$ by \cite[(5.2.1)]{CamargoSolid}, proving the claim.   
\end{proof}
With this idempotent algebra in hand, the morphism to the free algebra factors through the open immersion it defines.
\begin{proposition}\label{prop: explicit solid affine line}
    Let $f:\Sph_{\sld}\to\Sph\{T\}_{\sld}$ be the morphism of discrete solid $\EE_{\infty}$-rings to the free $\EE_{\infty}$-algebra. 
    \begin{enumerate}
        \item $f$ admits a factorization as 
        $$\Sph_{\sld}\xrightarrow{p}(\Sph\{T\},\Sph)_{\sld}\xrightarrow{j}\Sph\{T\}_{\sld}$$
        where $p$ is analytically proper. 
        \item There is an equivalence of categories $\ker(j^{*})\simeq\Mod_{\Qcal}(\Mod_{(\Sph\{T\},\Sph)_{\sld}})$. In particular, $j$ is an open immersion of analytic $\EE_{\infty}$-rings. 
    \end{enumerate}
\end{proposition}
\begin{proof}[Proof of (1)]
    The existence of the factorization follows from \Cref{prop: factorization of Huber ring maps} (2). 
\end{proof}
\begin{proof}[Proof of (2)]
    We need to verify that for $\Ical:=\fib_{(\Sph\{T\},\Sph)_{\sld}}(\Sph\{T\}\xrightarrow{e}\Qcal)$, the object $\intHom_{(\Sph\{T\},\Sph)_{\sld}}(\Ical,M)$ is in the essential image of the forgetful functor $j_{*}:\Mod_{\Sph\{T\}_{\sld}}\to\Mod_{(\Sph\{T\},\Sph)_{\sld}}$. By \Cref{cor: hom membership} this can be reduced to a check on homotopy, where it follows from \cite[Prop. 5.2.2]{CamargoSolid}. 
\end{proof}
\subsection{Exceptional functors for morphisms of finite expansion}\label{subsec: exceptional functors finite expansion}
In the animated setting \cite[\S 2.7]{HansenMannCLL}, the morphisms of finite expansion of Hamacher \cite{FiniteExpansion} form the exceptional class of the solid six-functor formalism. They adapt to connective $\EE_{\infty}$-rings as follows.
\begin{definition}\label{def: finite expansion Einfty}
    Let $f:A\to B$ be a morphism in $\CAlg^{\cn}$. The morphism $f$ is \emph{of finite expansion} if it admits a factorization as $A\to A\{T_{1},\dots,T_{n}\}\to B$ such that $\pi_{0}(A\{T_{1},\dots,T_{n}\})\to\pi_{0}(B)$ is integral. 
\end{definition}
\begin{remark}\label{rmk: finite expansion Einfty is static finite expansion}
    By \Cref{lem: finite expansion static detection} this is equivalent to $\pi_{0}(f):\pi_{0}(A)\to\pi_{0}(B)$ admitting a factorization as $\pi_{0}(A)\to\pi_{0}(A)[T_{1},\dots,T_{n}]\to\pi_{0}(B)$ where $\pi_{0}(A)[T_{1},\dots,T_{n}]\to\pi_{0}(B)$ is an integral morphism. 
\end{remark}
The exceptional functors of a finite-expansion morphism are assembled from those of the two factors: the free-algebra morphism $A_{\sld}\to A\{T_{1},\dots,T_{n}\}_{\sld}$, treated below, and the morphism $A\{T_{1},\dots,T_{n}\}_{\sld}\to B_{\sld}$, which has integral $\pi_{0}$ and is analytically proper by the following lemma.
\begin{lemma}\label{lem: integral implies analytically proper}
    Let $A\to B$ be a morphism in $\CAlg^{\cn}$ such that $\pi_{0}(A)\to\pi_{0}(B)$ is integral. Then the morphism of discrete solid $\EE_{\infty}$-rings $A_{\sld}\to B_{\sld}$ is analytically proper. 
\end{lemma}
\begin{proof}
    We want to show that a module $M$ is $B_{\sld}$-analytic if and only if its underlying $A$-module is $A_{\sld}$-analytic. By \Cref{prop: t-structure on A-analytic modules} (vi) we can reduce this to the case where $A,B,M$ are static, where it is \cite[Prop. 3.22]{Andreychev}. 
\end{proof}
It remains to construct the exceptional functors of the free-algebra morphisms, beginning with the solid affine line and the factorization of \Cref{prop: explicit solid affine line}. That $j$ is an open immersion was shown there through the idempotent algebra $\Qcal$. We reprove it below by the recognition criterion of \Cref{prop: open immersion detected on static rings}, exhibiting the requisite finite analytic Tor-amplitude directly.
\begin{proposition}\label{prop: exceptional functors for the affine line}
    Let $f:\Sph_{\sld}\to\Sph\{T\}_{\sld}$ be the morphism of discrete solid $\EE_{\infty}$-rings to the free $\EE_{\infty}$-algebra.
    \begin{enumerate}
        \item $f$ admits a factorization $\Sph_{\sld}\xrightarrow{p}(\Sph\{T\},\Sph)_{\sld}\xrightarrow{j}\Sph\{T\}_{\sld}$ in which $p$ is analytically proper and $j$ is an open immersion, so that $j^{*}$ admits a left adjoint $j_{!}$ with $j_{!}j^{*}M\simeq M\otimes_{(\Sph\{T\},\Sph)_{\sld}}j_{!}\Sph\{T\}$.
        \item The functor $f_{!}:=p_{*}j_{!}:\Mod_{\Sph\{T\}_{\sld}}\to\Mod_{\Sph_{\sld}}$ preserves small colimits and satisfies the projection formula: the natural morphism $f_{!}(f^{*}M\otimes_{\Sph\{T\}_{\sld}}N)\to M\otimes_{\Sph_{\sld}}f_{!}N$ is an equivalence in $\Mod_{\Sph_{\sld}}$.
    \end{enumerate}
\end{proposition}
\begin{proof}[Proof of (1)]
    That $p$ is analytically proper and $j$ an open immersion is \Cref{prop: explicit solid affine line}, and $j_{!}$ with the stated formula is then part of the open-immersion data of \Cref{def: open immersion}.

    We give a second proof that $j$ is an open immersion by verifying the hypotheses of \Cref{prop: open immersion detected on static rings} for $j:(\Sph\{T\},\Sph)_{\sld}\to\Sph\{T\}_{\sld}$. Both rings have underlying $\EE_{\infty}$-ring $\Sph\{T\}$, with $\pi_{0}(\Sph\{T\})\simeq\ZZ[T]$. The base-change condition holds since $\ZZ[T]\simeq\pi_{0}(\Sph\{T\})$ already lies in $\Mod_{\Sph\{T\}_{\sld}}$, so that $\Sph\{T\}_{\sld}\otimes_{(\Sph\{T\},\Sph)_{\sld}}\ZZ[T]\simeq\ZZ[T]$. The restriction $(-)|_{\Sph\{T\}}:\Mod_{\Sph\{T\}_{\sld}}\to\Mod_{(\Sph\{T\},\Sph)_{\sld}}$ is fully faithful, both categories being full subcategories of $\Mod_{\Sph\{T\}}$, which is (i). The $t$-structure on $\Mod_{(\Sph\{T\},\Sph)_{\sld}}$ is compatible with small products by \Cref{prop: product compatibility rings under solid sphere}, as $(\Sph\{T\},\Sph)_{\sld}$ receives the morphism $p$ from $\Sph_{\sld}$. Finally, the induced scalar extension on static rings is
    $$\pi_{0}(\Sph\{T\})_{\sld}\otimes_{(\pi_{0}(\Sph\{T\}),\pi_{0}(\Sph))_{\sld}}(-)\simeq\ZZ[T]_{\sld}\otimes_{(\ZZ[T],\ZZ)_{\sld}}(-),$$
    which is of analytic $\Tor$-amplitude $\leq 1$ by the proof of \cite[Prop. 2.9.36]{MannRigid6FF}, giving (ii), and an open immersion by \cite[Prop. 5.2.2]{CamargoSolid}, completing (iii). Thus $j$ is an open immersion.
\end{proof}
\begin{proof}[Proof of (2)]
    This follows from the six-functor formalism for analytic $\EE_{\infty}$-rings. Explicitly, $p_{*}$ preserves small colimits as the forgetful functor of an analytically proper map, and $j_{!}$ does so as a left adjoint, so $f_{!}$ preserves colimits as well. The projection formula is obtained by composing those for open immersions \Cref{def: open immersion} (a) and analytically proper morphisms \Cref{prop: analytically proper base change properties}. 
\end{proof}
By base change and the six-functor formalism for analytic $\EE_{\infty}$-rings established in \Cref{subsec: six operations on analytic rings}, we extend this to morphisms of finite expansion. 
\begin{theorem}\label{thm: exceptional functor for finite expansion}
    Let $f:A_{\sld}\to B_{\sld}$ be a morphism of discrete solid $\EE_{\infty}$-rings of finite expansion. Then there exists a small colimit-preserving functor $f_{!}:\Mod_{B_{\sld}}\to\Mod_{A_{\sld}}$ satisfying the projection formula and which admits a right adjoint $f^{!}:\Mod_{A_{\sld}}\to\Mod_{B_{\sld}}$. 
\end{theorem}
\begin{proof}
    Since $f$ is of finite expansion, it admits a factorization as $$A_{\sld}\xrightarrow{g}A\{T_{1},\dots,T_{n}\}_{\sld}\xrightarrow{h} B_{\sld}$$ where $h: A\{T_{1},\dots,T_{n}\}_{\sld}\to B_{\sld}$ is analytically proper by \Cref{lem: integral implies analytically proper}. 

    We first treat the case of $g$. The functor to discrete solid $\EE_{\infty}$-rings preserves nonempty small colimits by \Cref{lem: functoriality of solid rings}, and we implicitly identify base change of discrete solid and analytic $\EE_{\infty}$-rings in what follows. The morphism $\Sph_{\sld}\to\Sph\{T_{1},\dots,T_{n}\}_{\sld}$ is the iterated base change of $\Sph_{\sld}\to\Sph\{T\}_{\sld}$ hence admits exceptional functors as it lies in the homotopy class $E_{!}$ of \Cref{lem: geometric setup on analytic E-infinity rings}. The morphism $g$ is the base change of $\Sph_{\sld}\to\Sph\{T_{1},\dots,T_{n}\}_{\sld}$ along $\Sph_{\sld}\to A_{\sld}$ hence in $E_{!}$. Since $h$ is analytically proper, it lies in $E_{!}$. By closure under composition, $f$ lies in $E_{!}$ as well. The projection formula and existence of a right adjoint are formal from this being a presentable six-functor formalism as in \Cref{prop: six-functor formalism on analytic E-infinity rings}. 
\end{proof}
\subsection{Flat morphisms of discrete solid \texorpdfstring{$\EE_{\infty}$}{E-infinity}-rings}\label{subsec: flat morphisms of solid rings}
We now study the morphisms of discrete solid $\EE_{\infty}$-rings that are flat and almost of finite presentation in the sense of \Cref{def: finiteness conditions of morphisms}, following Mikami \cite[\S 4]{MikamiFPPF} and \cite[\S 2]{MikamiFFDescent}. 

This structure theory underlies the descent of \Cref{subsec: descent for solid rings} and enters the duality theory of the six-functor formalism for spectral algebraic stacks. For a morphism of discrete solid $\EE_{\infty}$-rings $A_{\sld}\to B_{\sld}$, write $\pi_{0}(A)_{\sld}\to\pi_{0}(B)_{\sld}$ for the induced morphism of discrete solid static rings. We begin with the static case.
\begin{proposition}[{\cite[Cor. 2.11 \& Prop. 2.20]{MikamiFFDescent}}]\label{prop: flat static representability}
    Let $\pi_{0}(A)_{\sld}\to \pi_{0}(B)_{\sld}$ be a finitely presented flat morphism of discrete solid static rings. 
    \begin{enumerate}
        \item The functor $\pi_{0}(B)_{\sld}\otimes_{\pi_{0}(A)_{\sld}}-:\Mod_{\pi_{0}(A)_{\sld}}\to\Mod_{\pi_{0}(B)_{\sld}}$ is of $\Tor$-amplitude $\leq 1$ and preserves small limits. 
        \item There exists a compact object $N_{\pi_{0}(B)/\pi_{0}(A)}^{\sld}\in\Mod_{\pi_{0}(A)_{\sld}}^{\omega}$ which corepresents the functor 
        \begin{equation}\label{eqn: static representability}
            \Mod_{\pi_{0}(A)_{\sld}}\xrightarrow{\pi_{0}(B)_{\sld}\otimes_{\pi_{0}(A)_{\sld}}-}\Mod_{\pi_{0}(B)_{\sld}}\xrightarrow{\Mor_{\pi_{0}(B)_{\sld}}(\pi_{0}(B),-)}\Ani
        \end{equation}
        and an equivalence of functors 
        $$(\pi_{0}(B)_{\sld}\otimes_{\pi_{0}(A)_{\sld}}-)|_{\pi_{0}(A)}\simeq\intHom_{\pi_{0}(A)_{\sld}}(N_{\pi_{0}(B)/\pi_{0}(A)}^{\sld},-)$$ 
        in the enriched endofunctor category $\End_{\pi_{0}(A)_{\sld}}$. 
    \end{enumerate}
\end{proposition}
This lifts to discrete solid $\EE_{\infty}$-rings.
\begin{proposition}[Cf. {\cite[\S 4]{MikamiFPPF}}]\label{prop: flat spectral representability}
    Let $A_{\sld}\to B_{\sld}$ be a morphism of discrete solid $\EE_{\infty}$-rings that is flat and almost of finite presentation.
    \begin{enumerate}
        \item The scalar extension functor $B_{\sld}\otimes_{A_{\sld}}-:\Mod_{A_{\sld}}\to\Mod_{B_{\sld}}$ is of analytic $\Tor$-amplitude $\leq 1$ and preserves small limits.
        \item There exists a compact object $N_{B/A}^{\sld}\in\Mod_{A_{\sld}}^{\omega}$ which corepresents the functor 
        \begin{equation}\label{eqn: spectral representability}
            \Mod_{A_{\sld}}\xrightarrow{B_{\sld}\otimes_{A_{\sld}}-}\Mod_{B_{\sld}}\xrightarrow{\Mor_{B_{\sld}}(B,-)}\Ani
        \end{equation}
        satisfying $N_{B/A}^{\sld}\otimes_{A_{\sld}}\pi_{0}(A)\simeq N_{\pi_{0}(B)/\pi_{0}(A)}^{\sld}$ and an equivalence of functors 
        $$(B_{\sld}\otimes_{A_{\sld}}-)|_{A}\simeq\intHom_{A_{\sld}}(N_{B/A}^{\sld},-)$$
        in the enriched endofunctor category $\End_{A_{\sld}}$. 
    \end{enumerate}
\end{proposition}
\begin{proof}[Proof of (1)]
    We show the flat morphism $A_{\sld}\to B_{\sld}$ satisfies the hypotheses of \Cref{lem: tor amplitude and limits} (2). Flatness gives $B\otimes_{A}\pi_{0}(A)\simeq\pi_{0}(B)$. The $t$-structure on $\Mod_{A_{\sld}}$ is compatible with small products by \Cref{prop: product compatibility rings under solid sphere}, as $A_{\sld}$ receives a morphism from $\Sph_{\sld}$. Finally, $\pi_{0}(B)_{\sld}\otimes_{\pi_{0}(A)_{\sld}}-$ is of analytic $\Tor$-amplitude $\leq 1$ and preserves small limits by \Cref{prop: flat static representability} (1). Hence $B_{\sld}\otimes_{A_{\sld}}-$ is of analytic $\Tor$-amplitude $\leq 1$ and preserves small limits.
\end{proof}
\begin{proof}[Proof of (2)]
    By (1), $B_{\sld}\otimes_{A_{\sld}}-$ preserves small limits and colimits, and $\Mod_{B_{\sld}}$ is $\omega$-presentable with compact unit, so $\Mor_{B_{\sld}}(B,-)$ preserves small limits and filtered colimits. This implies that the composite (\ref{eqn: spectral representability}) preserves small limits and filtered colimits. 

    By \cite[Prop. 5.5.2.7]{HTT}, such a functor is corepresentable by an object of $\Mod_{A_{\sld}}$ and it is compact as the functor furthermore preserves filtered colimits. Denote this object by $N_{B/A}^{\sld}$. By adjunction, $\Mor_{\pi_{0}(A)_{\sld}}(N_{B/A}^{\sld}\otimes_{A_{\sld}}\pi_{0}(A),-)$ corepresents the composite 
    $$\Mod_{\pi_{0}(A)_{\sld}}\xrightarrow{(-)|_{A}}\Mod_{A_{\sld}}\xrightarrow{\Mor_{A_{\sld}}(N_{B/A}^{\sld},-)}\Ani$$
    but this functor identifies with (\ref{eqn: static representability}) of \Cref{prop: flat static representability} (2) and we conclude their identification $N_{B/A}^{\sld}\otimes_{A_{\sld}}\pi_{0}(A)\simeq N_{\pi_{0}(B)/\pi_{0}(A)}^{\sld}$ by the Yoneda lemma. 

    We have equivalences of animae 
    \begin{align*}
        \Mor_{A_{\sld}}(N_{B/A}^{\sld},N_{B/A}^{\sld}) &\simeq\Mor_{B_{\sld}}(B,B_{\sld}\otimes_{A_{\sld}}N_{B/A}^{\sld}) && \text{representability}\\
        &\simeq \Mor_{A_{\sld}}\left(A,(B_{\sld}\otimes_{A_{\sld}}N_{B/A}^{\sld})|_{A}\right) && B_{\sld}\otimes_{A_{\sld}}-\dashv(-)|_{A}.
    \end{align*}
    Let $f:A\to (B_{\sld}\otimes_{A_{\sld}}N_{B/A}^{\sld})|_{A}$ be the image of $\id_{N_{B/A}^{\sld}}:N_{B/A}^{\sld}\to N_{B/A}^{\sld}$ under this chain of equivalences. Observing that $(B_{\sld}\otimes_{A_{\sld}}-)|_{A}\in\End_{A_{\sld}}$, \cite[Lem. A.4.8]{MannRigid6FF} implies that there is an equivalence of animae 
    $$\Mor_{\End_{A_{\sld}}}\left(\intHom_{A_{\sld}}(N_{B/A}^{\sld},-),(B_{\sld}\otimes_{A_{\sld}}-)|_{A}\right)\simeq\Mor_{A_{\sld}}\left(A,(B_{\sld}\otimes_{A_{\sld}}N_{B/A}^{\sld})|_{A}\right).$$
    Let $\eta:\intHom_{A_{\sld}}(N_{B/A}^{\sld},-)\to (B_{\sld}\otimes_{A_{\sld}}-)|_{A}$ be the natural transformation of $\Mod_{A_{\sld}}$-enriched endofunctors corresponding to $f:A\to (B_{\sld}\otimes_{A_{\sld}}N_{B/A}^{\sld})|_{A}$ under this equivalence. The functor $(B_{\sld}\otimes_{A_{\sld}}-)|_{A}$ preserves small limits and colimits as $B_{\sld}\otimes_{A_{\sld}}-$ preserves both by (1) and forgetful functors from module categories $(-)|_{A}$ generally do so. Similarly $\intHom_{A_{\sld}}(N_{B/A}^{\sld},-)$ preserves limits as a right adjoint and preserves colimits by compactness, using that compactness and internal compactness coincide in $\Mod_{A_{\sld}}$ by \Cref{lem: compact implies internally compact}. We check that $\eta$ is an equivalence pointwise. Using that both functors are exact and Postnikov and Whitehead towers in $\Mod_{A_{\sld}}$ we can reduce to checking that $\eta_{M}:\intHom_{A_{\sld}}(N_{B/A}^{\sld},M)\to(B_{\sld}\otimes_{A_{\sld}}M)|_{A}$ is an equivalence for $M\in\Mod_{A_{\sld}}^{\heartsuit}$. By adjunction $\pi_{0}(A)_{\sld}\otimes_{A_{\sld}}-\dashv(-)|_{A}$ and \Cref{prop: relationship with heart} (1), $\intHom_{A_{\sld}}(N_{B/A}^{\sld},M)\simeq\intHom_{\pi_{0}(A)_{\sld}}(N_{\pi_{0}(B)/\pi_{0}(A)}^{\sld},M)$ and $(B_{\sld}\otimes_{A_{\sld}}M)|_{A}\simeq(\pi_{0}(B)_{\sld}\otimes_{\pi_{0}(A)_{\sld}}M)|_{A}$ which in conjunction with conservativity of the forgetful functor $\Mod_{\pi_{0}(A)_{\sld}}\to\Mod_{A_{\sld}}$ is precisely \Cref{prop: flat static representability} (2).  
\end{proof}
\subsection{Descent for discrete solid \texorpdfstring{$\EE_{\infty}$}{E-infinity}-rings}\label{subsec: descent for solid rings}
Solid modules satisfy descent for the faithfully flat and finitely presented topology, following Mikami \cite[\S 4]{MikamiFPPF}. The argument rests on the structure theory of \Cref{subsec: flat morphisms of solid rings} together with the steadiness of morphisms of discrete solid $\EE_{\infty}$-rings, which we record first.
\begin{lemma}\label{lem: any finite expansion map of solid rings is steady}
    Any morphism of discrete solid $\EE_{\infty}$-rings $A_{\sld}\to B_{\sld}$ of finite expansion is steady as a morphism of complete analytic $\EE_{\infty}$-rings. 
\end{lemma}
\begin{proof}
    We argue as in \Cref{thm: exceptional functor for finite expansion}. The free algebra morphism $\Sph_{\sld}\to\Sph\{T\}_{\sld}$ factors as $\Sph_{\sld}\to(\Sph\{T\},\Sph)_{\sld}\to\Sph\{T\}_{\sld}$ where $\Sph_{\sld}\to(\Sph\{T\},\Sph)_{\sld}$ is analytically proper. Using that $\Sph\{T\}$ is discrete hence nuclear by \Cref{cor: discrete implies nuclear}, \Cref{prop: nuclear implies steady} implies $\Sph_{\sld}\to(\Sph\{T\},\Sph)_{\sld}$ is steady. The morphism $(\Sph\{T\},\Sph)_{\sld}\to\Sph\{T\}_{\sld}$ is an open immersion by \Cref{prop: exceptional functors for the affine line} hence steady by \Cref{cor: open immersions are steady}, showing $\Sph_{\sld}\to\Sph\{T\}_{\sld}$ is steady. Using the permanence properties of steady morphisms established in \Cref{prop: permanence properties of steady morphisms}, it follows that $A_{\sld}\to A\{T_{1},\dots,T_{n}\}_{\sld}$ is steady, being obtained by iterated base change from $\Sph_{\sld}\to\Sph\{T\}_{\sld}$. Similarly for $A\{T_{1},\dots,T_{n}\}_{\sld}\to B_{\sld}$, this is analytically proper by \Cref{lem: integral implies analytically proper} with $B$ discrete and hence nuclear in $A\{T_{1},\dots,T_{n}\}_{\sld}$-modules by \Cref{cor: discrete implies nuclear} so it is steady by \Cref{prop: nuclear implies steady}. Since steady morphisms are preserved under composition, $A_{\sld}\to B_{\sld}$ is steady. 
\end{proof}
\begin{remark}\label{rmk: steadiness not optimal}
    L. Mann shows that any morphism of discrete solid animated Huber rings is steady \cite[Prop. 2.9.7]{MannRigid6FF}, so \Cref{lem: any finite expansion map of solid rings is steady} and \Cref{cor: fppf is steady} are not optimal. The restricted statements suffice for the descent of \Cref{thm: solid fppf descent}, and we do not reprove the spectral analogue of Mann's result here (cf. \Cref{rmk: Huber to analytic ring preserves colimits}). 
\end{remark}
\begin{corollary}\label{cor: fppf is steady}
    Any faithfully flat and finitely presented morphism of discrete solid $\EE_{\infty}$-rings $A_{\sld}\to B_{\sld}$ is steady as a morphism of complete analytic $\EE_{\infty}$-rings. 
\end{corollary}
\begin{proof}
    By \Cref{lem: finiteness implies finite expansion}, any finitely presented morphism, and in particular any faithfully flat finitely presented morphism is of finite expansion. Thus it is steady by \Cref{lem: any finite expansion map of solid rings is steady}. 
\end{proof}
The finitely presented faithfully flat morphisms are descendable. 
\begin{proposition}\label{prop: solid descendable realizations}
    Let $A_{\sld}\to B_{\sld}$ be a finitely presented and faithfully flat morphism of discrete solid $\EE_{\infty}$-rings.
    \begin{enumerate}
        \item For all $r\geq 1$, $(N_{B/A}^{\sld})^{\otimes r}\in\Mod_{A_{\sld},\geq -1}$. 
        \item The natural map $|(N_{B/A}^{\sld})^{\otimes \bullet}|\to A$ is an equivalence.
        \item The map $A_{\sld}\to B_{\sld}$ is descendable. 
    \end{enumerate}
\end{proposition}
\begin{proof}[Proof of (1)]
    Note that by repeating the constructions and arguments of \Cref{prop: flat spectral representability}, $(N_{B/A}^{\sld})^{\otimes r}\simeq N_{B^{r}/A}^{\sld}$ where 
    $$B^{r}:=\underbrace{B_{\sld}\otimes_{A_{\sld}}\dots\otimes_{A_{\sld}}B_{\sld}}_{r\text{-times}}$$
    so it suffices to treat the case of $r=1$. Suppose there is $k>1$ such that $N_{B/A}^{\sld}\to\tau_{\leq -k}N_{B/A}^{\sld}$ is nonzero. Using the equivalence of functors in \Cref{prop: flat spectral representability} (2), we get $\intHom_{A_{\sld}}(N_{B/A}^{\sld},\tau_{\leq -k}N_{B/A}^{\sld})\simeq (B_{\sld}\otimes_{A_{\sld}}\tau_{\leq -k}N_{B/A}^{\sld})|_{A}$ which is nonvanishing on $\pi_{0}$ by the existence of the nontrivial map $N_{B/A}^{\sld}\to\tau_{\leq -k}N_{B/A}^{\sld}$. But this is a contradiction as $B_{\sld}\otimes_{A_{\sld}}-$ is of $\Tor$-amplitude $\leq 1$. 
\end{proof}
\begin{proof}[Proof of (2)]
    Observe that the unit of the extension-restriction adjunction furnishes a natural map $\id_{\Mod_{A_{\sld}}}\to\intHom_{A_{\sld}}(N_{B/A}^{\sld},-)$ which under the contravariant fully faithful embedding $\Mod_{A_{\sld}}^{\Opp}\to\End_{A_{\sld}}$ given on objects by $P\mapsto\intHom_{A_{\sld}}(P, -)$ gives a map $N_{B/A}^{\sld}\to A$ \cite[Cor. A.4.9]{HeyerMann6FF}. From (1), the realization $|(N_{B/A}^{\sld})^{\otimes \bullet}|$ is bounded below so the natural morphism $|(N_{B/A}^{\sld})^{\otimes \bullet}|\to A$ is one of bounded below objects. Applying \Cref{lem: conservativity of bounded below}, it suffices to demonstrate the equivalence after applying $\pi_{0}(A)\otimes_{A_{\sld}}-$, namely that 
    $$|(N_{\pi_{0}(B)/\pi_{0}(A)}^{\sld})^{\otimes\bullet}|\longrightarrow\pi_{0}(A)$$
    is an equivalence, which is the content of \cite[Thm. 2.15]{MikamiFFDescent}. 
\end{proof}
\begin{proof}[Proof of (3)]
    Let $(N_{B/A}^{\sld})^{r}:=\colim_{[m]\in\Delta^{\Opp}_{s,\leq r}}(N_{B/A}^{\sld})^{\otimes (m+1)}$ be the partial geometric realizations and $C_{r}:=\cofib((N_{B/A}^{\sld})^{r}\to A)$ where both terms lie in $\Mod_{A_{\sld},\geq-1}$ which is closed under colimits computed in $\Mod_{A_{\sld}}$. The partial geometric realizations assemble into (co)fiber sequences $(N_{B/A}^{\sld})^{r}\to (N_{B/A}^{\sld})^{r+1}\to (N_{B/A}^{\sld})^{\otimes (r+2)}[r+1]$ \cite[Lem. 4.12]{MikamiFFDescent}. By (1), $(N^{\sld}_{B/A})^{\otimes (r+2)}\in\Mod_{A_{\sld},\geq-1}$ so $(N^{\sld}_{B/A})^{\otimes(r+2)}[r+1]\in\Mod_{A_{\sld},\geq r}$. Noting that finite colimits of uniformly bounded below objects are uniformly bounded below, we have that $(N_{B/A}^{\sld})^{r},(N_{B/A}^{\sld})^{r+1}\in\Mod_{A_{\sld},\geq-1}$. Examining the long exact sequence on homotopy, we obtain $\pi_{k}((N_{B/A}^{\sld})^{r})\to\pi_{k}((N_{B/A}^{\sld})^{r+1})$ is an equivalence for $k\leq r-2$ and surjective for $k=r-1$. Writing $|(N_{B/A}^{\sld})^{\otimes\bullet}|$ as a sequential colimit of partial geometric realizations, we obtain from (2) an equivalence $\colim_{r}\tau_{\leq k}(N_{B/A}^{\sld})^{r}\simeq\tau_{\leq k}\colim_{r}(N_{B/A}^{\sld})^{r}\simeq\tau_{\leq k}A$ for all $k$. 
    
    For $r=1$ the transition maps $\pi_{-1}((N_{B/A}^{\sld})^{1})\to\pi_{-1}((N_{B/A}^{\sld})^{t})$ and $\pi_{0}((N_{B/A}^{\sld})^{1})\to\pi_{0}((N_{B/A}^{\sld})^{t})$ are, respectively, an equivalence and a surjection for $t\geq 2$. Since the sequential colimit of the partial realizations is equivalent to $A$, it follows that $\pi_{-1}((N_{B/A}^{\sld})^{1})\to\pi_{-1}(A)\simeq0$ is an equivalence and $\pi_{0}((N_{B/A}^{\sld})^{1})\to\pi_{0}(A)$ is surjective. Examining the long exact sequence on homotopy, we conclude $\pi_{0}(C_{1})\simeq0$ and thus $C_{1}\in\Mod_{A_{\sld},\geq1}$. Consider the commutative diagram 
    $$% https://q.uiver.app/#q=WzAsNixbMCwwLCIoTl97Qi9BfV57XFxzbGR9KV57MX0iXSxbMiwwLCJBIl0sWzQsMCwiQ197MX0iXSxbNCwxLCJDX3swfSJdLFsyLDEsIkEiXSxbMCwxLCIoTl97Qi9BfV57XFxzbGR9KV57MH09Tl97Qi9BfV57XFxzbGR9Il0sWzAsMV0sWzEsMl0sWzMsMl0sWzQsM10sWzQsMSwiXFx3clxcaWRfe0F9IiwyXSxbNSw0XSxbNSwwXV0=
    \begin{tikzcd}
        {(N_{B/A}^{\sld})^{1}} && A && {C_{1}} \\
        {(N_{B/A}^{\sld})^{0}=N_{B/A}^{\sld}} && A && {C_{0}}
        \arrow[from=1-1, to=1-3]
        \arrow[from=1-3, to=1-5]
        \arrow[from=2-1, to=1-1]
        \arrow[from=2-1, to=2-3]
        \arrow["{\wr\id_{A}}"', from=2-3, to=1-3]
        \arrow[from=2-3, to=2-5]
        \arrow[from=2-5, to=1-5]
    \end{tikzcd}$$
    with horizontal (co)fiber sequences. Write $(N_{B/A}^{\sld})^{1}\xrightarrow{u} A\xrightarrow{v}C_{1}$ for the morphisms in the (co)fiber sequence. $\pi_{0}(C_{1})\simeq0$ implies $v$ is zero. Applying $\Hom_{A_{\sld}}(A,-)$ to the top cofiber sequence, we obtain a (co)fiber sequence $\Hom_{A_{\sld}}(A,(N_{B/A}^{\sld})^{1})\to\Hom_{A_{\sld}}(A,A)\to\Hom_{A_{\sld}}(A,C_{1})$ in $\Sp$ where the final term has vanishing $\pi_{0}$ by internal compact projectivity of the unit established in \Cref{prop: A-analytic tensor product}. Thus $\id_{A}$ lifts to a morphism $h:A\to (N_{B/A}^{\sld})^{1}$ so that $u\circ h\simeq\id_{A}$ exhibiting $A$ as a retract of $(N_{B/A}^{\sld})^{1}$. 

    Let $D$ be the minimal full subcategory of $\Mod_{A_{\sld}}$ containing $N_{B/A}^{\sld}$ and closed under tensor products, retracts, and shifts. Let $D'$ be the minimal full subcategory of $\Mod_{A_{\sld}}$ containing the fibers of morphisms in $D$ and closed under tensor products, retracts, and shifts. Rotating the (co)fiber sequence $N_{B/A}^{\sld}\to(N_{B/A}^{\sld})^{1}\to(N_{B/A}^{\sld})^{\otimes2}[1]$ gives a (co)fiber sequence $(N_{B/A}^{\sld})^{1}\to(N_{B/A}^{\sld})^{\otimes 2}[1]\to N_{B/A}^{\sld}[1]$. The last two objects belong to $D$, so $(N_{B/A}^{\sld})^{1}\in D'$. Since $D'$ is closed under retracts and $A$ is a retract of $(N_{B/A}^{\sld})^{1}$, we obtain $A\in D'$. In particular, under the exact, monoidal, fully faithful embedding of \cite[Rmk. 2.5.15 \& Cor. A.4.9]{MannRigid6FF}, the identity endofunctor is contained in the smallest full subcategory of $\End_{A_{\sld}}$ containing $(B_{\sld}\otimes_{A_{\sld}}-)|_{A}$ closed under finite limits, finite colimits, retracts, and composition. Together with \Cref{cor: fppf is steady}, this shows that $A_{\sld}\to B_{\sld}$ is descendable by \Cref{lem: descendable equivalent}.  
\end{proof}
Descendability and steadiness yield descent.
\begin{theorem}\label{thm: solid fppf descent}
    Let $A_{\sld}\to B_{\sld}$ be a finitely presented and faithfully flat morphism of discrete solid $\EE_{\infty}$-rings. Then the functor 
    $$\Mod_{A_{\sld}}\longrightarrow\lim_{[n]\in\Delta}\Mod_{B_{\tsld}^{\otimes_{A_{\tsld}}(n+1)}}$$
    is an equivalence.
\end{theorem}
\begin{proof}
    This follows from \Cref{thm: descent for descendable morphisms} in light of \Cref{prop: solid descendable realizations} (3). 
\end{proof}
\begin{remark}\label{rmk: descendability index at most two}
    The proof of \Cref{prop: solid descendable realizations} (3) exhibits $A$ as a retract of $(N_{B/A}^{\sld})^{1}$, so a finitely presented faithfully flat morphism of discrete solid $\EE_{\infty}$-rings is descendable of index at most $2$. The bound rests on \cite[Lem. 4.5]{MikamiFPPF}, which is stated incorrectly there and has yet to be corrected, so we do not prove it.
\end{remark}
\section{Six operations for solid quasicoherent sheaves}\label{sec: six functors for solid sheaves}
With the analytic and solid theories of the preceding sections in place, we construct the six-functor formalism for solid quasicoherent sheaves that motivates this work. On affine spectral schemes it is obtained by restricting the six-functor formalism on complete analytic $\EE_{\infty}$-rings of \Cref{subsec: six operations on analytic rings} along the embedding of discrete solid $\EE_{\infty}$-rings. Descent along the faithfully flat and finitely presented topology globalizes these to spectral algebraic stacks through the machinery of Heyer--L. Mann \cite[\S 3.4]{HeyerMann6FF}. Together with the duality it supports, the construction recovers the formalism announced in the introduction.

One point of divergence from the animated theory governs the treatment of suaveness. There a morphism of solid affine schemes is suave exactly when it is almost of finite presentation and of finite $\Tor$-amplitude, the sufficiency resting on a factorization $A\to A[T_{1},\dots,T_{n}]\to B$ with $B$ perfect over the flat algebra $A[T_{1},\dots,T_{n}]$ \cite[Lem. 2.7.12]{HansenMannCLL}. Over the sphere this fails, the free $\EE_{\infty}$-algebra $\Sph\{T\}\simeq\bigoplus_{n\geq0}\Sigma^{\infty}_{+}B\Sigma_{n}$ being of infinite $\Tor$-amplitude. We impose the flat factorization as a hypothesis instead, working with the pseudo-lisse morphisms of \Cref{subsec: pseudo-lisse morphisms} and proving that these are suave. It is not clear that every morphism almost of finite presentation and of finite $\Tor$-amplitude is pseudo-lisse, or even that it lies in the closure of the pseudo-lisse morphisms under composition and base change.
\subsection{Six operations on affine spectral schemes}\label{subsec: six functors on affines}
Throughout this subsection we fix an uncountable regular cardinal $\kappa$ and restrict to $\kappa$-small affine spectral schemes, sidestepping the set-theoretic difficulties attached to sheaves on large sites. Much of the following discussion holds without the smallness restriction, but the $\kappa$-small setting suffices for our use and eschews the need to develop a theory of macrotopoi (cf. \cite[App. A]{DiracII}).
\begin{definition}\label{def: kappa small rings}
    For an uncountable regular cardinal $\kappa$, let the \emph{category of $\kappa$-small affine spectral schemes} $\Aff_{\kappa}$ be the opposite category of $\An\CAlg_{\kappa}^{\sld}$, the full subcategory of $\An\CAlg^{\sld}$ spanned by the essential image of the $\kappa$-compact connective $\EE_{\infty}$-algebras in $\Sp$ under the fully faithful functor $\CAlg^{\cn}\to\An\CAlg^{\sld}$ of \Cref{lem: functoriality of solid rings}. 
\end{definition}
\begin{remark}\label{rmk: affine spectral schemes essentially small}
    The category $\CAlg^{\cn}$ is $\omega$-presentable by \Cref{prop: approximation by finitely presented algebras}, so its full subcategory $\CAlg^{\cn}_{\kappa}$ of $\kappa$-compact objects is essentially small. The embedding of \Cref{lem: functoriality of solid rings} restricts to an equivalence $\CAlg^{\cn}_{\kappa}\xrightarrow{\sim}\An\CAlg^{\sld}_{\kappa}$ onto its essential image, so $\An\CAlg^{\sld}_{\kappa}$ and its opposite $\Aff_{\kappa}$ are essentially small as well.
\end{remark}
The class of morphisms along which the exceptional functors are defined is the one detected on $\pi_{0}$ by finite expansion. Transporting the permanence properties of \Cref{prop: finite expansion permanence} to the opposite category exhibits it as a geometric setup.
\begin{lemma}\label{lem: morphism of geometric setups}
    Let $\kappa$ be an uncountable regular cardinal. 
    \begin{enumerate}
        \item The pair $(\Aff_{\kappa},\mathsf{FE})$, where $\mathsf{FE}$ is the homotopy class of morphisms of finite expansion, is a geometric setup. 
        \item The restriction of the fully faithful embedding $\CAlg^{\cn}\to\An\CAlg^{\wedge}$ to $\CAlg^{\cn}_{\kappa}$ gives on opposite categories a morphism of geometric setups $(\Aff_{\kappa},\mathsf{FE})\to(\An\Aff,E_{!})$ where $(\An\Aff,E_{!})$ is the geometric setup on the opposite category of $\An\CAlg^{\wedge}$ constructed in \Cref{lem: geometric setup on analytic E-infinity rings}. 
    \end{enumerate}
\end{lemma}
\begin{proof}[Proof of (1)]
    By \Cref{prop: finite expansion permanence}, the morphisms of finite expansion contain all equivalences and are closed under composition, base change, and left cancellation among connective $\EE_{\infty}$-algebras. Passing to opposite categories, $\mathsf{FE}$ contains all equivalences and is closed under composition, pullback, and right cancellation in $\Aff_{\kappa}$, which are the defining conditions of \Cref{def: geometric setup and suitable decomposition} (1). Pullbacks are computed by the relative tensor products $B\otimes_{A}C$, finite colimits in $\CAlg^{\cn}$ and hence internal to $\An\CAlg^{\sld}_{\kappa}\simeq\CAlg^{\cn}_{\kappa}$ by closure of $\kappa$-compact objects under $\kappa$-small colimits.
\end{proof}
\begin{proof}[Proof of (2)]
    The embedding preserves nonempty small colimits by \Cref{lem: functoriality of solid rings}, so it preserves the pushouts computing base changes of morphisms of finite expansion, and on opposite categories $\Aff_{\kappa}\hookrightarrow\An\Aff$ therefore preserves the corresponding fiber products, giving the compatibility required by \Cref{def: geometric setup and suitable decomposition} (2).
    
    That every morphism of finite expansion of discrete solid $\EE_{\infty}$-rings lands in the exceptional class $E_{!}$ is furnished by its factorization into analytically proper morphisms and open immersions in \Cref{thm: exceptional functor for finite expansion}, each factor being exceptional, and $E_{!}$ closed under composition by \Cref{lem: geometric setup on analytic E-infinity rings}.
\end{proof}
Restriction along this morphism of geometric setups carries the six operations on complete analytic $\EE_{\infty}$-rings to $\kappa$-small affine spectral schemes.
\begin{proposition}\label{prop: six-functor formalism affine spectral schemes}
    Let $\kappa$ be an uncountable regular cardinal. There is a presentable six-functor formalism
    $$\QCoh_{\sld}:\Corr(\Aff_{\kappa},\mathsf{FE})\longrightarrow\PrL \hspace{1cm} \spec(A)\mapsto\Mod_{A_{\sld}}$$
    obtained by restricting the six-functor formalism of \Cref{prop: six-functor formalism on analytic E-infinity rings} along the morphism of geometric setups of \Cref{lem: morphism of geometric setups}.
\end{proposition}
\begin{proof}
    The morphism of geometric setups of \Cref{lem: morphism of geometric setups} induces a functor on correspondence operads $\Corr(\Aff_{\kappa},\mathsf{FE})\to\Corr(\An\Aff,E_{!})$ \cite[Prop. 2.3.5 \& 2.3.6]{HeyerMann6FF}. Precomposing the six-functor formalism of \Cref{prop: six-functor formalism on analytic E-infinity rings} with this functor is again a lax symmetric monoidal functor into $\PrL$, hence a presentable six-functor formalism, and its value on $\spec(A)$ is $\Mod_{A_{\sld}}$ by construction of the embedding in \Cref{lem: functoriality of solid rings}.
\end{proof}
This is the six-functor formalism for solid quasicoherent sheaves on $\kappa$-small affine spectral schemes. 
\begin{definition}\label{def: solid six-functor formalism affine}
    Let $\kappa$ be an uncountable regular cardinal. The \emph{solid six-functor formalism on $\kappa$-small affine spectral schemes} is the presentable six-functor formalism $\QCoh_{\sld}$ of \Cref{prop: six-functor formalism affine spectral schemes}.
\end{definition}
The six-functor formalism satisfies the categorical K\"{u}nneth formula in the sense of \Cref{def: categorical Kunneth formula}. 
\begin{corollary}\label{cor: solid six functor formalism affine satisfies categorical kunneth}
    The solid six-functor formalism on $\kappa$-small affine spectral schemes satisfies the categorical K\"{u}nneth formula. 
\end{corollary}
\begin{proof}
    The functor $\CAlg_{\kappa}^{\cn}\hookrightarrow\An\CAlg^{\wedge}$ preserves nonempty small colimits by \Cref{lem: functoriality of solid rings}, so the composite $\CAlg^{\cn}_{\kappa}\to\PrL$ preserves pushouts by the categorical K\"{u}nneth formula for the six-functor formalism for complete analytic $\EE_{\infty}$-rings established in \Cref{prop: categorical kunneth for analytic rings}. 
\end{proof}
\subsection{Primness and suaveness in the affine case}\label{subsec: affine prim and suave}
With the affine formalism in hand, we take up the study of its prim, suave, and smooth morphisms, together with the related properties of properness and \'{e}taleness relative to a six-functor formalism. The definitions are those of Heyer--L. Mann \cite[\S 4.4-4.6]{HeyerMann6FF}, and the discussion parallels the solid formalism for animated algebraic geometry of Hansen--L. Mann \cite[\S 2.7]{HansenMannCLL}.

The proper morphisms are $\msld$-prim without further hypotheses.
\begin{lemma}\label{lem: proper morphisms of affines}
    Let $f:\spec(B)\to\spec(A)$ be a morphism of $\kappa$-small affine spectral schemes.
    \begin{enumerate}
        \item If $f$ is an integral morphism on $\pi_{0}$, then $f$ is $\msld$-prim. 
        \item If $f$ is a proper morphism, then $f$ is $\msld$-prim. 
    \end{enumerate}
    In both cases, there is a natural equivalence of functors $f_{*}\simeq f_{!}:\Mod_{B_{\sld}}\to\Mod_{A_{\sld}}$. 
\end{lemma}
\begin{proof}
    \Cref{lem: integral implies analytically proper} furnishes the natural identification of functors in both cases since proper morphisms are in particular integral on $\pi_{0}$ by \cite[Prop. 5.2.1.1]{SAG}. This also implies that (2) follows from (1). 

    We now show (1). By \cite[\href{https://stacks.math.columbia.edu/tag/02JM}{Tag 02JM}]{stacks-project}, integrality on $\pi_{0}$ is left cancellative, hence closed under diagonals by \cite[Lem. 2.1.5]{HeyerMann6FF}. One then verifies $\msld$-primness of $f$ using \cite[Lem. 4.5.5]{HeyerMann6FF}: the natural map $f_{!}\pr_{2*}\Delta_{f!}B\to f_{*}B$ is an equivalence under the equivalences $\Delta_{f*}\simeq\Delta_{f!}, f_{!}\simeq f_{*}$, and $\pr_{2*}\Delta_{f*}\simeq\id_{\Mod_{B_{\sld}}}$. 
\end{proof}
\begin{remark}\label{rmk: primness and properness}
    A proper morphism of $\kappa$-small affine spectral schemes need not be $\msld$-proper in the sense of the six-functor formalism since proper morphisms need not be truncated. Compare the discussion in \cite[Def. 2.7.11]{HansenMannCLL} where proper morphisms in the solid animated setting are not necessarily proper in the sense of the six-functor formalism. 
\end{remark}
A large collection of suave morphisms for the affine formalism is provided by the pseudo-lisse morphisms of \Cref{subsec: pseudo-lisse morphisms}, recalled below, whose flat factorization stands in for the structure theory unavailable in the spectral setting.
\begin{definition}\label{def: pseudo-lisse body}
    Let $A\to B$ be a morphism in $\CAlg^{\cn}$.
    \begin{enumerate}
        \item $A\to B$ is \emph{pseudo-lisse} if it admits a factorization
        $$A\to C\to B$$
        in which $A\to C$ is flat and almost of finite presentation and $B$ is perfect as a $C$-module.
        \item $A\to B$ is \emph{elementary pseudo-lisse} if it admits a factorization
        $$A\to A[T_{1},\dots,T_{n}]\to B$$
        for some $n\geq0$ in which $B$ is perfect as an $A[T_{1},\dots,T_{n}]$-module.
    \end{enumerate}
\end{definition}
\begin{remark}\label{rmk: pseudo-lisse properties recalled}
    The permanence of such morphisms is recorded in \Cref{prop: properties of pseudo-lisse morphisms}: every elementary pseudo-lisse morphism is pseudo-lisse, the elementary pseudo-lisse morphisms are closed under composition and base change, and the pseudo-lisse morphisms are closed under base change. A pseudo-lisse morphism is moreover almost of finite presentation and of finite $\Tor$-amplitude, hence of finite expansion by \Cref{lem: finiteness implies finite expansion}, so it lies in $\mathsf{FE}$ and carries the exceptional functors of $\QCoh_{\sld}$.
\end{remark}
As a special case of pseudo-lisse morphisms, we may consider the fiber-smooth morphisms in the sense of \cite[Def. 11.2.3.1]{SAG}. 
\begin{definition}\label{def: fiber-smooth morphism}
    A morphism $A\to B$ in $\CAlg^{\cn}$ is \emph{fiber-smooth} if it satisfies the following conditions: 
    \begin{enumerate}[label=(\roman*)]
        \item $A\to B$ is flat and almost of finite presentation. 
        \item For every field $k$ and morphism of connective $\EE_{\infty}$-rings $B\to k$, the static $k$-vector space $\pi_{1}(k\otimes_{B}\LL_{B/A})$ vanishes, where $\LL_{B/A}$ is the $\EE_{\infty}$-cotangent complex of $A\to B$. 
    \end{enumerate}
\end{definition}
We show pseudo-lisse morphisms are $\msld$-suave. The pseudo-lisse morphisms subsume those for which Grothendieck duality is considered in \cite[\S 6.4.2]{SAG}.  
\begin{proposition}\label{prop: affine suave morphisms}
    Let $f:\spec(B)\to\spec(A)$ be a morphism of $\kappa$-small affine spectral schemes.
    \begin{enumerate}
        \item If $f$ is proper, almost of finite presentation, and of finite $\Tor$-amplitude, then $B$ is perfect as an $A$-module and $f$ is $\msld$-suave. Furthermore, there is an equivalence $f^{!}A\simeq\omega_{f}^{\circ}$ in $\Mod_{B_{\sld}}$, where $\omega_{f}^{\circ}$ is the dualizing sheaf of \cite[Def. 6.4.2.4]{SAG}. 
        \item If $f$ is pseudo-lisse, then $f$ is $\msld$-suave. In particular, every finitely presented and faithfully flat morphism is $\msld$-suave.
        \item If $f$ is fiber-smooth in the sense of \Cref{def: fiber-smooth morphism}, then $f$ is $\msld$-suave and its diagonal $\Delta_{f}:\spec(B)\to\spec(B\otimes_{A}B)$ is both $\msld$-prim and $\msld$-suave. In particular, every fiber-smooth morphism is $\msld$-smooth. 
        \item If $f$ is \'{e}tale, then $f$ is $\msld$-\'{e}tale. 
    \end{enumerate}
\end{proposition}
\begin{proof}[Proof of (1)]
    Properness makes $A\to B$ finite \cite[Prop. 5.2.1.1]{SAG}, so almost finite presentation renders $B$ pseudocoherent over $A$ \cite[Cor. 5.2.2.2]{SAG}, and finite $\Tor$-amplitude in addition makes $B$ perfect over $A$ \cite[Prop. 7.2.4.23]{HA}. By \Cref{lem: proper morphisms of affines}, $f$ is $\msld$-prim with $f_{!}\simeq f_{*}$ and \cite[Lem. 4.5.10]{HeyerMann6FF} shows that perfectness of $B$ over $A$ implies suaveness of $f$. 

    We have an identification $\omega_{f}\simeq f^{!}A\simeq\intHom_{A_{\sld}}(B, A)$. Now $B$ is dualizable in $\Mod_{A}(\Sp)$, being perfect \cite[Prop. 7.2.4.4]{HA}, with dual $\intHom_{A}(B,A)$, and the symmetric monoidal functor $\Mod_{A}(\Sp)\to\Mod_{A_{\sld}}$ preserves duality, hence takes $\intHom_{A}(B,A)$ to $\intHom_{A_{\sld}}(B, A)$. Since the dualizing sheaf of a finite morphism is the $A$-linear dual $\intHom_{A}(B,A)$ \cite[Ex. 6.0.4.3]{SAG}, we get $f^{!}A\simeq\omega_{f}^{\circ}$, as desired. 
\end{proof}
\begin{proof}[Proof of (2)]
    Let $A_{\sld}\to C_{\sld}\to B_{\sld}$ be a pseudo-lisse factorization, with $A\to C$ flat and almost of finite presentation and $B$ perfect over $C$. Since $\msld$-suave morphisms are closed under composition \cite[Lem. 4.5.9]{HeyerMann6FF}, it suffices to the first factor, the second already having been treated in (1). 

    For $A_{\sld}\to C_{\sld}$, the underlying functor $\Aff_{\kappa}^{\Opp}\simeq\An\CAlg^{\sld}_{\kappa}\to\CAlg(\PrL)$ preserves pushouts by \Cref{cor: solid six functor formalism affine satisfies categorical kunneth}, so the morphism is $\msld$-suave once scalar extension $C_{\sld}\otimes_{A_{\sld}}(-)$ admits a $\Mod_{A_{\sld}}$-linear left adjoint \cite[Lem. 6.3.6]{CamargoSolid}. For a flat morphism almost of finite presentation, such a left adjoint is constructed in \Cref{prop: flat spectral representability}. 
\end{proof}
\begin{proof}[Proof of (3)]
    Suaveness of $f$ follows from (2). By \cite[Prop. 11.3.3.1]{SAG}, $B$ is perfect as a $B\otimes_{A}B$-module in $\Sp$, hence $\pi_{0}(B)$ is finitely generated as a $\pi_{0}(B\otimes_{A}B)$-module \cite[Cor. 7.2.4.5]{HA}, so the diagonal $\Delta_{f}$ is finite and $\msld$-prim by \Cref{lem: proper morphisms of affines}. This exhibits $\Delta_{f}$ as $\msld$-suave by (1). 

    To show $\msld$-smoothness, we need to show $\omega_{f}$ is $\otimes_{B_{\sld}}$-invertible. $\omega_{f}$ is $f$-suave, being the suave dual of the unit, so \cite[Cor. 4.5.18 (ii)]{HeyerMann6FF} implies that $\omega_{f}$ is dualizable, and therefore $\otimes_{B_{\sld}}$-invertible by \cite[Rmk. 4.5.12]{HeyerMann6FF}, showing $f$ is $\msld$-smooth. 
\end{proof}
\begin{proof}[Proof of (4)]
    An \'{e}tale morphism is flat and almost of finite presentation by \Cref{def: finiteness conditions of morphisms} (8), hence pseudo-lisse and $\msld$-suave by (2). Its diagonal is a localization at an idempotent element \cite[Rmk. 7.5.0.2, Thm. 7.5.0.6 \& Ex. 7.5.0.7]{HA}, hence is \'{e}tale and $\msld$-suave by (2). Since the diagonal is also a monomorphism, it follows from \cite[Def. 4.6.1]{HeyerMann6FF} that $f$ is $\msld$-\'{e}tale. 
\end{proof}
\begin{definition}\label{def: relative dualizing sheaf}
    Let $f:\spec(B)\to\spec(A)$ be a $\msld$-suave morphism of $\kappa$-small affine spectral schemes. The \emph{relative dualizing sheaf} $\omega_{f}$ is the object $f^{!}A\in\Mod_{B_{\sld}}$.
\end{definition}
\begin{remark}\label{rmk: comparison to lurie jiang}
    The pseudo-lisse hypothesis is related to other treatments of Grothendieck duality in spectral algebraic geometry as follows. By \Cref{prop: affine suave morphisms} (1), the pseudo-lisse morphisms contain the affine morphisms treated in \cite[\S 6]{SAG}. The pseudo-lisse hypothesis imposes neither the Noetherian nor the homological boundedness conditions appearing in \cite[\S 2]{JiangGD}, where Grothendieck duality is considered for morphisms almost of finite presentation and of finite $\Tor$-amplitude under these additional hypotheses. Every pseudo-lisse morphism is almost of finite presentation and of finite $\Tor$-amplitude by \Cref{rmk: pseudo-lisse properties recalled}. Whether every morphism of connective $\EE_{\infty}$-rings almost of finite presentation and of finite $\Tor$-amplitude is suave, or contained in the closure of pseudo-lisse morphisms under composition and base change, is open. 
\end{remark}
\subsection{Perfect and coherent sheaves in the affine case}\label{subsec: affine perfect and coherent}
The morphism classes of the preceding subsection have object-level counterparts. A sheaf is perfect when it is dualizable for the solid tensor product, and $f$-coherent when it is suave for a morphism $f$ of finite expansion, in the sense of the prim and suave objects of Heyer--L. Mann \cite[\S 4.4 \& 4.5]{HeyerMann6FF}. On these subcategories, the duality of a suave morphism restricts to an anti-autoequivalence, the affine form of the Grothendieck--Verdier duality globalized in \Cref{subsec: duality theory for spectral algebraic stacks}. Compare the animated affine case of \cite[Prop. 2.8.3]{HansenMannCLL}.
\begin{definition}\label{def: perfect and relatively coherent objects}
    Let $\kappa$ be an uncountable regular cardinal. 
    \begin{enumerate}
        \item For $\spec(A)\in\Aff_{\kappa}$ an object of $\QCoh_{\sld}(A):=\Mod_{A_{\sld}}$ is \emph{perfect} if it is dualizable. Denote the full subcategory of \emph{perfect sheaves} by $\Perf_{\sld}(A)$.  
        \item Let $f:\spec(B)\to\spec(A)$ be a morphism of finite expansion. An object of $\QCoh_{\sld}(B):=\Mod_{B_{\sld}}$ is \emph{$f$-coherent} if it is an $f$-suave object in $\QCoh_{\sld}$. Denote the full subcategory of \emph{$f$-coherent sheaves} by $\Coh_{\sld}(B,f)$. 
    \end{enumerate}
\end{definition}
\begin{remark}\label{rmk: perfect sheaves discrete}
    By \Cref{prop: solid rings are Fredholm}, every dualizable object of $\Mod_{A_{\sld}}$ is discrete, so a perfect sheaf is discrete and $\Perf_{\sld}(A)$ is identified with the perfect $A$-modules in $\Sp$. Namely, we have an equivalence of categories $\Perf_{\sld}(A)\simeq\Perf(A)$, where the latter is the full subcategory of $\Mod_{A}(\Sp)$ spanned by the compact, or equivalently dualizable, objects \cite[Prop. 7.2.4.4]{HA}. 
\end{remark}
\begin{remark}\label{rmk: coherent sheaves}
    Lurie defines a coherent sheaf on a Noetherian affine spectral scheme $\spec(B)$ to be a (discrete) $B$-module in $\Sp$ such that each homotopy group is finitely generated as a $\pi_{0}(B)$-module in static Abelian groups \cite[\S 6.4.3]{SAG}. This does not agree with \Cref{def: perfect and relatively coherent objects} (2) even for $A=B=\Sph$ where the $\id_{\Sph}$-coherent sheaves are the compact spectra, but a spectrum with each homotopy group finitely generated need not be compact, or even pseudocoherent. We have elected to refer to the $f$-suave sheaves as $f$-coherent in line with \cite[\S 2.8]{HansenMannCLL} (cf. \cite[Cor. 2.8.4 (iii)]{HansenMannCLL}).  
\end{remark}
Unlike the perfect sheaves, the relatively coherent sheaves are characterized less directly, through the six-functor formalism rather than a finiteness condition in $\Sp$. Both subcategories are nonetheless stable and carry a duality, recorded next.
\begin{proposition}\label{prop: f-coherent and perfect sheaves affine}
    Let $f:\spec(B)\to\spec(A)$ be a morphism of finite expansion between $\kappa$-small affine spectral schemes. 
    \begin{enumerate}
        \item $\Perf_{\sld}(B)$ and $\Coh_{\sld}(B,f)$ are stable subcategories of $\Mod_{B_{\sld}}$. 
        \item The functor given on objects by $M\mapsto \intHom_{B_{\sld}}(M, f^{!}A)$ defines an anti-autoequivalence of $f$-coherent sheaves $\Coh_{\sld}(B,f)^{\Opp}\xrightarrow{\sim}\Coh_{\sld}(B,f)$. 
        \item The functor given on objects by $M\mapsto \intHom_{B_{\sld}}(M, B)$ defines an anti-autoequivalence of perfect sheaves $\Perf_{\sld}(B)^{\Opp}\xrightarrow{\sim}\Perf_{\sld}(B)$. 
        \item If $f$ is $\msld$-suave, then $\Perf_{\sld}(B)\subseteq\Coh_{\sld}(B,f)$ and the anti-autoequivalence of (2) is given on objects by $M\mapsto\intHom_{B_{\sld}}(M,\omega_{f})$. If the diagonal $\Delta_{f}$ is $\msld$-suave, then $\Coh_{\sld}(B,f)\subseteq\Perf_{\sld}(B)$.
        \item If $f$ is fiber-smooth, then there is an equivalence of categories $\Perf_{\sld}(B)\simeq\Coh_{\sld}(B,f)$. 
    \end{enumerate}
\end{proposition}
\begin{proof}[Proof of (1)]
    The identification of $\Perf_{\sld}(B)$ with the perfect $B$-modules in $\Sp$ \Cref{rmk: perfect sheaves discrete} exhibits it as a stable subcategory \cite[Prop. 7.2.4.2]{HA}. That $\Coh_{\sld}(B,f)$ is a stable subcategory is \cite[Cor. 4.4.13]{HeyerMann6FF}.
\end{proof}
\begin{proof}[Proof of (2)]
    This is \cite[Lem. 4.4.4 \& 4.4.17]{HeyerMann6FF}.
\end{proof}
\begin{proof}[Proof of (3)]
    The identity identifies $\Perf_{\sld}(B)\simeq\Coh_{\sld}(B,\id_{B_{\sld}})$ \cite[Ex. 4.4.3]{HeyerMann6FF}, so this is the case $f=\id_{B_{\sld}}$ of (2).
\end{proof}
\begin{proof}[Proof of (4)]
    The containment $\Perf_{\sld}(B)\subseteq\Coh_{\sld}(B,f)$ is \cite[Cor. 4.5.18]{HeyerMann6FF} (i), and $f^{!}A=\omega_{f}$ by \Cref{def: relative dualizing sheaf} puts the anti-autoequivalence of (2) in the stated form. When $\Delta_{f}$ is $\msld$-suave, the reverse containment $\Coh_{\sld}(B,f)\subseteq\Perf_{\sld}(B)$ is \cite[Cor. 4.5.18]{HeyerMann6FF} (ii).
\end{proof}
\begin{proof}[Proof of (5)]
    By \Cref{prop: affine suave morphisms} (3), $f$ is $\msld$-suave with $\msld$-prim and $\msld$-suave diagonal, so both containments of (4) apply.
\end{proof}
\subsection{Six operations on spectral algebraic stacks: construction}\label{subsec: construction on stacks}
The affine formalism and its sheaves now established, we globalize the six-functor formalism to spectral algebraic stacks. Solid modules satisfy descent for the faithfully flat and finitely presented topology by \Cref{thm: solid fppf descent}, and the arguments of Heyer--L. Mann \cite[\S 3.4]{HeyerMann6FF} promotes the affine formalism to one on stacks. The morphisms carrying exceptional functors are those locally of finite expansion, characterized through the affine case by descent.
\begin{definition}\label{def: kappa-small spectral stacks}
    Let $\kappa$ be an uncountable regular cardinal. The \emph{category of $\kappa$-small spectral algebraic stacks} $\Stk_{\kappa}\subseteq\PSh(\Aff_{\kappa};\Ani)$ is the full subcategory of presheaves of animae on $\Aff_{\kappa}$ that are sheaves for the faithfully flat and finitely presented topology.
\end{definition}
Globalizing the six operations to these stacks turns on the descent that solid modules enjoy along faithfully flat and finitely presented covers and provides the class of stacky morphisms along which exceptional functors are defined, and the six-functor formalism carrying them.
\begin{theorem}\label{thm: six-functor formalism on stacks}
    Let $\kappa$ be an uncountable regular cardinal. There is a homotopy class $\mathsf{LFE}$ of morphisms of $\Stk_{\kappa}$ containing those whose base change along every morphism from a $\kappa$-small affine spectral scheme to the target is affine and of finite expansion. The pair $(\Stk_{\kappa},\mathsf{LFE})$ is a geometric setup, the inclusion $(\Aff_{\kappa},\mathsf{FE})\to(\Stk_{\kappa},\mathsf{LFE})$ a morphism of geometric setups, and there is a six-functor formalism
    $$\QCoh_{\sld}:\Corr(\Stk_{\kappa},\mathsf{LFE})\longrightarrow\widehat{\Cat},\qquad X\mapsto\QCoh_{\sld}(X),$$
    whose restriction along $\Aff_{\kappa}\hookrightarrow\Stk_{\kappa}$ is the solid six-functor formalism of \Cref{def: solid six-functor formalism affine}.
\end{theorem}
\begin{proof}
    The affine formalism of \Cref{prop: six-functor formalism affine spectral schemes} descends along the faithfully flat and finitely presented topology: solid modules satisfy descent for it as shown in \Cref{thm: solid fppf descent}, and faithfully flat and finitely presented morphisms are $\msld$-suave by \Cref{prop: affine suave morphisms} (2). These verify the hypotheses of the globalization theorem \cite[Thm. 3.4.11]{HeyerMann6FF}, which produces $\mathsf{LFE}$, the geometric setup, the morphism of geometric setups, and the six-functor formalism restricting to the affine one.
\end{proof}
The construction produces a six-functor formalism and a distinguished class of morphisms, named as follows.
\begin{definition}\label{def: solid six-functor formalism on stacks}
    Let $\kappa$ be an uncountable regular cardinal.
    \begin{enumerate}
        \item The \emph{solid six-functor formalism on $\kappa$-small spectral algebraic stacks} is the six-functor formalism $\QCoh_{\sld}$ of \Cref{thm: six-functor formalism on stacks}.
        \item A morphism of $\kappa$-small spectral algebraic stacks is \emph{locally of finite expansion} if it lies in the class $\mathsf{LFE}$ of \Cref{thm: six-functor formalism on stacks}.
    \end{enumerate}
\end{definition}
The class $\mathsf{LFE}$ is abstractly characterized by minimality and locality. For computations, however, one wants a concrete handle reducing statements about a stack to a covering family of affine schemes. Such reductions run along the covers for which solid quasicoherent sheaves satisfy exceptional descent, which we isolate below. 
\begin{definition}\label{def: universal QCoh-cover}
    Let $\{V_{i}\to Y\}_{i\in I}$ be a small family of morphisms in $\mathsf{LFE}$ with common target $Y\in\Stk_{\kappa}$, and for a morphism $X\to Y$ write $U_{i}:=V_{i}\times_{Y}X$. The family $\{V_{i}\to Y\}_{i\in I}$ is a \emph{universal $\QCoh^{!}_{\sld}$-cover} if for every morphism $X\to Y$ of $\kappa$-small spectral algebraic stacks the augmented cosimplicial diagram
    $$\QCoh_{\sld}(X)\to\prod_{i_{0}\in I}\QCoh_{\sld}(U_{i_{0}})\rightrightarrows\prod_{i_{0},i_{1}\in I}\QCoh_{\sld}(U_{i_{0}}\times_{X}U_{i_{1}})\to\cdots$$
    with transition functors the exceptional pullbacks along the projections is a limit diagram in $\widehat{\Cat}$.
\end{definition}
The prototypical such cover is a faithfully flat and finitely presented morphism of affine spectral schemes (cf. \Cref{thm: solid fppf descent}).
\begin{lemma}\label{lem: ffp is universal cover}
    Let $f:\spec(B)\to\spec(A)$ be a faithfully flat and finitely presented morphism. Then $f$ is a $\msld$-suave cover and a universal $\QCoh^{!}_{\sld}$-cover.
\end{lemma}
\begin{proof}
    By \Cref{prop: affine suave morphisms} (2) the morphism $f$ is $\msld$-suave, and faithful flatness makes it a cover. It is a universal $\QCoh^{!}_{\sld}$-cover by \cite[Lem. 4.7.1]{HeyerMann6FF}.
\end{proof}
Using these covers, the morphisms locally of finite expansion admit the anticipated local characterization: the condition is detected on affine targets by the affine notion, and are local on source and target.
\begin{proposition}\label{prop: locally of finite expansion characterization}
    The class $\mathsf{LFE}$ of morphisms of $\kappa$-small spectral algebraic stacks locally of finite expansion satisfies the following.
    \begin{enumerate}[label=(\roman*)]
        \item $\mathsf{LFE}$ contains every morphism $f: X\to Y$ such that for each $\kappa$-small affine spectral scheme $V\to Y$ the base change $X\times_{Y}V\to V$ is locally of finite expansion.
        \item A morphism $f: X\to Y$ with $Y$ a $\kappa$-small affine spectral scheme is locally of finite expansion if and only if there is a universal $\QCoh^{!}_{\sld}$-cover $\{U_{i}\to X\}_{i\in I}$ by $\kappa$-small affine spectral schemes with each composite $U_{i}\to X\to Y$ of finite expansion.
        \item $\mathsf{LFE}$ contains every morphism $f: X\to Y$ admitting a universal $\QCoh^{!}_{\sld}$-cover $\{U_{i}\to X\}_{i\in I}$ with each composite $U_{i}\to X\to Y$ locally of finite expansion.
        \item $\mathsf{LFE}$ contains every morphism $f: X\to Y$ admitting a universal $\QCoh_{\sld}^{!}$-cover $\{V_{j}\to Y\}_{j\in J}$ with each base change $X\times_{Y}V_{j}\to V_{j}$ locally of finite expansion.
        \item $\mathsf{LFE}$ contains every morphism $f: X\to Y$ admitting a small family $\{V_{j}\to Y\}_{j\in J}$ with $\coprod_{j\in J}V_{j}\to Y$ an effective epimorphism in $\Stk_{\kappa}$ and each base change $X\times_{Y}V_{j}\to V_{j}$ locally of finite expansion.
    \end{enumerate}
    Moreover, $\mathsf{LFE}$ is minimal among the homotopy classes of morphisms of $\Stk_{\kappa}$ satisfying (i) to (iv) and the conclusions of \Cref{thm: six-functor formalism on stacks}. 
\end{proposition}
\begin{proof}
    The statements (i) to (iv) and minimality follow from \cite[Thm. 3.4.11]{HeyerMann6FF}, whose hypotheses are the faithfully flat and finitely presented descent of \Cref{thm: solid fppf descent} and the affine six-functor formalism of \Cref{prop: six-functor formalism affine spectral schemes}.

    For general covers in (v), we follow the strategy of \cite[Thm. 2.7.8 (iii) (b)]{HansenMannCLL}. Let $f:X\to Y$ be a morphism in $\Stk_{\kappa}$ and assume there is a jointly effective epimorphic family $\{V_{j}\to Y\}_{j\in J}$ such that for all $j\in J$ the morphism $V_{j}\times_{Y}X\to V_{j}$ is locally of finite expansion. Being locally of finite expansion can be checked after pullback along all morphisms of $\kappa$-small affine spectral schemes to $Y$ so we may, without loss of generality, take $Y$ to itself be a $\kappa$-small affine spectral scheme. By \cite[Lem. A.3.9]{MannRigid6FF}, there exists a faithfully flat and finitely presented cover $\{W_{i}\to Y\}_{i\in I}$ by $\kappa$-small affine spectral schemes such that each $W_{i}\to Y$ admits a factorization as $W_{i}\to V_{j}\to Y$ for some $j\in J$. By the pasting lemma and closure of morphisms locally of finite expansion under base change, it follows that each $W_{i}\times_{Y}X\to W_{i}$ is locally of finite expansion. Since faithfully flat and finitely presented morphisms of $\kappa$-small affine spectral schemes are universal $\QCoh^{!}_{\sld}$-covers by \Cref{lem: ffp is universal cover}, (iv) applies to $\{W_{i}\to Y\}_{i\in I}$ and $f$ is locally of finite expansion.
\end{proof}
By \Cref{lem: ffp is universal cover} every faithfully flat and finitely presented morphism is a suave cover, so many questions about stacks reduce to the affine schemes covering them. This reduction underlies the duality theory of the next subsection.
\subsection{Duality theory for spectral algebraic stacks}\label{subsec: duality theory for spectral algebraic stacks}
The construction complete, we now consider duality. The perfect and coherent sheaves of the affine case have stacky counterparts, defined via descent. The subsection closes with the Grothendieck--Verdier duality for suave morphisms of stacks and a restatement of the main theorem. Throughout, compare the animated development of Hansen--L. Mann \cite[\S 2.8]{HansenMannCLL}.
\begin{definition}\label{def: perfect and coherent sheaves on stacks}
    Let $f: X\to Y$ be a morphism of $\kappa$-small spectral algebraic stacks locally of finite expansion.
    \begin{enumerate}
        \item An object $M\in\QCoh_{\sld}(X)$ is \emph{perfect} if it is dualizable. Denote the full subcategory of perfect sheaves by $\Perf_{\sld}(X)$.
        \item An object $M\in\QCoh_{\sld}(X)$ is \emph{$f$-coherent} if it is an $f$-suave object of $\QCoh_{\sld}$. Denote the full subcategory of $f$-coherent sheaves by $\Coh_{\sld}(X,f)$.
    \end{enumerate}
\end{definition}
Both subcategories are governed by their affine restrictions. Dualizability can be checked on an affine cover and descends along all covers. 
\begin{proposition}\label{prop: perfect sheaves on stacks}
    Let $X$ be a $\kappa$-small spectral algebraic stack.
    \begin{enumerate}
        \item Perfect sheaves are preserved by arbitrary pullback and descend along all covers.
        \item An object $M\in\QCoh_{\sld}(X)$ is perfect if and only if for every $\kappa$-small affine spectral scheme $g:\spec(A)\to X$ the pullback $g^{*}M$ is perfect in $\Mod_{A_{\sld}}$. In particular, every perfect sheaf is discrete.
    \end{enumerate}
\end{proposition}
\begin{proof}[Proof of (1)]
    Dualizable objects are preserved by symmetric monoidal functors, hence by pullback, and admit descent.
\end{proof}
\begin{proof}[Proof of (2)]
    Passing to an affine cover reduces the assertion to the affine case, where $\Perf_{\sld}(A)$ is identified with the perfect $A$-modules in $\Sp$ via \Cref{rmk: perfect sheaves discrete}. Discreteness is affine-local, hence descends, giving the final claim.
\end{proof}
Relative coherence behaves similarly, but its descent is confined to the suave covers.
\begin{proposition}\label{prop: coherent sheaves on stacks}
    Let $f: X\to Y$ be a morphism of $\kappa$-small spectral algebraic stacks locally of finite expansion. The $f$-coherent sheaves are preserved by pullback along $\msld$-suave morphisms and descend along $\msld$-suave covers. In particular, they descend along covers in the faithfully flat and finitely presented topology.
\end{proposition}
\begin{proof}
    This is \cite[Lem. 4.4.8 \& 4.5.16]{HeyerMann6FF}. The final assertion follows by \Cref{lem: ffp is universal cover} using suaveness of such covers.
\end{proof}
With coherence available on stacks, the duality of a suave morphism can be stated globally. As in the affine case, its dualizing sheaf is the exceptional pullback of the unit.
\begin{definition}\label{def: dualizing sheaf on stacks}
    Let $f: X\to Y$ be a $\msld$-suave morphism of $\kappa$-small spectral algebraic stacks. The \emph{relative dualizing sheaf} $\omega_{f}$ is the object $f^{!}1_{Y}\in\QCoh_{\sld}(X)$, agreeing with \Cref{def: relative dualizing sheaf} when $f$ is a morphism of affine spectral schemes.
\end{definition}
\begin{theorem}[Grothendieck--Verdier duality]\label{thm: Grothendieck Verdier duality stacks}
    Let $f: X\to Y$ be a $\msld$-suave morphism of $\kappa$-small spectral algebraic stacks.
    \begin{enumerate}
        \item The natural morphism $\intHom_{Y}(f_{!}M,N)\to f_{*}\intHom_{X}(M,f^{!}N)$ is an equivalence in $\QCoh_{\sld}(Y)$ for all $M\in\QCoh_{\sld}(X)$ and $N\in\QCoh_{\sld}(Y)$.
        \item The functor given on objects by $M\mapsto\intHom_{X}(M,\omega_{f})$ is an anti-autoequivalence $\Coh_{\sld}(X,f)^{\Opp}\xrightarrow{\sim}\Coh_{\sld}(X,f)$.
    \end{enumerate}
\end{theorem}
\begin{proof}[Proof of (1)]
    This is \cite[Prop. 3.2.2]{HeyerMann6FF} (iii).
\end{proof}
\begin{proof}[Proof of (2)]
    This is \cite[Lem. 4.4.4 \& 4.4.17]{HeyerMann6FF}.
\end{proof}
Combining the construction, the morphism classes, and the duality just obtained yields the theorem announced at the outset. We restate it here with some terms unpacked. 
\begin{theorem}\label{thm: main six-functor formalism}
    Let $\kappa$ be an uncountable regular cardinal and $\Stk_{\kappa}$ the category of $\kappa$-small spectral algebraic stacks. There is a six-functor formalism
    $$\QCoh_{\sld}:\Corr(\Stk_{\kappa},\mathsf{LFE})\longrightarrow\widehat{\Cat} \hspace{1cm} X\mapsto\QCoh_{\sld}(X)$$
    satisfying the following.
    \begin{enumerate}
        \item For every morphism $f: X\to Y$ there is an adjoint pair 
        $$f^{*}:\QCoh_{\sld}(Y)\rightleftarrows\QCoh_{\sld}(X): f_{*}$$ 
        with $f^{*}$ symmetric monoidal.
        \item For every morphism $f: X\to Y$ locally of finite expansion there is an adjoint pair 
        $$f_{!}:\QCoh_{\sld}(X)\rightleftarrows\QCoh_{\sld}(Y): f^{!}$$ 
        satisfying proper base change and the projection formula.
        \item If $f: X\to Y$ is $\msld$-suave-locally pseudo-lisse (ie. pseudo-lisse on a $\msld$-suave cover of source and target) then $f$ is $\msld$-suave and there is a natural equivalence $f^{!}(-)\simeq\omega_{f}\otimes_{X}^{\sld}f^{*}(-)$. This holds in particular when $f$ is pseudo-lisse locally in the faithfully flat and finitely presented topology.
        \begin{itemize}
            \item [(3.1)] If $f$ is a fiber-smooth map of $\kappa$-small affine spectral schemes, then $f$ is $\msld$-smooth and $\omega_{f}$ is $\otimes_{X}^{\sld}$-invertible. 
            \item [(3.2)] If $f$ is \'{e}tale-locally an \'{e}tale map, then $f$ is $\msld$-\'{e}tale and $\omega_{f}\simeq1_{X}$.
        \end{itemize}
        \item If $f$ is an integral map of $\kappa$-small affine spectral schemes, then $f$ is $\msld$-prim and there is an equivalence $f_{!}(-)\simeq f_{*}(-)$. In particular, this holds if $f$ is proper. 
        \item If $f$ is $\msld$-suave, then the following variant of Grothendieck--Verdier duality holds: 
        \begin{itemize}
            \item [(5.1)] The natural morphism 
            $$\intHom_{Y}(f_{!}M,N)\to f_{*}\intHom_{X}(M,f^{!}N)$$ 
            is an equivalence in $\QCoh_{\sld}(Y)$ for all $M\in\QCoh_{\sld}(X),N\in\QCoh_{\sld}(Y)$. 
            \item [(5.2)] The functor given on objects by $M\mapsto\intHom_{X}(M,\omega_{f})$ is an anti-autoequivalence of $\Coh_{\sld}(X,f)$.
        \end{itemize}
    \end{enumerate}
\end{theorem}
\begin{proof}
    The six-functor formalism together with (1) and (2) is \Cref{thm: six-functor formalism on stacks}, drawing on \cite[Prop. 3.1.8 \& 3.2.2]{HeyerMann6FF}. In (3) membership in $\mathsf{LFE}$ is furnished by \Cref{prop: locally of finite expansion characterization}, and the suaveness is \Cref{prop: affine suave morphisms} (2) carried along suave covers by the locality of suave morphisms on source and target \cite[\S 4.5]{HeyerMann6FF}, the dualizing-sheaf formula holding by definition of $\msld$-suaveness. The fiber-smooth case (3.1) is \Cref{prop: affine suave morphisms} (3). The \'{e}tale case (3.2) follows from \Cref{prop: affine suave morphisms} (4) and \'{e}tale-locality of \'{e}taleness on source and target \cite[Lem. 4.6.3]{HeyerMann6FF}. Property (4) is \Cref{lem: proper morphisms of affines}, and (5) is \Cref{thm: Grothendieck Verdier duality stacks}.
\end{proof}
\appendix
\crefalias{section}{appendix}
\crefalias{subsection}{appendix}
\section{\texorpdfstring{$t$}{t}-bistable categories and their localizations}\label{app: categories}
In this appendix we collect the category-theoretic results used in the preceding text. Our main objects of study are stable presentable categories satisfying mild additional conditions on their $t$-structure and their symmetric monoidal counterparts. Much of the underlying formalism is standard, developed in the works of Aambø \cite{LocalizingTStruct}, Antieau \cite{AntieauEnhancements}, and Lurie \cite{HA, SAG}. Our primary original contribution is in \Cref{subsec: detecting smashing kernel} where we prove a recognition criterion for when the kernel of a symmetric monoidal localization is the category of modules over an idempotent algebra. 
\subsection{Stable categories and their \texorpdfstring{$t$}{t}-structures}\label{subsec: t-structures}
We recall the conditions that a $t$-structure on a stable category may satisfy. Throughout, we work with stable presentable categories. 
\begin{definition}\label{def: t-structure adjectives}
    Let $\Cscr$ be a stable presentable category with $t$-structure $(\Cscr_{\geq0},\Cscr_{\leq0})$. 
    \begin{enumerate}
        \item The $t$-structure on $\Cscr$ is \emph{compatible with filtered colimits} if $\Cscr_{\leq0}$ is closed under small filtered colimits in $\Cscr$. 
        \item The $t$-structure on $\Cscr$ is \emph{compatible with countable (resp. small) products} if $\Cscr_{\geq0}$ is closed under countable (resp. small) products computed in $\Cscr$.
        \item The $t$-structure on $\Cscr$ is \emph{right complete} if the natural functor 
        $$\colim\left(\Cscr_{\geq0}\hookrightarrow\Cscr_{\geq-1}\hookrightarrow\Cscr_{\geq-2}\hookrightarrow\dots\right)\to\Cscr,$$
        with the colimit taken in $\PrL$, is an equivalence. In particular, all Whitehead towers are convergent.  
        \item The $t$-structure on $\Cscr$ is \emph{left complete} if the natural functor 
        $$\Cscr\to\lim\left(\dots\to\Cscr_{\leq2}\xrightarrow{\tau_{\leq 1}}\Cscr_{\leq 1}\xrightarrow{\tau_{\leq0}}\Cscr_{\leq0}\right)$$
        is an equivalence. In particular, all Postnikov towers are convergent. 
        \item The $t$-structure on $\Cscr$ is \emph{right separated} if 
        $$\bigcap_{n\in \ZZ}\Cscr_{\leq n}\simeq0.$$
        \item The $t$-structure on $\Cscr$ is \emph{left separated} if 
        $$\bigcap_{n\in\ZZ}\Cscr_{\geq n}\simeq0.$$
    \end{enumerate}
\end{definition}
\begin{remark}\label{rmk: filtered colimits implies accessible}
    If $\Cscr$ is a stable presentable category with a $t$-structure compatible with filtered colimits, the functor $\tau_{\leq0}$ is accessible, implying accessibility of the $t$-structure \cite[Prop. 1.4.4.13]{HA}. 
\end{remark}
We will primarily be concerned with a variant of $t$-stable categories as defined by Aambø \cite[Def. 2.22]{LocalizingTStruct}, subject to the further requirement of left completeness. 
\begin{definition}\label{def: t-bistable category}
    Let $\Cscr$ be a stable presentable category with $t$-structure $(\Cscr_{\geq0},\Cscr_{\leq0})$. 
    \begin{enumerate}
        \item $\Cscr$ is \emph{$t$-bistable} if the $t$-structure is compatible with filtered colimits and both left and right complete. 
        \item $\Cscr$ is \emph{$t^{+}$-bistable} if it is $t$-bistable and the $t$-structure on $\Cscr$ is furthermore compatible with countable products.  
    \end{enumerate} 
\end{definition}
\begin{remark}\label{rmk: hearts are Grothendieck}
    By \cite[Rmk. 1.3.5.23]{HA}, the heart of a stable presentable category with accessible $t$-structure and compatible with filtered colimits is Grothendieck Abelian. In particular, the hearts of $t$-bistable categories are Grothendieck Abelian categories and are in particular presentable (cf. \cite[Prop. 3.10]{BekeModel}). If the category is furthermore $t^{+}$-bistable, then countable products are exact in the heart, as can be observed by countable products in the connective and coconnective subcategories agreeing.  
\end{remark}
$t$-bistable categories satisfy Whitehead's theorem. 
\begin{lemma}\label{lem: separatedness and prestable}
    Let $\Cscr$ be a $t$-bistable category. 
    \begin{enumerate}
        \item $\Cscr$ is both left and right separated. 
        \item Whitehead's theorem holds in the stable category $\Cscr$: $f:M\to N$ is an equivalence in $\Cscr$ if and only if $\pi_{*}(f):\pi_{*}(M)\to\pi_{*}(N)$ is an equivalence of graded objects in $\Cscr^{\heartsuit}$. 
    \end{enumerate} 
\end{lemma}
\begin{proof}[Proof of (1)]
    This is \cite[Lem. 2.15]{AntieauEnhancements}. 
\end{proof}
\begin{proof}[Proof of (2)]
    $(\Rightarrow)$ If $f:M\to N$ is an equivalence in $\Cscr$ then we can take the trivial cofiber sequence $M\xrightarrow{f}N\to0$ and consider the long exact sequence on homotopy objects to obtain the claim. 

    $(\Leftarrow)$ Suppose $f:M\to N$ is a morphism in $\Cscr$ which induces an equivalence of graded homotopy objects in the heart and set $C:=\cofib(f)$. Since $\pi_{*}(f)$ is an equivalence, the long exact sequence gives $\pi_{k}(C)\simeq0$ for all $k\in\ZZ$. Fix $r\leq0$ and consider $C_{r}:=\tau_{\geq r}C\in\Cscr_{\geq r}$. We have (co)fiber sequences $\pi_{k}(C_{r})[k]\to\tau_{\leq k}C_{r}\to\tau_{\leq k-1}C_{r}$ so induction with base case $k=r$ where $\tau_{\leq r}C_{r}\simeq\pi_{r}(C_{r})[r]$ exhibits each Postnikov truncation $\tau_{\leq k}C_{r}$ as an extension of zero objects, hence zero. By left completeness, $C_{r}\simeq\lim_{k}\tau_{\leq k}C_{r}\simeq0$ and $C$ is such that $C_{r}\simeq0$ for all $r\leq0$ so writing $C$ as the colimit of its Whitehead tower yields $C\simeq0$, as desired. 
\end{proof}
We isolate a further condition on a $t$-bistable category, not automatic but satisfied in cases of interest, which strengthens the exactness of products considered in \Cref{rmk: hearts are Grothendieck}. It rests on the following notions.
\begin{definition}\label{def: compact projective objects}
    Let $\Cscr$ be a $t$-bistable category. 
    \begin{enumerate}
        \item An object $P\in\Cscr_{\geq0}$ is \emph{compact projective} if $\Hom_{\Cscr}(P,-):\Cscr\to\Sp$ is $t$-exact and preserves small colimits. 
        \item $\Cscr$ is \emph{generated by compact projective objects} if there exists a set of compact projective objects $\{P_{i}\}_{i\in I}$ such that the functors 
        $$\{\Hom_{\Cscr}(P_{i},-):\Cscr\to\Sp\}_{i\in I}$$ 
        are jointly conservative.
    \end{enumerate}
\end{definition}
The $t$-structure of a $t$-bistable category generated by compact projective objects is compatible with products, a stable counterpart of the fact that a Grothendieck Abelian category generated by compact projective objects has exact products \cite[Lem. A.3.15]{NeemanTriangulated}. 
\begin{proposition}\label{prop: compact projectively generated implies product compatibility}
    Let $\Cscr$ be a $t$-bistable category generated by compact projective objects. Then the $t$-structure on $\Cscr$ is compatible with small products. 
\end{proposition}
\begin{proof}
    Let $\{M_{i}\}_{i\in I}$ be a small collection of connective objects with product $M:=\prod_{i\in I}M_{i}$. For each compact projective generator $P$, using that $\Hom_{\Cscr}(-,-)$ commutes with limits in the second argument, we have an equivalence $\Hom_{\Cscr}(P,M)\simeq\prod_{i\in I}\Hom_{\Cscr}(P,M_{i})$ which is connective by right $t$-exactness of $\Hom_{\Cscr}(P,-)$, connectivity of $M_{i}$, and compatibility of the $t$-structure on $\Sp$ with small products, as the formation of homotopy groups commutes with products in $\Sp$. Consider the (co)fiber sequence $\tau_{\geq 0}M\to M\to\tau_{\leq -1}M$. For any compact projective generator $P$, we obtain a (co)fiber sequence $\Hom_{\Cscr}(P,\tau_{\geq0}M)\to \Hom_{\Cscr}(P,M)\to\Hom_{\Cscr}(P, \tau_{\leq -1}M)$ where $\Hom_{\Cscr}(P,\tau_{\geq0}M)$ is connective by $P$ being compact projective and $\Hom_{\Cscr}(P,M)$ is connective by the previous argument. Since connective objects are closed under cofibers, it follows that $\Hom_{\Cscr}(P, \tau_{\leq -1}M)$ is connective as well. But compact projectivity of $P$ also implies that $\Hom_{\Cscr}(P, \tau_{\leq -1}M)\in\Sp_{\leq-1}$ showing that $\Hom_{\Cscr}(P, \tau_{\leq -1}M)\simeq0$. By joint conservativity of the compact projective generators, we conclude $\tau_{\leq-1}M\simeq0$ and $\tau_{\geq0}M\to M$ an equivalence. 
\end{proof}
When $\Cscr$ is $t^{+}$-bistable, the exactness properties of the heart do not by themselves recover the homotopy of a sequential limit from that of its terms: the relationship is mediated by a Milnor sequence, whose theory is the concern of the next section.
\subsection{Milnor sequences in the heart}\label{subsec: Milnor sequences in the heart}
For a sequential diagram 
$$\dots\to X_{2}\to X_{1}\to X_{0}$$
in the category $\Ani_{*}$ of pointed animae with limit $X$, the Milnor sequence 
$$0\to \lim_{n}^{1}\pi_{k+1}(X_{n})\to\pi_{k}(X)\to\lim_{n}\pi_{k}(X_{n})\to0$$
recovers the homotopy groups of $X$ in degrees $k\geq 2$ as an extension. In the category of Abelian groups, the Mittag--Leffler condition gives a convenient criterion for the derived limit term $\lim^{1}_{n}\pi_{k+1}(X_{n})$ to vanish so that there is an isomorphism $\pi_{k}(X)\simeq\lim_{n}\pi_{k}(X_{n})$. 
\begin{definition}\label{def: tower}
    \begin{enumerate}
        \item Let $\Cscr$ be a $t^{+}$-bistable category. The \emph{category of towers} $\Tow(\Cscr)$ is the functor category $\Fun(\NN^{\Opp},\Cscr)$.
        \item Let $\Asf$ be a Grothendieck Abelian category in which countable products are exact. The \emph{category of towers} $\Tow(\Asf)$ is the functor category $\Fun(\NN^{\Opp},\Asf)$.
    \end{enumerate} 
\end{definition}
As in the case of the category of Abelian groups, the higher derived limits of sequential diagrams vanish by a result of Roos.  
\begin{lemma}[{\cite[Thm. 2.1 \& Rmk. 2.3]{RoosLimits}}]\label{lem: vanishing of higher derived limits}
    Let $\Asf$ be a Grothendieck Abelian category in which countable products are exact. Then the functor $\lim_{n}:\Tow(\Asf)\to\Asf$ given on objects by $(\dots\to A_{2}\to A_{1}\to A_{0})\mapsto\lim_{n}A_{n}$ admits higher derived functors $\lim_{n}^{m}$ which vanish for $m\geq 2$. 
\end{lemma}
Only the limit and its first derived functor contribute, and Milnor sequences take the familiar form.
\begin{proposition}\label{prop: Milnor sequence in t-plus bistable}
    Let $\Cscr$ be a $t^{+}$-bistable category and $\{M_{n}\}_{n}\in\Tow(\Cscr)$ a tower with limit $M$. Then for all $k\in\ZZ$ there is a short exact sequence in $\Cscr^{\heartsuit}$
    \begin{equation}\label{eqn: t plus bistable Milnor sequence}
        0\to\lim_{n}^{1}\pi_{k+1}(M_{n})\to\pi_{k}(M)\to\lim_{n}\pi_{k}(M_{n})\to0.
    \end{equation}
\end{proposition}
\begin{proof}
    Writing $M$ as an equalizer of a product, we obtain a (co)fiber sequence 
    $$M\xrightarrow{e}\prod_{\NN}M_{n}\xrightarrow{f}\prod_{\NN}M_{n}$$
    where we have regarded $\NN$ as a set and $f$ is the endomorphism induced by the transition maps $\rho_{n+1}:M_{n+1}\to M_{n}$ of the sequential diagram by $(m_{n})_{n\in\NN}\mapsto(m_{n}-\rho_{n+1}(m_{n+1}))_{n\in\NN}$. Passing to the long exact sequence on homotopy objects 
    $$% https://q.uiver.app/#q=WzAsNixbMCwwLCJcXGRvdHMiXSxbMSwwLCJcXHByb2Rfe1xcTk59XFxwaV97aysxfShNX3tufSkiXSxbMywwLCJcXHByb2Rfe1xcTk59XFxwaV97aysxfShNX3tufSkiXSxbMiwxLCJcXHBpX3trfShNKSJdLFszLDEsIlxccHJvZF97XFxOTn1cXHBpX3trfShNX3tufSkiXSxbNCwxLCJcXGRvdHMiXSxbMCwxXSxbMSwyLCJcXHBpX3trKzF9KGYpIl0sWzIsMywiXFxwYXJ0aWFsIiwxXSxbMyw0LCJcXHBpX3trfShlKSIsMl0sWzQsNV1d
    \begin{tikzcd}
        \dots & {\prod_{\NN}\pi_{k+1}(M_{n})} && {\prod_{\NN}\pi_{k+1}(M_{n})} & \\
        && {\pi_{k}(M)} & {\prod_{\NN}\pi_{k}(M_{n})} & \dots
        \arrow[from=1-1, to=1-2]
        \arrow["{\pi_{k+1}(f)}", from=1-2, to=1-4]
        \arrow["\partial"{description}, from=1-4, to=2-3]
        \arrow["{\pi_{k}(e)}"', from=2-3, to=2-4]
        \arrow[from=2-4, to=2-5]
    \end{tikzcd}$$
    we obtain short exact sequences 
    \begin{equation}\label{eqn: Milnor SES}
        0\to\coker(\pi_{k+1}(f))\to\pi_{k}(M)\to\img(\pi_{k}(e))\simeq\ker(\pi_{k}(f))\to0.
    \end{equation}
    By exactness of countable products and the formation of homotopy objects preserving countable products, we get that $\ker(\pi_{k}(f))\simeq\lim_{n}\pi_{k}(M_{n})$. The identification of the cokernel with the higher derived limit $\coker(\pi_{k+1}(f))\simeq\lim_{n}^{1}\pi_{k+1}(M_{n})$ follows from \cite[Lem. 2.2]{RoosLimits}. 
\end{proof}
\begin{remark}\label{rmk: Mittag-Leffler under exact products}
    The Mittag--Leffler criterion persists under stronger exactness hypotheses. If small products are exact in the Grothendieck Abelian category $\Asf$, then $\lim_{n}^{1}$ vanishes on every tower in $\Asf$ with epimorphic transition maps \cite[Thm. 3.1]{RoosLimits}. In particular, if the $t$-structure on the $t^{+}$-bistable category $\Cscr$ is furthermore compatible with small products, so that small products are exact in $\Cscr^{\heartsuit}$ by the argument of \Cref{rmk: hearts are Grothendieck}, then the third term of (\ref{eqn: t plus bistable Milnor sequence}) computes $\pi_{k}(M)$ whenever the tower $\{\pi_{k+1}(M_{n})\}_{n}$ has epimorphic transition maps. 
\end{remark}
The derived-limit term in \Cref{prop: Milnor sequence in t-plus bistable} vanishes under a hypothesis weaker than exactness of small products.
\begin{definition}\label{def: lim1 acyclic categories}
    A $t^{+}$-bistable category $\Cscr$ is \emph{$\lim^{1}$-epi-acyclic} if every tower $\{A_{n}\}_{n\geq0}$ in $\Cscr^{\heartsuit}$ with epimorphic transition maps satisfies $\lim^{1}_{n}A_{n}\simeq0$. 
\end{definition}
\begin{remark}\label{rmk: 1-localic is lim1 epi acyclic}
    If $\Xscr$ is a replete 1-topos, then any $t^{+}$-bistable category whose heart is equivalent to Abelian sheaves on $\Xscr$ is $\lim^{1}$-epi-acyclic \cite[Prop. 3.1.10]{Proetale}. 
\end{remark}
Under this condition, a heart-preserving functor that is left $t$-exact is automatically $t$-exact.
\begin{proposition}\label{prop: lim1 acyclic t-exact}
    Let $\Cscr$ be a $t$-bistable category, $\Escr$ a $\lim^{1}$-epi-acyclic category, and $F:\Cscr\to\Escr$ an exact, limit-preserving, left $t$-exact functor. If $F(\Cscr^{\heartsuit})\subseteq\Escr^{\heartsuit}$ then $F$ is $t$-exact. 
\end{proposition}
\begin{proof}
    We need to show that $F$ is right $t$-exact. 

    Let $M\in\Cscr_{[0,n]}$ be connective and bounded above. Whitehead truncations yield (co)fiber sequences $\tau_{\geq k}M\to\tau_{\geq k-1}M\to\pi_{k-1}(M)[k-1]$, giving a (co)fiber sequence 
    $$F(\tau_{\geq k}M)\to F(\tau_{\geq k-1}M)\to F(\pi_{k-1}(M)[k-1])\simeq F(\pi_{k-1}(M))[k-1]$$
    where the final equivalence is from exactness and heart-preservation of $F$. Decreasing induction with base case $k=n$ exhibits each Whitehead truncation for $k\geq0$ as an extension of connective objects, the connective objects are closed under extensions, showing each $F(\tau_{\geq k}M)$ is connective, and in particular $F(M)\simeq F(\tau_{\geq0}M)$ is connective. 

    Now let $M\in\Cscr_{\geq0}$ be connective. By left completeness, we may write $M$ as its Postnikov tower $M\simeq\lim_{n}\tau_{\leq n}M$ and use the (co)fiber sequences from Postnikov truncations $\pi_{k}(M)[k]\to\tau_{\leq k}M\to\tau_{\leq k-1}M$. From the resulting (co)fiber sequence 
    $$F(\pi_{k}(M))[k]\to F(\tau_{\leq k}M)\to F(\tau_{\leq k-1}M)$$
    in $\Escr$, we note that $F(\tau_{\leq k}M)\to F(\tau_{\leq k-1}M)$ induces an equivalence on $\pi_{r}$ in $\Escr^{\heartsuit}$ for all $r\leq0$ and $k\geq 1$. As $F$ preserves limits, we have $F(M)\simeq\lim_{n}F(\tau_{\leq n}M)$. Using the Milnor sequence constructed in \Cref{prop: Milnor sequence in t-plus bistable}, we have a short exact sequence in $\Escr^{\heartsuit}$
    $$0\to\lim^{1}_{n}\pi_{r+1}(F(\tau_{\leq n}M))\to\pi_{r}(F(M))\to\lim_{n}\pi_{r}(F(\tau_{\leq n}M))\to0.$$
    Since $\Escr$ is $\lim^{1}$-epi-acyclic, the equivalences $F(\tau_{\leq k}M)\to F(\tau_{\leq k-1}M)$ on non-positive homotopy groups imply that $\pi_{r}(F(M))\simeq\lim_{n}\pi_{r}(F(\tau_{\leq n}M))$ for $r\leq-1$. $F$ sends bounded above connective objects to connective ones $\pi_{r}(F(\tau_{\leq n}M))\simeq0$ for $r\leq -1$ and thus $\pi_{r}(F(M))\simeq0$ for $r\leq -1$ showing that $F(M)$ is connective in $\Escr$, as desired. 
\end{proof}
\subsection{Symmetric monoidal structures}\label{subsec: symmetric monoidal structures}
The $t$-bistable categories of interest are frequently equipped with a symmetric monoidal structure compatible with the $t$-structure. We record that compatibility and its elementary consequences. 
\begin{definition}\label{def: symmetric monoidal t-bistable category}
    Let $\Cscr$ be a stable presentable category equipped with a $t$-structure and symmetric monoidal structure. 
    \begin{enumerate}
        \item $\Cscr$ is a \emph{presentably symmetric monoidal $t$-bistable category} if it is presentably symmetric monoidal with underlying category $t$-bistable and the symmetric monoidal structure on $\Cscr$ restricts to one on $\Cscr_{\geq0}$.
        \item $\Cscr$ is a \emph{presentably symmetric monoidal $t^{+}$-bistable category} if it is presentably symmetric monoidal $t$-bistable and its underlying category is $t^{+}$-bistable. 
    \end{enumerate}  
\end{definition}
\begin{remark}\label{rmk: t-bistable sym mon depends is structure}
    Relative to a fixed symmetric monoidal structure and $t$-structure, the conditions of \Cref{def: symmetric monoidal t-bistable category} -- that the tensor product preserve colimits in each variable, that the $t$-structure be both left and right complete, and that $\Cscr_{\geq0}$ be closed under tensor products with connective unit -- are properties.
\end{remark}
\begin{remark}\label{rmk: Postnikov is symmetric monoidal}
    In the language of \cite[Ex. 2.2.1.3]{HA}, the $t$-structure on the $t$-bistable (resp. $t^{+}$-bistable) category $\Cscr$ is compatible with the symmetric monoidal structure. Under this hypothesis, the Postnikov truncation functors $\tau_{\leq n}$ become symmetric monoidal functors $\Cscr_{\geq0}\to\Cscr_{[0,n]}$ \cite[Ex. 2.2.1.10]{HA}.
\end{remark}
A recurring device compares such a category with modules over the static commutative algebra object $\pi_{0}(1_{\Cscr})$. We begin with the case of a general connective commutative algebra object $A$. 
\begin{lemma}\label{lem: modules inherit bistability}
    Let $\Cscr$ be a presentably symmetric monoidal $t$-bistable category, $A\in\CAlg(\Cscr_{\geq0})$, and $\Escr:=\Mod_{A}(\Cscr)$. Then:
    \begin{enumerate}
        \item Denote by $f^{*}:\Cscr\to\Escr$ the scalar extension functor and by $f_{*}$ its right adjoint. The category $\Escr$ is itself a presentably symmetric monoidal $t$-bistable category with respect to which the forgetful functor $f_{*}:\Escr\to\Cscr$ is conservative, $t$-exact, and preserves small colimits. If $\Cscr$ is $t^{+}$-bistable, then the same holds for $\Escr$. 
        \item For all $M\in\Cscr,N\in\Escr$ there is an equivalence 
        $$f_{*}\intHom_{\Escr}(f^{*}M,N)\simeq\intHom_{\Cscr}(M, f_{*}N).$$
    \end{enumerate}
\end{lemma}
\begin{proof}[Proof of (1)]
    Since the $t$-structure on $\Cscr$ is accessible, \cite[Prop. A.15]{CartierModules} endows $\Escr$ with the induced accessible $t$-structure compatible with the symmetric monoidal structure, in the sense that $\Escr_{\geq0}\hookrightarrow\Escr$ is a symmetric monoidal functor. Moreover, $f_{*}$ is conservative \cite[Prop. 4.8.5.8 (5)]{HA} and $t$-exact and, by \cite[Cor. 4.2.3.7]{HA}, preserves small limits and colimits. 

    For any filtered diagram of coconnective objects, coconnectivity of the colimit can be checked after applying the $t$-exact functor $f_{*}$. Since the functor is $t$-exact and colimit-preserving, compatibility of the $t$-structure on $\Escr$ with filtered colimits follows from compatibility of the $t$-structure on $\Cscr$ with filtered colimits. Similarly, $t$-exactness implies $f_{*}$ preserves Whitehead and Postnikov towers in $\Escr$ so by conservativity their convergence can be checked in $\Cscr$ after applying $f_{*}$. Thus left and right completeness of $\Cscr$ implies the same for $\Escr$. 

    Finally, if $\Cscr$ is $t^{+}$-bistable, then so is $\Escr$ as compatibility with countable products holds in $\Escr$ by \cite[Prop. A.15 (2)]{CartierModules}. 
\end{proof}
\begin{proof}[Proof of (2)]
    For $P\in\Cscr$ we compute 
    \begin{align*}
        \Hom_{\Cscr}(P, \intHom_{\Cscr}(M, f_{*}N)) &\simeq \Hom_{\Cscr}(P\otimes_{\Cscr}M,f_{*}N) && \otimes_{\Cscr}\dashv\intHom_{\Cscr} \\
        &\simeq \Hom_{\Escr}(f^{*}(P\otimes_{\Cscr}M),N) && f^{*}\dashv f_{*} \\
        &\simeq \Hom_{\Escr}(f^{*}P\otimes_{\Escr}f^{*}M,N) && f^{*}\text{ sym. mon.} \\
        &\simeq \Hom_{\Escr}(f^{*}P,\intHom_{\Escr}(f^{*}M,N)) && \otimes_{\Escr}\dashv \intHom_{\Escr} \\
        &\simeq \Hom_{\Cscr}(P, f_{*}\intHom_{\Escr}(f^{*}M, N)) && f^{*}\dashv f_{*}
    \end{align*}
    and we conclude by the Yoneda lemma. 
\end{proof}
In the case of the connective commutative algebra object $\pi_{0}(1_{\Cscr})$, the comparison provides a convenient tool to reduce questions about $\Cscr$ to the module category over $\pi_{0}(1_{\Cscr})$. In particular, scalar extension induces an equivalence on hearts, is well behaved on internal homs, and is conservative on bounded below objects. 
\begin{proposition}\label{prop: relationship with heart}
    Let $\Cscr$ be a presentably symmetric monoidal $t$-bistable category with unit $1_{\Cscr}$ and $\Escr:=\Mod_{\pi_{0}(1_{\Cscr})}(\Cscr)$. Then:
    \begin{enumerate}
        \item Denote by $f^{*}:\Cscr\to\Escr$ the scalar extension functor and by $f_{*}$ its right adjoint. The functor $f^{*}$ induces an equivalence $f^{\heartsuit,*}:\Cscr^{\heartsuit}\to\Escr^{\heartsuit}$ given on objects by $M\mapsto\pi_{0}^{\Escr}(M\otimes_{\Cscr}\pi_{0}^{\Cscr}(1_{\Cscr}))$, with inverse the restriction of $f_{*}$. 
        \item The functor $f^{*}:\Cscr\to\Escr$ is conservative on bounded below objects: if $g:M\to N$ is a morphism of bounded below objects in $\Cscr$ such that $f^{*}(g):f^{*}M\to f^{*}N$ is an equivalence in $\Escr$ then $g$ is an equivalence in $\Cscr$. 
        \item For $M,N\in\Cscr^{\heartsuit}$, $\intHom_{\Cscr}(M,N)\in\Cscr_{\leq0}$.
    \end{enumerate}
\end{proposition}
\begin{proof}[Proof of (1)]
    This is \cite[Lem. A.13 \& Prop. A.15 (3)]{CartierModules}. 
\end{proof}
\begin{proof}[Proof of (2)]
    By shifting we may take $M$ and $N$ to be connective. Let $M\xrightarrow{g}N\to C$ be the associated (co)fiber sequence in $\Cscr$, so $C\in\Cscr_{\geq0}$ as a cofiber of connective objects. Exactness of the tensor product identifies $\pi_{0}^{\Cscr}(1_{\Cscr})\otimes_{\Cscr}C$ with the cofiber of $f^{*}(g)$ computed in $\Cscr$, which vanishes as $f^{*}(g)$ is an equivalence. 

    We claim $C\in\Cscr_{\geq n}$ for all $n\geq0$, arguing by induction with base case $n=0$ established above. Suppose $C\in\Cscr_{\geq n}$, so $C[-n]$ is connective and $\pi_{0}^{\Cscr}(1_{\Cscr})\otimes_{\Cscr}C[-n]\simeq(\pi_{0}^{\Cscr}(1_{\Cscr})\otimes_{\Cscr}C)[-n]\simeq0$. Since $\tau_{\leq0}:\Cscr_{\geq0}\to\Cscr^{\heartsuit}$ is symmetric monoidal by \Cref{rmk: Postnikov is symmetric monoidal}, we compute in $\Cscr^{\heartsuit}$
    $$0\simeq\pi_{0}(\pi_{0}^{\Cscr}(1_{\Cscr})\otimes_{\Cscr}C[-n])\simeq\pi_{0}^{\Cscr}(1_{\Cscr})\otimes_{\Cscr^{\heartsuit}}\pi_{0}^{\Cscr}(C[-n])\simeq\pi_{0}^{\Cscr}(C[-n])\simeq\pi_{n}^{\Cscr}(C)$$
    where the third equivalence holds as $\pi_{0}^{\Cscr}(1_{\Cscr})$ is the unit of $\Cscr^{\heartsuit}$. So $\tau_{\leq n}C\simeq\pi_{n}(C)[n]\simeq0$ and $C\in\Cscr_{\geq n+1}$ for all $n\geq0$. Left separatedness -- which holds by \Cref{lem: separatedness and prestable} (1) -- gives $C\simeq0$, so $g$ is an equivalence.
\end{proof}
\begin{proof}[Proof of (3)]
    As $\intHom_{\Cscr}(M,-)$ is right adjoint to the right $t$-exact $M\otimes_{\Cscr}-$ it is left $t$-exact, so $\intHom_{\Cscr}(M,N)\in\Cscr_{\leq0}$ for $N\in\Cscr^{\heartsuit}\subseteq\Cscr_{\leq0}$.
\end{proof}
An object was called compact projective in \Cref{def: compact projective objects} when the spectrum-valued $\Hom_{\Cscr}(P,-)$ is $t$-exact and preserves small colimits. With a closed symmetric monoidal structure now at hand, the same two conditions may be asked of the internal hom $\intHom_{\Cscr}(M,-):\Cscr\to\Cscr$. It is convenient to name these separately.
\begin{definition}\label{def: internally compact projective}
    Let $\Cscr$ be a presentably symmetric monoidal $t$-bistable category. 
    \begin{enumerate}
        \item An object $M\in\Cscr$ is \emph{internally compact} if $\intHom_{\Cscr}(M,-)$ preserves small colimits. 
        \item An object $M\in\Cscr_{\geq0}$ is \emph{internally projective} if the left $t$-exact functor $\intHom_{\Cscr}(M,-)$ is $t$-exact as an endofunctor of $\Cscr$. 
        \item An object $M\in\Cscr_{\geq0}$ is \emph{internally compact projective} if it is both internally compact and internally projective. 
    \end{enumerate}
\end{definition}
Internal compact projectivity is inherited by the unit and preserved under finite (co)products and tensor products, and implies ordinary compact projectivity when the unit is compact projective.
\begin{lemma}\label{lem: permanence of internal compact projectivity}
    Let $\Cscr$ be a presentably symmetric monoidal $t$-bistable category. 
    \begin{enumerate}
        \item The unit $1_{\Cscr}$ is internally compact projective, and the internally compact projective objects are closed under finite (co)products and tensor products. 
        \item If the unit $1_{\Cscr}$ is compact projective, then every internally compact projective object is compact projective. 
    \end{enumerate}
\end{lemma}
\begin{proof}[Proof of (1)]
    $1_{\Cscr}$ is internally compact projective as $\intHom_{\Cscr}(1_{\Cscr},-)\simeq\id_{\Cscr}$ which is $t$-exact and preserves small colimits. In the case of finite (co)products we can reduce to the case of binary ones, where it follows $\intHom_{\Cscr}(-,-)$ preserving finite (co)products in the first argument and $t$-exactness of finite (co)products. Now let $M,N$ be internally compact projective. Then $\intHom_{\Cscr}(M\otimes_{\Cscr}N,-)\simeq\intHom_{\Cscr}(M,\intHom_{\Cscr}(N,-))$ is the composite of two $t$-exact functors which preserve small colimits, hence is $t$-exact and preserves small colimits.
\end{proof}
\begin{proof}[Proof of (2)]
    Let $M$ be internally compact projective. The mapping spectrum factors as $\Hom_{\Cscr}(M,-)\simeq\Hom_{\Cscr}(1_{\Cscr},\intHom_{\Cscr}(M,-))$. Now $\intHom_{\Cscr}(M,-)$ is $t$-exact and preserves small colimits as $M$ is internally compact projective, and $\Hom_{\Cscr}(1_{\Cscr},-)$ is $t$-exact and preserves small colimits as $1_{\Cscr}$ is compact projective, so their composite $\Hom_{\Cscr}(M,-)$ is $t$-exact and preserves small colimits. Thus $M$ is compact projective.
\end{proof}
Internal compactness is moreover implied by ordinary compactness whenever the compact objects are closed under tensoring. This uses neither the $t$-structure nor a compact unit, so we state it for a stable $\omega$-presentable category. 
\begin{lemma}\label{lem: compact implies internally compact}
   Let $\Cscr$ be a stable $\omega$-presentable presentably symmetric monoidal category with compact unit in which the compact objects are closed under tensor product. Then every compact object is internally compact. 
\end{lemma}
\begin{proof}
    For each compact generator $M$ and compact object $N$ we have a natural equivalence $\Hom_{\Cscr}(M,\intHom_{\Cscr}(N,-))\simeq\Hom_{\Cscr}(M\otimes_{\Cscr}N,-)$ which preserves small colimits as $M\otimes_{\Cscr}N$ is compact so $\intHom_{\Cscr}(N,-)$ commutes with small colimits, as desired. 
\end{proof}
\begin{remark}\label{rmk: internally compact notation clarification}
    Equivalently, \Cref{lem: compact implies internally compact} holds for stable $\omega$-presentable presentably symmetric monoidal category in which binary tensor products of compact objects are compact, and in particular for stable $\omega$-presentably symmetric monoidal categories. See \Cref{subsec: notation and conventions} for the conventions we adopt. 
\end{remark}
This relationship holds not only within a fixed category, but along free functors to the module category over some connective commutative algebra object in $\Cscr$. 
\begin{proposition}\label{prop: scalar extension preserves internal compact projectives}
    Let $\Cscr$ be a presentably symmetric monoidal $t$-bistable category and $A\in\CAlg(\Cscr_{\geq0})$.
    \begin{enumerate}
        \item If $M$ is internally compact projective in $\Cscr$ then $A\otimes_{\Cscr}M$ is internally compact projective in $\Mod_{A}(\Cscr)$. 
        \item If $\Cscr$ is generated under small colimits and shifts by a set $\{P_{i}\}_{i\in I}$ of internally compact projective objects, then $\Mod_{A}(\Cscr)$ is generated under small colimits and shifts by $\{A\otimes_{\Cscr}P_{i}\}_{i\in I}$. 
    \end{enumerate}
\end{proposition}
\begin{proof}[Proof of (1)]
    Using \Cref{lem: modules inherit bistability} (2), this follows from $f_{*}$ preserving small colimits and being $t$-exact. 
\end{proof}
\begin{proof}[Proof of (2)]
    The objects $A\otimes_{\Cscr}P_{i}$ for varying $i\in I$ are internally compact projective by (1). These objects generate because, by adjunction, $\Hom_{\Cscr}(P_{i},f_{*}(-))\simeq\Hom_{\Mod_{A}(\Cscr)}(A\otimes_{\Cscr}P_{i},-)$ are jointly conservative as functors to $\Sp$ by conservativity of $f_{*}$ and $\Hom_{\Cscr}(P_{i},-)$ being jointly conservative in $\Cscr$. 
\end{proof}
This completes the framework of presentably symmetric monoidal $t$-bistable categories and their internal notions of compactness and projectivity. We turn now to the localizations of such categories.
\subsection{Localizations of stable presentable categories}\label{subsec: localizations of stable categories}
Localizations of $t$-bistable categories arise repeatedly in the main text. We isolate the class relevant to our purposes, whose defining feature is that the reflection is compatible with the connective part of the $t$-structure.
\begin{definition}\label{def: pi stably reflective}
    Let $\Cscr$ be a $t$-bistable category. A full subcategory $\Cscr^{L}\subseteq\Cscr$ is \emph{$\pi$-stably reflective} if $\Cscr^{L}$ is closed under small limits and small colimits in $\Cscr$ and the inclusion $\Cscr^{L}\hookrightarrow\Cscr$ admits a left adjoint $L:\Cscr\to\Cscr^{L}$ which is right $t$-exact regarded as an endofunctor of $\Cscr$.  
\end{definition}
Before turning to the induced $t$-structure on the local objects, we record the closure property that makes the argument go through in its natural generality. Membership in a full subcategory is detected by retracts as soon as the subcategory is idempotent complete, with no reference to any additional structure.
\begin{lemma}\label{lem: retracts in full subcategories}
    Let $\Cscr$ be a category and $\Escr\subseteq\Cscr$ an idempotent complete full subcategory closed under equivalences. If an object $C\in\Cscr$ admits a retraction $C\xrightarrow{i}E\xrightarrow{r}C$ in $\Cscr$ with $E\in\Escr$, then $C\in\Escr$. 
\end{lemma}
\begin{proof}
    Recall the simplicial sets $\Idem\subseteq\Idem^{+}\supseteq\Ret$ of \cite[Def. 4.4.5.2]{HTT}. The simplicial set $\Idem^{+}$ has exactly two vertices $x$ and $y$, of which $x$ is the unique vertex of $\Idem$, and $\Ret$ is the quotient of $\Delta^{2}$ collapsing the edge $\Delta^{\{0,2\}}$ to $y$. A map $\Ret\to\Cscr$ is then the datum of a $2$-simplex witnessing a retraction, so the given retraction determines a weak retraction diagram $\overline{F}:\Ret\to\Cscr$ with $\overline{F}(x)=E$ and $\overline{F}(y)=C$ in the sense of \cite[Def. 4.4.5.4]{HTT}. The inclusion $\Ret\subseteq\Idem^{+}$ is inner anodyne by \cite[Prop. 4.4.5.6]{HTT}, so the restriction map $\Fun(\Idem^{+},\Cscr)\to\Fun(\Ret,\Cscr)$ is a trivial fibration by \cite[Cor. 4.4.5.7]{HTT}. Choose an extension $F:\Idem^{+}\to\Cscr$ of $\overline{F}$ and set
    $$\widetilde{F}:=F|_{\Idem}:\Idem\to\Cscr,$$
    a coherent idempotent carried by $E$ with underlying endomorphism $i\circ r:E\to E$. 

    We claim $\widetilde{F}$ factors through $\Escr$. As $\Escr\subseteq\Cscr$ is full, a simplex of $\Cscr$ lies in $\Escr$ if and only if each of its vertices does. The simplicial set $\Idem$ has the single vertex $x$ by \cite[Rem. 4.4.5.3]{HTT}, and $\widetilde{F}(x)=E\in\Escr$, so every simplex in the image of $\widetilde{F}$ has all its vertices equal to $E$. Thus $\widetilde{F}=\iota\circ G$ for a unique $G:\Idem\to\Escr$, where $\iota:\Escr\to\Cscr$ is the inclusion. Since $\Escr$ is idempotent complete, $G$ is effective, so extends to a strong retraction diagram $H:\Idem^{+}\to\Escr$ \cite[Cor. 4.4.5.14 \& 4.4.5.15]{HTT}. 

    Set $E':=H(y)\in\Escr$ and let $\Dscr\subseteq\Fun(\Idem,\Cscr)$ be the full subcategory spanned by the effective idempotents. The idempotent $\widetilde{F}$ is effective, being the restriction of $F$, so lies in $\Dscr$, as do the restrictions of $F$ and $\iota\circ H$ to $\Idem$, both of which recover $\widetilde{F}$. The restriction map $\Fun(\Idem^{+},\Cscr)\to\Dscr$ is a trivial Kan fibration by \cite[Cor. 4.4.5.16]{HTT}, so its fiber over $\widetilde{F}$ is a contractible Kan complex, and $F\simeq\iota\circ H$ in $\Fun(\Idem^{+},\Cscr)$. Evaluating at $y$ gives
    $$C=F(y)\simeq(\iota\circ H)(y)=E'$$
    in $\Cscr$. As $\Escr$ is closed under equivalences and $E'\in\Escr$, we conclude $C\in\Escr$. 
\end{proof}
This is slightly more than the retract closure recorded in \cite[\S 4.4.5]{HTT}, but it is elementary and streamlines the verification that $\Cscr^{L}$ is closed under retracts, reducing it to idempotent completeness of the subcategory of local objects $\Cscr^{L}$. With it in hand, the $\pi$-stably reflective condition shows that the local objects admit the induced $t$-structure.
\begin{proposition}\label{prop: properties of pi stably reflective localizations}
    Let $\Cscr$ be a $t$-bistable category and $\Cscr^{L}\subseteq\Cscr$ a $\pi$-stably reflective subcategory. 
    \begin{enumerate}
        \item $\Cscr^{L}$ is itself a stable presentable category closed under extensions and retracts in $\Cscr$. 
        \item $\Cscr^{L}$ admits the induced $t$-structure from $\Cscr$ with respect to which $\Cscr^{L}$ is itself $t$-bistable. 
        \item $M\in\Cscr^{L}$ if and only if $\pi_{k}(M)\in(\Cscr^{L})^{\heartsuit}$ for all $k\in\ZZ$. 
        \item If $\Cscr$ is furthermore $t^{+}$-bistable then $\Cscr^{L}$ is itself $t^{+}$-bistable with respect to the induced $t$-structure. 
    \end{enumerate}
\end{proposition}
\begin{proof}[Proof of (1)]
    Closure under small limits and colimits implies $\Cscr^{L}$ is a stable category \cite[Lem. 1.1.3.3]{HA} and presentability follows from \cite[Thm. 1.1]{RagimovReflection}. 

    For closure under extensions, we can contemplate the diagram 
    \begin{equation}\label{eqn: closed under extensions diagram}
        % https://q.uiver.app/#q=WzAsNixbMCwwLCJNX3sxfSJdLFsyLDAsIk1fezJ9Il0sWzAsMSwiMCJdLFsyLDEsIk1fezN9Il0sWzQsMCwiMCJdLFs0LDEsIk1fezF9WzFdIl0sWzAsMV0sWzEsNF0sWzQsNV0sWzEsM10sWzMsNV0sWzAsMl0sWzIsM11d
        \begin{tikzcd}
            {M_{1}} && {M_{2}} && 0 \\
            0 && {M_{3}} && {M_{1}[1]}
            \arrow[from=1-1, to=1-3]
            \arrow[from=1-1, to=2-1]
            \arrow[from=1-3, to=1-5]
            \arrow[from=1-3, to=2-3]
            \arrow[from=1-5, to=2-5]
            \arrow[from=2-1, to=2-3]
            \arrow[from=2-3, to=2-5]
        \end{tikzcd}
    \end{equation}
    in $\Cscr$ where all squares are biCartesian. Since $M_{1},M_{3}\in\Cscr^{L}$ and $\Cscr^{L}$ is closed under colimits, $M_{1}[1]\in\Cscr^{L}$. Closure under limits then places $M_{2}$, the fiber of $M_{3}\to M_{1}[1]$ in $\Cscr$, in $\Cscr^{L}$.

    Cocompleteness of $\Cscr$ and $\Cscr^{L}$ imply they are idempotent complete \cite[Rmk. 4.4.5.3 \& Cor. 4.4.5.15]{HTT}. Thus $\Cscr^{L}$ is closed under retracts by \Cref{lem: retracts in full subcategories}. 
\end{proof}
\begin{proof}[Proof of (2)]
    To show that $\Cscr^{L}$ admits the induced $t$-structure, it suffices to show that $\Cscr^{L}$ is closed under the Postnikov truncation functor $\tau_{\leq0}$ in $\Cscr$ \cite[\S 1.3.19 \& \S 2]{Perverse}. Fix $M\in\Cscr^{L}$. Since $\tau_{\leq0}M\simeq\cofib(\tau_{\geq 1}M\to M)$ and $\Cscr^{L}$ is closed under colimits, it suffices to show that $\tau_{\geq 1}M\in\Cscr^{L}$. We exhibit $\tau_{\geq 1}M$ as a retract of $L(\tau_{\geq 1}M)$ in which case (1) implies $\tau_{\geq 1}M\in\Cscr^{L}$. 

    Write $c:\tau_{\geq 1}M\to M$ for the counit of the Whitehead truncation with coconnective fiber $(\tau_{\leq 0}M)[-1]$. Since $L$ is exact and right $t$-exact, $L(\tau_{\geq 1}M)\in\Cscr_{\geq 1}$. Applying $\Hom_{\Cscr}(L(\tau_{\geq 1}M),-)$ to the (co)fiber sequence $(\tau_{\leq 0}M)[-1]\to\tau_{\geq 1}M\to M$ gives an isomorphism on $\pi_{0}$ from $\Hom_{\Cscr}(L(\tau_{\geq 1}M),\tau_{\geq 1}M)$ to $\Hom_{\Cscr}(L(\tau_{\geq 1}M),M)$, as the fiber $\Hom_{\Cscr}(L(\tau_{\geq 1}M),(\tau_{\leq0}M)[-1])$ lies in $\Sp_{\leq -2}$. Indeed, $\Hom_{\Cscr}(L(\tau_{\geq 1}M),\tau_{\leq 0}M)$ is coconnective as $L(\tau_{\geq 1}M)\in\Cscr_{\geq 0}$ and $\tau_{\leq 0}M\in\Cscr_{\leq 0}$, and vanishes in degree $0$ as $L(\tau_{\geq 1}M)\in\Cscr_{\geq 1}$, so $\tau_{\leq0}M$ lies in $\Sp_{\leq -1}$, and shifting gives $(\tau_{\leq0}M)[-1]\in\Sp_{\leq -2}$. There exists a unique morphism $\beta:L(\tau_{\geq 1}M)\to \tau_{\geq 1}M$ such that $c\circ\beta\simeq\alpha$ where $\alpha:L(\tau_{\geq 1}M)\to M$ is the composite of $L(c)$ with the inverse of the unit $\eta_{M}:M\to L(M)$, which is invertible since $M\in\Cscr^{L}$. 

    We claim $\beta\circ\eta_{\tau_{\geq 1}M}\simeq\id_{\tau_{\geq 1}M}$. Using the naturality square involving the downward-pointing vertical arrows
    \begin{equation}\label{eqn: retraction square}
        % https://q.uiver.app/#q=WzAsNCxbMCwwLCJcXHRhdV97XFxnZXEgMX1NIl0sWzIsMCwiTSJdLFsyLDEsIkwoTSkiXSxbMCwxLCJMKFxcdGF1X3tcXGdlcSAxfU0pIl0sWzAsMSwiYyJdLFszLDIsIkwoYykiLDJdLFsxLDIsIlxcZXRhX3tNfSIsMix7ImN1cnZlIjoxfV0sWzAsMywiXFxldGFfe1xcdGF1X3tcXGdlcSAxfU19IiwyXSxbMiwxLCJcXGV0YV97TX1eey0xfSIsMix7ImN1cnZlIjoxfV1d
        \begin{tikzcd}
            {\tau_{\geq 1}M} && M \\
            {L(\tau_{\geq 1}M)} && {L(M).}
            \arrow["c", from=1-1, to=1-3]
            \arrow["{\eta_{\tau_{\geq 1}M}}"', from=1-1, to=2-1]
            \arrow["{\eta_{M}}"', curve={height=6pt}, from=1-3, to=2-3]
            \arrow["{L(c)}"', from=2-1, to=2-3]
            \arrow["{\eta_{M}^{-1}}"', curve={height=6pt}, from=2-3, to=1-3]
        \end{tikzcd}
    \end{equation}
    we compute 
    $$c\circ \beta\circ\eta_{\tau_{\geq 1}M}\simeq\alpha\circ\eta_{\tau_{\geq 1}M}\simeq\eta_{M}^{-1}\circ L(c)\circ\eta_{\tau_{\geq 1}M}\simeq\eta_{M}^{-1}\circ\eta_{M}\circ c\simeq c.$$
    Since $\tau_{\geq 1}M\in\Cscr_{\geq1}$, the same argument applied to the (co)fiber sequence $(\tau_{\leq0}M)[-1]\to\tau_{\geq 1}M\xrightarrow{c}M$ with $\Hom_{\Cscr}(\tau_{\geq 1}M,-)$ in place of $\Hom_{\Cscr}(L(\tau_{\geq 1}M),-)$ shows that the map
    $$\Hom_{\Cscr}(\tau_{\geq 1}M,\tau_{\geq 1}M)\to\Hom_{\Cscr}(\tau_{\geq 1}M,M),\qquad g\mapsto c\circ g,$$
    is an isomorphism on $\pi_{0}$, the fiber $\Hom_{\Cscr}(\tau_{\geq 1}M,(\tau_{\leq0}M)[-1])$ lying in $\Sp_{\leq -2}$ as before. In particular, it is injective on $\pi_{0}$, so $\beta\circ\eta_{\tau_{\geq 1}M}$ and $\id_{\tau_{\geq 1}M}$, which the preceding computation sends to the common value $c$, are homotopic. This exhibits $\tau_{\geq 1}M$ as a retract of $L(\tau_{\geq 1}M)\in\Cscr^{L}$, so $\tau_{\geq 1}M\in\Cscr^{L}$ by (1). Using the (co)fiber sequence $\tau_{\geq 1}M\to M\to\tau_{\leq 0}M$ and closure of $\Cscr^{L}$ under colimits, we conclude $\tau_{\leq 0}M\in\Cscr^{L}$, which is what we wanted to show.

    Since $\Cscr^{L}$ is equipped with the induced $t$-structure and the inclusion $\Cscr^{L}\hookrightarrow\Cscr$ preserves and reflects small limits and colimits, each defining property of a $t$-bistable category is inherited by $\Cscr^{L}$ from $\Cscr$.
\end{proof}
\begin{proof}[Proof of (3)]
    $(\Rightarrow)$ If $M\in\Cscr^{L}$ then $\pi_{k}(M)\in(\Cscr^{L})^{\heartsuit}$ for all $k\in\ZZ$, since $\Cscr^{L}$ admits the induced $t$-structure.

    $(\Leftarrow)$ By right completeness of $\Cscr$, we can write $M$ as the colimit of its Whitehead tower, and each Whitehead truncation again satisfies the hypothesis, its homotopy objects being among those of $M$. Since $\Cscr^{L}$ is closed under small colimits in $\Cscr$, it suffices to treat $M$ bounded below, say $M\in\Cscr_{\geq r}$. By left completeness, we have $M\simeq\lim_{k}\tau_{\leq k}M$, and since $\Cscr^{L}$ is closed under limits it suffices to show that each Postnikov truncation $\tau_{\leq k}M$ lies in $\Cscr^{L}$. We argue by induction on $k\geq r$ using the (co)fiber sequences 
    $$\pi_{k}(M)[k]\to\tau_{\leq k}M\to\tau_{\leq k-1}M,$$
    which exhibit $\tau_{\leq k}M$ as an extension of $\tau_{\leq k-1}M$ by $\pi_{k}(M)[k]$. In the base case $k=r$, the boundedness $M\in\Cscr_{\geq r}$ gives $\tau_{\leq r}M\simeq\pi_{r}(M)[r]$ which lies in $\Cscr^{L}$ because $\pi_{r}(M)\in(\Cscr^{L})^{\heartsuit}$ and $\Cscr^{L}$ is closed under shifts. For the inductive step, $\pi_{k}(M)[k]\in\Cscr^{L}$ by hypothesis and $\tau_{\leq k-1}M\in\Cscr^{L}$ by the inductive hypothesis, so $\tau_{\leq k}M\in\Cscr^{L}$ by closure under extensions in (1). Hence each $\tau_{\leq k}M$ lies in $\Cscr^{L}$ and therefore so does $M\simeq\lim_{k}\tau_{\leq k}M$.
\end{proof}
\begin{proof}[Proof of (4)]
    The only condition left to verify atop those shown in (2) is compatibility with countable products. Since $\Cscr^{L}$ admits the induced $t$-structure and the inclusion $\Cscr^{L}\hookrightarrow\Cscr$ reflects limits, compatibility of the induced $t$-structure on $\Cscr^{L}$ with countable products follows from compatibility of the $t$-structure on $\Cscr$ with countable products. 
\end{proof}
We extend these considerations to the symmetric monoidal setting. The following criterion of Lurie gives conditions under which a localization is compatible with the tensor product.
\begin{proposition}[{\cite[Prop. 2.2.1.9]{HA}}]\label{prop: sym mon localizations}
    Let $\Cscr^{\otimes}$ be a symmetric monoidal category and $L:\Cscr\to\Cscr^{L}$ a localization of the underlying category such that for all $L$-equivalences $f:M\to N$ and all $P\in\Cscr$, the induced morphism $P\otimes_{\Cscr}M\xrightarrow{\id_{P}\otimes f}P\otimes_{\Cscr}N$ is an $L$-equivalence. Then $\Cscr^{L}$ admits a unique symmetric monoidal structure with respect to which $L$ assembles into a symmetric monoidal functor $L^{\otimes}:\Cscr^{\otimes}\to(\Cscr^{L})^{\otimes}$. 
\end{proposition}
\begin{remark}\label{rmk: inclusion is lax}
    In the situation of \Cref{prop: sym mon localizations}, the inclusion $\Cscr^{L}\hookrightarrow\Cscr$ is lax symmetric monoidal with respect to the respective symmetric monoidal structures. 
\end{remark}
We name the localizations meeting this hypothesis in the setting of interest.
\begin{definition}\label{def: compatible localization}
    Let $\Cscr$ be a presentably symmetric monoidal $t$-bistable category and $\Cscr^{L}$ a $\pi$-stably reflective subcategory. The localization functor $L:\Cscr\to\Cscr^{L}$ is \emph{$\otimes_{\Cscr}$-compatible} if it satisfies the hypothesis of \Cref{prop: sym mon localizations}: $L$-equivalences are preserved by tensoring. 
\end{definition}
Compatibility promotes the induced structure on $\Cscr^{L}$ to a monoidal one.
\begin{proposition}\label{prop: compatible localization implies pres sym mon t-bistable}
    Let $\Cscr$ be a presentably symmetric monoidal $t$-bistable category and $\Cscr^{L}$ a $\pi$-stably reflective subcategory with $\otimes_{\Cscr}$-compatible localization $L:\Cscr\to\Cscr^{L}$. 
    \begin{enumerate}
        \item $\Cscr^{L}$ is a presentably symmetric monoidal $t$-bistable category. 
        \item If $\Cscr$ was furthermore $t^{+}$-bistable, then $\Cscr^{L}$ is a presentably symmetric monoidal $t^{+}$-bistable category. 
    \end{enumerate}
\end{proposition}
\begin{proof}[Proof of (1)]
    Note that the underlying category $\Cscr^{L}$ is $t$-bistable with respect to the induced $t$-structure by \Cref{prop: properties of pi stably reflective localizations}. The symmetric monoidal structure on $\Cscr^{L}$ is given by $M\otimes_{\Cscr^{L}}N:=L(M\otimes_{\Cscr}N)$ \cite[Prop. 2.2.1.9]{HA}, and $L$ is right $t$-exact, so the tensor product on $\Cscr^{L}$ restricts to one on $\Cscr^{L}_{\geq0}$. 
\end{proof}
\begin{proof}[Proof of (2)]
    This follows readily from (1) since $t^{+}$-bistability is a condition on the underlying category which is satisfied by \Cref{prop: properties of pi stably reflective localizations} (4). 
\end{proof}
\subsection{Recognizing symmetric monoidal localizations with smashing kernel}\label{subsec: detecting smashing kernel}
We characterize when a $\otimes_{\Cscr}$-compatible localization of a $t$-bistable category is an open immersion, that is, has kernel the category of modules over an idempotent algebra \cite[Prop. 6.5]{Complex}.
\begin{definition}\label{def: categorical open immersion}
    A morphism $f^{*}:\Cscr\to\Escr$ in $\CAlg(\PrL_{\St})$ with a fully faithful right adjoint is an \emph{open immersion} if it satisfies either of the following equivalent conditions: 
    \begin{enumerate}[label=(\alph*)]
        \item $f^{*}$ admits a fully faithful left adjoint $f_{\sharp}$ that satisfies the projection formula: for all $M\in\Cscr$ and $N\in\Escr$ the natural morphism $f_{\sharp}(f^{*}M\otimes_{\Escr}N)\to M\otimes_{\Cscr}f_{\sharp}N$ is an equivalence in $\Cscr$.
        \item The kernel $\ker(f^{*}):=\{K\in\Cscr:f^{*}K\simeq0\}\subseteq\Cscr$ is equivalent, as a full subcategory of $\Cscr$, to $\Mod_{Q}(\Cscr)$ for some idempotent algebra $Q\in\Idem(\Cscr)$.
    \end{enumerate}
\end{definition}
The two formulations are connected by the adjoint functor theorem: an open immersion $f^{*}$ admits a left adjoint and so preserves limits. Preservation of limits is one of the two structural inputs the kernel analysis requires -- the other being homotopy-detection which we define below. 
\begin{definition}\label{def: kernel detected homotopy}
    Let $L:\Cscr\to\Cscr^{L}$ be a $\pi$-stably reflective localization of a $t$-bistable category. Its kernel is \emph{homotopy detected} if it admits the induced $t$-structure: for each $K\in\ker(L)$ and every $n\in\ZZ$, both $\tau_{\leq n}K$ and $\tau_{\geq n}K$ lie in $\ker(L)$. 
\end{definition}
\begin{remark}\label{rmk: homotopy detection of kernel}
    Homotopy detection of $\ker(L)$ is implied by, but weaker than $t$-exactness of $L$. If $L$ is $t$-exact, then it commutes with the truncation functors, so $L(K)\simeq0$ forces $L(\tau_{\leq n}K)\simeq\tau_{\leq n}L(K)\simeq0$ and $\ker(L)$ is closed under truncation. The converse need not hold: we only require the comparison map for truncations to be an equivalence for all $K\in\ker(L)$, but not necessarily for all objects of $\Cscr$. 
\end{remark}
We consider the structure of the kernel of an $\otimes_{\Cscr}$-compatible $\pi$-stably reflective localization under the additional assumption that $L$ preserves small limits. 
\begin{lemma}\label{lem: kernel structure L limit preserving}
    Let $L:\Cscr\to\Cscr^{L}$ be a $\otimes_{\Cscr}$-compatible $\pi$-stably reflective localization of a presentably symmetric monoidal $t$-bistable category with $L$ preserving small limits. Then: 
    \begin{enumerate}
        \item $\ker(L):=\{K\in\Cscr:L(K)\simeq0\}\subseteq\Cscr$ is a stable tensor ideal closed under small limits, colimits, and extensions in $\Cscr$. 
        \item If moreover $\ker(L)$ is homotopy detected, $K\in\ker(L)$ if and only if $\pi_{k}(K)\in\ker(L)\cap\Cscr^{\heartsuit}$ for all $k\in\ZZ$. 
    \end{enumerate}
\end{lemma}
\begin{proof}[Proof of (1)]
    By \cite[Def. 2.26 \& Lem. 2.27]{LocalizingTStruct}, $\ker(L)$ is a stable subcategory of $\Cscr$ closed under small colimits and extensions. Since $L$ preserves limits, for $M:=\lim_{i\in I}M_{i}$ a limit of objects $M_{i}\in\ker(L)$, we have $L(\lim_{i\in I}M_{i})\simeq \lim_{i\in I}L(M_{i})\simeq\lim_{i\in I}0\simeq0$ showing $L(M)\simeq0$ and $M\in\ker(L)$. To see $\ker(L)$ is a tensor ideal, we note that for $K\in\ker(L)$ that $K\to 0$ is an $L$-equivalence, so by \Cref{prop: sym mon localizations}, $M\otimes_{\Cscr}K\to M\otimes_{\Cscr}0\simeq 0$ is an $L$-equivalence showing $M\otimes_{\Cscr}K\in\ker(L)$ for all $M\in\Cscr$. 
\end{proof}
\begin{proof}[Proof of (2)]
    $(\Rightarrow)$ Let $K\in\ker(L)$. Since $\ker(L)$ is homotopy detected, each homotopy group $\pi_{k}(K)$ lies in $\ker(L)\cap\Cscr^{\heartsuit}$. 

    $(\Leftarrow)$ By right completeness we can write $K$ as the colimit of its Whitehead tower, and each Whitehead truncation again satisfies the hypothesis. Since $\ker(L)$ is closed under small colimits, we may assume $K$ to be bounded below. For $K\in\Cscr_{\geq r}$, we appeal to left completeness to write $K\simeq\lim_{k}\tau_{\leq k}K$ as the limit of its Postnikov tower, and further reduce to showing each term of the Postnikov tower lies in $\ker(L)$. Induction with base case 
    $$\pi_{r}(K)[r]\to\tau_{\leq r}K\to\tau_{\leq r-1}K\simeq0$$
    exhibits each term of the Postnikov tower as an extension of two objects in $\ker(L)$, and thus lies in $\ker(L)$ by closure of $\ker(L)$ under extensions. We then conclude that $K\in\ker(L)$ for $\ker(L)$ is closed under limits. 
\end{proof}
We use the following recognition criterion for the category of modules over an idempotent algebra. 
\begin{proposition}[{\cite[Thm. 7.7]{RagimovReflection}}]\label{prop: idempotent recognition}
    Let $\Cscr$ be a presentably symmetric monoidal category and $\Escr\subseteq\Cscr$ a full subcategory. The following are equivalent: 
    \begin{enumerate}[label=(\alph*)]
        \item $\Escr\simeq\Mod_{Q}(\Cscr)$ for some idempotent algebra $Q\in\Idem(\Cscr)$. 
        \item The inclusion $\iota:\Escr\hookrightarrow\Cscr$ admits a left adjoint given on objects by $M\mapsto M\otimes_{\Cscr}Q$ for some object $Q\in\Cscr$. 
        \item $\Escr$ is closed under limits and colimits in $\Cscr$ and $E\otimes_{\Cscr}M,\intHom_{\Cscr}(M,E)\in\Escr$ for $E\in\Escr,M\in\Cscr$.
    \end{enumerate}
\end{proposition}
The internal hom condition of \Cref{prop: idempotent recognition} (c) is the substantive one to verify. The next two results reduce it to a condition on static objects: first by commuting the internal hom past Whitehead and Postnikov towers, then by descending membership to homotopy objects.
\begin{lemma}\label{lem: internal hom and towers}
    Let $\Cscr$ be a presentably symmetric monoidal $t$-bistable category. 
    \begin{enumerate}
        \item If $M\in\Cscr$ is bounded below, then $\intHom_{\Cscr}(M,-)$ commutes with Whitehead towers of bounded above objects: $\intHom_{\Cscr}(M, N)\simeq\colim_{k}\intHom_{\Cscr}(M,\tau_{\geq k}N)$. 
        \item If $N\in\Cscr$ is bounded above, then $\intHom_{\Cscr}(-,N)$ commutes with Postnikov towers of bounded below objects: $\intHom_{\Cscr}(M, N)\simeq\colim_{k}\intHom_{\Cscr}(\tau_{\leq k}M, N)$. 
    \end{enumerate}
\end{lemma}
\begin{proof}
    By exactness of $\intHom_{\Cscr}(-,-)$ and shifting, we may reduce to the case of $M$ connective and $N$ coconnective. 

    We first show (1). As the right adjoint of the right $t$-exact functor $M\otimes_{\Cscr}-$, the functor $\intHom_{\Cscr}(M,-)$ is left $t$-exact. Writing $N\simeq\colim_{k}\tau_{\geq k}N$ for the Whitehead tower, which converges by right completeness, we obtain a comparison morphism
    \begin{equation}\label{eqn: Whitehead comparison map}
        \colim_{k}\intHom_{\Cscr}(M,\tau_{\geq k}N)\longrightarrow\intHom_{\Cscr}(M, N).
    \end{equation}
    Applying $\intHom_{\Cscr}(M,-)$ to the (co)fiber sequences $\tau_{\geq k}N\to N\to\tau_{\leq k-1}N$ identifies the fiber of $\intHom_{\Cscr}(M,\tau_{\geq k}N)\to\intHom_{\Cscr}(M,N)$ with $\intHom_{\Cscr}(M,\tau_{\leq k-1}N)[-1]$. Since finite limits commute with filtered colimits, the fiber $F$ of \eqref{eqn: Whitehead comparison map} is
    $$F\simeq\colim_{k}\intHom_{\Cscr}(M,\tau_{\leq k-1}N)[-1].$$
    By left $t$-exactness of $\intHom_{\Cscr}(M,-)$, each $\intHom_{\Cscr}(M,\tau_{\leq k-1}N)[-1]$ lies in $\Cscr_{\leq k}$, so for every $r\in\ZZ$ the homotopy object $\pi_{r}$ of the $k$-th term vanishes once $k<r$. This holds cofinally in the filtered colimit computing $\pi_{r}(F)$, so $\pi_{r}(F)\simeq0$ for all $r$, hence $F\simeq0$ and \eqref{eqn: Whitehead comparison map} is an equivalence. 

    The proof of (2) is dual. Applying the contravariant $\intHom_{\Cscr}(-,N)$ to the (co)fiber sequences $\tau_{\geq k+1}M\to M\to\tau_{\leq k}M$ identifies the cofiber of the comparison morphism $\colim_{k}\intHom_{\Cscr}(\tau_{\leq k}M,N)\to\intHom_{\Cscr}(M,N)$ with $\colim_{k}\intHom_{\Cscr}(\tau_{\geq k+1}M,N)$. Coconnectivity of $N$ and adjunction means $\intHom_{\Cscr}(-,N)$ carries $\Cscr_{\geq k+1}$ into $\Cscr_{\leq -k-1}$, so these terms grow arbitrarily coconnective and the colimit vanishes, whence the comparison is an equivalence.
\end{proof}
Iterating along both towers reduces an internal hom membership to its values on the homotopy objects of the two arguments.
\begin{corollary}\label{cor: hom membership}
    Let $\Cscr$ be a presentably symmetric monoidal $t$-bistable category and $\Escr\subseteq\Cscr$ a subcategory closed under small colimits, countable limits, and extensions in $\Cscr$. Let $M\in\Cscr,E\in\Escr$ and suppose $\intHom_{\Cscr}(\pi_{a}(M),\pi_{b}(E))\in\Escr$ for all $a,b\in\ZZ$. Then $\intHom_{\Cscr}(M, E)\in\Escr$. 
\end{corollary}
\begin{proof}
    By closure of $\Escr$ under colimits and countable limits, we may use Whitehead towers in $M$ and Postnikov towers in $E$ to reduce to the case of $M$ bounded below and $E$ bounded above. Then, by shifting and exactness of $\intHom_{\Cscr}(-,-)$, we can take $M$ to be connective and $E$ to be coconnective.

    Now suppose $\intHom_{\Cscr}(\pi_{a}(M),\pi_{b}(E))\in\Escr$ for all $a,b\in\ZZ$. Since $E$ is coconnective and $\pi_{a}(M)$ is in the heart -- and hence connective -- we have by \Cref{lem: internal hom and towers} (1) that $\colim_{k}\intHom_{\Cscr}(\pi_{a}(M),\tau_{\geq k}E)\simeq\intHom_{\Cscr}(\pi_{a}(M),E)$. We show each term of the colimit lies in $\Escr$. Using the (co)fiber sequences $\tau_{\geq k}E\to\tau_{\geq k-1}E\to\pi_{k-1}(E)[k-1]$ in $\Cscr$, induction with base case 
    $$\intHom_{\Cscr}(\pi_{a}(M),\tau_{\geq 0}E)\to\intHom_{\Cscr}(\pi_{a}(M),\tau_{\geq -1}E)\to\intHom_{\Cscr}(\pi_{a}(M),\pi_{-1}(E))[-1]$$
    exhibits each term of the colimit as an extension of two objects that lie in $\Escr$, hence each lies in $\Escr$ by closure under extensions. Here the base term satisfies $\intHom_{\Cscr}(\pi_{a}(M),\tau_{\geq 0}E)\simeq\intHom_{\Cscr}(\pi_{a}(M),\pi_{0}(E))\in\Escr$ as $\tau_{\geq0}E\simeq\pi_{0}(E)$ by coconnectivity of $E$. We thus obtain $\intHom_{\Cscr}(\pi_{a}(M),E)\in\Escr$ for all $a\in\ZZ$.

    Now using the equivalence $\colim_{k}\intHom_{\Cscr}(\tau_{\leq k}M, E)\simeq\intHom_{\Cscr}(M, E)$ of \Cref{lem: internal hom and towers} (2), it remains to show each term of the colimit lies in $\Escr$. Applying the exact contravariant $\intHom_{\Cscr}(-,E)$ to the (co)fiber sequences $\pi_{k}(M)[k]\to\tau_{\leq k}M\to\tau_{\leq k-1}M$ produces (co)fiber sequences 
    $$\intHom_{\Cscr}\left(\tau_{\leq k-1}M,E\right)\to\intHom_{\Cscr}\left(\tau_{\leq k}M,E\right)\to\intHom_{\Cscr}\left(\pi_{k}(M),E\right)[-k]$$
    in $\Cscr$, exhibiting $\intHom_{\Cscr}(\tau_{\leq k}M,E)$ as an extension of $\intHom_{\Cscr}(\pi_{k}(M),E)[-k]$ by $\intHom_{\Cscr}(\tau_{\leq k-1}M,E)$. The base case $k=0$ gives $\intHom_{\Cscr}(\tau_{\leq 0}M,E)\simeq\intHom_{\Cscr}(\pi_{0}(M),E)\in\Escr$ by the previous paragraph, and the third term lies in $\Escr$ by the same together with stability of $\Escr$. Induction on $k$ and closure under extensions then give $\intHom_{\Cscr}(\tau_{\leq k}M,E)\in\Escr$ for every $k$, and we conclude $\intHom_{\Cscr}(M, E)\in\Escr$ as desired. 
\end{proof}
We now prove the recognition criterion in $\Cscr$.
\begin{theorem}\label{thm: smashing kernel recognition}
    Let $L:\Cscr\to\Cscr^{L}$ be a $\otimes_{\Cscr}$-compatible $\pi$-stably reflective localization of a presentably symmetric monoidal $t$-bistable category with $L$ preserving small limits and $\ker(L)$ homotopy detected. The following conditions are equivalent:
    \begin{enumerate}[label=(\alph*)]
        \item $L$ is an open immersion with $\ker(L)\simeq\Mod_{Q}(\Cscr)$ for some idempotent algebra $Q$ in $\Cscr$. 
        \item $\intHom_{\Cscr}(M,K)\in\ker(L)$ for all $M\in\Cscr^{\heartsuit},K\in\ker(L)\cap\Cscr^{\heartsuit}$. 
    \end{enumerate}
\end{theorem}
\begin{proof}
    ($\Rightarrow$) If $L$ is an open immersion, $\intHom_{\Cscr}(M, K)\in\ker(L)$ for all $M\in\Cscr,K\in\ker(L)$ by \Cref{prop: idempotent recognition}, of which (b) is a special case. 

    ($\Leftarrow$) By \Cref{lem: kernel structure L limit preserving} (1), $\ker(L)$ is a stable tensor ideal closed under small limits, small colimits, and extensions in $\Cscr$. To conclude $\ker(L)\simeq\Mod_{Q}(\Cscr)$ for an idempotent algebra $Q\in\Idem(\Cscr)$ it therefore suffices, by \Cref{prop: idempotent recognition} (c), to show that $\intHom_{\Cscr}(M,K)\in\ker(L)$ for all $M\in\Cscr$ and $K\in\ker(L)$. 

    Fix such $M$ and $K$. Apply \Cref{cor: hom membership} to $\ker(L)$, whose closure under small colimits, small limits, and extensions is furnished by \Cref{lem: kernel structure L limit preserving} (1). We show $\intHom_{\Cscr}(\pi_{a}(M),\pi_{b}(K))\in\ker(L)$ for all $a,b\in\ZZ$. Here $\pi_{a}(M)\in\Cscr^{\heartsuit}$ and $\pi_{b}(K)\in\ker(L)\cap\Cscr^{\heartsuit}$ by \Cref{lem: kernel structure L limit preserving} (2), so the hypothesis of \Cref{cor: hom membership} holds by (b). 

    As $M\in\Cscr$ and $K\in\ker(L)$ were arbitrary, the conditions of \Cref{prop: idempotent recognition} (c) hold, so there is an idempotent algebra $Q\in\Idem(\Cscr)$ with $\ker(L)\simeq\Mod_{Q}(\Cscr)$, exhibiting $L$ as an open immersion.
\end{proof}
It remains to record a practically verifiable sufficient condition for homotopy detection. In particular, we show the hypothesis holds when the localization has $t$-amplitude at most 1. 
\begin{definition}\label{def: t-amplitude}
    Let $F:\Cscr\to\Escr$ be an exact, right $t$-exact, small colimit-preserving functor between $t$-bistable categories, and let $n\geq0$. $F$ is \emph{of $t$-amplitude $\leq n$} if $F(\Cscr^{\heartsuit})\subseteq\Escr_{\leq n}$. 
\end{definition}
\begin{lemma}\label{lem: homotopy detected kernels of t-amplitude less than 1 functors}
    If $F:\Cscr\to\Escr$ is an exact, right $t$-exact, small colimit-preserving functor of $t$-amplitude $\leq 1$ between $t$-bistable categories, then $F(\Cscr_{\leq m})\subseteq\Escr_{\leq m+1}$ and 
    $$\ker(F):=\{K\in\Cscr:F(K)\simeq0\}\subseteq\Cscr$$ 
    admits the induced $t$-structure.  
\end{lemma}
\begin{proof}
    First observe that for $M\in\Cscr^{\heartsuit}$, exactness of $F$ gives $F(M[k])\simeq F(M)[k]\in\Escr_{\leq k+1}$ for $k\in\ZZ$. 

    Now assume $M\in\Cscr_{[a,b]}$ is bounded. We can write it as a Whitehead tower 
    $$0\simeq\tau_{\geq b+1}M\to \tau_{\geq b}M\to\dots\to\tau_{\geq a+1}M\to\tau_{\geq a}M\simeq M$$
    with (co)fiber sequences $\tau_{\geq k+1}M\to\tau_{\geq k}M\to \pi_{k}(M)[k]$. Applying $F$ we get a diagram 
    $$0\simeq F(\tau_{\geq b+1}M)\to F(\tau_{\geq b}M)\to\dots\to F(\tau_{\geq a+1}M)\to F(\tau_{\geq a}M)\simeq F(M)$$
    with (co)fiber sequences $F(\tau_{\geq k+1}M)\to F(\tau_{\geq k}M)\to F(\pi_{k}(M))[k]$. Each $F(\tau_{\geq k}M)$ is an extension of objects in $\Escr_{\leq b+1}$ and hence lies in $\Escr_{\leq b+1}$ showing $F(M)\in\Escr_{\leq b+1}$ as well. 

    Now let $M\in\Cscr_{\leq m}$. Writing $M$ as its Whitehead tower, we have $F(M)\simeq\colim_{k}F(\tau_{\geq k}M)$ in $\Escr$ as $F$ preserves small colimits. Since each $\tau_{\geq k}M$ term is bounded, $F(M)$ is a filtered colimit of objects in $\Escr_{\leq m+1}$. Since passage to homotopy objects in $\Escr$ commutes with filtered colimits, we compute $\pi_{r}(F(M))\simeq\colim_{k}(\pi_{r}(F(\tau_{\geq k}M)))$ which vanishes for $r\geq m+2$ as each term does. This shows $F(M)\in\Escr_{\leq m+1}$ by \Cref{lem: separatedness and prestable} (2), proving the first claim.

    For the second claim, let $K\in\ker(F)$ and $m\in\ZZ$. By stability, $\ker(F)$ is closed under fibers and cofibers. Using the (co)fiber sequences $\tau_{\geq m+1}K\to K\to\tau_{\leq m}K$, it suffices to show that $\tau_{\leq m}K\in\ker(F)$. Applying $F$ to the (co)fiber sequence, we get a (co)fiber sequence $F(\tau_{\geq m+1}K)\to F(K)\simeq0\to F(\tau_{\leq m}K)$ in $\Escr$ showing $F(\tau_{\leq m}K)\simeq\cofib(F(\tau_{\geq m+1}K)\to 0)\simeq F(\tau_{\geq m+1}K)[1]$. We have that $F(\tau_{\geq m+1}K)[1]\in\Escr_{\geq m+2}$ by right $t$-exactness of $F$ but $F(\tau_{\leq m}K)\in\Escr_{\leq m+1}$ by the preceding argument. So $F(\tau_{\leq m}K)\in\Escr_{\leq m+1}\cap\Escr_{\geq m+2}\simeq0$, showing that $\tau_{\leq m}K\in\ker(F)$, as desired. 
\end{proof}
\begin{remark}\label{rmk: sharpness of t-amplitude bound}
    The argument does not carry over for functors of $t$-amplitude $n\geq 2$. For $F(\Cscr^{\heartsuit})\subseteq\Escr_{\leq n}$ the same computation only places $F(\tau_{\leq m}K)$ in $\Escr_{\leq m+n}\cap\Escr_{\geq m+2}$, which need not be zero. 
\end{remark}
We obtain the following corollary. 
\begin{corollary}\label{cor: recognition at t-amplitude 1}
    Let $L:\Cscr\to\Cscr^{L}$ be an $\otimes_{\Cscr}$-compatible $\pi$-stably reflective localization of a presentably symmetric monoidal $t^{+}$-bistable category such that $L$ preserves small limits and is of $t$-amplitude $\leq 1$. The following are equivalent: 
    \begin{enumerate}[label=(\alph*)]
        \item $L$ is an open immersion with $\ker(L)\simeq\Mod_{Q}(\Cscr)$ for some idempotent algebra $Q\in\Idem(\Cscr)$. 
        \item $\intHom_{\Cscr}(M, K)\in\ker(L)$ for all $M\in \Cscr^{\heartsuit},K\in\ker(L)\cap\Cscr^{\heartsuit}$. 
    \end{enumerate}
\end{corollary}
\begin{proof}
    $L$ is exact, right $t$-exact, small colimit-preserving, and of $t$-amplitude $\leq 1$, so \Cref{lem: homotopy detected kernels of t-amplitude less than 1 functors} gives that $\ker(L)$ is homotopy detected. The hypotheses of \Cref{thm: smashing kernel recognition} apply, giving the desired equivalence. 
\end{proof}
\section{Nuclearity in the \texorpdfstring{$\kappa$}{kappa}-presentable setting}\label{app: nuclearity}
In \cite[Lecture VIII]{Complex}, Clausen and Scholze introduce trace-class morphisms and nuclear objects in a stable $\omega$-presentable symmetric monoidal category with compact unit (though not necessarily $\omega$-presentably symmetric monoidal -- see \Cref{subsec: notation and conventions}) and relate them to dualizability. 

The categories of interest in the preceding text are not necessarily $\omega$-presentable. Assuming that the tensor unit is ($\omega$-)compact, we introduce a $\kappa$-presentable version of nuclearity adapted from Ramzi's rigidification. These nuclear objects form a full subcategory closed under small colimits, and we prove that an object is dualizable if and only if it is both nuclear and compact. Much of the discussion below builds on work of Aoki \cite{AokiSchwartz} and Ramzi \cite{RamziLocRigid,RamziDualizable}.
\subsection{Dualizability and trace-class morphisms}\label{subsec: dualizability and trace-class morphisms}
We begin with a set-theoretic preliminary. 
\begin{lemma}\label{lem: kappa-presentably symmetric monoidal}
    Let $\Cscr$ be a presentably symmetric monoidal category such that the underlying category is $\kappa$-presentable for an uncountable regular cardinal $\kappa$. Then $\Cscr\in\CAlg(\PrL_{\kappa})$ if and only if the unit $1_{\Cscr}$ is $\kappa$-compact and the $\kappa$-compact objects are preserved by the tensor product. 
\end{lemma}
\begin{proof}
    This follows from \cite[Not. 5.3.2.8 \& Lem. 5.3.2.11]{HA}.
\end{proof}
\begin{remark}\label{rmk: large kappa}
    By \cite[Lem. 5.3.2.12]{HA}, for every presentably symmetric monoidal category $\Cscr$ there exists an uncountable regular cardinal $\kappa$ such that $\Cscr\in\CAlg(\PrL_{\kappa})$. Thus one may always choose $\kappa$ such that this hypothesis is satisfied. The substantive additional requirement in what follows is that the unit be ($\omega$-)compact, in contrast to the $\omega$-presentability assumed in \cite[Lecture VIII]{Complex}.
\end{remark}
Trace-class morphisms, and the nuclear objects of the next subsection, are both closely tied to dualizability. We first recall that notion, in the slightly restricted generality of a closed symmetric monoidal category.
\begin{definition}\label{def: dualizability in closed symmetric monoidal categories}
    Let $\Cscr$ be a closed symmetric monoidal category. 
    \begin{enumerate}
        \item Let $M\in\Cscr$. The \emph{pre-dual} of $M$ is the object $M^{\vee}:=\intHom_{\Cscr}(M,1_{\Cscr})$. Write $\ev_{M}:M\otimes_{\Cscr}M^{\vee}\to 1_{\Cscr}$ for the counit morphism of the adjunction $M\otimes_{\Cscr}(-)\dashv\intHom_{\Cscr}(M,-)$ at $1_{\Cscr}$. 
        \item An object $M\in\Cscr$ is \emph{dualizable} if there exists a morphism $\coev_{M}:1_{\Cscr}\to M^{\vee}\otimes_{\Cscr}M$ such that the composites 
        $$M\xrightarrow{\id_{M}\otimes_{\Cscr}\coev_{M}}M\otimes_{\Cscr}M^{\vee}\otimes_{\Cscr}M\xrightarrow{\ev_{M}\otimes_{\Cscr}\id_{M}}M$$
        $$M^{\vee}\xrightarrow{\coev_{M}\otimes_{\Cscr}\id_{M^{\vee}}}M^{\vee}\otimes_{\Cscr}M\otimes_{\Cscr}M^{\vee}\xrightarrow{\id_{M^{\vee}}\otimes_{\Cscr}\ev_{M}}M^{\vee}$$
        are identities. In this case, we say the pre-dual $M^{\vee}$ of $M$ is the \emph{dual} of $M$.  
    \end{enumerate}
\end{definition}
Under the set-theoretic assumptions of \Cref{lem: kappa-presentably symmetric monoidal}, we have the following definition of trace-class morphisms. 
\begin{definition}\label{def: trace-class}
    Let $\kappa$ be an uncountable regular cardinal and $\Cscr\in\CAlg(\PrL_{\kappa})$ such that $1_{\Cscr}$ is compact (ie. $\omega$-compact). A morphism $f:M\to N$ in $\Cscr$ is \emph{trace-class} if the morphism $1_{\Cscr}\to\intHom_{\Cscr}(M, N)$ classifying $f$ admits a factorization as $1_{\Cscr}\to M^{\vee}\otimes_{\Cscr}N\to\intHom_{\Cscr}(M, N)$ where $M^{\vee}\otimes_{\Cscr}N\to \intHom_{\Cscr}(M, N)$ is the adjoint morphism to $M\otimes_{\Cscr}M^{\vee}\otimes_{\Cscr}N\xrightarrow{\ev_{M}\otimes_{\Cscr}\id_{N}}N$. 
\end{definition}
\begin{remark}[{\cite[\S 7.1]{AokiSchwartz}}]\label{rmk: trace-class factorization}
    In the situation of \Cref{def: trace-class}, write $\eta_{f}:1_{\Cscr}\to M^{\vee}\otimes_{\Cscr}N$ for a lift of the classifying morphism. Then $f$ itself admits a factorization
    \begin{equation}\label{eqn: trace-class factorization}
        M\xrightarrow{\id_{M}\otimes_{\Cscr}\eta_{f}}M\otimes_{\Cscr}M^{\vee}\otimes_{\Cscr}N\xrightarrow{\ev_{M}\otimes_{\Cscr}\id_{N}}N
    \end{equation}
    where $\ev_{M}$ is the evaluation morphism for the pre-dual in \Cref{def: dualizability in closed symmetric monoidal categories}.
\end{remark}
\subsection{Rigidification and (basic) nuclear objects}\label{subsec: rigidification and basic nuclear objects}
Throughout this subsection, fix an uncountable regular cardinal $\kappa$ and a stable $\kappa$-presentably symmetric monoidal category $\Cscr\in\CAlg(\PrL_{\St,\kappa})$ with compact unit. We define the (basic) nuclear objects of $\Cscr$ directly, and then identify them with the essential image of the rigidification of \cite{RamziLocRigid}. The comparison functor is not in general fully faithful, but is so under this hypothesis.  
\begin{definition}\label{def: basic nuclear in category}
    Let $\kappa$ be an uncountable regular cardinal and $\Cscr\in\CAlg(\PrL_{\St,\kappa})$ be a stable $\kappa$-presentably symmetric monoidal category with compact unit.
    \begin{enumerate}
        \item An object of $\Cscr$ is \emph{basic nuclear} if it is a $\QQ_{\geq0}$-indexed colimit of $\kappa$-compact objects along trace-class morphisms. Write $\BNuc(\Cscr)$ for the full subcategory of basic nuclear objects.  
        \item An object of $\Cscr$ is \emph{nuclear} if it is in the full subcategory of $\Cscr$ generated under small colimits by basic nuclear objects. Write $\Nuc(\Cscr)$ for the full subcategory of nuclear objects. 
    \end{enumerate}
\end{definition}
\begin{remark}\label{rmk: nuclear is kappa-independent}
    Under the identification of nuclear objects with the essential image of the rigidification (\Cref{prop: nuclearity and rigidity}), \cite[Thm. 4.81]{RamziLocRigid} shows that nuclearity is independent of the cardinal $\kappa$. 
\end{remark}
When the underlying category is $\omega$-presentable, every basic nuclear object in the sense above is basic nuclear in the sense of Clausen and Scholze, as recorded in the following lemma.
\begin{lemma}\label{lem: CS nuclear comparison}
    Let $\kappa$ be an uncountable regular cardinal and $\Cscr\in\CAlg(\PrL_{\St,\kappa})$ be a stable $\kappa$-presentably symmetric monoidal category with compact unit, and suppose furthermore that the underlying category is $\omega$-presentable. Then every basic nuclear object of $\Cscr$ in the sense of \Cref{def: basic nuclear in category} (1) is basic nuclear in the sense of \cite[Def. 8.5]{Complex}.
\end{lemma}
\begin{proof}
    Suppose $M$ is basic nuclear in the sense of \Cref{def: basic nuclear in category} (1), and write $M\simeq\colim_{q\in\QQ_{\geq0}}M_{q}$ where each $M_{q}$ is $\kappa$-compact and each transition morphism is trace-class. The inclusion of posets $\NN\hookrightarrow\QQ_{\geq0}$ is cofinal \cite[Rmk. 2.59]{RamziDualizable}, so $M\simeq\colim_{n\in\NN}M_{n}$. The transition morphisms $M_{n}\to M_{n+1}$ are trace-class, being composites of trace-class morphisms in the original diagram, hence this expresses $M$ as a sequential colimit along trace-class morphisms. Thus $M$ is basic nuclear in the sense of \cite[Def. 8.5]{Complex}.
\end{proof}
To establish the properties of nuclear objects, we identify them with the essential image of the rigidification of $\Cscr$, following Ramzi \cite[Const. 4.79 \& Thm. 4.81]{RamziLocRigid}.
\begin{definition}\label{def: rigidification}
    Let $\kappa$ be an uncountable regular cardinal, $\Cscr\in\CAlg(\PrL_{\St,\kappa})$ be a stable $\kappa$-presentably symmetric monoidal category, and denote by $\ysf:\Cscr_{\kappa}\hookrightarrow\Ind(\Cscr_{\kappa})$ the Yoneda embedding. The \emph{rigidification} $\Rig(\Cscr)$ of $\Cscr$ is the full subcategory of $\Ind(\Cscr_{\kappa})$ generated under small colimits by objects of the form $\colim_{q\in\QQ_{\geq0}}\ysf(M_{q})$ where $\{M_{q}\}_{q\in\QQ_{\geq0}}$ is a diagram of $\kappa$-compact objects with trace-class transition morphisms. 
\end{definition}
\begin{remark}\label{rmk: computation of trace-class morphisms}
    A morphism between $\kappa$-compact objects is trace-class in $\Cscr_{\kappa}$ if and only if it is trace-class in $\Cscr$ \cite[Rmk. 3.2]{RamziLocRigid}. By symmetric monoidality of the Yoneda embedding, the corresponding morphisms are trace-class in $\Ind(\Cscr_{\kappa})$.
\end{remark}
We record the properties of the comparison functor that we will need.
\begin{lemma}\label{lem: comparison functor from rigidification}
    Let $\kappa$ be an uncountable regular cardinal and $\Cscr\in\CAlg(\PrL_{\St,\kappa})$ be a stable $\kappa$-presentably symmetric monoidal category with compact unit.
    \begin{enumerate}
        \item The category $\Rig(\Cscr)$ admits a symmetric monoidal structure for which the composite 
        $$\Rig(\Cscr)\hookrightarrow\Ind(\Cscr_{\kappa})\xrightarrow{\colim}\Cscr$$
        is symmetric monoidal. 
        \item $\Rig(\Cscr)$ is $\omega_{1}$-presentable. Moreover, every object of the form $\colim_{q\in\QQ_{\geq0}}\ysf(M_{q})$ occurring in \Cref{def: rigidification} is $\omega_{1}$-compact.
        \item The functor in (1) is fully faithful, induces an equivalence on the dualizable objects, and its essential image is closed under small colimits computed in $\Cscr$. 
    \end{enumerate}
\end{lemma}
\begin{proof}[Proof of (1)]
    The symmetric monoidal structure on $\Rig(\Cscr)$ is constructed in \cite[Const. 4.79]{RamziLocRigid} and symmetric monoidality of the composite is verified in \cite[Thm. 4.81]{RamziLocRigid}. 
\end{proof}
\begin{proof}[Proof of (2)]
    This is shown as part of the proof of \cite[Thm. 4.81]{RamziLocRigid}. 
\end{proof}
\begin{proof}[Proof of (3)]
    The functor in (1) is fully faithful when the unit is compact \cite[Cor. 4.75 \& Lem. 4.84]{RamziLocRigid}, and restricts to an equivalence on dualizable objects by \cite[Rmk. 7.21]{AokiSchwartz}: the rigidifications of Aoki and Ramzi agree by their universal properties stated in \cite[Cor. 7.20]{AokiSchwartz} and \cite[Thm. 4.81]{RamziLocRigid}. 
    
    As a fully faithful morphism in $\CAlg(\PrL_{\St})$, it exhibits $\Rig(\Cscr)$ as a coreflective subcategory of $\Cscr$, so by (the dual of) \cite[\href{https://kerodon.net/tag/04JX}{Tag 04JX}]{kerodon} its essential image is closed under small colimits computed in $\Cscr$.
\end{proof}
Together, these properties identify nuclear objects with the essential image of the rigidification. 
\begin{proposition}\label{prop: nuclearity and rigidity}
    Let $\kappa$ be an uncountable regular cardinal and $\Cscr\in\CAlg(\PrL_{\St,\kappa})$ be a stable $\kappa$-presentably symmetric monoidal category with compact unit. An object of $\Cscr$ is nuclear if and only if it is in the essential image of the comparison functor $\iota:\Rig(\Cscr)\to\Cscr$ of \Cref{lem: comparison functor from rigidification}. 
\end{proposition}
\begin{proof}
    $(\Rightarrow)$ Suppose $M$ is nuclear. First suppose $M$ is basic nuclear and choose a presentation $M\simeq\colim_{q\in\QQ_{\geq0}}M_{q}$ by $\kappa$-compact objects along trace-class morphisms. Applying the Yoneda embedding gives an object $\colim_{q\in\QQ_{\geq0}}\ysf(M_{q})$ of $\Rig(\Cscr)\subseteq\Ind(\Cscr_{\kappa})$ whose image under $\iota$ is $M$. Thus every basic nuclear object lies in the essential image of $\iota$. Since the essential image of $\iota$ is closed under small colimits by \Cref{lem: comparison functor from rigidification} (3), every nuclear object lies in the essential image of $\iota$.

    $(\Leftarrow)$ Conversely, by definition $\Rig(\Cscr)$ is generated under small colimits by objects of the form $\colim_{q\in\QQ_{\geq0}}\ysf(M_{q})$ in $\Ind(\Cscr_{\kappa})$, where each $M_{q}$ is $\kappa$-compact and the transition morphisms are trace-class. Since $\iota$ is the restriction of the colimit functor, it preserves small colimits and carries each such generator to $\colim_{q\in\QQ_{\geq0}}M_{q}$, which is basic nuclear by definition. The nuclear objects are closed under small colimits by definition, so every object in the essential image of $\iota$ is nuclear. Hence the essential image of $\iota$ is precisely $\Nuc(\Cscr)$.
\end{proof}
Nuclearity is preserved by symmetric monoidal functors and satisfies a projection formula when the right adjoint is sufficiently well-behaved. 
\begin{proposition}\label{prop: preservation of nuclearity}
    Let $\kappa$ be an uncountable regular cardinal and $F:\Cscr\to\Escr$ be a morphism in $\CAlg(\PrL_{\St,\kappa})$ with compact units. Then:
    \begin{enumerate}
        \item $F$ preserves nuclear objects. 
        \item Suppose the right adjoint $G:\Escr\to\Cscr$ preserves small colimits. For all $M\in\Nuc(\Cscr)$ and $N\in\Escr$, the natural morphism
        $$M\otimes_{\Cscr}G(N)\longrightarrow G(F(M)\otimes_{\Escr}N)$$
        is an equivalence.
    \end{enumerate}
\end{proposition}
\begin{proof}[Proof of (1)]
    As a morphism in $\CAlg(\PrL_{\St,\kappa})$, $F$ is symmetric monoidal, preserves small colimits, and preserves $\kappa$-compact objects. Being symmetric monoidal, it preserves trace-class morphisms \cite[Rmk. 3.2]{RamziLocRigid}. Thus $F$ carries basic nuclear objects of $\Cscr$ to basic nuclear objects of $\Escr$. The full subcategory of objects of $\Cscr$ carried into $\Nuc(\Escr)$ therefore contains the basic nuclear objects, and it is closed under small colimits since $F$ preserves small colimits and $\Nuc(\Escr)$ is closed under them, so it contains all of $\Nuc(\Cscr)$.
\end{proof}
\begin{proof}[Proof of (2)]
    The functors $(-)\otimes_{\Cscr}G(N)$ and $G(F(-)\otimes_{\Escr}N)$ preserve small colimits, so the objects of $\Cscr$ on which the natural morphism is an equivalence form a subcategory closed under small colimits. Since the nuclear objects are generated under small colimits by the basic nuclear objects, it suffices to treat basic nuclear objects. Thus suppose $M\simeq\colim_{q\in\QQ_{\geq0}}M_{q}$ with the $M_{q}$ $\kappa$-compact and the transition morphisms trace-class. For each $q'>q$ in $\QQ_{\geq0}$ with $M_{q}\to M_{q'}$ trace-class, \cite[Prop. 3.4]{RamziLocRigid} supplies a diagonal filler 
    $$% https://q.uiver.app/#q=WzAsNCxbMCwwLCJNX3txfVxcb3RpbWVzX3tcXENzY3J9RyhOKSJdLFsyLDAsIk1fe3EnfVxcb3RpbWVzX3tcXENzY3J9RyhOKSJdLFswLDEsIkcoRihNX3txfSlcXG90aW1lc197XFxFc2NyfU4pIl0sWzIsMSwiRyhGKE1fe3EnfSlcXG90aW1lc197XFxFc2NyfU4pIl0sWzAsMl0sWzEsM10sWzAsMV0sWzIsM10sWzIsMSwiXFxleGlzdHMiLDEseyJzdHlsZSI6eyJib2R5Ijp7Im5hbWUiOiJkYXNoZWQifX19XV0=
    \begin{tikzcd}
        {M_{q}\otimes_{\Cscr}G(N)} && {M_{q'}\otimes_{\Cscr}G(N)} \\
        {G(F(M_{q})\otimes_{\Escr}N)} && {G(F(M_{q'})\otimes_{\Escr}N)}
        \arrow[from=1-1, to=1-3]
        \arrow[from=1-1, to=2-1]
        \arrow[from=1-3, to=2-3]
        \arrow["\exists"{description}, dashed, from=2-1, to=1-3]
        \arrow[from=2-1, to=2-3]
    \end{tikzcd}$$
    inducing an equivalence on the filtered colimits by \cite[Lem. 2.16]{AokiSchwartz}. 
\end{proof}
Finally, we obtain the identification of the dualizable objects with the compact nuclear objects \cite[Prop. 9.3]{Complex}.
\begin{proposition}\label{prop: dualizable iff nuclear and compact category}
    Let $\kappa$ be an uncountable regular cardinal and $\Cscr\in\CAlg(\PrL_{\St,\kappa})$ be a stable $\kappa$-presentably symmetric monoidal category with compact unit. An object $M\in\Cscr$ is dualizable if and only if it is compact and nuclear.
\end{proposition}
\begin{proof}
    $(\Rightarrow)$ Suppose $M$ is dualizable. A dualizable object is $\Cscr$-atomic, and since $1_{\Cscr}$ is compact the $\Cscr$-atomic objects are compact \cite[Ex. 1.23 \& Cor. 1.26]{RamziDualizable}, so $M$ is compact. Moreover its coevaluation exhibits $\id_{M}$ as trace-class, so the constant $\QQ_{\geq0}$-indexed diagram at $M$, with every transition morphism $\id_{M}$, realizes $M$ as a $\QQ_{\geq0}$-indexed colimit of the compact, in particular $\kappa$-compact, object $M$ along trace-class morphisms, exhibiting $M$ as basic nuclear and hence nuclear. 

    $(\Leftarrow)$ Conversely, suppose $M$ is compact and nuclear. By \Cref{prop: nuclearity and rigidity}, $M\simeq\iota(\widetilde{M})$ for some $\widetilde{M}\in\Rig(\Cscr)$, where $\iota:\Rig(\Cscr)\to\Cscr$ is the comparison functor of \Cref{lem: comparison functor from rigidification}. For a filtered diagram $\{Y_{i}\}_{i\in I}$ in $\Rig(\Cscr)$ we compute
    \begin{align*}
        \Hom_{\Rig(\Cscr)}(\widetilde{M},\colim_{i\in I}Y_{i}) &\simeq \Hom_{\Cscr}(M,\iota(\colim_{i\in I}Y_{i})) && \iota \text{ fully faithful}\\
        &\simeq \Hom_{\Cscr}(M,\colim_{i\in I}\iota(Y_{i})) && \iota \text{ colimit-preserving}\\
        &\simeq \colim_{i\in I}\Hom_{\Cscr}(M,\iota(Y_{i})) && M \text{ compact}\\
        &\simeq \colim_{i\in I}\Hom_{\Rig(\Cscr)}(\widetilde{M},Y_{i}) && \iota \text{ fully faithful}
    \end{align*}
    so that $\widetilde{M}$ too is compact. Now $\Rig(\Cscr)$ is rigid \cite[Thm. 4.81]{RamziLocRigid}, and in a rigid category a compact object is dualizable: the identity of a compact object is a compact morphism \cite[Ex. 3.7]{AokiSchwartz}, a compact morphism in a rigid category is trace-class \cite[Prop. 7.17]{AokiSchwartz}, and an object whose identity is trace-class is dualizable by \cite[Ex. 2.9 \& 2.15]{RamziLocRigid} and \cite[Ex. 1.23]{RamziDualizable}. So $\widetilde{M}$ is dualizable, and by \Cref{lem: comparison functor from rigidification} so is $M\simeq\iota(\widetilde{M})$.
\end{proof}

\section{Commutative algebra of \texorpdfstring{$\EE_{\infty}$}{E-infinity}-rings}\label{app: commutative algebra}
The preceding text relies on several facts about the commutative algebra of connective $\EE_{\infty}$-rings, which this appendix collects. The finiteness conditions on modules and morphisms used in the main text are recalled first, together with some structural properties of the category of connective $\EE_{\infty}$-rings. 

We then turn our attention to two classes of morphisms. The morphisms of finite expansion are a natural enlargement of the collection of morphisms of finite type, while the pseudo-lisse morphisms carve out a generality in which a morphism almost of finite presentation and of finite $\Tor$-amplitude retains a workable structure theory, one that the animated setting supplies automatically but that the spectral setting does not.
\subsection{Finiteness properties of modules and morphisms}\label{subsec: finiteness properties}
The finiteness conditions used below follow \cite[\S 7.2]{HA} and \cite[\S 4]{SAG}. 
\begin{definition}\label{def: finiteness properties of modules}
    Let $A\in\CAlg^{\cn}$ and $M\in\Mod_{A}(\Sp)$. 
    \begin{enumerate}
        \item $M$ is \emph{pseudocoherent} if there exists $k\in\ZZ$ such that $M\in\Mod_{A}(\Sp)_{\geq k}$ and $\tau_{\leq n}M$ is compact in $\Mod_{A}(\Sp)_{[k,n]}$ for all $n\geq0$. 
        \item $M$ is \emph{perfect} if it is compact as an object of $\Mod_{A}(\Sp)$. 
    \end{enumerate}
\end{definition} 
\begin{remark}\label{rmk: pseudocoherent and almost perfect}
    In \cite[\S 7]{HA} and \cite{SAG}, pseudocoherent modules are referred to as almost perfect modules. 
\end{remark}
Passing from modules to morphisms, one meets a familiar obstruction. As Lurie points out, many conditions on morphisms of static commutative rings do not generalize directly to connective $\EE_{\infty}$-rings, since the higher coherences force one to prescribe relations among relations without end. The conditions recorded below appear in \cite[\S 7.2.4]{HA} and \cite[\S 4.1 \& \S 5.2]{SAG} (cf.\ \cite[Rmk. 5.2.0.2]{SAG}) and are formulated with that constraint in mind.
\begin{definition}\label{def: finiteness conditions of morphisms}
    Let $f:A\to B$ be a morphism in $\CAlg^{\cn}$. 
    \begin{enumerate}
        \item $f$ is \emph{of finite presentation} if $B$ is a compact object of the category of connective $A$-algebras $\CAlg^{\cn}_{A}$. 
        \item $f$ is \emph{almost of finite presentation} if $\tau_{\leq n}B$ is a compact object of the category of connective $n$-truncated $A$-algebras for all $n\geq0$. 
        \item $f$ is \emph{of finite type} if $\pi_{0}(f):\pi_{0}(A)\to\pi_{0}(B)$ is a finite type morphism of static commutative rings. 
        \item $f$ is \emph{finite} if $\pi_{0}(B)$ is finitely generated as a module over $\pi_{0}(A)$. 
        \item $f$ is \emph{of finite $\Tor$-amplitude} if $B\otimes_{A}(-):\Mod_{A}(\Sp)\to\Mod_{B}(\Sp)$ carries $\Mod_{A}(\Sp)^{\heartsuit}$ into $\Mod_{B}(\Sp)_{\leq n}$ for some $n$.
        \item $f$ is \emph{flat} if $\pi_{0}(f):\pi_{0}(A)\to\pi_{0}(B)$ is a flat morphism of static commutative rings and the map $\pi_{0}(B)\otimes_{\pi_{0}(A)}\pi_{k}(A)\to\pi_{k}(B)$ is an isomorphism of $\pi_{0}(B)$-modules for all $k\geq0$. 
        \item $f$ is \emph{faithfully flat} if $f$ is flat and $\pi_{0}(f):\pi_{0}(A)\to\pi_{0}(B)$ is faithfully flat as a morphism of static commutative rings. 
        \item $f$ is \emph{\'{e}tale} if $f$ is flat and $\pi_{0}(f):\pi_{0}(A)\to\pi_{0}(B)$ is \'{e}tale as a morphism of static commutative rings. 
    \end{enumerate}
\end{definition}
\begin{remark}\label{rmk: terminology for finite presentation}
    Our terminology differs from that of \cite[\S 7.2.4]{HA}, where morphisms of finite presentation in the sense of \Cref{def: finiteness conditions of morphisms} are called morphisms locally of finite presentation.
\end{remark}
Two families of algebras over a connective $\EE_{\infty}$-ring recur throughout this appendix, the free $\EE_{\infty}$-algebras and the polynomial algebras.
\begin{definition}\label{def: free and flat algebras}
    Let $A$ be a connective $\EE_{\infty}$-algebra in $\Sp$ and let $n\geq 0$.
    \begin{enumerate}
        \item The \emph{free $\EE_{\infty}$-algebra} on generators $T_{1},\dots,T_{n}$ over $A$ is $A\{T_{1},\dots,T_{n}\}:=\mathsf{Sym}_{A}(A^{\oplus n})$.
        \item The \emph{polynomial algebra} in the variables $T_{1},\dots,T_{n}$ over $A$ is $A[T_{1},\dots,T_{n}]:=A\otimes_{\Sph}\Sigma^{\infty}_{+}\NN^{n}$, the monoid algebra of the free commutative monoid $\NN^{n}$.
    \end{enumerate}
\end{definition}
\begin{remark}\label{rmk: free and flat algebras}
    Let $A$ be a connective $\EE_{\infty}$-algebra in $\Sp$ and $n\geq 1$. By \cite[Rmk. 5.4.1.2 \& Not. B.1.1.2]{SAG}, we have equivalences
    $$\pi_{0}(A\{T_{1},\dots,T_{n}\})\simeq\pi_{0}(A[T_{1},\dots,T_{n}])\simeq\pi_{0}(A)[T_{1},\dots,T_{n}],$$
    and $A[T_{1},\dots,T_{n}]$ is flat over $A$. Over the sphere, $\Sph\{T_{1},\dots,T_{n}\}$ carries the homology of the symmetric groups and is not of finite $\Tor$-amplitude, while $\Sph[T_{1},\dots,T_{n}]$ is flat. This underlies the contrast between the morphisms of finite expansion of \Cref{subsec: morphisms of finite expansion} and the pseudo-lisse morphisms of \Cref{subsec: pseudo-lisse morphisms}.
\end{remark}
Many statements about algebras over a connective $\EE_{\infty}$-ring reduce to the finitely presented case, by writing a general algebra as a filtered colimit of finitely presented ones as in the static setting. The relevant approximation is the following, an analogue of \cite[\href{https://stacks.math.columbia.edu/tag/00QL}{Tag 00QL}]{stacks-project}.
\begin{proposition}[{\cite[Prop. 7.2.4.27]{HA}}]\label{prop: approximation by finitely presented algebras}
    Let $A$ be a connective $\EE_{\infty}$-algebra in $\Sp$.
    \begin{enumerate}
        \item The category $\CAlg^{\cn}_{A}$ is generated under sifted colimits by retracts of finitely generated free algebras (\Cref{def: free and flat algebras} (1)). In particular, it is $\omega$-presentable.
        \item Write $\CAlg^{\mathsf{fp}}_{A}\subseteq\CAlg^{\cn}_{A}$ for the full subcategory spanned by the finitely presented algebras. The inclusion induces an equivalence $\Ind(\CAlg^{\mathsf{fp}}_{A})\simeq\CAlg^{\cn}_{A}$.
        \item A connective $A$-algebra $B$ is almost of finite presentation if and only if for every $n\geq 0$ there is a finitely presented $A$-algebra $C$ such that $\tau_{\leq n}B$ is a retract of $\tau_{\leq n}C$.
    \end{enumerate}
\end{proposition}
\subsection{Morphisms of finite expansion}\label{subsec: morphisms of finite expansion}
The morphisms of finite expansion were introduced by Hamacher \cite{FiniteExpansion} to carry the exceptional pushforward and pullback functors of \'{e}tale cohomology beyond the case of morphisms of finite type.

The notion is recalled first for static commutative rings, then transported to connective $\EE_{\infty}$-rings through its effect on $\pi_{0}$, before its permanence properties are taken up.
\begin{definition}[{\cite[Def. 1.1 \& Rmk. 1.2]{FiniteExpansion}}]\label{def: static finite expansion}
    Let $A\to B$ be a morphism of static commutative rings. Then $A\to B$ is \emph{of finite expansion} if it admits a factorization
    $$A\to A[T_{1},\dots,T_{n}]\to B$$
    where $A[T_{1},\dots,T_{n}]\to B$ is integral.
\end{definition}
The finite expansion condition is phrased using the free algebra $A\{T_{1},\dots,T_{n}\}$ of \Cref{def: free and flat algebras}. That it depends only on $\pi_{0}$, where it reduces to the static notion, is the content of the following lemma.
\begin{lemma}\label{lem: finite expansion static detection}
    Let $f:A\to B$ be a morphism of connective $\EE_{\infty}$-algebras in $\Sp$. The following are equivalent.
    \begin{enumerate}[label=(\alph*)]
        \item The morphism $f$ admits a factorization $A\to A\{T_{1},\dots,T_{n}\}\to B$ for which $\pi_{0}(A\{T_{1},\dots,T_{n}\})\to\pi_{0}(B)$ is integral.
        \item The induced morphism of static rings $\pi_{0}(f):\pi_{0}(A)\to\pi_{0}(B)$ is of finite expansion in the sense of \Cref{def: static finite expansion}.
    \end{enumerate}
\end{lemma}
\begin{proof}
    $(\Rightarrow)$ Connectivity of $A$ and $B$ gives $\pi_{0}(A\{T_{1},\dots,T_{n}\})\simeq\pi_{0}(A)[T_{1},\dots,T_{n}]$ \cite[Not. B.1.1.2]{SAG}, so applying $\pi_{0}(-)$ to the factorization of (a) presents $\pi_{0}(A)\to\pi_{0}(B)$ as $\pi_{0}(A)\to\pi_{0}(A)[T_{1},\dots,T_{n}]\to\pi_{0}(B)$ with the second map integral.

    $(\Leftarrow)$ The classifying map $\pi_{0}(A)[T_{1},\dots,T_{n}]\to\pi_{0}(B)$ of (b) is the datum of $n$ elements of $\pi_{0}(B)$. Lifting these to points of $\Omega^{\infty} B$ and invoking the universal property of the free algebra of \Cref{def: free and flat algebras}, furnished by the free--forgetful adjunction, 
    $$\Mor_{\CAlg_{A}}(A\{T_{1},\dots,T_{n}\},B)\simeq(\Omega^{\infty} B)^{n},$$ 
    produces a factorization $A\to A\{T_{1},\dots,T_{n}\}\to B$. On $\pi_{0}$ this recovers the map $\pi_{0}(A)[T_{1},\dots,T_{n}]\to\pi_{0}(B)$, which is integral, so $\pi_{0}(A\{T_{1},\dots,T_{n}\})\to\pi_{0}(B)$ is integral.
\end{proof}
With the equivalence in hand, we define the notion of a morphism of finite expansion between connective $\EE_{\infty}$-rings.
\begin{definition}\label{def: finite expansion appendix}
    Let $A\to B$ be a morphism of connective $\EE_{\infty}$-algebras in $\Sp$. Then $A\to B$ is \emph{of finite expansion} if it satisfies the equivalent conditions of \Cref{lem: finite expansion static detection}.
\end{definition}
\begin{proposition}\label{prop: finite expansion permanence}
    The morphisms of finite expansion between connective $\EE_{\infty}$-algebras in $\Sp$ satisfy the following permanence properties.
    \begin{enumerate}
        \item (Composition) If $A\to B$ and $B\to C$ are of finite expansion, then so is the composite $A\to C$.
        \item (Base Change) If $A\to B$ is of finite expansion and $A\to C$ is arbitrary, then $C\to B\otimes_{A} C$ is of finite expansion.
        \item (Left Cancellation) If the composite $A\to B\to C$ and $A\to B$ are of finite expansion, then so is $B\to C$.
    \end{enumerate}
\end{proposition}
\begin{proof}
    By \Cref{lem: finite expansion static detection}, each assertion reduces to the corresponding statement for the induced morphism on static rings. For base change, $\pi_{0}(B\otimes_{A}C)\simeq\pi_{0}(B)\otimes_{\pi_{0}(A)}\pi_{0}(C)$ for connective $\EE_{\infty}$-rings $A,B,C$, so that $\pi_{0}(C)\to\pi_{0}(B\otimes_{A}C)$ is the base change of $\pi_{0}(A)\to\pi_{0}(B)$ along $\pi_{0}(A)\to\pi_{0}(C)$ \cite[Cor. 7.2.1.23]{HA}. Composition and left cancellation are verified in \cite[Lem. 1.3 (1)]{FiniteExpansion} and base change in \cite[Lem. 1.3 (3)]{FiniteExpansion}. 
\end{proof}
The finiteness conditions of \Cref{def: finiteness conditions of morphisms} relate to finite expansion as follows.
\begin{lemma}\label{lem: finiteness implies finite expansion}
    Let $f:A\to B$ be a morphism in $\CAlg^{\cn}$. If $f$ is of finite presentation, almost of finite presentation, or of finite type, then $f$ is of finite expansion.
\end{lemma}
\begin{proof}
    By \Cref{lem: finite expansion static detection}, it suffices to show that $\pi_{0}(f)$ is of finite expansion in each case. Finite presentation implies almost of finite presentation, which in turn implies finite type \cite[\S 4.1]{SAG}, so under any of the three hypotheses $\pi_{0}(f):\pi_{0}(A)\to\pi_{0}(B)$ is a finite type morphism of static rings. Such a morphism factors as a surjection $\pi_{0}(A)[T_{1},\dots,T_{n}]\twoheadrightarrow\pi_{0}(B)$ out of a polynomial algebra, and a surjection is integral, so $\pi_{0}(f)$ is of finite expansion.
\end{proof}
\subsection{Pseudo-lisse morphisms}\label{subsec: pseudo-lisse morphisms}
Finite expansion constrains a morphism only through its effect on $\pi_{0}$ and leaves its $\Tor$-amplitude uncontrolled -- the morphism $A\to A\{T_{1},\dots,T_{n}\}$ is of finite expansion yet need not be of finite $\Tor$-amplitude when $n\geq 1$. The last class considered in this appendix restores that control. 

In the animated setting, a morphism of animated rings almost of finite presentation and of finite $\Tor$-amplitude admits a factorization $A\to A[T_{1},\dots,T_{n}]\to B$ with $B$ perfect over $A[T_{1},\dots,T_{n}]$ \cite[Lem. 2.7.12 (iv)]{HansenMannCLL}. In the spectral setting, choosing representatives in $\Omega^{\infty}B$ of $n$ elements of $\pi_{0}(B)$ produces a morphism of $A$-algebras $A\{T_{1},\dots,T_{n}\}\to B$ but need not produce a morphism of $A$-algebras $A[T_{1},\dots,T_{n}]\to B$. Since the free $\EE_{\infty}$-algebra may fail this $\Tor$-amplitude condition, the animated factorization theorem does not directly carry over. The definition below instead imposes a factorization through a flat algebra with perfect second leg, at the cost of a slightly more restricted generality than that of morphisms almost of finite presentation and of finite $\Tor$-amplitude.
\begin{definition}\label{def: pseudo-lisse appendix}
    Let $A\to B$ be a morphism in $\CAlg^{\cn}$.
    \begin{enumerate}
        \item $A\to B$ is \emph{pseudo-lisse} if it admits a factorization
        $$A\to C\to B$$
        in which $A\to C$ is flat and almost of finite presentation, and $B$ is perfect as a $C$-module.
        \item $A\to B$ is \emph{elementary pseudo-lisse} if it admits a factorization
        $$A\to A[T_{1},\dots,T_{n}]\to B$$
        for some $n\geq 0$ in which $B$ is perfect as an $A[T_{1},\dots,T_{n}]$-module.
    \end{enumerate}
\end{definition}
Elementary pseudo-lisse and pseudo-lisse morphisms satisfy the following. 
\begin{proposition}\label{prop: properties of pseudo-lisse morphisms}
    \begin{enumerate}
        \item Every elementary pseudo-lisse morphism is pseudo-lisse. 
        \item Every pseudo-lisse morphism is almost of finite presentation and of finite $\Tor$-amplitude. In particular, every pseudo-lisse morphism is of finite expansion. 
        \item The elementary pseudo-lisse morphisms satisfy the following permanence properties: 
        \begin{itemize}
            \item [(3.1)] (Composition) If $A\to B$ and $B\to C$ are elementary pseudo-lisse, then so is the composite $A\to C$. 
            \item [(3.2)] (Base Change) If $A\to B$ is elementary pseudo-lisse, then for any morphism $A\to C$ the induced morphism $C\to B\otimes_{A}C$ is elementary pseudo-lisse. 
        \end{itemize}
        \item The pseudo-lisse morphisms satisfy the following permanence property: 
        \begin{itemize}
            \item [(4.1)] (Base Change) If $A\to B$ is pseudo-lisse, then for any morphism $A\to C$ the induced morphism $C\to B\otimes_{A}C$ is pseudo-lisse. 
        \end{itemize}
    \end{enumerate}
\end{proposition}
\begin{proof}[Proof of (1)]
    Flatness follows from \cite[Rmk. 5.4.1.2]{SAG}. Being almost of finite presentation is preserved by base change \cite[Rmk. 7.2.4.28]{HA}, which reduces the claim to the universal case $\Sph\to\Sph[T_{1},\dots,T_{n}]$, and this follows from the Hilbert Basis Theorem for connective $\EE_{\infty}$-rings \cite[Prop. 7.2.4.31]{HA}: $\Sph$ and $\Sph[T_{1},\dots,T_{n}]$ are Noetherian in the sense of \cite[Def. 7.2.4.30]{HA} as both are connective, on $\pi_{0}$ these are $\ZZ$ and $\ZZ[T_{1},\dots,T_{n}]$, which are Noetherian (and in particular coherent) as static commutative rings, and each higher homotopy group of either is finitely generated -- equivalently, as $\pi_{0}$ is Noetherian, finitely presented -- over $\pi_{0}$. For $\Sph$ this is Serre's celebrated finiteness of the stable stems in positive degrees. For $\Sph[T_{1},\dots,T_{n}]=\Sigma^{\infty}_{+}\NN^{n}$ one has $\pi_{k}(\Sph[T_{1},\dots,T_{n}])\simeq\pi_{k}(\Sph)\otimes_{\ZZ}\ZZ[T_{1},\dots,T_{n}]$, which is finitely generated over $\ZZ[T_{1},\dots,T_{n}]$ because $\pi_{k}(\Sph)$ is finite when $k\geq1$.
\end{proof}
\begin{proof}[Proof of (2)]
    We first show $A\to B$ is almost of finite presentation. By \cite[Cor. 7.4.3.19]{HA}, being almost of finite presentation is preserved by composition and $A\to C$ is almost of finite presentation by assumption, so it suffices to show that $C\to B$ is almost of finite presentation. Since $B$ is perfect over $C$, the morphism $C\to B$ is finite and $B$ is pseudocoherent as a $C$-module. A finite morphism is almost of finite presentation if and only if its target is pseudocoherent over the source \cite[Cor. 5.2.2.2]{SAG}, so $C\to B$ is almost of finite presentation. 

    For the finite $\Tor$-amplitude, $A\to C$ is flat and $B$ is of finite $\Tor$-amplitude over $C$, being perfect \cite[Prop. 7.2.4.23]{HA}, so the composite $A\to B$ is of finite $\Tor$-amplitude \cite[Lem. 6.1.1.6]{SAG}. 
\end{proof}
\begin{proof}[Proof of (3)]
    Let $A\to A[T_{1},\dots,T_{n}]\to B$ and $B\to B[T_{1},\dots,T_{m}]\to C$ be factorizations of $A\to B$ and $B\to C$ as elementary pseudo-lisse morphisms. We claim that the composite factors as $A\to A[T_{1},\dots,T_{m+n}]\to C$ where $C$ is perfect over $A[T_{1},\dots,T_{m+n}]$. First note that $A[T_{1},\dots,T_{m+n}]\otimes_{A[T_{1},\dots,T_{n}]}B\simeq B[T_{1},\dots,T_{m}]$, providing a factorization of $A$-algebras $A\to A[T_{1},\dots,T_{m+n}]\to C$. Since scalar extension preserves compact objects, $B[T_{1},\dots,T_{m}]$ is perfect over $A[T_{1},\dots,T_{m+n}]$, hence $B[T_{1},\dots,T_{m}]$ is contained in the minimal full subcategory of $A[T_{1},\dots,T_{m+n}]$-modules containing the unit and closed under finite colimits and retracts. 
    
    Since $C$ is perfect over $B[T_{1},\dots,T_{m}]$, restriction of scalars exhibits it as an object of the minimal full subcategory of $A[T_{1},\dots,T_{m+n}]$-modules containing $B[T_{1},\dots,T_{m}]$ and closed under finite colimits and retracts. As $B[T_{1},\dots,T_{m}]$ is perfect over $A[T_{1},\dots,T_{m+n}]$, this subcategory is contained in the perfect $A[T_{1},\dots,T_{m+n}]$-modules, so $C$ is perfect over $A[T_{1},\dots,T_{m+n}]$. 

    For closure under base change, applying $C\otimes_{A}(-)$ to a factorization of $A\to B$ as an elementary pseudo-lisse morphism yields a factorization of $C\to B\otimes_{A}C$ as an elementary pseudo-lisse morphism by functoriality of the polynomial algebra and the permanence of perfectness under base change.  
\end{proof}
\begin{proof}[Proof of (4)]
    This follows from the base change argument of (3), instead using that flatness and being almost of finite presentation are preserved by base change along general maps \cite[Prop. 7.2.2.16 \& Rmk. 7.2.4.28]{HA}. 
\end{proof}
\printbibliography
\end{document}